\chapter{流形与微分形式}
在第 VIII 章中，我们了解了普法夫形式，并看到一次微分形式与线积分理论密切相关。在本章中，我们将讨论线积分的高维类似物，即更高次微分形式在 $ \mathbb{R}^n $ 的某些子流形上的积分。因此本章将讨论微分形式的理论。

在第 1 节中，我们推广已有的关于流形的知识。特别地，我们探讨给定流形的子流形这一概念，并引入带边流形。

在第 2 节中，我们汇集多重线性代数中所需的结果。它们构成微分形式理论的代数基础：在第 3 节中，我们讨论 $ \mathbb{R}^n $ 的开子集上的微分形式。在第 4 节中，我们使这一理论整体化，然后讨论流形的可定向性。

因为我们始终考虑欧几里得空间的子流形，所以可以自然地赋予它们黎曼度量。在第 5 节中，我们将更仔细地考察这一附加结构，并说明黎曼几何的若干基本事实。为适应物理学的需要，我们也讨论半黎曼度量；在例子中，我们将始终局限于闵可夫斯基空间。

作为本章的结尾，第 6 节建立微分形式理论与经典向量分析之间的联系。特别地，我们研究梯度、散度与旋度这些算子，并导出它们的基本性质。我们给出它们的局部坐标表示，并在几个重要的例子中显式计算这些表示。

在第 2 节中，该节的其他内容属于线性代数，我们还引入霍奇星算子，后面几节需要用它来定义余微分。然后我们就能把向量分析的各种算子统一到霍奇演算的语言之中。初次阅读时可以跳过这部分材料：出于这一原因，我们把任何用到霍奇理论的材料都留到每节末尾再讨论。

在全书中，我们都限于 $ \mathbb{R}^n $ 的子流形。不过，除切空间的定义之外，我们组织所有证明的方式使它们对抽象流形仍然成立，或者稍加修改即可成立。因此第 XI 章与第 XII 章给出了微分拓扑与微分几何的初步介绍；尽管它们有时缺少一般理论那样的完整优美，我们所考虑的众多例子确实为这一课题的进一步学习奠定了坚实的基础。

\section{子流形}

本节中，

- $ M $ 是 $ m $ 维流形，$ N $ 是 $ n $ 维流形。

更确切地说，这意味着 $ M $ 是某个 $ \mathbb{R}^{\overline{m}} $ 的 $ m $ 维 $ C^\infty $ 子流形，其中 $ \overline{m}\ge m $；对 $ N $ 有类似的陈述。

为简单起见并突出本质，我们限于研究光滑映射。特别地，我们总把\textbf{微分同胚}理解为 $ C^\infty $ 微分同胚，并置

\[
\operatorname{Diff}(M,N):=\operatorname{Diff}^\infty(M,N).
\]

然而，下文中凡有证明之处，结论对 $ C^k $ 流形与 $ C^k $ 映射同样成立，其中必要时须对 $ k\in\mathbb{N}^\times $ 作适当限制。我们通常把这些调整放在注\footnote{在小字排印的题为“正则性”的各节中。}中，并把它们的验证留给你。

\needspace{3\baselineskip}\subsection*{定义与基本性质}

设 $ 0\le\ell\le m $。$ M $ 的子集 $ L $ 称为 $ M $ 的（$ \ell $ 维）子流形，若对每个 $ p\in L $，存在 $ M $ 在 $ p $ 附近的坐标卡 $ (\varphi,U) $，使得\footnote{为避免麻烦的特殊情形，我们把空集看作 $ \ell $ 维子流形，对每个 $ \ell\in\{0,\ldots,m\} $ 皆然（见第 VII.9 节）。}

\[
\varphi(U\cap L)=\varphi(U)\cap(\mathbb{R}^\ell\times\{0\}).
\]

每个这样的坐标卡称为 $ M $ 关于 $ L $ 的\textbf{子流形坐标卡}。数 $ m-\ell $ 称为 $ L $ 在 $ M $ 中的\textbf{余维数}。

\bookfigure[0.68]{figures/16c48a90a5be9eee.png}

显然，这个定义直接推广了 $ \mathbb{R}^m $ 的子流形的概念。

在子流形的语境中，浸入起着重要作用。它们将类比第 VII.9 节给出的定义来引入。

设 $ k\in\mathbb{N}^\times\cup\{\infty\} $。则 $ f\in C^k(M,N) $ 称为 \textbf{$ C^k $ 浸入}，若对每个 $ p\in M $，$ T_p f:T_pM\to T_{f(p)}N $ 是单射。我们称 $ C^k $ 浸入 $ f $ 为 $ M $ 在 $ N $ 中的 \textbf{$ C^k $ 嵌入}，若 $ f $ 是从 $ M $ 到 $ f(M) $ 的同胚（其中 $ f(M) $ 自然地赋予 $ N $ 的相对拓扑）。代替 $ C^\infty $ 浸入［或 $ C^\infty $ 嵌入］，我们简称\textbf{浸入}［或\textbf{嵌入}］。

\begin{remark}

\begin{items}
\item

若 $ L $ 是 $ M $ 的 $ \ell $ 维子流形，且 $ M $ 是 $ N $ 的子流形，则 $ L $ 是 $ N $ 的 $ \ell $ 维子流形。

\proofhead
设 $ p\in L $，并设 $ (\varphi,U) $ 是 $ M $ 关于 $ L $ 在 $ p $ 附近的子流形坐标卡。又设 $ (\psi,V) $ 是 $ N $ 关于 $ M $ 在 $ p $ 附近的子流形坐标卡。我们还可假设 $ U=V\cap M $。令 $ X:=\varphi(U)\subset\mathbb{R}^m $，$ Y:=\operatorname{pr}\circ\psi(V)\subset\mathbb{R}^m $，其中 $ \operatorname{pr}:\mathbb{R}^m\times\mathbb{R}^{n-m}\to\mathbb{R}^m $ 表示标准投影，则
\[
\chi:=\operatorname{pr}\circ\psi\circ\varphi^{-1}\in\operatorname{Diff}(X,Y).
\]

现由

\[
\Phi(y,z):=(\chi^{-1}(y),z)\qquad \text{对 } (y,z)\in Y\times\mathbb{R}^{n-m}
\]

定义 $ \Phi\in\operatorname{Diff}(Y\times\mathbb{R}^{n-m},X\times\mathbb{R}^{n-m}) $，并置 $ \Psi:=\Phi\circ\psi $。则 $ \Psi(V) $ 在 $ \mathbb{R}^n $ 中开，且 $ \Psi\in\operatorname{Diff}(V,\Psi(V)) $，并有

\[
\Psi(V\cap L)=(\varphi(U\cap L)\times\{0\})\cap(\mathbb{R}^\ell\times\{0\})=\Psi(V)\cap(\mathbb{R}^\ell\times\{0\})\subset\mathbb{R}^n,
\]

正如容易验证的那样。因此 $ (\Psi,V) $ 是 $ N $ 关于 $ L $ 在 $ p $ 附近的子流形坐标卡。\nobreak\hspace{0.35em}■

\item

因为 $ \mathbb{R}^{\overline{m}}=\mathbb{R}^{\overline{m}}\times\{0\}\subset\mathbb{R}^n $ 对 $ n\ge\overline{m} $ 是 $ \mathbb{R}^n $ 的子流形，故由 (a) 推出：对每个 $ n\ge\overline{m} $，$ M $ 是 $ \mathbb{R}^n $ 的 $ m $ 维子流形。这表明，只要只关心 $ M $ 的“内部性质”，即只借助 $ M $ 的坐标卡与切空间描述、不依赖于 $ M $ 如何“处于”周围空间中的那些性质，$ M $ 的“周围空间” $ \mathbb{R}^{\overline{m}} $ 就不起重要作用。\footnote{在第 4 节将会清楚，切空间也有一种“内部”刻画。} 然而，$ M $ 如何处于 $ \mathbb{R}^{\overline{m}} $ 中确实有影响，例如在定义法丛 $ T^\perp M $ 时。

\item

设 $ L $ 是 $ M $ 的子流形。对 $ M $ 关于 $ L $ 的子流形坐标卡 $ (\varphi,U) $，置

\[
(\varphi_L,U_L):=(\varphi\mid U\cap L,U\cap L).
\]

则 $ (\varphi_L,U_L) $ 是 $ L $ 的坐标卡，其中 $ \varphi(U_L) $ 理解为 $ \mathbb{R}^\ell $ 的开子集，即 $ \mathbb{R}^\ell\times\{0\}\subset\mathbb{R}^m $ 与 $ \mathbb{R}^\ell $ 等同。
若 $ \mathcal{A}:=\{(\varphi_\lambda,U_\lambda);\lambda\in\Lambda\} $ 是 $ M $ 关于 $ L $ 的一族子流形坐标卡，使得 $ L $ 被坐标片（坐标区域）$ \{U_\lambda;\lambda\in\Lambda\} $ 覆盖，则 $ \{(\varphi_{\lambda,L},U_{\lambda,L});\lambda\in\Lambda\} $ 是 $ L $ 的图册，称为\textbf{由 $ \mathcal{A} $ 诱导的图册}。

\proofhead
简单的验证留给你。\nobreak\hspace{0.35em}■
\item

设 $ L $ 与 $ K $ 分别是 $ M $ 与 $ N $ 的 $ \ell $ 维与 $ k $ 维子流形。则 $ L\times K $ 是流形 $ M\times N $ 的 $ (\ell+k) $ 维子流形，而 $ M\times N $ 是 $ (m+n) $ 维的。

\proofhead
这直接由定义推出。我们再次把证明留给你。\footnote{见习题 VII.9.4。}\nobreak\hspace{0.35em}■
\item

设 $ L $ 是 $ M $ 的子流形。则

\[
i:L\to M,\qquad p\mapsto p
\]

是嵌入，称为 $ L $ 在 $ M $ 中的\textbf{自然嵌入}；我们把它写作 $ i:L\hookrightarrow M $。对 $ p\in L $，我们在单射 $ T_p i $ 下把 $ T_p L $ 与其在 $ T_p M $ 中的像等同，即把 $ T_p L $ 看作 $ T_p M $ 的向量子空间：$ T_p L\subset T_p M $。

\proofhead
设 $ (\varphi,U) $ 是 $ M $ 关于 $ L $ 的子流形坐标卡。则 $ i $ 有局部表示
\[
\varphi\circ i\circ\varphi_L^{-1}:\varphi_L(U_L)\to\varphi(U),\qquad x\mapsto(x,0).
\]

现在结论显然。\nobreak\hspace{0.35em}■

\item

若 $ f:M\to N $ 是浸入，则 $ m\le n $。

\item

设 $ L $ 是 $ M $ 的 $ \ell $ 维子流形，并设 $ f $ 属于 $ \operatorname{Diff}(M,N) $。则 $ f(L) $ 是 $ N $ 的 $ \ell $ 维子流形。

\proofhead
简单的验证留给你。\nobreak\hspace{0.35em}■
\item

$ M $ 的每个开子集都是 $ M $ 的 $ m $ 维子流形。

\item

若 $ (\varphi,U) $ 是 $ M $ 的坐标卡，则 $ \varphi:U\to\mathbb{R}^m $ 是嵌入，且 $ \varphi $ 是从 $ U $ 到 $ \varphi(U) $ 的微分同胚。

\item

设 $ L $ 与 $ K $ 分别是 $ M $ 与 $ N $ 的子流形，$ i_L:L\hookrightarrow M $ 与 $ i_K:K\hookrightarrow N $ 是它们各自的自然嵌入。设 $ k\in\mathbb{N}\cup\{\infty\} $，$ f\in C^k(M,N) $ 且 $ f(L)\subset K $。则 $ f $ 在 $ L $ 上的限制满足

\[
f\mid L:=f\circ i_L\in C^k(L,K),
\]

并且下面的图交换。

\end{items}

\bookfigure[0.95]{figures/f888b417202e7f1d.png}

在 $ T_p i_L $ 下把 $ T_p L $ 与其在 $ T_p M $ 中的像等同，即以典范方式把 $ T_p L $ 看作 $ T_p M $ 的向量子空间，我们特别有 $ T_p(f\mid L)=(T_p f)\mid T_p L $。

\proofhead
这由对例 VII.10.10(b) 的证明作显然的修改得到，而本陈述推广了该例。\nobreak\hspace{0.35em}■

\begin{items}
\item[（k）]

（正则性）当 $ M $ 是 $ C^k $ 流形（$ k\in\mathbb{N}^\times $）时，类似的定义与陈述成立。此时 $ L $ 也是 $ C^k $ 流形，且自然包含 $ i:L\hookrightarrow M $ 属于 $ C^k $ 类。\nobreak\hspace{0.35em}■

\end{items}

下一个定理是命题 VII.9.10 的推广，它表明我们可以用嵌入生成子流形。

\end{remark}

\begin{theorem}

\begin{subitems}
\item

设 $ f:M\to N $ 是浸入。则 $ f $ 局部地是嵌入，即对 $ M $ 中的每个 $ p $，存在邻域 $ U $，使得 $ f\mid U $ 是嵌入。

\item

若 $ f:M\to N $ 是嵌入，则 $ f(M) $ 是 $ N $ 的 $ m $ 维子流形，且 $ f $ 是从 $ M $ 到 $ f(M) $ 的微分同胚。

\proofhead
(i) 设 $ p\in M $，并设 $ (\varphi,U_0) $ 与 $ (\psi,V) $ 分别是 $ M $ 在 $ p $ 附近与 $ N $ 在 $ f(p) $ 附近的坐标卡，使得 $ f(U_0)\subset V $。则
\[
f_{\varphi,\psi}:=\psi\circ f\circ\varphi^{-1}:\varphi(U_0)\to\psi(V)
\]

由注 1.1(i) 是浸入。由浸入定理（定理 VII.9.7），存在 $ \varphi(p) $ 在 $ \varphi(U_0) $ 中的开邻域 $ X $，使得 $ f_{\varphi,\psi}(X) $ 是 $ \mathbb{R}^n $ 的 $ m $ 维子流形。于是 $ \psi\in\operatorname{Diff}(V,\psi(V)) $ 与注 1.1(g) 推出：$ f(U) $（其中 $ U:=\varphi^{-1}(X) $）是 $ N $ 的 $ m $ 维子流形。
通过适当缩小 $ X $，注 VII.9.9(d) 表明 $ f_{\varphi,\psi} $ 是从 $ X=\varphi(U) $ 到 $ f_{\varphi,\psi}(X)=\psi\circ f(U) $ 的微分同胚。因此 $ f $ 是从 $ U $ 到 $ f(U) $ 的微分同胚，其中 $ f(U) $ 赋予由 $ N $ 诱导的拓扑。因此 $ f\mid U $ 是嵌入。

\end{subitems}

(ii) 设 $ f $ 是嵌入。对 $ q\in f(M) $，设 $ (\psi,V) $ 是 $ N $ 在 $ q $ 附近的坐标卡，$ (\varphi,U) $ 是 $ M $ 在 $ p:=f^{-1}(q) $ 附近的坐标卡，且 $ f(U)\subset V $。因为 $ f $ 是从 $ M $ 到 $ f(M) $ 的同胚，我们知道 $ f(U) $ 在 $ f(M) $ 中开。因此可假设 $ f(U)=f(M)\cap V $。现在由 (i) 的证明推出 $ f(M)\cap V $ 是 $ N $ 的 $ m $ 维子流形。由于这对每个 $ q\in f(M) $ 成立，我们得出 $ f(M) $ 是 $ N $ 的 $ m $ 维子流形。

由 (i)，$ f $ 是从 $ M $ 到 $ f(M) $ 的局部微分同胚。因为 $ f $ 是同胚，推出 $ f\in\operatorname{Diff}(M,f(M)) $。\nobreak\hspace{0.35em}■

\end{theorem}

由注 VII.9.9(c) 我们知道，单浸入的像一般不是子流形。下面的定理给出一个简单的充分条件，用以判断单浸入是否为嵌入。

\begin{theorem}

设 $ M $ 紧致，$ f:M\to N $ 是单浸入。则 $ f $ 是嵌入，$ f(M) $ 是 $ N $ 的 $ m $ 维子流形，且 $ f\in\operatorname{Diff}(M,f(M)) $。

\proofhead
因为 $ M $ 紧致且 $ f(M) $ 是度量空间，双射连续映射 $ f:M\to f(M) $ 是同胚（见习题 III.3.3）。现在结论由定理 1.2 推出。\nobreak\hspace{0.35em}■
\end{theorem}

\begin{remark}

（正则性）设 $ k\in\mathbb{N}^\times $。则当 $ M $ 与 $ N $ 是 $ C^k $ 流形且 $ f $ 属于 $ C^k $ 类时，定理 1.2 与 1.3 的相应版本仍然成立。\nobreak\hspace{0.35em}■

\end{remark}

\begin{example}

\begin{items}
\item

设 $ 1\le\ell<m $，并令 $ (x,y) $ 表示 $ \mathbb{R}^{\ell+1}\times\mathbb{R}^{m-\ell}=\mathbb{R}^{m+1} $ 的一般点。则

\[
L_y:=\sqrt{1-|y|^2}\,S^\ell\times\{y\}
\]

是 $ m $ 维球面 $ S^m $ 的 $ \ell $ 维子流形，对每个 $ y\in\mathbb{B}^{m-\ell} $ 成立。它与 $ S^\ell $ 微分同胚。点 $ p\in L_y $ 处的切空间满足

\begin{equation}
T_p L_y=T_p S^m\cap(p,\mathbb{R}^{\ell+1}\times\{0\})\subset T_p\mathbb{R}^{m+1}.\tag{1.1}
\end{equation}
\end{items}

\bookfigure[0.28]{figures/ca183af51d821ba6.png}

\proofhead
对 $ y\in\mathbb{B}^{m-\ell} $，映射
\begin{equation}
F_y:\mathbb{R}^{\ell+1}\to\mathbb{R}^{m+1},\qquad x\mapsto(\sqrt{1-|y|^2}\,x,y)\tag{1.2}
\end{equation}

是光滑浸入。因为 $ S^\ell $ 与 $ S^m $ 分别是 $ \mathbb{R}^{\ell+1} $ 与 $ \mathbb{R}^{m+1} $ 的子流形，且 $ F_y(S^\ell)\subset S^m $，注 1.1(j) 取 $ i_\ell:S^\ell\hookrightarrow\mathbb{R}^{\ell+1} $ 给出

\begin{equation}
f_y:=F_y\mid S^\ell=F_y\circ i_\ell\in C^\infty(S^\ell,S^m).\tag{1.3}
\end{equation}

显然 $ f_y $ 是单射，且注 VII.10.9(b) 的链式法则推出

\[
T_p f_y=T_p F_y\circ T_p i_\ell\qquad \text{对 } p\in S^\ell.
\]

因此 $ T_p f_y $ 是单射（见习题 I.3.3），即 $ f_y $ 是浸入。因为 $ S^\ell $ 紧致，定理 1.3 表明 $ L_y=f_y(S^\ell) $ 是 $ S^m $ 的 $ \ell $ 维子流形，且与 $ S^\ell $ 微分同胚。于是 (1.1) 是 (1.2) 与 (1.3) 的简单推论。\nobreak\hspace{0.35em}■

\begin{items}
\item[（b）]

（\textbf{旋转环面型超曲面}）设

\[
\gamma:S^1\to(0,\infty)\times\mathbb{R},\qquad t\mapsto(\rho(t),\sigma(t))
\]

是单浸入，从而由定理 1.3 是嵌入。又设 $ i:S^m\hookrightarrow\mathbb{R}^{m+1} $，并定义

\[
f:S^m\times S^1\to\mathbb{R}^{m+1}\times\mathbb{R},\qquad (q,t)\mapsto(\rho(t)i(q),\sigma(t)).
\]

则 $ f $ 是嵌入，且

\[
T^{m+1}:=f(S^m\times S^1)
\]

是 $ \mathbb{R}^{m+2} $ 中的超曲面，它与 $ S^m\times S^1 $ 微分同胚。
当 $ m=0 $ 时，集合 $ T^1 $ 由闭光滑曲线 $ \gamma(S^1) $ 的两份拷贝组成，该曲线没有自交点\footnote{这里及以下，“曲线”指一维流形（见注 1.19(a)）。}，并关于 $ y $ 轴对称。

\end{items}

\bookfigure[0.40]{figures/10e6c2e9a5f96b76.png}

当 $ m=1 $ 时，$ T^2 $ 是 $ \mathbb{R}^3 $ 中由子午曲线

\[
\Gamma:=\{(\rho(t),0,\sigma(t));t\in S^1\}
\]

绕 $ z $ 轴旋转生成的旋转曲面（见例 VII.9.11(e)）。$ T^2 $“是一个 2 环面”，即它与 $ \mathsf{T}^2:=S^1\times S^1 $ 微分同胚。特别地，例 VII.9.11(f) 中的 2 环面 $ \mathsf{T}^2_{a,r} $ 与 $ \mathsf{T}^2 $ 微分同胚。

\bookfigure[0.28]{figures/9e60a8d3a3249c83.png}

在一般情形，我们称 $ T^{m+1} $ 为\textbf{旋转环面型超曲面}。

\proofhead
由例 VII.9.5(b)，$ S^m $ 与 $ S^1 $ 分别是 $ m $ 维与 1 维流形。因此 $ S^m\times S^1 $ 是 $ (m+1) $ 维流形。

设 $ (\varphi\times\psi,U\times V) $ 是 $ S^m\times S^1 $ 的乘积坐标卡\footnote{即 $ (\varphi,U) $ 与 $ (\psi,V) $ 分别是 $ S^m $ 与 $ S^1 $ 的坐标卡，且 $ \varphi\times\psi(q,t):=(\varphi(q),\psi(t)) $。}。因为 $ \gamma $ 是浸入，它关于 $ \psi $（以及 $ \mathbb{R}^2 $ 的平凡坐标卡 $ \operatorname{id}_{\mathbb{R}^2} $）的局部表示，即 $ \gamma_\psi=(r,s) $，其中 $ r:=\rho\circ\psi^{-1} $、$ s:=\sigma\circ\psi^{-1} $，满足

\begin{equation}
(\dot{r}(y),\dot{s}(y))\ne(0,0)\qquad \text{对 } y\in\psi(V).\tag{1.4}
\end{equation}

此外，$ f $ 关于 $ \varphi\times\psi $ 的局部表示具有形式

\[
f_{\varphi\times\psi}(x,y)=(r(y)g(x),s(y))\qquad \text{对 } (x,y)\in\varphi(U)\times\psi(V),
\]

其中 $ g:=i\circ\varphi^{-1} $ 是属于 $ \varphi $ 的 $ S^m $ 的参数化。由此推出

\begin{equation*}
[\partial f_{\varphi\times\psi}(x,y)]=\begin{bmatrix} r(y)\partial g(x) & \dot{r}(y)g(x)\\ \cdots\cdots\cdots\cdots & \cdots\cdots\cdots\cdots\\ 0 & \dot{s}(y)\end{bmatrix}\in\mathbb{R}^{(m+2)\times(m+1)}.
\end{equation*}

因为 $ r(y)>0 $ 且 $ \partial g(x) $ 是单射，该矩阵的前 $ m $ 列线性无关。若 $ \dot{s}(y)\ne0 $，则矩阵的秩为 $ m+1 $。若 $ \dot{s}(y)=0 $，则由 (1.4) 得 $ \dot{r}(y)\ne0 $。由 $ |g(x)|^2=(g(x)\mid g(x))=1 $（$ x\in\varphi(U) $）推出 $ (g(x)\mid\partial_j g(x))=0 $，其中 $ 1\le j\le m $、$ x\in\varphi(U) $。这表明此时矩阵的秩仍为 $ m+1 $。因此 $ f $ 是浸入。

现在对某个 $ (y,s)\in T^{m+1} $ 考虑方程 $ f(q,t)=(y,s) $。由关系 $ \rho(t)i(q)=y $ 与 $ |i(q)|=1 $ 推出 $ \rho(t)=|y| $。因为 $ \gamma $ 是单射，恰有一个 $ t\in S^1 $ 使 $ (\rho(t),\sigma(t))=(|y|,s) $。同样，恰有一个 $ q\in S^m $ 使 $ i(q)=y/|y| $。因此方程 $ (\rho(t)i(q),\sigma(t))=(y,s) $（其中 $ (y,s) $ 如上）有唯一解（因为 $ y=|y|(y/|y|) $）。因此 $ f $ 是单浸入，

……$ S^m \times S^1 $ 到 $ \mathbb{R}^{m+2} $ 中。现在所有断言都可由定理 1.3 推出，因为 $ S^m \times S^1 $ 是紧的。 \nobreak\hspace{0.35em}■

\begin{items}
\item[（c）]

设 $ L $ 与 $ M $ 是 $ N $ 的子流形，且 $ L \subset M $。则 $ L $ 是 $ M $ 的子流形。

\proofhead
因为 $ \operatorname{id}_N \in \operatorname{Diff}(N, N) $，我们知道 $ i := \operatorname{id}_N|_L $ 是 $ L $ 在 $ N $ 中的浸入，且 $ i(L) \subset M $。因此由注 1.1(j) 可知，$ i $ 是 $ L $ 在 $ M $ 中的双射浸入。因为 $ L $ 与 $ M $ 都带有由 $ N $ 诱导的拓扑，并且 $ M $ 在 $ L $ 上诱导出相同的拓扑，所以作为微分同胚的限制，$ i $ 是拓扑的。因此 $ i $ 是嵌入，断言由定理 1.2 推出。 \nobreak\hspace{0.35em}■
\item[（d）]

设 (b) 的假设在 $ m = 1 $ 时得到满足。则对每个 $ (q_0, t_0) \in S^1 \times S^1 $，映射

\[
f(\cdot, t_0): S^1 \to \mathbb{R}^3
\]

与

\[
f(q_0, \cdot): S^1 \to \mathbb{R}^3
\]

的像都是 $ T^2 $ 的一维子流形，并且微分同胚于 $ S^1 $（从而是“圆”）。

\end{items}

\bookfigure[0.40]{figures/d057479a9d232051.png}

\proofhead
因为作为嵌入的限制，$ f(\cdot, t_0) $ 与 $ f(q_0, \cdot) $ 本身都是嵌入，所以 $ f(S^1, t_0) $ 与 $ f(q_0, S^1) $ 是 $ \mathbb{R}^3 $ 中微分同胚于 $ S^1 $ 的子流形，并且它们位于 $ T^2 $ 中。断言现在由 (c) 推出。 \nobreak\hspace{0.35em}■
\end{example}

\needspace{3\baselineskip}\subsection*{淹没}

设 $ f \in C^1(M, N) $。若 $ T_p f $ 是满射，则称 $ p \in M $ 是 $ f $ 的\textbf{正则点}。否则称 $ p $ 是\textbf{奇点}。若每个 $ p \in f^{-1}(q) $ 都是正则点，则称点 $ q \in N $ 是 $ f $ 的\textbf{正则值}。若 $ M $ 的每个点都是正则的，则称 $ f $ 是\textbf{正则映射}或\textbf{淹没}。

这些定义推广了第 VII.8 节引入的概念。

\begin{remark}

\begin{items}
\item

若 $ p $ 是 $ f $ 的正则点，则 $ m \ge n $。每个 $ q \in N \setminus f(M) $ 都是 $ f $ 的正则值。

\item

点 $ p \in M $ 是 $ f = (f^1, \ldots, f^n) \in C^1(M, \mathbb{R}^n) $ 的正则点，当且仅当余切向量\footnote{见第 VIII.3 节。}

\[
df^j(p) := d_p f^j = \operatorname{pr}_2 \circ T_p f^j \in T_p^* M \quad \text{对 } 1 \le j \le n
\]

线性无关。

\item

$ f \in C^1(M, \mathbb{R}) $ 的奇点也称为\textbf{临界点}。因此 $ p \in M $ 是 $ f $ 的临界点，当且仅当 $ df(p) = 0 $。\footnote{见注 VII.3.14(a)。} \nobreak\hspace{0.35em}■

\end{items}

下面这个定理把正则值定理推广到流形之间映射的情形。

\end{remark}

\begin{theorem}

（正则值） 设 $ q \in N $ 是映射 $ f \in C^\infty(M, N) $ 的正则值。则 $ L := f^{-1}(q) $ 是 $ M $ 的余维数为 $ n $ 的子流形。对 $ p \in L $，$ T_p f $ 的核是 $ T_p L $。

\proofhead
设 $ p_0 \in f^{-1}(q) $。设 $ (\varphi, U) $ 是 $ M $ 在 $ p_0 $ 附近的坐标卡，$ (\psi, V) $ 是 $ N $ 在 $ q $ 附近的坐标卡，且 $ f(U) \subset V $。则由链式法则可知，对每个 $ p \in U \cap f^{-1}(q) $，点 $ \varphi(p) $ 是局部表示
\[
f_{\varphi,\psi} := \psi \circ f \circ \varphi^{-1} \in C^\infty(\varphi(U), \mathbb{R}^n)
\]

的正则点。

换言之，$ y := \psi(q) $ 是 $ f_{\varphi,\psi} $ 的正则值。因此定理 VII.9.3 保证 $ (f_{\varphi,\psi})^{-1}(y) $ 是 $ \mathbb{R}^m $ 的 $ (m-n) $ 维子流形。从而存在 $ \mathbb{R}^m $ 的开集 $ X $ 与 $ Y $，以及 $ \Phi \in \operatorname{Diff}(X, Y) $，使得

\[
\Phi(X \cap (f_{\varphi,\psi})^{-1}(y)) = Y \cap (\mathbb{R}^{m-n} \times \{0\}).
\]

把 $ \varphi(U) $ 与 $ X $ 换成它们的交集，可以假设 $ \varphi(U) = X $。但此时 $ \varphi_1 := \Phi \circ \varphi $ 是 $ M $ 在 $ p $ 附近的坐标卡，满足

\[
\varphi_1(f^{-1}(q) \cap U) = \Phi \circ \varphi(f^{-1} \circ \psi^{-1}(y) \cap U) = \Phi((f_{\varphi,\psi})^{-1}(y) \cap X) = Y \cap (\mathbb{R}^{m-n} \times \{0\}),
\]

因此它是 $ M $ 关于 $ f^{-1}(q) $ 的子流形坐标卡。第二个断言现在可由对定理 VII.10.7 证明的显然修改得到。 \nobreak\hspace{0.35em}■

\end{theorem}

\begin{remark}

\begin{items}
\item

定理 1.7 有一个逆命题，它说 $ M $ 的每个子流形都可以局部表示为某个正则映射的纤维。更确切地说，它说：若 $ L $ 是 $ M $ 的 $ \ell $ 维子流形，则对每个 $ p \in L $，存在 $ M $ 中的邻域 $ U $ 与 $ f \in C^\infty(U, \mathbb{R}^{m-\ell}) $，使得 $ f^{-1}(0) = U \cap L $，且 $ 0 $ 是 $ f $ 的正则值。

\proofhead
设 $ (\varphi, U) $ 是 $ M $ 在 $ p $ 附近关于 $ L $ 的子流形坐标卡。则由 $ f(q) := (\varphi^{\ell+1}(q), \ldots, \varphi^m(q)) $（$ q \in U $）定义的函数属于 $ C^\infty(U, \mathbb{R}^{m-\ell}) $，并满足 $ f^{-1}(0) = U \cap L $。因为 $ \varphi $ 是微分同胚，$ 0 $ 是 $ f $ 的正则值。 \nobreak\hspace{0.35em}■
\item

（正则性） 若 $ q $ 是 $ f \in C^k(M, N) $ 的正则值（对某个 $ k \in \mathbb{N}^\times $），则 $ f^{-1}(q) $ 是 $ M $ 的 $ C^k $ 子流形。此时只需假设 $ M $ 本身是 $ C^k $ 流形。 \nobreak\hspace{0.35em}■

\end{items}

\end{remark}

\begin{example}

\begin{items}
\item

设 $ X $ 是 $ \mathbb{R}^m \times \mathbb{R}^n $ 中的开集，$ q \in \mathbb{R}^n $ 是 $ f \in C^\infty(X, \mathbb{R}^n) $ 的正则值，且 $ M := f^{-1}(q) \ne \varnothing $。则 $ M $ 是 $ X $ 的 $ m $ 维子流形。对

\[
\pi := \operatorname{pr}|_M : M \to \mathbb{R}^m
\]

其中

\[
\operatorname{pr}: \mathbb{R}^m \times \mathbb{R}^n \to \mathbb{R}^m, \quad (x, y) \mapsto x,
\]

我们有 $ \pi \in C^\infty(M, \mathbb{R}^m) $。最后设 $ p \in M $，并假设 $ D_1 f(p) \in \mathcal{L}(\mathbb{R}^m, \mathbb{R}^n) $ 是满射。\footnote{我们采用第 VII.8 节的记号。} 则 $ p $ 是 $ \pi $ 的正则点，当且仅当 $ D_2 f(p) $ 是双射。

\proofhead
正则值定理保证 $ M $ 是 $ X $ 的 $ m $ 维子流形，且对 $ p \in M $ 有 $ T_p M = \ker(T_p f) $。因为 $ \pi $ 是线性映射（从而是光滑映射）的限制，由注 1.1(j) 可知 $ \pi \in C^\infty(M, \mathbb{R}^m) $ 且 $ T_p \pi = T_p \operatorname{pr}|_{T_p M} $。
由 $ T_p \operatorname{pr} = (p, \partial \operatorname{pr}(p)) $ 以及 $ \partial \operatorname{pr}(p)(h, k) = h $（$ (h, k) \in \mathbb{R}^m \times \mathbb{R}^n $）可知，$ T_p \pi $ 是满射，当且仅当对每个 $ y \in \mathbb{R}^m $ 都存在 $ (h, k) \in \mathbb{R}^m \times \mathbb{R}^n $ 使得
\[
\partial f(p)(h, k) = D_1 f(p)h + D_2 f(p)k = 0
\]

且 $ h = y $。这是因为 $ D_1 f(p) $ 是满射，当且仅当对每个 $ z \in \mathbb{R}^n $ 都存在 $ k \in \mathbb{R}^n $ 使 $ D_2 f(p)k = z $；等价地，当且仅当 $ D_2 f(p) $ 本身是满射。由于 $ D_2 f(p) \in \mathcal{L}(\mathbb{R}^n) $，证明完毕。 \nobreak\hspace{0.35em}■

\item

（“尖点灾变”） 对

\[
f: \mathbb{R}^2 \times \mathbb{R} \to \mathbb{R}, \quad ((u, v), x) \mapsto u + vx + x^3,
\]

我们有

\[
[D_1 f(w, x)] = [1, x] \in \mathbb{R}^{1 \times 2}, \quad \text{其中 } w := (u, v).
\]

因此 $ 0 $ 是 $ f $ 的正则值，且 $ M := f^{-1}(0) $ 是 $ \mathbb{R}^3 $ 中的曲面。因为 $ D_2 f(w, x) = v + 3x^2 $，由 (a) 可知

\[
K := \{((u, v), x) \in M ; v + 3x^2 = 0\}
\]

是投影 $ \pi: M \to \mathbb{R}^2 $ 的奇点集。它满足

\begin{equation}
K = \gamma(\mathbb{R}) \quad \text{其中 } \gamma: \mathbb{R} \to \mathbb{R}^3, \ t \mapsto (2t^3, -3t^2, t).\tag{1.5}
\end{equation}

特别地，$ K $ 是 $ M $ 的一维子流形，是一条光滑嵌入的曲线。它的投影 $ B := \pi(K) $ 是

\[
\sigma: \mathbb{R} \to \mathbb{R}^2, \quad t \mapsto (2t^3, -3t^2)
\]

的像，即一条尼尔抛物线。\footnote{见注 VII.9.9(a)。} 它是 0 维流形 $ P := \{(0, 0)\} \in \mathbb{R}^2 $（即“尖点”）与两条一维流形 $ B_1 := \sigma((-\infty, 0)) $ 和 $ B_2 := \sigma((0, \infty)) $ 的并。

\end{items}

\bookfigure[0.42]{figures/ada05689643058f1.png}

\proofhead
点 $ (u, v, x) \in \mathbb{R}^3 $ 属于 $ K $，当且仅当它满足方程
\begin{equation}
u + vx + x^3 = 0 \quad \text{且} \quad v + 3x^2 = 0.\tag{1.6}
\end{equation}

由第一个方程消去 $ v $，我们看到 (1.6) 等价于

\[
2x^3 = u \quad \text{且} \quad 3x^2 = -v.
\]

这证明了 (1.5)。对映射

\[
g: \mathbb{R}^3 \to \mathbb{R}^2, \quad (u, v, x) \mapsto (u - 2x^3, v + 3x^2)
\]

的导数，我们求得

\begin{equation*}
[\partial g(u, v, x)] = \begin{bmatrix} 1 & 0 & -6x^2 \\ 0 & 1 & 6x \end{bmatrix} \in \mathbb{R}^{2 \times 3}.
\end{equation*}

这个矩阵的秩为 2，这表明 $ 0 $ 是 $ g $ 的正则值。因此由正则值定理，$ K = g^{-1}(0) $ 是 $ \mathbb{R}^3 $ 的一维子流形。因为 $ K \subset M $，由注 1.5(c) 可知 $ K $ 是 $ M $ 的子流形。其余部分是显然的。 \nobreak\hspace{0.35em}■

\end{example}

\begin{remark}

（灾变理论） 我们现在考虑一个质量为 1 的质点沿实轴运动，其势能为 $ U $，总能量为

\[
E(\dot{x}, x) = \frac{\dot{x}^2}{2} + U(x) \quad \text{对 } x \in \mathbb{R}.
\]

根据例 VII.6.14(a)，牛顿运动方程是

\[
\ddot{x} = -U'(x).
\]

由例 VII.8.17(b) 与 (c) 可知，能量 $ E $ 的临界点恰好是满足 $ U'(x_0) = 0 $ 的点 $ (0, x_0) $。因为 $ E $ 的海森矩阵在 $ (0, x_0) $ 处具有形式

\begin{equation*}
\begin{bmatrix} 1 & 0 \\ 0 & U''(x_0) \end{bmatrix},
\end{equation*}

所以它正定，当且仅当 $ U''(x_0) > 0 $。因此由定理 VII.5.14 可知，$ (0, x_0) $ 是总能量的孤立极小值点，当且仅当 $ x_0 $ 是势能的孤立极小值点。\footnote{我们只考虑“通有”情形，即当 $ U'(x_0) = 0 $ 时满足 $ U''(x_0) \ne 0 $。} 从图形上看很清楚，总能量的孤立极小值点是“稳定的”，其意义是：若在运动开始时（即 $ t = 0 $ 时）这一性质成立，则对一切 $ t \in \mathbb{R}^+ $，$ (\dot{x}(t), x(t)) $ 都停留在 $ (0, x_0) $ 的“邻域”内。

直观上，可以通过想象一个小球沿 $ U $ 的图像无摩擦地滚动、同时受重力作用，来理解 $ x $ 将如何沿轴 $ \mathbb{R} $ 运动。如果小球位于“势阱底部”，即位于局部极小值点，那么它不会运动，因为 $ \dot{x}(t) = U'(x_0) = 0 $。如果小球在局部极小值点附近被释放，那么小球将滚下坡、越过极小值点，再沿另一侧“山谷的斜坡”向上滚，直到它在原来的高度上“耗尽”动能。然后它将反向运动并滚动，直到再次在最初被释放的位置瞬时静止。因此小球将绕 $ x_0 $ 作周期振动。\footnote{这个看似合理的图景可以用常微分方程理论证明；例如见 [Ama95]。}

\bookfigure[0.45]{figures/e04ef29a714ec21c.png}

我们现在假设 $ U $ 连续依赖于额外的“控制参数” $ u, v, \ldots $。通过改变这些参数，我们可以连续地改变 $ U $ 的图像。这样，可能出现局部极小值先并入鞍点、然后不再是临界点的情形。先前被限制在局部极小值点邻域内的小球，此时就会离开这个邻域，并绕另一个静止点振动。

\bookfigure[0.80]{figures/20634627367798c0.png}

现在考虑一位能看到小球运动、却不知道过程背后机制的观察者。她将看到，原先平静地静止在某一位置的小球，会突然“无缘无故地”开始滚动，并绕另一个（虚拟的）中心周期振动。这看起来像是情形的突然而剧烈的变化，即一场“灾变”。

为了理解这类灾变（并在必要时避免它们），必须理解它们发生的机制。在上述情形中，这归结为理解势的临界点（特别是相对极小值点）如何依赖于控制参数。

为说明这一点，我们考虑势

\[
U_{(u, v)}: \mathbb{R} \to \mathbb{R}, \quad x \mapsto ux + vx^2/2 + x^4/4
\]

（$ (u, v) \in \mathbb{R}^2 $）。$ U_{(u, v)} $ 的临界点恰好是例 1.9(b) 中函数 $ f $ 的零点。因此，流形 $ M $，即\textbf{灾变流形}，描述了势的两参数族 $ \{U_{(u, v)} ; (u, v) \in \mathbb{R}^2\} $ 的所有临界点。特别令人感兴趣的是 $ M $ 的那个子集，即\textbf{灾变集} $ K $，它由从 $ M $ 到参数空间的投影 $ \pi $ 的所有奇点组成。在我们的例子中，$ K $ 是光滑嵌入 $ M $ 中的一条曲线，即\textbf{折叠曲线}，因为灾变流形沿 $ K $ 被“折叠”。$ K $ 在 $ \pi $ 下的像，即折叠曲线到参数平面上的投影，是\textbf{分岔集} $ B $。$ \mathbb{R}^2 \setminus B $ 的每个点都是 $ \pi $ 的正则点。纤维 $ \pi^{-1}(u, v) $ 对 $ (u, v) \in A \cup P $ 恰含一个点，对 $ (u, v) \in B_1 \cup B_2 $ 恰含两个点，对 $ (u, v) \in I $ 恰含三个点，其中 $ A $ 与 $ I $ 在例 1.9(b) 的插图中绘出。下面这些图显示了当 $ (u, v) $ 属于这些集合时势 $ U_{(u, v)} $ 的定性形式。

\bookfigure[0.95]{figures/e4dcd3109dbec91f.png}

现在考虑参数空间中的一条连续曲线 $ C $，它从 $ A $ 出发、终止于 $ I $（或反向），并且始终停留在 $ B_1\cup B_2 $ 中。沿这条曲线连续移动时，$ \pi $ 的原像中点的个数会突然从 1 变成 3（或从 3 变成 1）。如右图所示，这样一条曲线 $ C $ 可由投影灾变流形 $ M $ 上的一条曲线 $ \Gamma $ 得到，该曲线在穿过折叠曲线时发生“跳变”。简言之，$ x $ 的值经历了一次“灾变”。

这些事实导致了对“灾变理论”的许多解释；这些解释——尤其是因为其名称——被用来获得极大的名声，并且在科普文献中尤其激起了过高的期望，以为这门学科会以某种方式解释或帮助防止现实世界中的灾变。关于灾变理论的批判性、非技术性导论，我们参考 [Arn84]；我们还推荐 [PS78]，其中详细介绍了奇点理论的数学理论及其若干应用，而灾变理论正是其中的一部分。

\bookfigure[0.30]{figures/1441e043718626f3.png}
\end{remark}

\needspace{3\baselineskip}\subsection*{带边子流形}

我们知道，开单位球 $ \mathbb{B}^m $ 及其边界，即 $ (m-1) $ 维球面 $ S^{m-1} $，分别是 $ \mathbb{R}^m $ 的 $ m $ 维与 $ (m-1) $ 维子流形。然而，闭球 $ \overline{\mathbb{B}}^m = \mathbb{B}^m\cup S^{m-1} $ 不是流形，因为点 $ p \in \partial\mathbb{B}^m = S^{m-1} $ 在 $ \overline{\mathbb{B}}^m $ 中没有邻域 $ U $ 能被拓扑地映成 $ \mathbb{R}^m $ 的开集 $ V $；这样的邻域 $ U $ 作为开集 $ V $ 的同胚像，同样须在 $ \mathbb{R}^m $ 中为开，而这是不对的。在 $ p $ 的邻域中，也就是“用很强的显微镜观察”时，$ \overline{\mathbb{B}}^m $ 看起来不像 $ \mathbb{R}^m $，而更像一个半空间。为了也涵盖这类情形，我们必须推广流形的概念，允许把半空间的子集作为参数集。

以下设 $ m \in \mathbb{N}^{\times} $，并令

\[
\mathbb{H}^m := \mathbb{R}^{m-1}\times(0,\infty)
\]

为 $ \mathbb{R}^m $ 的\textbf{开上半空间}。若不会引起误解，我们将其边界 $ \partial\mathbb{H}^m = \mathbb{R}^{m-1}\times\{0\} $ 与 $ \mathbb{R}^{m-1} $ 等同。若 $ U $ 是 $ \overline{\mathbb{H}}^m := \overline{\mathbb{H}}^m = \mathbb{R}^{m-1}\times\mathbb{R}^{+} $ 的开子集，则称 $ \operatorname{int}(U) := U\cap\mathbb{H}^m $ 为 $ U $ 的\textbf{内部}，称 $ \partial U := U\cap\partial\mathbb{H}^m $ 为 $ U $ 的\textbf{边界}。注意边界 $ \partial U $ 既不是 $ U $ 在 $ \overline{\mathbb{H}}^m $ 中的拓扑边界\footnote{从今往后，我们只用符号 $ \partial M $ 表示边界；为清楚起见，对拓扑空间子集 $ M $ 的拓扑边界，我们写作 $ \operatorname{Rd}(M) $，即令 $ \operatorname{Rd}(M) := \overline{M}\setminus\mathring{M} $。}，也不是 $ U $ 在 $ \mathbb{R}^m $ 中的拓扑边界（在后一情形除非 $ U = \mathbb{H}^m $）。

\bookfigure[0.50]{figures/a0569b8109ce828a.png}

设 $ X $ 在 $ \overline{\mathbb{H}}^m $ 中开，$ E $ 是巴拿赫空间。若存在 $ x_0 $ 在 $ \mathbb{R}^m $ 中的邻域 $ U $ 以及可微函数 $ f_U : U \to E $，使得 $ f_U $ 在 $ U\cap X $ 上与 $ f $ 一致，则称 $ f : X \to E $ 在边界点 $ x_0 \in \partial X $ 处\textbf{可微}。于是由命题 VII.2.5 可得

\[
\partial_j f_U(x_0) = \lim_{t \to 0+} \left(f_U(x_0 + te_j) - f_U(x_0)\right)/t = \lim_{t \to 0+} \left(f(x_0 + te_j) - f(x_0)\right)/t
\]

其中 $ 1 \le j \le m $，$ (e_1,\ldots,e_m) $ 是 $ \mathbb{R}^m $ 的标准基。这与命题 VII.2.8 表明 $ \partial f_U(x_0) $ 已由 $ f $ 确定。因此 $ f $ 的\textbf{导数}

\[
\partial f(x_0) := \partial f_U(x_0) \in \mathcal{L}(\mathbb{R}^m,E)
\]

在 $ x_0 $ 处是良定义的，即与 $ f $ 的局部延拓 $ f_U $ 的选择无关。

若映射 $ f : X \to E $ 在 $ X $ 的每一点处可微，且映射

\[
\partial f : X \to \mathcal{L}(\mathbb{R}^m,E) ,\quad x \mapsto \partial f(x)
\]

连续，则称 $ f $ 是\textbf{连续可微的}\footnote{自然，若 $ f|\operatorname{int}(X) $ 在 $ x_0 $ 处可微，则称 $ f $ 在 $ x_0 \in \operatorname{int}(X) $ 处可微。}。

$ f $ 的高阶导数类似地定义，它们也与具体的局部延拓无关。对 $ k \in \mathbb{N}^{\times}\cup\{\infty\} $，从 $ X $ 到 $ E $ 的 $ C^k $ 映射构成一个向量空间，与 $ \mathbb{R}^m $ 的开子集的情形一样，我们把它记作 $ C^k(X,E) $。

设 $ Y $ 在 $ \overline{\mathbb{H}}^m $ 中开。若 $ f $ 是双射，且 $ f $ 与 $ f^{-1} $ 都属于类 $ C^k $，则也称 $ f : X \to Y $ 为 $ C^k $ \textbf{微分同胚}，并记 $ f \in \operatorname{Diff}^k(X,Y) $。特别地，$ \operatorname{Diff}(X,Y) := \operatorname{Diff}^{\infty}(X,Y) $ 是从 $ X $ 到 $ Y $ 的所有光滑（即 $ C^{\infty} $）微分同胚的集合。

\begin{remark}

设 $ X $ 与 $ Y $ 在 $ \overline{\mathbb{H}}^m $ 中开，且对某个 $ k \in \mathbb{N}^{\times}\cup\{\infty\} $，$ f : X \to Y $ 是 $ C^k $ 微分同胚。

\begin{items}
\item

若 $ \partial X $ 非空，则 $ \partial Y \neq \varnothing $，且 $ f|\partial X $ 是从 $ \partial X $ 到 $ \partial Y $ 的 $ C^k $ 微分同胚。又 $ f|\operatorname{int}(X) $ 属于 $ \operatorname{Diff}^k(\operatorname{int}(X),\operatorname{int}(Y)) $。

\proofhead
设 $ p \in \partial X $，且 $ q := f(p) $ 属于 $ \operatorname{int}(Y) $。则由反函数定理（定理 VII.7.3，施用于 $ f $ 的一个局部延拓）可知 $ \partial f^{-1}(q) $ 是 $ \mathbb{R}^m $ 的自同构。于是再由定理 VII.7.3 可知，$ f^{-1} $ 把 $ q $ 在 $ \operatorname{int}(Y) $ 中的某个适当邻域 $ V $ 映为 $ p $ 在 $ \mathbb{R}^m $ 中的开邻域 $ U $。但由于 $ f^{-1}(V) \subset X \subset \overline{\mathbb{H}}^m $ 且 $ p = f^{-1}(q) \in \partial X $，这是不可能的。因此 $ f(\partial X) \subset \partial Y $。类似地可得 $ f^{-1}(\partial Y) \subset \partial X $。这就表明 $ f(\partial X) = \partial Y $。
由于 $ X $ 与 $ Y $ 在 $ \overline{\mathbb{H}}^m $ 中开，$ \partial X $ 与 $ \partial Y $ 都在 $ \partial\mathbb{H}^m = \mathbb{R}^{m-1} $ 中开，且 $ f|\partial X $ 是从 $ \partial X $ 到 $ \partial Y $ 的双射。因为 $ f|\partial X $ 与 $ f^{-1}|\partial Y $ 显然属于类 $ C^k $，所以 $ f|\partial X $ 是从 $ \partial X $ 到 $ \partial Y $ 的 $ C^k $ 微分同胚。最后一个断言现在清楚了。\nobreak\hspace{0.35em}■
\item

对 $ p \in \partial X $，有 $ \partial f(p)(\partial\mathbb{H}^m) \subset \partial\mathbb{H}^m $ 以及 $ \partial f(p)(\pm\mathbb{H}^m) \subset \pm\mathbb{H}^m $。

\proofhead
由 $ f(\partial X) = \partial Y $ 可知，对 $ f $ 的第 $ m $ 个坐标函数 $ f^m $ 有 $ f^m|\partial X = 0 $。由此得到 $ \partial_j f^m(p) = 0 $，$ 1 \le j \le m-1 $。因此 $ f $ 在 $ p $ 处的雅可比矩阵形如
\begin{equation}
[\partial f(p)] = \begin{bmatrix} & \partial(f\mid\partial\mathbb{H}^m)(p) & \vdots & \partial_m f^1(p) \\ & & \vdots & \vdots \\ & & \vdots & \partial_m f^{m-1}(p) \\ \cdots & \cdots\cdots\cdots\cdots & \cdots & \cdots\cdots\cdots \\ 0 & \cdots & 0 & \partial_m f^m(p) \end{bmatrix}\tag{1.7}
\end{equation}

因为 $ f(X) \subset Y \subset \overline{\mathbb{H}}^m $，所以对 $ q \in X $ 有不等式 $ f^m(q) \ge 0 $。于是

\[
\partial_m f^m(p) = \lim_{t \to 0+} t^{-1}\left(f^m(p + te_m) - f^m(p)\right) = \lim_{t \to 0+} t^{-1} f^m(p + te_m) \ge 0 .
\]

由于 $ \partial f(p) \in \mathcal{L}\mathrm{aut}(\mathbb{R}^m) $（见注 VII.7.4(d)）且 $ \partial_m f^m(p) \ge 0 $，我们有 $ \partial_m f^m(p) > 0 $。由 (1.7) 可读出

\[
(\partial f(p)x)^m = \partial_m f^m(p)t \quad \text{对 } x := (y,t) \in \mathbb{R}^{m-1}\times\mathbb{R} .
\]

因此 $ \partial f(p)x $ 的第 $ m $ 个坐标的符号与 $ \operatorname{sign}(t) $ 一致，证毕。\nobreak\hspace{0.35em}■

\end{items}

现在我们可以定义带边子流形的概念。若对每个 $ p \in B $，都存在 $ N $ 的一个围绕 $ p $ 的坐标卡 $ (\psi,V) $——即 $ N $ 关于 $ B $ 的一个围绕 $ p $ 的\textbf{子流形坐标卡}——使得

\begin{equation}
\psi(V\cap B) = \psi(V)\cap(\overline{\mathbb{H}}^b\times\{0\}) \subset \mathbb{R}^n ,\tag{1.8}
\end{equation}

则称 $ n $ 维流形 $ N $ 的子集 $ B $ 为 $ N $ 的\textbf{带边 $ b $ 维子流形}。

这里，若 $ \psi(p) $ 落在 $ \partial\mathbb{H}^b := \partial\mathbb{H}^b\times\{0\} $ 中，则称 $ p $ 是 $ B $ 的\textbf{边界点}。

\bookfigure[0.85]{figures/ff1a74dfc31b8a7d.png}

所有边界点的集合构成 $ B $ 的\textbf{边界}\footnote{注意边界 $ \partial B $ 与内部 $ \operatorname{int}(B) $ 通常不同于 $ B $ 的拓扑边界 $ \operatorname{Rd}(B) $ 与拓扑内部 $ \mathring{B} $。在关于流形的陈述中，我们将在上述定义的意义下理解“边界”与“内部”。} $ \partial B $。集合 $ \operatorname{int}(B) := B\setminus\partial B $ 称为带边子流形 $ B $ 的\textbf{内部}。最后，若 $ b = n-1 $，则称 $ B $ 是 $ N $ 中的\textbf{带边超曲面}。

\end{remark}

\begin{remark}

\begin{items}
\item

本节开头意义下的 $ N $ 的每个子流形 $ M $ 都是带边子流形，但其边界为空。我们称这样的对象为\textbf{无边（子）流形}。

\item

边界 $ \partial B $ 与内部 $ \operatorname{int}(B) $ 是良定义的，即与坐标卡无关。

\proofhead
设 $ (\chi,W) $ 是 $ N $ 的另一个关于 $ B $ 的、围绕 $ p $ 的子流形坐标卡。又设 $ f $ 是转移函数 $ \chi\circ\psi^{-1} $ 在 $ \psi(V\cap W)\cap(\overline{\mathbb{H}}^b\times\{0\}) $ 上的限制，把它看作 $ \overline{\mathbb{H}}^b $ 的开子集。则由注 1.11(a) 可知，$ \chi(p) $ 属于 $ \partial\mathbb{H}^b $ 当且仅当 $ \psi(p) $ 属于 $ \partial\mathbb{H}^b $。\nobreak\hspace{0.35em}■
\item

设 $ p \in \operatorname{int}(B) $。则 (1.8) 蕴含

\[
\psi(V\cap\operatorname{int}(B)) = \psi(V)\cap(\mathbb{H}^b\times\{0\}) .
\]

因为 $ \mathbb{H}^b $ 微分同胚于 $ \mathbb{R}^b $，这表明 $ \operatorname{int}(B) $ 是 $ N $ 的无边 $ b $ 维子流形。

\item

在 $ p \in \partial B $ 的情形，由 (1.8) 可得

\[
\psi(V\cap\partial B) = \psi(V)\cap(\mathbb{R}^{b-1}\times\{0\}) .
\]

因此 $ \partial B $ 是 $ N $ 的无边 $ (b-1) $ 维子流形。

\item

$ N $ 的每个带边 $ b $ 维子流形都是 $ \mathbb{R}^n $ 的带边 $ b $ 维子流形。

\proofhead
这类似于注 1.1(a) 的证明。\nobreak\hspace{0.35em}■
\item

（正则性）对 $ k \in \mathbb{N}^{\times} $，$ C^k $ 带边子流形如何定义是清楚的，且 (a)–(c) 的类似结论仍然成立。\nobreak\hspace{0.35em}■

\end{items}

\end{remark}

\needspace{3\baselineskip}\subsection*{局部坐标卡}

设 $ B $ 是 $ N $ 的带边 $ b $ 维子流形。若以下条件成立，则称映射 $ \varphi $ 为 $ B $ 在 $ p $ 处的（$ b $ 维局部）\textbf{坐标卡}：

- $ U := \operatorname{dom}(\varphi) $ 在 $ B $ 中是开集，其中 $ B $ 带有由 $ N $（从而由 $ \mathbb{R}^n $）诱导的拓扑；
- $ \varphi $ 是从 $ U $ 到 $ \overline{\mathbb{H}}^b $ 的开子集 $ X $ 的同胚；
- $ i_B\circ\varphi^{-1} : X \to N $ 是浸入，其中 $ i_B : B \to N $，$ p \mapsto p $ 表示单射。

注意，除了 $ \varphi(U) $ 在 $ \overline{\mathbb{H}}^b $ 中是开集、且 $ \mathbb{R}^n $ 被替换为 $ N $ 之外，这个定义与 $ \mathbb{R}^n $ 的子流形的 $ C^{\infty} $ 坐标卡的定义逐字一致（见第 VII.9 节）。

\bookfigure[0.80]{figures/cd416af3c78670db.png}

\begin{remark}

\begin{items}
\item

若 $ (\psi,V) $ 是 $ N $ 关于 $ B $ 的子流形坐标卡，则相交坐标卡 $ (\varphi,U) := (\psi|V\cap B, V\cap B) $ 是 $ B $ 的 $ b $ 维坐标卡。

\item

若 $ (\varphi_1,U_1) $ 与 $ (\varphi_2,U_2) $ 是 $ B $ 在 $ p \in B $ 处的坐标卡，则 $ \varphi_j(U_1\cap U_2) $ 在 $ \overline{\mathbb{H}}^b $ 中开，$ j = 1,2 $，且\textbf{转移函数} $ \varphi_2\circ\varphi_1^{-1} $ 满足

\[
\varphi_2\circ\varphi_1^{-1} \in \operatorname{Diff}(\varphi_1(U_1\cap U_2), \varphi_2(U_1\cap U_2)) .
\]
\item

设 $ (\varphi,U) $ 是 $ B $ 在 $ p \in \partial B $ 处的坐标卡。则

\[
(\varphi_{\partial B}, U_{\partial B}) := (\varphi|U\cap\partial B, U\cap\partial B)
\]

是 $ \partial B $（$ N $ 的无边 $ (b-1) $ 维子流形）的坐标卡。

\item

所有能用流形的坐标卡描述的概念与定义，例如映射的可微性与局部表示，都直截了当地推广到带边子流形。特别地，$ i_B : B \hookrightarrow N $，即 $ B $ 在 $ N $ 中的自然嵌入 $ p \mapsto p $，是光滑映射。

\item

若 $ C $ 是 $ M $ 的带边子流形且 $ f \in \operatorname{Diff}(B,C) $，则 $ f(\partial B) = \partial C $，且 $ f|\partial B $ 是从 $ \partial B $ 到 $ \partial C $ 的微分同胚。

\proofhead
这由注 1.11(a) 推出。\nobreak\hspace{0.35em}■
\item

设 $ B $ 是 $ N $ 的带边 $ b $ 维子流形，且 $ f \in C^{\infty}(B,M) $ 是\textbf{嵌入}，即 $ f $ 是双射浸入并且是从 $ B $ 到 $ f(B) $ 的同胚。则 $ f(B) $ 是 $ M $ 的带边 $ b $ 维子流形，满足 $ \partial f(B) = f(\partial B) $，且 $ f $ 是从 $ B $ 到 $ f(B) $ 的微分同胚。

\proofhead
定理 1.2(ii) 的证明同样适用。\nobreak\hspace{0.35em}■
\item

（正则性）前面所有结论都逐字转移到 $ C^k $ 带边子流形。\nobreak\hspace{0.35em}■

\end{items}

我们自然再次称 $ B $ 的一个坐标卡族 $ \{(\varphi_{\alpha},U_{\alpha}) ; \alpha \in A\} $（满足 $ B = \bigcup_{\alpha}U_{\alpha} $）为 $ B $ 的\textbf{图册}。

\end{remark}

\needspace{3\baselineskip}\subsection*{切空间与法空间}

设 $ B $ 是 $ N $ 的带边子流形，并设 $ p \in \partial B $。又设 $ (\varphi,U) $ 是 $ B $ 在 $ p $ 处的坐标卡。则定义 $ B $ 在点 $ p $ 处的\textbf{切空间} $ T_pB $ 为

\[
T_pB := T_{\varphi(p)}(i_B\circ\varphi^{-1})(T_{\varphi(p)}\mathbb{R}^b) ,
\]

其中 $ b := \dim(B) $。因此 $ T_pB $ 是 $ N $ 在 $ p $ 处的切空间 $ T_pN $ 的一个（“满”）$ b $ 维向量子空间（而不是，比如说，一个半空间）。对注 VII.10.3(a) 的证明作显然的修改即可表明 $ T_pB $ 是良定义的，即与所用的坐标卡无关。此时，我们通过 $ TB := \bigcup_{p\in B}T_pB $ 定义 $ B $ 的\textbf{切丛} $ TB $。

\bookfigure[0.30]{figures/36a9db60a295abb5.png}

\begin{remark}

\begin{items}
\item

对 $ p \in \partial B $，$ T_p\partial B $ 是 $ T_pB $ 的 $ (b-1) $ 维向量子空间。

\proofhead
这是注 1.12(d) 与 1.13(c) 的简单推论。\nobreak\hspace{0.35em}■
\item

设 $ p \in \partial B $，且 $ (\varphi,U) $ 是 $ B $ 在 $ p $ 处的坐标卡。令

\[
T_p^{\pm}B := T_{\varphi(p)}(i_B\circ\varphi^{-1})(\varphi(p), \pm\overline{\mathbb{H}}^b) ,
\]

则 $ T_pB = T_p^{+}B\cup T_p^{-}B $ 且 $ T_p^{+}B\cap T_p^{-}B = T_p(\partial B) $。向量 $ v $ 是\textbf{内指}［或\textbf{外指}］切向量，当且仅当 $ v $ 属于集合 $ T_p^{+}B\setminus T_p(\partial B) $［或 $ T_p^{-}B\setminus T_p(\partial B) $］。这成立当且仅当 $ (T_p\varphi)v $ 的第 $ b $ 个分量为正［或为负］。

\proofhead
由注 1.11(b) 与 1.13(b) 容易推出，$ T_p^{\pm}B $ 的定义与坐标无关。\nobreak\hspace{0.35em}■
\item

设 $ C $ 是 $ M $ 的带边或无边子流形。对 $ f \in C^1(C,N) $，$ f $ 在 $ p \in C $ 处的\textbf{切映射} $ T_pf $ 按无边流形的情形定义。于是注 VII.10.9 的类似结论仍然成立。\nobreak\hspace{0.35em}■

\end{items}

设 $ p \in \partial B $。则 $ T_{p}(\partial B) $ 是 $ b $ 维向量空间 $ T_{p}B $ 的 $ (b-1) $ 维向量子空间。作为 $ T_{p}N $（从而也是 $ T_{p}\mathbb{R}^{\overline{n}} $）的向量子空间，$ T_{p}B $ 以 $ \mathbb{R}^{\overline{n}} $ 上欧几里得标量积诱导的内积 $ (\cdot \mid \cdot)_{p} $ 构成内积空间。因此 $ T_{p}^{-}B $ 中恰有一个与 $ T_{p}(\partial B) $ 正交的单位向量 $ \nu(p) $，我们称它为 $ \partial B $ 在 $ p $ 处的\textbf{外向（单位）法向量}。显然 $ -\nu(p) \in T_{p}^{+}B $ 是 $ T_{p}B $ 中与 $ T_{p}(\partial B) $ 正交的唯一内指向量，我们称它为 $ \partial B $ 在 $ p $ 处的\textbf{内向（单位）法向量}。

\bookfigure[0.40]{figures/8fb811dbae7b9903.png}
\end{remark}

\needspace{3\baselineskip}\subsection*{正则值定理}

我们已经看到，无边界子流形在许多情形下（事实上局部总是如此）可以表示成正则映射的纤维。现在我们把这一重要而简单的判别准则推广到带边子流形的情形。

\begin{theorem}

（正则值） 设 $ c $ 是 $ f \in C^{\infty}(N, \mathbb{R}) $ 的正则值。则

\[
B := f^{-1}((-\infty, c]) = \{p \in N ; f(p) \le c\}
\]

是 $ N $ 的 $ n $ 维带边子流形，且 $ \partial B = f^{-1}(c) $、$ \operatorname{int}(B) = f^{-1}((-\infty, c)) $。对 $ p \in \partial B $，有 $ T_{p}(\partial B) = \ker(d_{p}f) $，并且 $ \partial B $ 上的外向单位法向量 $ \nu(p) $ 由 $ \nabla_{p}f(p)/|\nabla_{p}f|_{p} $ 给出。

\proofhead
由于 $ f^{-1}((-\infty, c)) $ 在 $ N $ 中是开的，从而是 $ N $ 的 $ n $ 维子流形，故只需考虑 $ p \in f^{-1}(c) $。

因此设 $ p \in f^{-1}(c) $，并设 $ (\psi, V) $ 是 $ N $ 在 $ p $ 附近的坐标卡，使得 $ \psi(p) = 0 $。则 $ g := c - f \circ \psi^{-1} $ 属于 $ C^{\infty}(\psi(V), \mathbb{R}) $，且满足 $ g(0) = 0 $，并且 $ g(x) \ge 0 $ 当且仅当 $ x $ 落在 $ \psi(V \cap B) $ 中。此外 $ 0 $ 是 $ g $ 的正则点。通过重新命名坐标（即用 $ \psi $ 与一个置换复合），我们可以假设 $ \partial_{n}g(0) \ne 0 $，从而 $ \partial_{n}g(0) > 0 $。

考虑由 $ \varphi(x) := (x^{1}, \ldots, x^{n-1}, g(x)) $ 定义的映射 $ \varphi \in C^{\infty}(\psi(V), \mathbb{R}^{n}) $。它满足 $ \varphi(0) = 0 $，并且
\begin{equation*}
\partial\varphi = \begin{bmatrix} 1_{n-1} & & & 0 \\ & & & \vdots \\ & & & 0 \\ \partial_{1}g & \cdots & \partial_{n-1}g & \partial_{n}g \end{bmatrix}
\end{equation*}

因此 $ \partial\varphi(0) $ 是 $ \mathbb{R}^{n} $ 的自同构，而定理 VII.7.3（反函数定理）保证存在 $ 0 $ 在 $ \psi(V) $ 中的开邻域 $ U $ 与 $ W $，使得 $ \varphi|U $ 是从 $ U $ 到 $ W $ 的微分同胚。

令 $ V_{0} := \psi^{-1}(U) $、$ \chi := \varphi \circ \psi|V_{0} $，我们看到 $ (\chi, V_{0}) $ 是 $ N $ 在 $ p $ 附近的坐标卡，且 $ \chi(p) = 0 $、$ \chi(B \cap V_{0}) = \chi(V_{0}) \cap \overline{\mathbb{H}}^{n} $。这表明 $ B $ 是 $ N $ 的带边子流形，且 $ \partial B = f^{-1}(c) $、$ \operatorname{int}(B) = f^{-1}((-\infty, c)) $。于是由定理 1.7 得到

\begin{equation}
T_{p}(\partial B) = \ker(T_{p}f) = \ker(d_{p}f) \quad \text{对 } p \in \partial B.\tag{1.9}
\end{equation}

由于对 $ v \in T_{p}N $ 有 $ \langle d_{p}f, v \rangle_{p} = (\nabla_{p}f \mid v)_{p} $，由 (1.9) 可得 $ \nabla_{p}f $ 与 $ T_{p}(\partial B) $ 正交。

最后，设 $ \lambda: (-\varepsilon, \varepsilon) \to N $ 是 $ N $ 中的 $ C^{1} $ 道路，满足 $ \lambda(0) = p $ 且 $ \dot{\lambda}(0) = \nabla_{p}f $（见定理 VII.10.6）。则

\[
(f \circ \lambda)^{\cdot}(0) = \langle d_{p}f, \nabla_{p}f \rangle = |\nabla_{p}f|_{p}^{2} > 0.
\]

因此由推论 IV.3.3 的泰勒公式可得

\[
f(\lambda(t)) = c + t|\nabla_{p}f|_{p}^{2} + o(t) \quad (t \to 0).
\]

从而对充分小的正 $ t $ 有 $ f(\lambda(t)) > c $，即 $ f(\lambda(t)) \notin B $。这表明 $ \nabla_{p}f $ 是 $ B $ 在 $ p $ 处的外向切向量。现在最后一个断言也清楚了。\nobreak\hspace{0.35em}■

\end{theorem}

\begin{remark}

\begin{items}
\item

由于我们可以局部地把子流形表示成正则映射的纤维（见注 1.8(a)），我们也可以局部地把带边子流形表示成半开区间的原像。更确切地说，设 $ B $ 是 $ N $ 的 $ n $ 维带边子流形。则对每一点 $ p \in B $，存在 $ N $ 中的邻域 $ U $ 及函数 $ f \in C^{\infty}(U, \mathbb{R}) $，使得当 $ p \in \operatorname{int}(B) $ 时 $ B \cap U = f^{-1}((-\infty, 1)) $，而当 $ p \in \partial B $ 时 $ f(p) = 0 $ 且 $ B \cap U = f^{-1}((-\infty, 0]) $，并且 $ 0 $ 是它的正则值。

\proofhead
设 $ (\varphi, U) $ 是 $ N $ 在 $ p $ 附近关于 $ B $ 的子流形坐标卡，且 $ \varphi(p) = 0 $。我们可以假设 $ \varphi(U) $ 包含在 $ \mathbb{B}_{\infty}^{n} $ 中。若 $ p $ 是 $ B $ 的内点，我们对 $ q \in U $ 令 $ f(q) := \varphi^{n}(q) $。则 $ f $ 属于 $ C^{\infty}(U, \mathbb{R}) $，且 $ f^{-1}((-\infty, 1)) = U $。若 $ p $ 属于 $ \partial B $，我们对 $ q \in U $ 令 $ f(q) := -\varphi^{n}(q) $。则 $ f(p) = 0 $，且 $ f^{-1}((-\infty, 0]) = U \cap B $。因为 $ \varphi \in \operatorname{Diff}(U, \varphi(U)) $，我们知道 $ f $ 是淹没。因此 $ 0 $ 是 $ f $ 的正则值。\nobreak\hspace{0.35em}■
\item

（正则性）设 $ c $ 是 $ f \in C^{k}(N, \mathbb{R}) $ 的正则值，其中 $ k \in \mathbb{N}^{\times} $。则 $ f^{-1}((-\infty, c]) $ 是 $ N $ 中的 $ n $ 维带边子流形。此时只需假设 $ N $ 是 $ C^{k} $ 流形。\nobreak\hspace{0.35em}■

\end{items}

\end{remark}

\begin{example}

\begin{items}
\item

对每个 $ r > 0 $，$ \overline{\mathbb{B}}_{r}^{n} := r\overline{\mathbb{B}}^{n} = \{x \in \mathbb{R}^{n} ; |x| \le r\} $ 是 $ \mathbb{R}^{n} $ 的 $ n $ 维带边子流形。它的边界与拓扑边界重合，从而与半径为 $ r $ 的 $ (n-1) $ 维球面重合，即 $ \partial\overline{\mathbb{B}}_{r}^{n} = rS^{n-1} $。在 $ p \in \partial\overline{\mathbb{B}}_{r}^{n} $ 处的外向法向量 $ \nu(p) $ 由 $ (p, p/|p|) $ 给出。
当 $ n = 1 $ 时，球 $ \overline{\mathbb{B}}_{r}^{1} $ 是 $ \mathbb{R} $ 中的闭区间 $ [-r, r] $，而半径为 $ r $ 的 0 维球面由 $ S_{r}^{0} = \{-r\} \cup \{r\} $ 给出。外向法向量为 $ \nu(-r) = (-r, -1) $ 和 $ \nu(r) = (r, 1) $。

\end{items}

\bookfigure[0.60]{figures/efa534af2190a9c6.png}

\proofhead
这由定理 1.15 取 $ N := \mathbb{R}^{n} $ 及 $ f(x) := |x|^{2} $（$ x \in \mathbb{R}^{n} $）即得。\nobreak\hspace{0.35em}■

\begin{items}
\item[（b）]

设 $ A \in \mathbb{R}_{\mathrm{Sym}}^{(n+1) \times (n+1)} $ 及 $ c \in \mathbb{R}^{\times} $。又设

\[
V_{c} := \{x \in \mathbb{R}^{n+1} ; (Ax \mid x) \le c\}
\]

非空。若 $ A $ 正定且 $ c > 0 $，则 $ V_{c} $ 是一个 $ (n+1) $ 维立体，其边界是 $ n $ 维椭球

\[
K_{c} := \{x \in \mathbb{R}^{n+1} ; (Ax \mid x) = c\}.
\]

若 $ A $ 负定且 $ c < 0 $，则 $ V_{c} $ 是 $ V_{-c} $ 内部的补集，而 $ V_{-c} $ 的边界是 $ n $ 维椭球 $ K_{-c} $。若 $ A $ 不定但可逆，则 $ V_{c} $ 是某个界定 $ V_{c} $ 的 $ n $ 维双曲面 $ K_{c} $ 的“内部”或“外部”。在每种情形下，$ Ax/|Ax| $ 都是 $ V_{c} $ 在 $ K_{c} $ 处的外向法向量。（将此与注 VII.10.18 比较，并相应地解释那里的图。）

\end{items}

\bookfigure[0.75]{figures/ce197bf4e24c2524.png}

\begin{items}
\item[（c）]

设 $ A \in \mathbb{R}^{(n+1) \times (n+1)} $ 是对称的，$ c \in \mathbb{R}^{\times} $ 且 $ K_{c} \ne \varnothing $。又设 $ v \in \mathbb{R}^{n+1} \setminus \{0\} $、$ \alpha, \beta \in \mathbb{R} $ 且 $ \alpha < \beta $。则

\[
B := \{x \in K_{c} ; \alpha \le (v \mid x) \le \beta\}
\]

是 $ K_{c} $ 中位于两个平行超平面

\[
H_{\gamma} := \{x \in \mathbb{R}^{n+1} ; (v \mid x) = \gamma\} \quad \text{对 } \gamma \in \{\alpha, \beta\}.
\]

之间的部分。若 $ H_{\alpha} $ 与 $ H_{\beta} $ 不是 $ K_{c} $ 的切超平面，则 $ B $ 是 $ K_{c} $ 的 $ n $ 维带边子流形，且

\[
\partial B = \{x \in K_{c} ; (v \mid x) \in \{\alpha, \beta\}\}.
\]

\proofhead
因为映射 $ g := (v \mid \cdot)|K_{c}: K_{c} \to \mathbb{R} $ 由注 1.1(j) 是光滑的，我们知道 $ g^{-1}((\alpha, \beta)) $ 在 $ K_{c} $ 中是开的。因此 $ g^{-1}((\alpha, \beta)) $ 是 $ K_{c} $ 的 $ n $ 维子流形。
因此只需证明每个 $ p \in g^{-1}(\{\alpha, \beta\}) $ 都是 $ B $ 的边界点。于是设 $ V $ 是 $ K_{c} $ 中包含 $ p \in g^{-1}(\beta) $ 的开邻域，使得 $ g^{-1}(\alpha) \cap V = \varnothing $。$ H_{\beta} $ 不是切超平面这一假设蕴涵 $ \beta $ 是 $ f := g|V $ 的正则值（请证明这一点！）。现在由定理 1.15 应用于流形 $ V $ 与函数 $ f $ 即得该断言。类似的论证表明每个 $ p \in g^{-1}(\alpha) $ 都是 $ B $ 的边界点。\nobreak\hspace{0.35em}■
\item[（d）]

（类圆柱旋转超曲面）设

\[
\gamma: [0, 1] \to (0, \infty) \times \mathbb{R}, \quad t \mapsto (\rho(t), \sigma(t))
\]

是光滑嵌入。又设 $ i: S^{m} \hookrightarrow \mathbb{R}^{m+1} $ 及

\[
f: S^{m} \times [0, 1] \to \mathbb{R}^{m+1} \times \mathbb{R}, \quad (q, t) \mapsto (\rho(t)i(q), \sigma(t)).
\]

则 $ f $ 是光滑嵌入，且

\[
Z^{m+1} := f(S^{m} \times [0, 1])
\]

是 $ \mathbb{R}^{m+2} $ 中的带边超曲面，它微分同胚于“球面柱面” $ S^{m} \times [0, 1] $。
当 $ m = 0 $ 时，$ Z^{1} $ 由两条光滑、不自交的紧曲线\footnote{即带边的一维流形。} $ \gamma([0, 1]) $ 组成，它们关于 $ y $ 轴对称。

\end{items}

\bookfigure[0.60]{figures/24e2351a3a6b4f37.png}

当 $ m = 1 $ 时，$ Z^{2} $ 是 $ \mathbb{R}^{3} $ 中把\textbf{子午曲线}

\[
\Gamma := \{(\rho(t), 0, \sigma(t)) ; t \in [0, 1]\}
\]

绕 $ z $ 轴旋转所得的旋转曲面。

在一般情形下，我们称 $ Z^{m+1} $ 为\textbf{带边类圆柱旋转曲面}。其边界满足

\[
\partial Z^{m+1} = f(S^{m} \times \{0\}) \cup f(S^{m} \times \{1\}),
\]

而其内部有

\[
\operatorname{int}(Z^{m+1}) = f(S^{m} \times (0, 1)).
\]

特别地，$ \operatorname{int}(Z^{m+1}) $ 是无边界的类圆柱旋转超曲面。当 $ m = 1 $ 时，它由把子午曲线

\[
\operatorname{int}(\Gamma) := \{(\rho(t), 0, \sigma(t)) ; 0 < t < 1\}
\]

绕 $ z $ 轴旋转生成。

\proofhead
容易看出\footnote{见习题 4。} $ S^{m} \times [0, 1] $ 是 $ \mathbb{R}^{m+2} $ 的带边子流形，且其边界为 $ (S^{m} \times \{0\}) \cup (S^{m} \times \{1\}) $。对例 1.5(b) 的证明作显然的修改可知 $ f $ 是嵌入。现在这些断言由注 1.13(f) 得到。\nobreak\hspace{0.35em}■
\end{example}

\needspace{3\baselineskip}\subsection*{一维流形}

显然 $ \mathbb{R} $ 中每个完美区间 $ J $ 都是 $ \mathbb{R}^{n} $ 的一维子流形，带或不带边界，取决于 $ J $ 是否开。又，我们已经知道 1 维球面 $ S^{1} $ 是 $ \mathbb{R}^{n} $ 的一维子流形，只要 $ n \ge 2 $。容易看出\footnote{见习题 7。}，非空完美区间在开时微分同胚于 $ (0, 1) $，在一侧闭时微分同胚于 $ [0, 1) $，在紧时微分同胚于 $ [0, 1] $。下面的重要分类定理表明，这些区间与 $ S^{1} $ 在微分同胚意义下是一维连通流形的全部。

\begin{theorem}

设 $ C $ 是 $ N $ 的连通一维子流形，它带 [或不带] 边界。则 $ C $ 微分同胚于 $ [0, 1] $ 或 $ [0, 1) $ [或 $ (0, 1) $] 或 $ S^{1} $。

\proofhead
证明可参看 [BG88] 的 3.4 节，那里处理的是无边界流形。对那里的论证作显然的修改也可覆盖带边流形的情形（见 [Mil65] 的附录）。\nobreak\hspace{0.35em}■
\end{theorem}

\begin{remark}

\begin{items}
\item

我们把嵌入 $ N $ 中的（光滑）\textbf{曲线} $ C $ 理解为完美区间或 $ S^{1} $ 在（光滑）嵌入下的像。在后一情形，我们也称 $ C $ 为\textbf{嵌入 $ N $ 中的 1 维球面}。于是定理 1.18 表明：$ N $ 的每个连通一维子流形（带或不带边界）都是嵌入曲线，反之亦然。

\item

（正则性）定理 1.18 对 $ C^{1} $ 流形仍成立。\nobreak\hspace{0.35em}■

\end{items}

\end{remark}

\needspace{3\baselineskip}\subsection*{单位分解}

我们以证明一个技术性结果来结束本节，它在从局部到整体（以及反之）的过渡中将特别有用。

设 $ X $ 是 $ \mathbb{R}^{\overline{n}} $ 的 $ n $ 维子流形，带或不带边界，其中 $ \overline{n} \in \mathbb{N}^{\times} $。又设 $ \{U_{\alpha} ; \alpha \in \mathsf{A}\} $ 是 $ X $ 的开覆盖。则若族 $ \{\pi_{\alpha} ; \alpha \in \mathsf{A}\} $ 满足下列性质，我们称它为从属于该覆盖的\textbf{光滑单位分解}：

\begin{subitems}
\item

$ \pi_{\alpha} \in C^{\infty}(X, [0, 1]) $，且 $ \operatorname{supp}(\pi_{\alpha}) \subset\subset U_{\alpha} $（$ \alpha \in \mathsf{A} $）；

\item

族 $ \{\pi_{\alpha} ; \alpha \in \mathsf{A}\} $ 是\textbf{局部有限的}，即对每个 $ p \in X $ 存在开邻域 $ V $，使得除有限多个 $ \alpha \in \mathsf{A} $ 外 $ \operatorname{supp}(\pi_{\alpha}) \cap V = \varnothing $；

\item

对每个 $ p \in X $ 有 $ \sum_{\alpha \in \mathsf{A}} \pi_{\alpha}(p) = 1 $。

\end{subitems}

\begin{proposition}

$ X $ 的每个开覆盖都有从属于它的光滑单位分解。

证明

\begin{subitems}
\item

设 $ (\varphi, U) $ 是 $ X $ 中 $ p $ 附近的坐标卡。则 $ \varphi(U) $ 在 $ \overline{\mathbb{H}}^{n} $ 中是开的。因此存在 $ \varphi(p) $ 在 $ \overline{\mathbb{H}}^{n} $ 中的紧邻域 $ K' $，使得 $ K' \subset \varphi(U) $。因为 $ \varphi $ 是同胚，$ K := \varphi^{-1}(K') $ 是 $ X $ 中 $ p $ 的紧邻域，且 $ K \subset U $，并且 $ (\varphi|\mathring{K}, \mathring{K}) $ 是 $ p $ 附近的坐标卡。特别地，$ X $ 是局部紧的。
命题 X.7.14 蕴涵存在 $ \chi' \in C^{\infty}(\varphi(U), [0, 1]) $，使得 $ \chi'|K' = 1 $ 且 $ \operatorname{supp}(\chi') \subset\subset \varphi(U) $。我们对 $ q \in U $ 令 $ \chi(q) := \varphi^{\ast}\chi'(q) $，而对 $ q \in X \setminus U $ 令 $ \chi(q) := 0 $。则 $ \chi $ 属于 $ C^{\infty}(X, [0, 1]) $ 且具有紧支撑，该支撑包含在 $ U $ 中。

\item

由推论 IX.1.9(ii) 与注 X.1.16(e)，存在 $ X $ 的可数覆盖 $ \{V_{j} ; j \in \mathbb{N}\} $，它由相对紧的开集组成。我们令 $ K_{0} := \overline{V}_{0} $。则存在 $ i_{0}, \ldots, i_{m} \in \mathbb{N} $，使得 $ K_{0} $ 被 $ \{V_{i_{0}}, \ldots, V_{i_{m}}\} $ 覆盖。此外，我们令 $ j_{1} := \max\{i_{0}, \ldots, i_{m}\} + 1 $ 及 $ K_{1} := \bigcup_{i=0}^{j_{1}} \overline{V}_{i} $。集合 $ K_{1} $ 是紧的，且 $ K_{0} \subset\subset K_{1} $。然后我们归纳地得到紧集序列 $ (K_{j}) $，满足 $ K_{j} \subset\subset K_{j+1} $，且 $ \bigcup_{j=0}^{\infty} K_{j} = \bigcup_{j=0}^{\infty} V_{j} = X $。

\item

我们首先假设对 $ j \in \mathbb{N} $ 有 $ K_{j} \ne K_{j+1} $，并对 $ j \in \mathbb{N} $ 令 $ W_{j} := K_{j} \setminus \mathring{K}_{j-1} $，其中 $ K_{-1} := \varnothing $。则 $ W_{j} $ 是紧的，且当 $ |j - k| \ge 2 $ 时 $ W_{j} \cap W_{k} = \varnothing $。我们还有 $ \bigcup_{j=0}^{\infty} W_{j} = X $。
设 $ \mathcal{U} := \{U_{\alpha} ; \alpha \in \mathsf{A}\} $ 是 $ X $ 的开覆盖。由 (i) 及 $ W_{j} $ 的紧性可得，对每个 $ j \in \mathbb{N} $ 都存在 $ W_{j} $ 的有限覆盖 $ \{\widetilde{U}_{j,i} \in \mathcal{U} ; 0 \le i \le m(j)\} $。

\end{subitems}

\bookfigure[0.35]{figures/53426faadea39d9d.png}

我们令

\[
U_{j,i} := \widetilde{U}_{j,i} \cap (\mathring{W}_{j-1} \cup W_{j} \cup \mathring{W}_{j+1})
\]

并选取函数 $ \chi_{j,i} \in C^{\infty}(U_{j,i}, [0, 1]) $，使得

\[
\operatorname{supp}(\chi_{j,i}) \subset\subset U_{j,i} \subset \mathring{W}_{j-1} \cup W_{j} \cup \mathring{W}_{j+1} \quad \text{对 } 0 \le i \le m(j),
\]

其中 $ W_{-1} := \varnothing $，并且

\[
\bigcup_{i=0}^{m(j)} [\chi_{j,i} > 0] \supset W_j
\]

对 $ j \in \mathbb{N} $ 成立。于是 $ \{\chi_{j,i}; 0 \le i \le m(j), j \in \mathbb{N}\} $ 是局部有限族。因此

\[
\chi := \sum_{j=0}^{\infty} \sum_{i=0}^{m(j)} \chi_{j,i}
\]

有定义，属于 $ C^{\infty}(X, [0,1]) $，且对 $ p \in X $ 满足 $ \chi(p) > 0 $。现在令

\[
\pi_{\alpha} := \sum_{\alpha} \chi_{j,i}/\chi \quad \text{对 } \alpha \in \mathsf{A},
\]

其中 $ \sum_{\alpha} $ 表示对使 $ U_{j,i} $ 包含于 $ U_{\alpha} $ 的所有指标对 $ (j,i) $ 求和。于是 $ \{\pi_{\alpha}; \alpha \in \mathsf{A}\} $ 是从属于覆盖 $ \{U_{\alpha}; \alpha \in \mathsf{A}\} $ 的光滑单位分解。

\begin{subitems}
\item[（iv）]

若存在 $ j \in \mathbb{N} $ 使 $ K_j = K_{j+1} $，则 $ X = K_j $。因此 $ X $ 是紧的。在这种情形下，只需对 (iii) 作简单修改即得结论（因为这时只需考虑单个紧集，即 $ X $）。\nobreak\hspace{0.35em}■

\end{subitems}

\end{proposition}

下面的注 (a) 表明，命题 1.20 是定理 X.7.16 的一个意义深远的推广。

\begin{remark}

\begin{items}
\item

设 $ K $ 是流形 $ X $ 的紧子集，并设 $ \{U_j; 1 \le j \le m\} $ 是 $ K $ 的开覆盖。则存在函数 $ \pi_j \in C^{\infty}(X, [0,1]) $，使得对 $ 1 \le j \le m $ 有 $ \operatorname{supp}(\pi_j) \subset\subset U_j $，且对 $ p \in K $ 有 $ \sum_{j=1}^{m} \pi_j(p) = 1 $。

\proofhead
设 $ U_0 := X \setminus K $。则 $ \{U_j; 0 \le j \le m\} $ 是 $ X $ 的开覆盖。现在由命题 1.20 容易得到结论。\nobreak\hspace{0.35em}■
\item

命题 1.20 的证明表明，$ \mathbb{R}^m $ 的每个子流形（无论是否带边）都是局部紧的、具有可数基，并且是 $ \sigma $ 紧的。

\item

（正则性）设 $ k \in \mathbb{N}^{\times} $。把上述定义第 (i) 部分中的 $ \pi_{\alpha} \in C^{\infty}(X, [0,1]) $ 换成 $ \pi_{\alpha} \in C^{k}(X, [0,1]) $，便得到从属于覆盖 $ \{U_{\alpha}; \alpha \in \mathsf{A}\} $ 的 $ C^{k} $ 单位分解。此时，若把“光滑单位分解”换成“$ C^{k} $ 单位分解”，命题 1.20 仍然成立。在这种情形下，只需假设 $ X $ 属于 $ C^{k} $ 类。\nobreak\hspace{0.35em}■

\end{items}

约定 在本书余下的部分中，我们把每个流形都理解成适当“环境空间” $ \mathbb{R}^m $ 中的光滑带边子流形。

\end{remark}

\exerciseshead

\begin{exercises}
\item

设 $ f: M \to N $ 是淹没。证明 $ f $“局部看起来像投影”，即：对每个 $ p \in M $，存在 $ M $ 在 $ p $ 处的坐标卡 $ (\varphi, U) $ 与 $ N $ 在 $ f(p) $ 处的坐标卡 $ (\psi, V) $，使得 $ f(U) \subset V $，并且

\[
f_{\varphi,\psi}: \mathbb{R}^n \times \mathbb{R}^{m-n} \to \mathbb{R}^n, \quad (x,y) \mapsto x.
\]
\item

设 $ f: M \to N $ 是浸入。证明 $ f $ 局部看起来像典范单射 $ \mathbb{R}^m \to \mathbb{R}^m \times \mathbb{R}^{n-m} $，$ x \mapsto (x,0) $。

\item

证明从 $ M $ 到 $ N $ 的每个微分同胚局部看起来都像 $ \mathbb{R}^m $ 中的恒等映射。

\item

设 $ B $ 是 $ N $ 的带边子流形。证明 $ M \times B $ 是 $ M \times N $ 的带边子流形，且 $ \partial(M \times B) = M \times \partial B $。

\item

证明“截面”为 $ M $ 的柱体 $ [0,1] \times M $ 与“实心”环面 $ S^1 \times \overline{\mathbb{B}}^2 $ 都是带边流形。确定它们各自的维数与边界。

\item

证明 $ \mathbb{R}^n $ 中的闭 $ r $ 球 $ r\overline{\mathbb{B}}^n $ 与闭单位球 $ \overline{\mathbb{B}}^n $ 微分同胚。

\item

证明 $ \mathbb{R} $ 中的完全区间微分同胚于 $ (0,1) $、$ [0,1) $ 或 $ [0,1] $。

\item

设 $ B $ 是 $ M $ 的非空 $ k $ 维子流形（带边与否均可）。证明 $ B $ 的豪斯多夫维数等于 $ k $。
（提示：IX.3 的习题 4–6 与注 1.21(b)。）

\item

设 $ B $ 是 $ M $ 的带边子流形，$ f \in C^{\infty}(B, N) $。证明 $ \operatorname{graph}(f) $ 是 $ M \times N $ 的带边子流形，并确定它的边界。

\item

设 $ X $ 是 $ \mathbb{R}^n $ 的 $ n $ 维子流形（带边与否均可），并设 $ U := \{U_{\alpha}; \alpha \in \mathsf{A}\} $ 与 $ V := \{V_{\beta}; \beta \in \mathsf{B}\} $ 表示 $ X $ 的开覆盖。若存在 $ j: \mathsf{B} \to \mathsf{A} $，使得对 $ \beta \in \mathsf{B} $ 有 $ V_{\beta} \subset U_{j(\beta)} $，则称 $ V $ 是 $ U $ 的\textbf{加细}。证明：从属于 $ V $ 的每个光滑单位分解都诱导出从属于 $ U $ 的光滑单位分解。

\end{exercises}

\section{多重线性代数}

为了构造并理解更高次的微分形式微积分，我们需要线性代数（更确切地说，多重线性代数）的若干结果，本节即给出这些结果。

\begin{remark}

设 $ V $ 是有限维向量空间。

\begin{items}
\item

$ V $ 可以赋以一个内积 $ (\cdot|\cdot)_{V} $，使得 $ (V,(\cdot|\cdot)_{V}) $ 是希尔伯特空间。$ V $ 上所有范数都等价。

\proofhead
由注 I.12.5，存在向量空间同构 $ T:\mathbb{K}^{m} \to V $，且 $ m := \dim(V) $。于是
\[
(v|w)_{V} := (T^{-1}v|T^{-1}w) \quad \text{对 } v,w \in V
\]

定义了 $ V $ 上的一个内积，其中 $ (\cdot|\cdot) $ 表示 $ \mathbb{K}^{m} $ 中的欧几里得内积。因此 $ (V,(\cdot|\cdot)_{V}) $ 是有限维内积空间，从而由注 VII.1.7(b) 可知它是希尔伯特空间。第二个断言由推论 VII.1.5 推出。\nobreak\hspace{0.35em}■

\item

正如（泛函分析中的）惯例，我们用 $ V^{*} $ 表示从 $ V $ 到 $ \mathbb{C} $ 的所有（连续）共轭线性映射构成的空间，而 $ V' $ 表示 $ V $ 的对偶空间，即 $ V $ 上所有（连续）线性形式构成的空间。于是由 (a) 及里斯表示定理（定理 VII.2.14）可知，映射

\begin{equation}
V \to V^{*}, \quad v \mapsto (v|\cdot)_{V}\tag{2.1}
\end{equation}

是等距同构，而

\begin{equation}
V \to V', \quad v \mapsto (\cdot|v)_{V}\tag{2.2}
\end{equation}

是共轭线性的。若 $ \mathbb{K} = \mathbb{R} $，则 $ V^{*} = V' $，且映射 (2.1) 与 (2.2) 相同，这是因为每个实内积都是对称的。在下文中我们只处理实情形，因此出于这一原因以及某些历史原因，我们将写 $ V^{*} $ 而不写 $ V' $。\nobreak\hspace{0.35em}■

\end{items}

本节中，设

- $ V $ 与 $ W $ 是有限维实向量空间。

\end{remark}

\needspace{3\baselineskip}\subsection*{外积}

对 $ r \in \mathbb{N} $，记 $ \mathcal{L}^{r}(V,\mathbb{R}) $ 为所有 $ r $ 线性映射 $ V^{r} \to \mathbb{R} $ 构成的向量空间。由注 2.1(b) 与定理 VII.4.2(iii)，此记号与第 VII.4 节引入的记号一致。特别地，有

\[
\mathcal{L}^{0}(V,\mathbb{R}) = \mathbb{R} \quad \text{与} \quad \mathcal{L}^{1}(V,\mathbb{R}) = V^{*}.
\]

称 $ r $ 线性映射 $ \alpha: V^{r} \to W $ 为\textbf{交错的}，若 $ r \ge 2 $ 且

\[
\alpha(v_{\sigma(1)},\ldots,v_{\sigma(r)}) = \operatorname{sign}(\sigma)\,\alpha(v_{1},\ldots,v_{r}) \quad \text{对 } v_{1},\ldots,v_{r} \in V
\]

且对每个置换 $ \sigma \in S_{r} $ 成立（见习题 I.9.6）。我们令

\[
\bigwedge^{0}V^{*} := \mathcal{L}^{0}(V,\mathbb{R}) = \mathbb{R} \quad \text{与} \quad \bigwedge^{1}V^{*} := \mathcal{L}^{1}(V,\mathbb{R}) = V^{*}
\]

以及

\[
\bigwedge^{r}V^{*} := \{\, \alpha \in \mathcal{L}^{r}(V,\mathbb{R}) ;\ \alpha \text{ 是交错的} \,\} \quad \text{对 } r \ge 2.
\]

这里 $ \bigwedge^{r}V^{*} $ 称为 $ V^{*} $ 的 \textbf{$ r $ 重外积}（$ r \in \mathbb{N} $），而 $ \alpha \in \bigwedge^{r}V^{*} $ 称为 $ V $ 上的\textbf{交错 $ r $ 形式}（或简称为 \textbf{$ r $ 形式}）。

\begin{remark}

\begin{items}
\item

$ \bigwedge^{r}V^{*} $ 是 $ \mathcal{L}^{r}(V,\mathbb{R}) $ 的向量子空间，即 $ V $ 上\textbf{交错 $ r $ 形式构成的向量空间}。

\item

设 $ r \ge 2 $ 且 $ \alpha \in \mathcal{L}^{r}(V,\mathbb{R}) $。以下四个断言等价：

\begin{subitems}
\item

$ \alpha \in \bigwedge^{r}V^{*} $。

\item

若对某对满足 $ j \ne k $ 的 $ (j,k) $ 有 $ v_{j} = v_{k} $，则 $ \alpha(v_{1},\ldots,v_{r}) = 0 $。

\item

对 $ j \ne k $，$ \alpha(\ldots,v_{j},\ldots,v_{k},\ldots) = -\alpha(\ldots,v_{k},\ldots,v_{j},\ldots) $，也就是说，若交换 $ \alpha(v_{1},\ldots,v_{r}) $ 中的两个位置，其符号反转。

\item

若 $ v_{1},\ldots,v_{r} \in V $ 线性无关，则 $ \alpha(v_{1},\ldots,v_{r}) = 0 $。

\proofhead
蕴含式 “$ \text{(i)}\Rightarrow\text{(iii)}\Rightarrow\text{(ii)} $” 是显然的。
“$ \text{(ii)}\Rightarrow\text{(iv)} $” 设 $ v_{1},\ldots,v_{r} \in V $ 线性无关。这意味着存在 $ k \in \{1,\ldots,r\} $ 与 $ \lambda_{1},\ldots,\lambda_{r} \in \mathbb{R} $ 使得 $ \lambda_{k} = 0 $ 且 $ v_{k} = \sum_{j=1}^{r}\lambda_{j}v_{j} $。现在由 $ \alpha $ 对其第 $ k $ 个变量的线性性以及由 (ii) 可得
\[
\alpha(v_{1},\ldots,v_{r}) = \sum_{j=1}^{r}\lambda_{j}\alpha(v_{1},\ldots,\underset{(k)}{v_{j}},\ldots,v_{r}) = 0.
\]

“$ \text{(iv)}\Rightarrow\text{(iii)} $” 由 (iv) 及多重线性性，得

\begin{equation*}
\begin{aligned}
0 &= \alpha(\ldots,v_{j}+v_{k},\ldots,v_{j}+v_{k},\ldots) \\ &= \alpha(\ldots,v_{j},\ldots,v_{j},\ldots) + \alpha(\ldots,v_{j},\ldots,v_{k},\ldots) \\ &\qquad + \alpha(\ldots,v_{k},\ldots,v_{j},\ldots) + \alpha(\ldots,v_{k},\ldots,v_{k},\ldots) \\ &= \alpha(\ldots,v_{j},\ldots,v_{k},\ldots) + \alpha(\ldots,v_{k},\ldots,v_{j},\ldots),
\end{aligned}
\end{equation*}

这便证明了该断言。
“$ \text{(iii)}\Rightarrow\text{(i)} $” 这由每个置换都可写成对换之积这一事实推出（见习题 I.9.6）。\nobreak\hspace{0.35em}■

\end{subitems}

\item

$ \bigwedge^{r}V^{*} = \{0\} $ 对 $ r > \dim(V) $。

\proofhead
这由 (b) 的 (iv) 推出。\nobreak\hspace{0.35em}■
\end{items}

对 $ r \in \mathbb{N}^{\times} $ 与 $ \varphi^{1},\ldots,\varphi^{r} \in V^{*} $，\textbf{外积}\footnote{外积也称为楔积。}

\[
\varphi^{1}\wedge\cdots\wedge\varphi^{r}
\]

的定义为

\begin{equation}
\varphi^{1}\wedge\cdots\wedge\varphi^{r}(v_{1},\ldots,v_{r}) := \det\left[\langle\varphi^{j},v_{k}\rangle\right] = \det\begin{bmatrix} \langle\varphi^{1},v_{1}\rangle & \cdots & \langle\varphi^{1},v_{r}\rangle \\ \vdots & & \vdots \\ \langle\varphi^{r},v_{1}\rangle & \cdots & \langle\varphi^{r},v_{r}\rangle \end{bmatrix}\tag{2.3}
\end{equation}

其中 $ v_{1},\ldots,v_{r} \in V $。由线性代数可知，$ (r\times r) $ 矩阵的行列式是其列向量的交错 $ r $ 形式。由此以及 $ \varphi^{1},\ldots,\varphi^{r} $ 的线性性立即推出 $ \varphi^{1}\wedge\cdots\wedge\varphi^{r} $ 属于 $ \bigwedge^{r}V^{*} $：外积 $ \varphi^{1}\wedge\cdots\wedge\varphi^{r} $ 是 $ V $ 上的交错 $ r $ 形式。

\end{remark}

\qyxffntextdef{1}{若 $ \{e_{1},\ldots,e_{m}\} $ 是有序基，即其中元素的顺序是固定的，我们写 $ (e_{1},\ldots,e_{m}) $。}
\begin{proposition}

\begin{subitems}
\item

设 $ m := \dim(V) > 0 $。若 $ (e_{1},\ldots,e_{m}) $ 是 $ V $ 的基\qyxffnmark{1}，且 $ (\varepsilon^{1},\ldots,\varepsilon^{m}) $ 是 $ V^{*} $ 的相伴对偶基，则

\[
\{\, \varepsilon^{j_{1}}\wedge\cdots\wedge\varepsilon^{j_{r}} ;\ 1 \le j_{1} < j_{2} < \cdots < j_{r} \le m \,\}
\]

是 $ \bigwedge^{r}V^{*} $ 的基（$ 1 \le r \le m $）。

\item

$ \dim\left(\bigwedge^{r}V^{*}\right) = \binom{m}{r} $ 对 $ r \in \mathbb{N} $。

\proofhead
为简短起见，我们令
\[
\mathbb{J}_{r} := \mathbb{J}_{r}^{m} := \{\, (j) := (j_{1},\ldots,j_{r}) \in \mathbb{N}^{r} ;\ 1 \le j_{1} < j_{2} < \cdots < j_{r} \le m \,\}.
\]

此外，对有序多重指标 $ (j) \in \mathbb{J}_{r} $，令

\[
\varepsilon^{(j)} := \varepsilon^{j_{1}}\wedge\cdots\wedge\varepsilon^{j_{r}}.
\]
\end{subitems}

(i) 设 $ \alpha $ 是交错 $ r $ 形式。由于每个向量 $ v \in V $ 都有基表示 $ v = \sum_{k=1}^{m}\langle\varepsilon^{k},v\rangle e_{k} $，故由注 2.2(b) 得

\begin{equation*}
\begin{aligned}
\alpha(v_{1},\ldots,v_{r}) &= \sum_{k_{1}=1}^{m}\cdots\sum_{k_{r}=1}^{m}\langle\varepsilon^{k_{1}},v_{1}\rangle\cdots\langle\varepsilon^{k_{r}},v_{r}\rangle\alpha(e_{k_{1}},\ldots,e_{k_{r}}) \\ &= \sum_{(j)\in\mathbb{J}_{r}} a_{(j)} \sum_{\sigma\in S_{r}} \operatorname{sign}(\sigma)\langle\varepsilon^{\sigma(j_{1})},v_{1}\rangle\cdots\langle\varepsilon^{\sigma(j_{r})},v_{r}\rangle,
\end{aligned}
\end{equation*}

其中

\begin{equation}
a_{(j)} := \alpha(e_{j_{1}},\ldots,e_{j_{r}}).\tag{2.4}
\end{equation}

由注 VII.1.19(a) 以及方阵的行列式在转置下不变这一事实，我们可以把上一表达式的内和改写为

\[
\det\left(\left[\langle\varepsilon^{j_{\mu}},v_{\nu}\rangle\right]_{1\le\mu,\nu\le r}\right) = \varepsilon^{(j)}(v_{1},\ldots,v_{r}).
\]

因此

\[
\alpha(v_{1},\ldots,v_{r}) = \sum_{(j)\in\mathbb{J}_{r}} a_{(j)}\varepsilon^{(j)}(v_{1},\ldots,v_{r}) \quad \text{对 } v_{1},\ldots,v_{r} \in V,
\]

从而

\begin{equation}
\alpha = \sum_{(j)\in\mathbb{J}_{r}} a_{(j)}\varepsilon^{(j)}.\tag{2.5}
\end{equation}

这表明集合 $ \{\, \varepsilon^{(j)} ;\ (j) \in \mathbb{J}_{r} \,\} $ 张成向量空间 $ \bigwedge^{r}V^{*} $。

现在设

\[
\alpha = \sum_{(j)\in\mathbb{J}_{r}} b_{(j)}\varepsilon^{(j)}
\]

是 $ \alpha $ 的另一个表示，其中 $ b_{(j)} \in \mathbb{R} $。于是特别有

\[
\alpha(e_{k_{1}},\ldots,e_{k_{r}}) = \sum_{(j)\in\mathbb{J}_{r}} b_{(j)}\varepsilon^{(j)}(e_{k_{1}},\ldots,e_{k_{r}}) \quad \text{对 } (k) \in \mathbb{J}_{r}.
\]

因为

\begin{equation*}
\varepsilon^{(j)}(e_{k_{1}},\ldots,e_{k_{r}}) = \det\left(\left[\delta^{j_{\mu}}_{k_{\nu}}\right]_{1\le\mu,\nu\le r}\right) = \begin{cases} 1 & \text{若 } (j) = (k), \\ 0 & \text{其他情形}, \end{cases}
\end{equation*}

所以 $ b_{(j)} = a_{(j)} $ 对 $ (j) \in \mathbb{J}_{r} $。因此表示 (2.5) 是唯一的。

(ii) 该断言现在很清楚，因为 $ m $ 元集合恰好含有 $ \binom{m}{r} $ 个含 $ r $ 个元素的子集（见习题 I.6.3）。\nobreak\hspace{0.35em}■

\end{proposition}

下文中，令

\[
\alpha^{1}\wedge\cdots\wedge\widehat{\alpha^{j}}\wedge\cdots\wedge\alpha^{r}
\]

（$ 1 \le j \le r $）表示从 $ \alpha^{1}\wedge\cdots\wedge\alpha^{r} $ 中删去线性形式 $ \alpha^{j} $ 所得的 $ (r-1) $ 形式。当删去更多线性形式时，我们使用类似的记号，例如

\[
\alpha^{1}\wedge\cdots\wedge\widehat{\alpha^{j}}\wedge\cdots\wedge\widehat{\alpha^{k}}\wedge\cdots\wedge\alpha^{r}.
\]

\begin{example}

\begin{items}
\item

一维向量空间 $ \bigwedge^{0}V^{*} = \mathbb{R} $ 与 $ \bigwedge^{m}V^{*} $ 分别以 $ 1 $ 与 $ \varepsilon^{1}\wedge\cdots\wedge\varepsilon^{m} $ 为基。

\item

$ \{\varepsilon^{1},\ldots,\varepsilon^{m}\} $ 是 $ \bigwedge^{1}V^{*} = V^{*} $ 的基。

\item

$ \{\, \varepsilon^{1}\wedge\cdots\wedge\widehat{\varepsilon^{j}}\wedge\cdots\wedge\varepsilon^{m} ;\ 1 \le j \le m \,\} $ 是 $ \bigwedge^{m-1}V^{*} $ 的基。

\item

对基表示

\[
\alpha^{j} = \sum_{k=1}^{m} a^{j}_{k}\varepsilon^{k} \in V^{*} \quad \text{对 } 1 \le j \le r,
\]

与

\[
\alpha^{1}\wedge\cdots\wedge\alpha^{r} = \sum_{(j)\in\mathbb{J}_{r}} a_{(j)}\varepsilon^{(j)} \in \bigwedge^{r}V^{*},
\]

我们有 $ a^{i}_{k} = \langle\alpha^{i},e_{k}\rangle $（$ 1 \le i \le r $，$ 1 \le k \le m $）。又有

\[
a_{(j)} = \det\left(\left[a^{i}_{j_{k}}\right]_{1\le i,k\le r}\right) \quad \text{对 } (j) = (j_{1},\ldots,j_{r}) \in \mathbb{J}_{r}.
\]

\proofhead
这由 (2.3) 与 (2.4) 推出。\nobreak\hspace{0.35em}■
\item

对 $ r \ge 1 $，有

\[
\bigwedge^{r}V^{*} = \operatorname{span}\{\, \varphi^{1}\wedge\cdots\wedge\varphi^{r} ;\ \varphi^{j} \in V^{*},\ 1 \le j \le r \,\}.
\]
\item

对 $ r \ge 2 $，$ \varphi^{1}\wedge\cdots\wedge\varphi^{r} = 0 $ 当且仅当 $ \varphi^{1},\ldots,\varphi^{r} $ 线性无关。

\proofhead
这由 (2.3) 推出。\nobreak\hspace{0.35em}■
\end{items}

正如下一个命题所示，我们可以利用基表示定义从 $ \bigwedge^{r}V^{*}\times\bigwedge^{s}V^{*} $ 到 $ \bigwedge^{r+s}V^{*} $ 的双线性映射。

\end{example}

\begin{proposition}

设 $ r,s \in \mathbb{N}^{\times} $。

\begin{subitems}
\item

恰有一个映射

\begin{equation}
\wedge:\ \bigwedge^{r}V^{*}\times\bigwedge^{s}V^{*} \to \bigwedge^{r+s}V^{*}, \quad (\alpha,\beta) \mapsto \alpha\wedge\beta,\tag{2.6}
\end{equation}

即\textbf{外积}，具有下述性质：
(α) $ \wedge $ 是双线性的；
(β) 对 $ \varphi^{1},\ldots,\varphi^{r},\psi^{1},\ldots,\psi^{s} \in V^{*} $，

\begin{equation}
(\varphi^{1}\wedge\cdots\wedge\varphi^{r})\wedge(\psi^{1}\wedge\cdots\wedge\psi^{s}) = \varphi^{1}\wedge\cdots\wedge\varphi^{r}\wedge\psi^{1}\wedge\cdots\wedge\psi^{s}.\tag{2.7}
\end{equation}
\item

给定基表示

\begin{equation}
\alpha = \sum_{(j)\in\mathbb{J}_{r}} a_{(j)}\varepsilon^{(j)} \quad \text{与} \quad \beta = \sum_{(k)\in\mathbb{J}_{s}} b_{(k)}\varepsilon^{(k)},\tag{2.8}
\end{equation}

则有

\begin{equation}
\alpha\wedge\beta = \sum_{\substack{(j)\in\mathbb{J}_{r}\\ (k)\in\mathbb{J}_{s}}} a_{(j)}b_{(k)}\varepsilon^{(j)}\wedge\varepsilon^{(k)}.\tag{2.9}
\end{equation}
\item

外积是结合且\textbf{分次反交换}的，也就是说，

\[
\alpha\wedge\beta = (-1)^{rs}\beta\wedge\alpha \quad \text{对 } \alpha \in \bigwedge^{r}V^{*} \text{ 与 } \beta \in \bigwedge^{s}V^{*}.
\]

\proofhead
若 $ \wedge $ 是某个从 $ \bigwedge^{r}V^{*}\times\bigwedge^{s}V^{*} $ 到 $ \bigwedge^{r+s}V^{*} $ 的线性映射并满足 (2.7)，则由 (2.8) 立即推出 (2.9) 成立。因此我们可以用 (2.9) 和一个给定的基来唯一地定义（即把基元素双线性延拓到整个空间）具有性质 (α) 与 (β) 的双线性映射 (2.6)。由 (2.3)、(2.7) 与例 2.4(e)，$ \wedge $ 与所选基无关。(iii) 此时是行列式性质的直接推论。\nobreak\hspace{0.35em}■
\end{subitems}

\end{proposition}

\begin{remark}

设 $ E_{k} $（$ k \in \mathbb{N} $）是同一域 $ K $ 上的向量空间。

\begin{items}
\item

\textbf{直和}

\[
E := \bigoplus_{k=0}^{\infty}E_{k} =: \bigoplus_{k\ge 0}E_{k}
\]

的定义如下：
$ E $ 是所有位于 $ \bigcup_{k=0}^{\infty}E_{k} $ 中的序列 $ (x_{k}) $ 构成的集合，其中 $ x_{k} \in E_{k} $（$ k \in \mathbb{N} $），且满足对几乎所有 $ k \in \mathbb{N} $ 有 $ x_{k} = 0 $。在 $ E $ 上，加法 $ + $ 与数乘由下式定义

\[
(x_{k}) + \lambda(y_{k}) := (x_{k} + \lambda y_{k}) \quad \text{对 } (x_{k}),(y_{k}) \in E \text{ 与 } \lambda \in \mathbb{K}.
\]

于是 $ E $ 是 $ K $ 向量空间。\footnote{$ E $ 是所有具有紧支集的映射 $ f:\mathbb{N} \to \bigcup_{k=0}^{\infty}E_{k} $（其中 $ f(k) \in E_{k} $，$ k \in \mathbb{N} $）构成的向量空间，赋以例 I.12.3(e) 的逐点乘积。} 此外，$ E_{k} $ 将通过线性映射

\[
E_{k} \to E, \quad x_{k} \mapsto (0,\ldots,0,x_{k},0,\ldots),
\]

与一个向量子空间等同，其中 $ x_{k} $ 占据右端序列的第 $ k $ 个位置。显然

\[
E = \operatorname{span}\{E_{k} ;\ k \in \mathbb{N}\} \quad \text{与} \quad E_{k}\cap E_{j} = \{0\} \quad \text{对 } k \ne j,
\]

这解释了“直和”这一名称（见例 I.12.3(l)）。

\item

令 $ E := \bigoplus_{k\ge 0}E_{k} $，我们定义一个乘法

\[
E\times E \to E, \quad (v,w) \mapsto v\odot w
\]

使得 $ (E,+,\odot) $ 是一个代数（见第 I.12 节）。我们称 $ E $ 为（$ K $ 上的）\textbf{分次代数}，并称该乘法是\textbf{分次的}，若

\[
E_{k}\odot E_{\ell} \subset E_{k+\ell} \quad \text{对 } k,\ell \in \mathbb{N}.
\]

若还满足关系

\[
v_{k}\odot v_{\ell} = (-1)^{k\ell}v_{\ell}\odot v_{k} \quad \text{对 } k,\ell \in \mathbb{N},
\]

则称该乘法与该代数都是\textbf{分次反交换}的。\nobreak\hspace{0.35em}■

\end{items}

我们令

\[
\bigwedge V^{*} := \bigoplus_{r\ge 0}\bigwedge^{r}V^{*}
\]

并通过定义

\begin{equation}
\alpha\wedge\beta := \beta\wedge\alpha := \alpha\beta \quad \text{对 } \alpha \in \bigwedge^{0}V^{*} = \mathbb{R},\ \beta \in \bigwedge V^{*}.\tag{2.10}
\end{equation}

来延拓外积的定义。我们还令 $ \mathbb{J}_{0} := \{0\} $。

\end{remark}

\begin{theorem}

\begin{subitems}
\item

恰有一个双线性、结合且分次反交换的映射

\[
\bigwedge V^{*}\times\bigwedge V^{*} \to \bigwedge V^{*}
\]

它把外积 (2.6) 与/或 (2.10) 延拓到整个 $ \bigwedge V^{*}\times\bigwedge V^{*} $ 上。它也将记作 $ \wedge $，并称为 $ \bigwedge V^{*} $ 上的\textbf{外积}。

\item

$ \dim\left(\bigwedge V^{*}\right) = 2^{\dim(V)} $。

\proofhead
(i) 这由命题 2.5 与定义 (2.10) 经自然的双线性延拓立即推出。
\end{subitems}

(ii) 因为 $ \bigwedge^{r}V^{*} = \{0\} $ 对 $ r > \dim(V) $，且 $ \bigwedge V^{*} $ 是各向量子空间 $ \bigwedge^{r}V^{*} $ 的直和，故由命题 2.3(ii) 与二项式定理得

\[
\dim\left(\bigwedge V^{*}\right) = \sum_{r=0}^{m}\binom{m}{r} = 2^{m}
\]

其中 $ m := \dim(V) $。\nobreak\hspace{0.35em}■

\end{theorem}

这个定理表明，当赋以自然的向量空间结构与外积时，$ \bigwedge V^{*} $ 是维数为 $ 2^{\dim(V)} $ 的结合、分次反交换实代数。它称为 $ V^{*} $ 的\textbf{格拉斯曼代数}（或\textbf{外代数}）。

\begin{remark}

因为 $ V $ 是有限维的，所以可以用\textbf{典范同构}

\[
\kappa : V \to V^{\ast\ast} := (V^{*})^{*}
\]

把 $ V $ 与 $ V^{\ast\ast} $ 等同起来，该同构由

\[
\langle \kappa(v), v^{*}\rangle_{V^{*}} := \langle v^{*}, v\rangle \qquad \text{对 } v \in V , \quad v^{*} \in V^{*}
\]

定义。因此，格拉斯曼代数

\[
\bigwedge V := \bigoplus_{r \ge 0} \bigwedge^{r} V
\]

在 $ V $ 上也是良定义的。

\proofhead
显然 $ \kappa : V \to V^{\ast\ast} $ 是线性的。设 $ \{e_{1},\ldots,e_{m}\} $ 是 $ V $ 的基，$ \{\varepsilon_{1},\ldots,\varepsilon_{m}\} $ 是 $ V^{*} $ 的相应对偶基。对 $ v \in \ker(\kappa) $，有
\[
\langle \varepsilon_{j}, v\rangle = \langle \kappa(v), \varepsilon_{j}\rangle_{V^{*}} = 0 \qquad \text{对 } j = 1,\ldots,m .
\]

于是 $ v = \sum_{j=1}^{m}\langle \varepsilon_{j}, v\rangle e_{j} $ 蕴涵 $ v = 0 $，故 $ \kappa $ 是单射。又 $ \dim(V^{\ast\ast}) = m $（见定理 VII.2.14）蕴涵 $ \kappa $ 是同构。 \nobreak\hspace{0.35em}■

\end{remark}

\needspace{3\baselineskip}\subsection*{拉回}

对 $ A \in \mathcal{L}(V,W) $ 与 $ \alpha \in \bigwedge^{r}W^{*} $，当 $ r \ge 1 $ 时定义 $ A^{*}\alpha $ 为

\[
A^{*}\alpha(v_{1},\ldots,v_{r}) := \alpha(Av_{1},\ldots,Av_{r}) \qquad \text{对 } v_{1},\ldots,v_{r} \in V ,
\]

而当 $ r = 0 $ 时，对 $ \alpha \in \bigwedge^{0}W^{*} = \mathbb{R} $ 定义

\[
A^{*}\alpha := \alpha .
\]

则称 $ A^{*}\alpha $ 为 $ \alpha $ 沿 $ A $ 在 $ V $ 上的\textbf{拉回}。

\begin{remark}

\begin{items}
\item

对 $ \alpha \in \bigwedge^{r}W^{*} $，拉回 $ A^{*}\alpha $ 属于 $ \bigwedge^{r}V^{*} $，且映射 $ A^{*} $ 是线性的：

\[
A^{*} \in \mathcal{L}(\bigwedge W^{*}, \bigwedge V^{*}) \qquad \text{且 } A(\bigwedge^{r}W^{*}) \subset \bigwedge^{r}V^{*} , \quad r \in \mathbb{N} .
\]

我们把 $ A^{*} $ 称为由 $ A $ 给出的\textbf{拉回变换}（或通常就称为\textbf{拉回}）。
当 $ r = 1 $ 时，$ A^{*} $ 是对偶于 $ A $ 的映射（在 §VIII.3 中记为 $ A^{\top} $）。还要注意，$ A $ 把向量空间 $ V $ 映到 $ W $，而 $ A^{*} $ 把 $ \bigwedge W^{*} $ 映到 $ \bigwedge V^{*} $，因而是“反方向”的：

\[
V \xrightarrow{A} W
\]
\[
\bigwedge V^{*} \xleftarrow{A^{*}} \bigwedge W^{*}
\]
\item

若 $ X $ 是另一个有限维实向量空间，且 $ B \in \mathcal{L}(W,X) $，则

\[
(BA)^{*} = A^{*}B^{*} \qquad \text{且 } (\mathrm{id}_{V})^{*} = \mathrm{id}_{\bigwedge V^{*}} .
\]

换言之，映射 $ A \mapsto A^{*} $ 是\textbf{反变}的。

\item

我们有

\[
A^{*}(\alpha \wedge \beta) = A^{*}\alpha \wedge A^{*}\beta \qquad \text{对 } \alpha,\beta \in \bigwedge W^{*} .
\]

因此 $ A^{*} $ 是从 $ \bigwedge W^{*} $ 到 $ \bigwedge V^{*} $ 的代数同态。

\proofhead
这些结论由拉回与外积的定义显然推出。 \nobreak\hspace{0.35em}■
\end{items}

设 $ m := \dim(V) $，$ W := V $，并设 $ \alpha \in \bigwedge^{m}V^{*} $。根据命题 2.3(ii)，$ \bigwedge^{m}V^{*} $ 是一维的，因此 $ A^{*}\alpha $ 必与 $ \alpha $ 成比例；下面我们来确定这个倍数。

\end{remark}

\begin{proposition}

对 $ m := \dim(V) $ 与 $ A \in \mathcal{L}(V) $，

\[
A^{*}\alpha = \det(A)\alpha \qquad \text{对 } \alpha \in \bigwedge^{m}V^{*} .
\]

\proofhead
设 $ \{e_{1},\ldots,e_{m}\} $ 是 $ V $ 的基，$ [a_{k}^{j}] \in \mathbb{R}^{m\times m} $ 是 $ A $ 在这组基下的矩阵（见 §VII.1）。则
\[
Ae_{k} = \sum_{j=1}^{m} a_{k}^{j}e_{j} \qquad \text{对 } 1 \le k \le m .
\]

由此以及 $ \alpha \in \bigwedge^{m}V^{*} $ 的性质可得

\begin{equation*}
\begin{aligned}
A^{*}\alpha(e_{1},\ldots,e_{m}) &= \alpha(Ae_{1},\ldots,Ae_{m}) \\ &= \sum_{j_{1}=1}^{m}\cdots\sum_{j_{m}=1}^{m} a_{1}^{j_{1}}\cdots a_{m}^{j_{m}}\alpha(e_{j_{1}},\ldots,e_{j_{m}}) \\ &= \sum_{\sigma \in S_{m}} \operatorname{sign}(\sigma)\, a_{1}^{\sigma(1)}\cdots a_{m}^{\sigma(m)}\alpha(e_{1},\ldots,e_{m}) \\ &= \det(A)\,\alpha(e_{1},\ldots,e_{m}) ,
\end{aligned}
\end{equation*}

其中最后一步用了注 VII.1.19 的符号公式以及 $ \det(A^{\top}) = \det(A) $ 这一事实。现由 $ \alpha $ 的多重线性性即得结论。 \nobreak\hspace{0.35em}■

\end{proposition}

\needspace{3\baselineskip}\subsection*{体积元}

现在设 $ \mathcal{O}r $ 是 $ V $ 的一个定向，即 $ V := (V, \mathcal{O}r) $ 是有向向量空间。为简便起见，我们把 $ V $ 的正向有序基（因而它是 $ \mathcal{O}r $ 的元素）称为\textbf{正基}（见注 VIII.2.4）。又令 $ m := \dim(V) $。

我们把每个 $ \alpha \in \bigwedge^{m}V^{*}\setminus\{0\} $ 称为 $ V $ 上的\textbf{体积形式}。两个体积形式 $ \alpha $ 与 $ \beta $ 称为\textbf{等价}的，若存在 $ \lambda > 0 $ 使得 $ \alpha = \lambda\beta $。容易验证，这个定义在所有体积形式构成的集合上诱导出一个等价关系 $ \sim $。因为 $ \dim(\bigwedge^{m}V^{*}) = 1 $，所以恰有两个等价类。

\begin{remark}

\begin{items}
\item

设 $ (e_{1},\ldots,e_{m}) $ 是 $ V $ 的正基，$ (\varepsilon^{1},\ldots,\varepsilon^{m}) $ 是相应对偶基。若 $ (\widetilde{e}_{1},\ldots,\widetilde{e}_{m}) $ 是 $ V $ 的基，且 $ \alpha \in \bigwedge^{m}V^{*}\setminus\{0\} $ 满足 $ \alpha \sim \varepsilon^{1}\wedge\cdots\wedge\varepsilon^{m} $，则 $ (\widetilde{e}_{1},\ldots,\widetilde{e}_{m}) $ 为正，当且仅当 $ \alpha(\widetilde{e}_{1},\ldots,\widetilde{e}_{m}) > 0 $。这意味着 $ \bigwedge^{m}V^{*}\setminus\{0\} $ 的两个等价类可以与 $ V $ 的两个定向等同起来。换言之，体积形式 $ \alpha $ 通过下述要求\textbf{确定}了 $ V $ 的定向 $ \mathcal{O}r $：

\[
\alpha(e_{1},\ldots,e_{m}) > 0 \iff (e_{1},\ldots,e_{m}) \in \mathcal{O}r .
\]

\proofhead
设 $ B \in \mathcal{L}(V) $ 是从 $ (e_{1},\ldots,e_{m}) $ 到 $ (\widetilde{e}_{1},\ldots,\widetilde{e}_{m}) $ 的基变换，即 $ \widetilde{e}_{j} = Be_{j} $，$ 1 \le j \le m $。则由命题 2.10 得
\[
\alpha(\widetilde{e}_{1},\ldots,\widetilde{e}_{m}) = \det(B)\alpha(e_{1},\ldots,e_{m}) = \det(B)\lambda
\]

其中 $ \alpha = \lambda\varepsilon^{1}\wedge\cdots\wedge\varepsilon^{m} $ 且 $ \lambda > 0 $。 \nobreak\hspace{0.35em}■

\item

若 $ \det(A) > 0 $（或 $ \det(A) < 0 $），则称 $ V $ 的自同构 $ A $ 是\textbf{保持定向}（或\textbf{反转定向}）的。令

\[
\mathcal{L}\mathrm{aut}^{+}(V) := GL^{+}(V) := \{\, A \in \mathcal{L}\mathrm{aut}(V) ;\ \det(A) > 0 \,\} .
\]
\end{items}

\begin{subitems}
\item

对 $ A \in \mathcal{L}\mathrm{aut}(V) $，下列陈述等价：
(α) $ A \in \mathcal{L}\mathrm{aut}^{+}(V) $。
(β) 对每个基 $ (b_{1},\ldots,b_{m}) $，基 $ (b_{1},\ldots,b_{m}) $ 与 $ (Ab_{1},\ldots,Ab_{m}) $ 具有相同的定向。
(γ) 对每个 $ \alpha \in \bigwedge^{m}V^{*}\setminus\{0\} $，体积形式 $ \alpha $ 与 $ A^{*}\alpha $ 确定了 $ V $ 的相同定向。

\item

$ \mathcal{L}\mathrm{aut}^{+}(V) $ 是 $ \mathcal{L}\mathrm{aut}(V) =: GL(V) $ 的子群。

\proofhead
(i) 这由 $ A^{*}\alpha = \det(A)\alpha $ 以及定向的定义推出。
\end{subitems}

(ii) 映射

\[
\mathcal{L}\mathrm{aut}(V) \to (\mathbb{R}^{\times}, \cdot) , \qquad A \mapsto \det(A)
\]

是同态。根据习题 I.7.5，$ \mathcal{L}\mathrm{aut}^{+}(V) $ 作为 $ (\mathbb{R}^{\times}, \cdot) $ 的子群 $ ((0,\infty), \cdot) $ 的原像，是 $ \mathcal{L}\mathrm{aut}(V) $ 的子群。 \nobreak\hspace{0.35em}■

现在设 $ (V, (\cdot\mid\cdot), \mathcal{O}r) $ 是有向内积空间。又设 $ (e_{1},\ldots,e_{m}) $ 是正标准正交基（ONB），$ (\varepsilon^{1},\ldots,\varepsilon^{m}) $ 是 $ V^{*} $ 的相应对偶基。则

\[
\omega := \omega_{V} := \varepsilon^{1}\wedge\cdots\wedge\varepsilon^{m}
\]

称为 $ V $ 的\textbf{体积元}。

\end{remark}

\begin{remark}

\begin{items}
\item

对 $ V $ 的每个正 ONB $ (\widetilde{e}_{1},\ldots,\widetilde{e}_{m}) $，有

\[
\omega(\widetilde{e}_{1},\ldots,\widetilde{e}_{m}) = 1 .
\]

\proofhead
设 $ B $ 是由 $ \widetilde{e}_{j} = Be_{j} $（$ 1 \le j \le m $）给出的基变换。则 $ B $ 属于 $ \mathcal{L}\mathrm{aut}^{+}(V) \cap O(m) $，从而 $ \det(B) = 1 $（见习题 VII.9.2）。因此由命题 2.10 得
\[
\omega(\widetilde{e}_{1},\ldots,\widetilde{e}_{m}) = B^{*}\omega(e_{1},\ldots,e_{m}) = \det(B)\varepsilon^{1}\wedge\cdots\wedge\varepsilon^{m}(e_{1},\ldots,e_{m}) = 1 ,
\]

这就证明了结论。 \nobreak\hspace{0.35em}■

\item

$ V $ 的体积元是唯一一个对任意（从而每个）正 ONB 都取值 $ 1 $ 的体积形式。

\proofhead
这由 (a) 推出。 \nobreak\hspace{0.35em}■
\item

对 $ v_{1},\ldots,v_{m} \in \mathbb{R}^{m} $，令

\[
P(v_{1},\ldots,v_{m}) := \left\{ \sum_{j=1}^{m} t^{j}v_{j} ;\ 0 \le t^{j} \le 1 \right\} ,
\]

即 $ P(v_{1},\ldots,v_{m}) $ 是由 $ v_{1},\ldots,v_{m} $ 张成的平行六面体。则

\[
|\omega_{\mathbb{R}^{m}}(v_{1},\ldots,v_{m})| = \mathrm{vol}_{m}(P(v_{1},\ldots,v_{m})) := \lambda_{m}(P(v_{1},\ldots,v_{m})) .
\]

换言之，体积元给每个 $ m $ 元向量组赋予它们张成的平行四边形的有向体积\footnote{有向体积为正，当且仅当 $ (v_{1},\ldots,v_{m}) $ 为正；它为负，当且仅当 $ (v_{1},\ldots,v_{m}) $ 属于 $ -\mathcal{O}r $。}。

\proofhead
我们由 $ v_{j} = Be_{j} $（$ 1 \le j \le m $）定义 $ B \in \mathcal{L}(\mathbb{R}^{m}) $。则
\[
P(v_{1},\ldots,v_{m}) = B([0,1]^{m}) .
\]

于是由命题 2.10 与 (a) 得

\[
\omega_{\mathbb{R}^{m}}(v_{1},\ldots,v_{m}) = B^{*}\omega_{\mathbb{R}^{m}}(e_{1},\ldots,e_{m}) = \det(B) .
\]

由定理 IX.5.25 知 $ \lambda_{m}(B([0,1]^{m})) = |\det(B)| $，这正是所要的。 \nobreak\hspace{0.35em}■

\end{items}

在下面的命题中，我们用 $ V $ 的任意正基来表示体积元 $ \omega $。

\end{remark}

\begin{proposition}

设 $ (b_{1},\ldots,b_{m}) $ 是 $ V $ 的正基，$ (\beta^{1},\ldots,\beta^{m}) $ 是它的对偶基。则

\[
\omega = \sqrt{G}\,\beta^{1}\wedge\cdots\wedge\beta^{m} ,
\]

其中 $ G := \det[(b_{j}\mid b_{k})] $ 是\textbf{格拉姆行列式}。特别地，

\[
\omega(b_{1},\ldots,b_{m}) = \sqrt{G} .
\]

\proofhead
设 $ (e_{1},\ldots,e_{m}) $ 是 $ V $ 的正 ONB，并由 $ b_{j} = Be_{j} $（$ 1 \le j \le m $）定义 $ B \in \mathcal{L}(V) $。根据注 2.11(b) 的 (i)，有 $ \det(B) > 0 $。由注 2.12(a) 与 (2.3) 得
\begin{equation*}
\begin{aligned}
\omega(e_{1},\ldots,e_{m}) = 1 &= \beta^{1}\wedge\cdots\wedge\beta^{m}(b_{1},\ldots,b_{m}) \\ &= B^{*}(\beta^{1}\wedge\cdots\wedge\beta^{m})(e_{1},\ldots,e_{m}) .
\end{aligned}
\end{equation*}

因此

\[
\omega = \det(B)\beta^{1}\wedge\cdots\wedge\beta^{m} ,
\]

因为 $ m $ 形式由它在 $ V $ 的一组基上的值确定，也依据命题 2.10。又

\begin{equation}
(b_{j}\mid b_{k}) = (Be_{j}\mid Be_{k}) = (B^{*}Be_{j}\mid e_{k}) \qquad \text{对 } 1 \le j,k \le m\tag{2.11}
\end{equation}

（见习题 VII.1.5）。因为 $ (e_{1},\ldots,e_{m}) $ 是 ONB，任意 $ v \in V $ 有表示 $ v = \sum_{k=1}^{m}(v\mid e_{k})e_{k} $。由此可得

\[
Te_{j} = \sum_{k=1}^{m}(Te_{j}\mid e_{k})e_{k} \qquad \text{对 } 1 \le j \le m \text{ 与 } T \in \mathcal{L}(V) .
\]

因此 (2.11) 表明 $ [(b_{j}\mid b_{k})] \in \mathbb{R}^{m\times m} $ 是 $ B^{*}B $ 在基 $ (e_{1},\ldots,e_{m}) $ 下的矩阵。所以

\[
G = \det[(b_{j}\mid b_{k})] = \det(B^{*}B) = (\det(B))^{2}
\]

因为 $ \det(B^{*}) = \det(B) $。结论得证。 \nobreak\hspace{0.35em}■

\end{proposition}

\needspace{3\baselineskip}\subsection*{里斯同构}

设 $ (V, (\cdot\mid\cdot)) $ 是内积空间，$ m := \dim(V) $。把里斯同构 (2.2) 记为

\[
\Theta := \Theta_{V} : V \to V^{*} , \qquad v \mapsto (\cdot\mid v) ,
\]

即

\begin{equation}
\langle \Theta v, w\rangle = (w\mid v) \qquad \text{对 } v,w \in V .\tag{2.12}
\end{equation}

则

\begin{equation}
(\alpha\mid\beta)_{*} := (\Theta^{-1}\alpha\mid\Theta^{-1}\beta) \qquad \text{对 } \alpha,\beta \in V^{*}\tag{2.13}
\end{equation}

定义了 $ V^{*} $ 上的内积，即对偶于 $ (\cdot\mid\cdot) $ 的\textbf{标量积}。以下我们总为 $ V^{*} $ 配备这个内积，于是 $ V^{*} := (V^{*}, (\cdot\mid\cdot)_{*}) $ 是内积空间。

\begin{remark}

设 $ \{e_{1},\ldots,e_{m}\} $ 是 $ V $ 中的基，$ \{\varepsilon^{1},\ldots,\varepsilon^{m}\} $ 是它在 $ V^{*} $ 中的对偶基。

\begin{items}
\item

令

\[
g_{jk} := (e_{j}\mid e_{k}) \qquad \text{对 } 1 \le j,k \le m , \quad \text{其中 } [g^{jk}] := [g_{jk}]^{-1} \in \mathbb{R}^{m\times m} .
\]

则

\[
\Theta e_{j} = \sum_{k=1}^{m} g_{jk}\varepsilon^{k} \quad \text{与} \quad \Theta^{-1}\varepsilon^{j} = \sum_{k=1}^{m} g^{jk}e_{k} \qquad \text{对 } 1 \le j \le m .
\]

\proofhead
由基展开 $ \Theta e_{j} = \sum_{k=1}^{m} a_{jk}\varepsilon^{k} $（$ 1 \le j \le m $）以及 (2.12) 得
\[
(e_{j}\mid v) = \langle \Theta e_{j}, v\rangle = \sum_{k=1}^{m} a_{jk}\langle \varepsilon^{k}, v\rangle \qquad \text{对 } v \in V \text{ 与 } 1 \le j \le m .
\]

用 $ e_{1},\ldots,e_{m} $ 逐个替换 $ v $，得 $ a_{jk} = (e_{j}\mid e_{k}) $，这证明了第一个结论。$ \Theta^{-1}\varepsilon^{j} $ 的表示是显然的。 \nobreak\hspace{0.35em}■

\item

对 $ v = \sum_{j=1}^{m}\xi^{j}e_{j} \in V $ 与 $ w = \sum_{j=1}^{m}\eta^{j}e_{j} \in V $，有

\[
(v\mid w) = \sum_{j,k=1}^{m} g_{jk}\xi^{j}\eta^{k} .
\]

对 $ \alpha = \sum_{j=1}^{m}a_{j}\varepsilon^{j} \in V^{*} $ 与 $ \beta = \sum_{j=1}^{m}b_{j}\varepsilon^{j} \in V^{*} $，有关系

\[
(\alpha\mid\beta)_{*} = \sum_{j,k=1}^{m} g^{jk}a_{j}b_{k} .
\]

\proofhead
第一个结论是显然的。由 (a) 与 (2.13) 可推出
\[
(\varepsilon^{i}\mid\varepsilon^{\ell})_{*} = (\Theta^{-1}\varepsilon^{i}\mid\Theta^{-1}\varepsilon^{\ell}) = \sum_{j,k=1}^{m} g^{ij}g^{\ell k}(e_{j}\mid e_{k}) = \sum_{j,k=1}^{m} g^{ij}g^{\ell k}g_{jk} = g^{i\ell}
\]

对 $ 1 \le i,\ell \le m $。现由 $ (\cdot\mid\cdot)_{*} $ 的双线性性即得第二个结论。 \nobreak\hspace{0.35em}■

\item

若 $ \{e_{1},\ldots,e_{m}\} $ 是 ONB，则 $ \Theta e_{j} = \varepsilon^{j} $（$ 1 \le j \le m $），且 $ \{\varepsilon^{1},\ldots,\varepsilon^{m}\} $ 同样是 ONB。

\item

你也许已经注意到，我们用上标来标记基表示中向量的系数，而用下标表示 1 形式的展开系数。也就是说，

\[
v = \sum_{j=1}^{m}\xi^{j}e_{j} \in V \quad \text{与} \quad \alpha = \sum_{j=1}^{m} a_{j}\varepsilon^{j} \in V^{*} .
\]

由 (a) 及 $ [g_{jk}] $ 的对称性可得

\[
\Theta v = \sum_{j=1}^{m} b_j\varepsilon^{j} \in V^{*} \quad \text{且} \quad \Theta^{-1}\alpha = \sum_{j=1}^{m} \eta^{j}e_j \in V
\]

其中

\[
b_j := \sum_{k=1}^{m} g_{jk}\xi^{k} \quad \text{且} \quad \eta^{j} := \sum_{k=1}^{m} g^{jk}a_k \quad \text{对 } 1 \le j \le m.
\]

$ \Theta $ [或 $ \Theta^{-1} $] 的作用在形式上实现指标的\textbf{下降} [或上升]。据此，我们可以借用音乐记号

\[
g^{\flat} := \Theta \quad \text{且} \quad g^{\sharp} := \Theta^{-1},
\]

或简写为对 $ v \in V $ 的 $ v^{\flat} := \Theta v $，以及对 $ \alpha \in V^{*} $ 的 $ \alpha^{\sharp} := \Theta^{-1}\alpha $。\nobreak\hspace{0.35em}■

\end{items}

\end{remark}

\needspace{3\baselineskip}\subsection*{霍奇星算子\footnote{本节与下一节初读时可以跳过。}}

设 $ (V, (\cdot \mid \cdot), \mathcal{O}r) $ 是定向内积空间，$ m := \dim(V) $，且 $ \omega $ 是 $ V $ 的体积元。又设 $ \{e_1,\ldots,e_m\} $ 是 $ V $ 的标准正交基，$ \{\varepsilon^1,\ldots,\varepsilon^m\} $ 是其对偶基。

现在我们如下在 $ \bigwedge^{r}V^{*} $ 上定义一个标量积 $ (\cdot \mid \cdot)_r $：

对 $ r = 0 $，令

\begin{equation}
(\alpha \mid \beta)_0 := \alpha\beta \quad \text{对 } \alpha, \beta \in \bigwedge^{0}V^{*} = \mathbb{R}.\tag{2.14}
\end{equation}

对 $ 1 \le r \le m $，令

\[
\alpha = \sum_{(j)\in\mathbb{J}_r} a_{(j)}\varepsilon^{(j)} \quad \text{且} \quad \beta = \sum_{(j)\in\mathbb{J}_r} b_{(j)}\varepsilon^{(j)},
\]

根据命题 2.3，它们是 $ \alpha, \beta \in \bigwedge^{r}V^{*} $ 的合法基表示。于是置

\begin{equation}
(\alpha \mid \beta)_r := \sum_{(j)\in\mathbb{J}_r} a_{(j)}b_{(j)}.\tag{2.15}
\end{equation}

显然 $ (\cdot \mid \cdot)_r $ 是 $ \bigwedge^{r}V^{*} $（$ 0 \le r \le m $）上的标量积。由注 2.14(b) 与 (c) 有 $ (\cdot \mid \cdot)_1 = (\cdot \mid \cdot)_{\ast} $。

\begin{remark}

\begin{items}
\item

基 $ \{\varepsilon^{(j)}; (j) \in \mathbb{J}_r\} $ 是 $ (\bigwedge^{r}V^{*}, (\cdot \mid \cdot)_r) $ 的标准正交基，$ 1 \le r \le m $。

\item

对 $ \alpha^1,\ldots,\alpha^r, \beta^1,\ldots,\beta^r \in V^{*} $，有

\[
(\alpha^1 \wedge \cdots \wedge \alpha^r \mid \beta^1 \wedge \cdots \wedge \beta^r)_r = \det[(\alpha^{j} \mid \beta^{k})_{\ast}].
\]

\proofhead
设
\[
\alpha^{j} = \sum_{i=1}^{m} a^{j}_{i}\varepsilon^{i} \quad \text{且} \quad \beta^{k} = \sum_{i=1}^{m} b^{k}_{i}\varepsilon^{i} \quad \text{对 } 1 \le j, k \le r.
\]

又设

\[
\alpha^1 \wedge \cdots \wedge \alpha^r = \sum_{(j)\in\mathbb{J}_r} a_{(j)}\varepsilon^{(j)} \quad \text{且} \quad \beta^1 \wedge \cdots \wedge \beta^r = \sum_{(k)\in\mathbb{J}_r} b_{(k)}\varepsilon^{(k)}
\]

为基表示。则依例 2.4(d) 有

\[
a_{(j)} = \det\left(\left[a^{i}_{j_k}\right]_{1\le i,k\le r}\right) \quad \text{且} \quad b_{(k)} = \det\left(\left[b^{i}_{k_\ell}\right]_{1\le i,\ell\le r}\right)
\]

其中 $ (j) = (j_1,\ldots,j_r) \in \mathbb{J}_r $，$ (k) = (k_1,\ldots,k_r) \in \mathbb{J}_r $。
由 $ (\cdot \mid \cdot)_{\ast} $ 的双线性与对称性，以及行列式是其行向量的交错 $ r $ 形式这一事实可得（见命题 2.3(i) 的证明）

\begin{equation*}
\begin{aligned}
\det[(\alpha^{j} \mid \beta^{k})_{\ast}] &= \sum_{(j)\in\mathbb{J}_r}\sum_{\sigma\in\mathbb{S}_r} \operatorname{sign}(\sigma) a^{1}_{j_{\sigma(1)}} \cdots a^{r}_{j_{\sigma(r)}} \det\left(\left[(\varepsilon^{j_k} \mid \beta^{\ell})_{\ast}\right]_{1\le k,\ell\le r}\right) \\ &= \sum_{(j)\in\mathbb{J}_r} \det\left(\left[a^{i}_{j_k}\right]_{1\le i,k\le r}\right) \det\left(\left[(\varepsilon^{j_k} \mid \beta^{\ell})_{\ast}\right]_{1\le k,\ell\le r}\right) \\ &= \sum_{(j)\in\mathbb{J}_r} a_{(j)} \det\left(\left[(\varepsilon^{j_k} \mid \beta^{\ell})_{\ast}\right]_{1\le k,\ell\le r}\right) \\ &= \sum_{(j)\in\mathbb{J}_r}\sum_{(k)\in\mathbb{J}_r} a_{(j)}b_{(k)} \det\left(\left[(\varepsilon^{j_i} \mid \varepsilon^{k_{\ell}})_{\ast}\right]_{1\le i,\ell\le r}\right).
\end{aligned}
\end{equation*}

由 (a) 有

\begin{equation*}
\det[(\varepsilon^{j_i} \mid \varepsilon^{k_{\ell}})_{\ast}] = \det\left(\left[\delta^{j_i,k_{\ell}}\right]_{1\le i,\ell\le r}\right) = \begin{cases} 1 & \text{若 } (j) = (k), \\ 0 & \text{其他情形}. \end{cases}
\end{equation*}

于是得到

\[
\det[(\alpha^{j} \mid \beta^{k})_{\ast}] = \sum_{(j)\in\mathbb{J}_r} a_{(j)}b_{(j)}.
\]

由 (2.15)，证明完成。\nobreak\hspace{0.35em}■

\item

$ \bigwedge^{r}V^{*} $ 上的标量积 $ (\cdot \mid \cdot)_r $ 不依赖于标准正交基或其定向的特殊选取，而只依赖于 $ V $ 上的内积 $ (\cdot \mid \cdot) $。

\proofhead
这由 (b)、例 2.4(e) 以及标量积的双线性推出。\nobreak\hspace{0.35em}■
\end{items}

因为

\begin{equation}
\dim\left(\bigwedge^{r}V^{*}\right) = \binom{m}{r} = \binom{m}{m-r} = \dim\left(\bigwedge^{m-r}V^{*}\right),\tag{2.16}
\end{equation}

所以对 $ 0 \le r \le m $，$ \bigwedge^{r}V^{*} $ 与 $ \bigwedge^{m-r}V^{*} $ 是同构的向量空间。现在我们引入从 $ \bigwedge^{r}V^{*} $ 到 $ \bigwedge^{m-r}V^{*} $ 的一个特殊（自然）同构：霍奇星算子。

我们首先注意，对每个 $ \alpha \in \bigwedge^{r}V^{*} $，命题 2.5 表明

\begin{equation}
(\beta \mapsto \alpha \wedge \beta) \in \mathcal{L}\left(\bigwedge^{m-r}V^{*}, \bigwedge^{m}V^{*}\right).\tag{2.17}
\end{equation}

因为 $ \bigwedge^{m}V^{*} $ 是一维的，所以存在唯一的 $ f_{\alpha}(\beta) \in \mathbb{R} $ 使得

\[
\alpha \wedge \beta = f_{\alpha}(\beta)\omega_V \quad \text{对 } \beta \in \bigwedge^{m-r}V^{*}.
\]

由 (2.17)，$ f_{\alpha} $ 属于 $ \mathcal{L}(\bigwedge^{m-r}V^{*}, \mathbb{R}) $。于是根据里斯表示定理，存在唯一的 $ \ast\alpha \in \bigwedge^{m-r}V^{*} $，使得对 $ \beta \in \bigwedge^{m-r}V^{*} $ 有 $ f_{\alpha}(\beta) = (\ast\alpha \mid \beta)_{m-r} $。换言之，每个 $ \alpha \in \bigwedge^{r}V^{*} $ 都有唯一的元素 $ \ast\alpha \in \bigwedge^{m-r}V^{*} $ 使得

\begin{equation}
\alpha \wedge \beta = (\ast\alpha \mid \beta)_{m-r}\omega_V \quad \text{对 } \beta \in \bigwedge^{m-r}V^{*}.\tag{2.18}
\end{equation}

因此 $ \ast\alpha = \Theta^{-1}f_{\alpha} $，其中 $ \Theta $ 表示空间 $ \bigwedge^{m-r}V^{*} $ 的里斯同构 $ \Theta $。于是

\begin{equation}
(\alpha \mapsto \ast\alpha) \in \mathcal{L}\left(\bigwedge^{r}V^{*}, \bigwedge^{m-r}V^{*}\right).\tag{2.19}
\end{equation}

这个映射称为\textbf{霍奇星算子}（或简称霍奇星）。

\end{remark}

\begin{remark}

\begin{items}
\item

霍奇星是同构。

\proofhead
由 $ \ast\alpha = 0 $ 与 (2.18) 可知，对每个 $ \beta \in \bigwedge^{m-r}V^{*} $ 都有 $ \alpha \wedge \beta = 0 $。对特殊选取 $ \beta := \varepsilon^{r+1} \wedge \cdots \wedge \varepsilon^{m} $，由
\[
\alpha = \sum_{(j)\in\mathbb{J}_r} a_{(j)}\varepsilon^{(j)} \quad \text{其中 } (j_0) := (1,\ldots,r)
\]

可得 $ 0 = \alpha \wedge \beta = a_{(j_0)}\omega_V $，从而 $ a_{(j_0)} = 0 $。类似地，对 $ (j) \in \mathbb{J}_r $ 可得 $ a_{(j)} = 0 $。因此 (2.19) 是单射。现在由 (2.16) 推出结论。\nobreak\hspace{0.35em}■

\item

霍奇星依赖于 $ V $ 的标量积与定向。\nobreak\hspace{0.35em}■

\end{items}

\end{remark}

\begin{example}

\begin{items}
\item

对 $ 1 \le j \le m $，有 $ \ast\varepsilon^{j} = (-1)^{j-1}\varepsilon^{1} \wedge \cdots \wedge \widehat{\varepsilon^{j}} \wedge \cdots \wedge \varepsilon^{m} $。

\proofhead
由交错性、外积的结合性以及例 2.4(f) 可得
\[
\varepsilon^{j} \wedge (\varepsilon^{1} \wedge \cdots \wedge \widehat{\varepsilon^{k}} \wedge \cdots \wedge \varepsilon^{m}) = (-1)^{j-1}\delta^{jk}\varepsilon^{1} \wedge \cdots \wedge \varepsilon^{m} = (-1)^{j-1}\delta^{jk}\omega
\]

对 $ 1 \le k \le m $。现在由 (2.18) 以及如下事实推出结论：根据注 2.15(a)，

\[
\{\varepsilon^{1} \wedge \cdots \wedge \widehat{\varepsilon^{k}} \wedge \cdots \wedge \varepsilon^{m}; 1 \le k \le m\}
\]

是 $ \bigwedge^{m-1}V^{*} $ 的标准正交基。\nobreak\hspace{0.35em}■

\item

对 $ 1 \le j \le m $，有

\[
\ast(\varepsilon^{1} \wedge \cdots \wedge \widehat{\varepsilon^{j}} \wedge \cdots \wedge \varepsilon^{r}) = (-1)^{m-j}\varepsilon^{j}.
\]

\proofhead
因为
\[
(\varepsilon^{1} \wedge \cdots \wedge \widehat{\varepsilon^{j}} \wedge \cdots \wedge \varepsilon^{m}) \wedge \varepsilon^{k} = (-1)^{m-j}\delta^{jk}\omega,
\]

所以如上一证明一样推出结论。\nobreak\hspace{0.35em}■

\item

$ \ast 1 = \omega $ 且 $ \ast\omega = 1 $。

\item

现在考虑同时涵盖 (a) 与 (b) 的一般情形。于是设 $ 1 \le r \le m-1 $ 且 $ (j) \in \mathbb{J}_r $。则存在唯一的 $ (j^{c}) \in \mathbb{J}_{m-r} $，使得 $ (j) \vee (j^{c}) := (j_1,\ldots,j_r,j_1^{c},\ldots,j_{m-r}^{c}) $ 是 $ \{1,\ldots,m\} $ 的一个排列。置 $ s(j) := \operatorname{sign}((j) \vee (j^{c})) $，我们便有

\begin{equation}
\ast\varepsilon^{(j)} = s(j)\varepsilon^{(j^{c})}.\tag{2.20}
\end{equation}

于是对 $ \alpha = \sum_{(j)\in\mathbb{J}_r} a_{(j)}\varepsilon^{(j)} \in \bigwedge^{r}V^{*} $ 有

\[
\ast\alpha = \sum_{(j)\in\mathbb{J}_r} s(j)a_{(j)}\varepsilon^{(j^{c})}.
\]

\proofhead
对 $ (k) \in \mathbb{J}_{m-r} $，当 $ (k) \ne (j^{c}) $ 时有 $ \varepsilon^{(j)} \wedge \varepsilon^{(k)} = 0 $，因为该乘积中至少有一个 $ \varepsilon^{j_i} $ 出现两次。当 $ (k) = (j^{c}) $ 时，由 (2.3) 推出
\begin{equation}
\varepsilon^{(j)} \wedge \varepsilon^{(j^{c})} = s(j)\omega.\tag{2.21}
\end{equation}

现在由 (2.18) 与注 2.15(a) 推出 (2.20)。\nobreak\hspace{0.35em}■

\item

对 $ \alpha \in \bigwedge^{r}V^{*} $，$ 0 \le r \le m $，有 $ \ast\ast\alpha := \ast(\ast\alpha) = (-1)^{r(m-r)}\alpha $。

\proofhead
对 $ (j), (k) \in \mathbb{J}_r $，由 (2.18)、(d) 与命题 2.5(iii) 可得
\[
(\ast\ast\varepsilon^{(j)} \mid \varepsilon^{(k)})_r\omega = (\ast\varepsilon^{(j)}) \wedge \varepsilon^{(k)} = s(j)\varepsilon^{(j^{c})} \wedge \varepsilon^{(k)} = (-1)^{r(m-r)}s(j)\varepsilon^{(k)} \wedge \varepsilon^{(j^{c})}.
\]

于是因为当 $ (k) \ne (j) $ 时 $ \varepsilon^{(k)} \wedge \varepsilon^{(j^{c})} = 0 $，并利用 (2.21) 可得

\[
(\ast\ast\varepsilon^{(j)} \mid \beta)_r = (-1)^{r(m-r)}(\varepsilon^{(j)} \mid \beta)_r \quad \text{对 } \beta \in \bigwedge^{r}V^{*}.
\]

因此对 $ (j) \in \mathbb{J}_r $ 有 $ \ast\ast\varepsilon^{(j)} = (-1)^{r(m-r)}\varepsilon^{(j)} $，由命题 2.3(i)，这证明了结论。\nobreak\hspace{0.35em}■

\item

对 $ \alpha, \beta \in \bigwedge^{r}V^{*} $，有关系

\begin{equation}
\alpha \wedge \ast\beta = \beta \wedge \ast\alpha = (\alpha \mid \beta)_r\omega.\tag{2.22}
\end{equation}

\proofhead
设 $ \alpha := \beta := \varepsilon^{(j)} $。则由 (d) 与 (2.21) 有
\[
\alpha \wedge \ast\beta = \alpha \wedge \ast\alpha = \beta \wedge \ast\alpha = s(j)\varepsilon^{(j)} \wedge \varepsilon^{(j^{c})} = \omega = (\varepsilon^{(j)} \mid \varepsilon^{(j)})_r\omega = (\alpha \mid \beta)_r\omega.
\]

再取 $ \alpha := \varepsilon^{(j)} $，$ \beta := \varepsilon^{(k)} $，$ (j) \ne (k) $，由 (d) 得到

\[
\alpha \wedge \ast\beta = s(k)\varepsilon^{(j)} \wedge \varepsilon^{(k^{c})} = 0 = s(j)\varepsilon^{(k)} \wedge \varepsilon^{(j^{c})} = \beta \wedge \ast\alpha.
\]

此外，由注 2.15(a) 有 $ (\alpha \mid \beta)_r = 0 $。因此 (2.22) 在此情形也成立。现在这一结论也由命题 2.3(i) 推出。\nobreak\hspace{0.35em}■

\end{items}

\end{example}

\needspace{3\baselineskip}\subsection*{不定内积}

对若干应用，特别是在物理学中，必须放弃标量积正定这一假设。因此，下面我们简短地过一遍处理这种情形所需的修改。

设 $ \mathbf{b}: V \times V \to \mathbb{R} $ 是双线性形式。若对每个 $ y \in V \setminus \{0\} $ 都存在 $ x \in V $ 使得 $ \mathbf{b}(x,y) \neq 0 $，则称 $ \mathbf{b} $ 是\textbf{非退化}的。若

\[
\mathbf{b}(x,y) = \mathbf{b}(y,x) \qquad \text{对 } x, y \in V,
\]

则称它是\textbf{对称}的。

设 $ \mathbf{b}: V \times V \to \mathbb{R} $ 是 $ V := (V, (\cdot \mid \cdot)) $ 上的非退化对称双线性形式。在下面的注中，我们列出 $ \mathbf{b} $ 的一些基本性质。

\begin{remark}

\begin{items}
\item

存在 $ V $ 的 $ \mathbf{b} $-规范正交基（$ \mathbf{b} $-ONB），即存在 $ V $ 的基 $ \{b_1, \ldots, b_m\} $ 使得 $ \mathbf{b}(b_j, b_k) = \pm \delta_{jk} $ 对 $ 1 \le j, k \le m $。若 $ r $ 是正号的个数，$ s $ 是负号的个数，则 $ r + s = m $。数 $ t := r - s $ 称为 $ \mathbf{b} $ 的\textbf{符号差}。符号差以及 $ r $ 和 $ s $ 都与 $ \mathbf{b} $-ONB 的选取无关。特别地，\footnote{不应当把 $ \operatorname{sign}(\mathbf{b}) $ 与符号差 $ t $ 混淆。显然 $ 2\operatorname{sign}(\mathbf{b}) = m - t $。}

\[
(-1)^s = \operatorname{sign}(\mathbf{b}) := \operatorname{sign}(\det[\mathbf{b}(b_j, b_k)]) .
\]

\proofhead
定理 VII.4.2(iii) 显然蕴含 $ \mathbf{b} $ 连续。于是 $ \mathbf{b}(x, \cdot): V \to \mathbb{R} $ 是 $ V $ 上的连续线性形式。因此，里斯表示定理（定理 VII.2.14）保证存在唯一的 $ \mathfrak{B}x \in V $ 使得
\[
\mathbf{b}(x,y) = (\mathfrak{B}x \mid y) \qquad \text{对 } y \in V.
\]

由 $ \mathbf{b}(\cdot, y) $ 的线性性可知 $ x \mapsto \mathfrak{B}x $ 是线性的。于是由定理 VII.1.6，$ \mathfrak{B} $ 属于 $ \mathcal{L}(V) $，且

\[
\mathbf{b}(x,y) = (\mathfrak{B}x \mid y) \qquad \text{对 } x, y \in V.
\]

映射 $ \mathfrak{B} $ 称为 $ \mathbf{b} $ 关于 $ (\cdot \mid \cdot) $ 的\textbf{表示算子}。因为 $ \mathbf{b} $ 非退化，$ \mathfrak{B} $ 是 $ V $ 的自同构（反之亦然）；又因为 $ \mathbf{b} $ 对称，$ \mathfrak{B} $ 也对称。
注 2.1(a) 允许我们把 $ V $ 与 $ \mathbb{R}^m $ 等同起来。因此主轴变换定理\footnote{见例 VII.10.17(b)。}保证存在 $ V $ 的 ONB $ \{v_1, \ldots, v_m\} $ 以及 $ \mathfrak{B} $ 的特征值 $ \lambda_1 \ge \cdots \ge \lambda_m $，使得

\begin{equation}
\mathfrak{B}v_j = \lambda_j v_j \qquad \text{对 } 1 \le j \le m.\tag{2.23}
\end{equation}

因为 $ \mathbf{b} $ 非退化，我们有 $ \lambda_j \neq 0 $ 对 $ 1 \le j \le m $。我们令 $ b_j := v_j / \sqrt{|\lambda_j|} $。则 $ \{b_1, \ldots, b_m\} $ 是 $ V $ 的基，且由 (2.23) 推出

\[
\mathbf{b}(b_j, b_k) = (\mathfrak{B}b_j \mid b_k) = (\mathfrak{B}v_j \mid v_k)/\sqrt{|\lambda_j \lambda_k|} = \lambda_j (v_j \mid v_k)/\sqrt{|\lambda_j \lambda_k|} = \operatorname{sign}(\lambda_j)\delta_{jk}
\]

对 $ 1 \le j, k \le m $。
要证明 $ t = r - s $ 与 $ \mathbf{b} $-ONB 的选取无关，由于 $ r + s = m $，只需证明 $ r $ 具有同样的性质。因此设 $ \{c_1, \ldots, c_m\} $ 是 $ V $ 的 $ \mathbf{b} $-ONB，使得 $ \mathbf{b}(c_j, c_j) = 1 $ 对 $ 1 \le j \le \rho $，且 $ \mathbf{b}(c_j, c_j) = -1 $ 对 $ \rho + 1 \le j \le m $。我们要证明向量 $ b_1, \ldots, b_r, c_{\rho+1}, \ldots, c_m $ 线性无关，因为这将蕴含 $ r + (m - \rho) \le m $，从而 $ r \le \rho $；此外，交换两个 $ \mathbf{b} $-ONB 会类似地得到 $ \rho \le r $，这便确定了 $ r $。
因此设

\[
\beta_1 b_1 + \cdots + \beta_r b_r = \gamma_{\rho+1} c_{\rho+1} + \cdots + \gamma_m c_m
\]

其中 $ \beta_1, \ldots, \beta_r, \gamma_{\rho+1}, \ldots, \gamma_m $ 为实数。集合 $ \{b_1, \ldots, b_r, c_{\rho+1}, \ldots, c_m\} $ 的每个线性相关关系都可写成这种形式。于是对 $ v := \beta_1 b_1 + \cdots + \beta_r b_r $，我们有

\[
\mathbf{b}(v,v) = \sum_{j=1}^{r} \beta_r^2 = - \sum_{j=\rho+1}^{m} \gamma_j^2 ,
\]

这蕴含 $ \beta_1 = \cdots = \beta_r = \gamma_{\rho+1} = \cdots = \gamma_m = 0 $。最后一个断言现在清楚了。 \nobreak\hspace{0.35em}■

\item

（里斯表示定理）对每个 $ v^* \in V^* $，恰有一个 $ v \in V $ 使得 $ \mathbf{b}(v,w) = \langle v^*, w \rangle $ 对 $ w \in W $。映射

\[
\Theta_{\mathbf{b}}: V \to V^*, \qquad v \mapsto \mathbf{b}(v, \cdot)
\]

是向量空间同构，即关于 $ \mathbf{b} $ 的\textbf{里斯同构}。注 2.14(a) 的结论在此情形也成立。

\proofhead
利用 $ \mathbf{b} $ 的表示算子 $ \mathfrak{B} $ 与 $ V $ 的里斯同构 $ \Theta $，定理 VII.2.14 蕴含
\[
\mathbf{b}(v,w) = (\mathfrak{B}v \mid w) = \langle \Theta\mathfrak{B}v, w \rangle \qquad \text{对 } v, w \in V.
\]

令 $ \Theta_{\mathbf{b}} := \Theta\mathfrak{B} $ 即得断言。 \nobreak\hspace{0.35em}■

\item

对 $ V $ 的每个基 $ \{v_1, \ldots, v_m\} $，关于 $ \mathbf{b} $ 的\textbf{格拉姆行列式}，即

\[
G_{\mathbf{b}} := \det([\mathbf{b}(v_j, v_k)]) ,
\]

非零。

\proofhead
行列式 $ G_{\mathbf{b}} $ 为零当且仅当线性方程组
\begin{equation}
\sum_{k=1}^{m} \mathbf{b}(v_j, v_k)\xi^k = 0 \qquad \text{对 } 1 \le j \le m ,\tag{2.24}
\end{equation}

有非平凡解。若 $ v := \sum_{k=1}^{m} \xi^k v_k $，则 (2.24) 等价于 $ \mathbf{b}(v_j, v) = 0 $ 对 $ 1 \le j \le m $。因为 $ \{v_1, \ldots, v_m\} $ 是 $ V $ 的基且 $ \mathbf{b} $ 非退化，推出 $ v = 0 $，证毕。 \nobreak\hspace{0.35em}■

\item

设 $ (b_1, \ldots, b_m) $ 是 $ V $ 的正定向基，$ (\beta^1, \ldots, \beta^m) $ 是其对偶基。又设 $ \mathfrak{B} $ 是酉的。则

\[
\varepsilon^1 \wedge \cdots \wedge \varepsilon^m = \sqrt{|G_{\mathbf{b}}|}\, \beta^1 \wedge \cdots \wedge \beta^m .
\]

\proofhead
命题 2.13 的证明的第一部分表明
\[
\varepsilon^1 \wedge \cdots \wedge \varepsilon^m = \det(B)\beta^1 \wedge \cdots \wedge \beta^m ,
\]

其中 $ B \in \mathcal{L}(V) $ 是从 $ (e_1, \ldots, e_m) $ 到 $ (b_1, \ldots, b_m) $ 的基变换。以 $ \mathbf{b} $ 的表示 $ \mathfrak{B} $，我们得到

\[
\mathbf{b}(b_j, b_k) = (\mathfrak{B}b_j \mid b_k) = (\mathfrak{B}Be_j \mid Be_k) = (B^*\mathfrak{B}Be_j \mid e_k) \qquad \text{对 } 1 \le j, k \le m .
\]

如同命题 2.13 的证明，这蕴含

\[
G_{\mathbf{b}} = \det[\mathbf{b}(b_j, b_k)] = \det(B^*\mathfrak{B}B) = \det(\mathfrak{B})(\det(B))^2 .
\]

因为 $ |\det(\mathfrak{B})| = 1 $，推出 $ (\det(B))^2 = |G_{\mathbf{b}}| $，这蕴含断言。 \nobreak\hspace{0.35em}■

\end{items}

现设 $ \mathcal{O}r $ 是 $ V $ 的一个定向，$ \mathbf{b} $ 是 $ V $ 上的非退化对称双线性形式。又设 $ (e_1, \ldots, e_m) $ 是 $ V $ 的正定向 $ \mathbf{b} $-ONB，其对偶基为 $ (\varepsilon^1, \ldots, \varepsilon^m) $。在 $ V^* $ 上，我们通过

\[
\mathbf{b}_*(v^*, w^*) := \mathbf{b}(\Theta_{\mathbf{b}}^{-1}v^*, \Theta_{\mathbf{b}}^{-1}w^*) \qquad \text{对 } v^*, w^* \in V^* .
\]

这就定义了非退化对称双线性形式 $ \mathbf{b}_* $。对 $ \alpha^1, \ldots, \alpha^r, \beta^1, \ldots, \beta^r \in V^* $，我们令

\[
\mathbf{b}_r(\alpha^1 \wedge \cdots \wedge \alpha^r, \beta^1 \wedge \cdots \wedge \beta^r) := \det[\mathbf{b}_*(\alpha^j, \beta^k)] ,
\]

从而 $ \mathbf{b}_1 = \mathbf{b}_* $；我们还借助命题 2.3(i) 的基表示，通过双线性延拓定义

\[
\mathbf{b}_r: \bigwedge^r V^* \times \bigwedge^r V^* \to \mathbb{R} \qquad \text{对 } r \ge 1 .
\]

如同 (2.16)–(2.19)（取 $ \mathbf{b}_0 := (\cdot \mid \cdot)_0 $），推出存在线性映射

\[
\bigwedge^r V^* \to \bigwedge^{m-r} V^*, \qquad \alpha \mapsto \ast\alpha
\]

（$ 0 \le r \le m $），称为\textbf{霍奇星算子}。它由下式刻画：

\begin{equation}
\alpha \wedge \beta = \mathbf{b}_{m-r}(\ast\alpha, \beta)\varepsilon^1 \wedge \cdots \wedge \varepsilon^m \qquad \text{对 } \beta \in \bigwedge^{m-r} V^* .\tag{2.25}
\end{equation}
\end{remark}

\begin{remark}

\begin{items}
\item

霍奇星是一个同构，它只依赖于双线性形式 $ \mathbf{b} $ 与定向，而不依赖于 $ \mathbf{b} $-ONB。

\item

对 $ 1 \le r \le m $，$ \{\varepsilon^{(j)} ;\ (j) \in \mathbb{J}_r\} $ 是 $ \mathbf{b}_r $-ONB；且对 $ \omega := \varepsilon^1 \wedge \cdots \wedge \varepsilon^m $，有 $ \mathbf{b}_m(\omega, \omega) = \operatorname{sign}(\mathbf{b}) $。

\proofhead
第一个断言由 $ \mathbf{b}_r $ 的定义立得。因为
\[
\mathbf{b}_m(\omega,\omega) = \det(\operatorname{diag}[\mathbf{b}_*(\varepsilon^1,\varepsilon^1),\ldots,\mathbf{b}_*(\varepsilon^m,\varepsilon^m)]) = \det(\operatorname{diag}[\mathbf{b}(e_1,e_1),\ldots,\mathbf{b}(e_m,e_m)]) ,
\]

第二个断言也成立。 \nobreak\hspace{0.35em}■

\item

我们有 $ \ast 1 = \operatorname{sign}(\mathbf{b})\omega $ 且 $ \ast\omega = 1 $。此外

\[
\ast\varepsilon^{(j)} = s(j)\mathbf{b}_{m-r}(\varepsilon^{(j^c)}, \varepsilon^{(j^c)})\varepsilon^{(j^c)} \qquad \text{对 } (j) \in \mathbb{J}_r ,
\]

其中 $ 1 \le r \le m - 1 $。

\proofhead
首先 $ \omega = 1 \wedge \omega = \mathbf{b}_m(\ast 1, \omega)\omega $ 蕴含 $ \mathbf{b}_m(\ast 1, \omega) = 1 $。其次 $ \dim(\bigwedge^m V^*) = 1 $ 给出 $ \ast 1 = a\omega $，其中 $ a \in \mathbb{R} $。由此并利用 (b) 得到
\[
1 = \mathbf{b}_m(\ast 1, \omega) = a\mathbf{b}_m(\omega, \omega) = a\operatorname{sign}(\mathbf{b}) ,
\]

故 $ a = \operatorname{sign}(\mathbf{b}) $。这就证明了第一个断言。类似地，我们得到 $ \ast\omega = 1 $。设 $ 1 \le r \le m - 1 $ 且 $ (j) \in \mathbb{J}_r $。则

\[
\omega = s(j)\varepsilon^{(j)} \wedge \varepsilon^{(j^c)} = s(j)\mathbf{b}_{m-r}(\ast\varepsilon^{(j)}, \varepsilon^{(j^c)})\omega ,
\]

故 $ \mathbf{b}_{m-r}(\ast\varepsilon^{(j)}, \varepsilon^{(j^c)}) = s(j) $。注意 $ \{\varepsilon^{(k)} ;\ (k) \in \mathbb{J}_{m-r}\} $ 是 $ \bigwedge^{m-r} V^* $ 的 $ \mathbf{b}_{m-r} $-ONB。又 $ \ast\varepsilon^{(j)} \in \bigwedge^{m-r} V^* $，且 $ \mathbf{b}_{m-r}(\ast\varepsilon^{(j)}, \varepsilon^{(k)}) = 0 $ 对 $ (k) \neq (j^c) $。由此推出 $ \ast\varepsilon^{(j)} = a\varepsilon^{(j^c)} $，其中 $ a \in \mathbb{R} $，从而

\[
a\mathbf{b}_{m-r}(\varepsilon^{(j^c)}, \varepsilon^{(j^c)}) = \mathbf{b}_{m-r}(\ast\varepsilon^{(j)}, \varepsilon^{(j^c)}) = s(j) .
\]

这蕴含 $ a = s(j)\mathbf{b}_{m-r}(\varepsilon^{(j^c)}, \varepsilon^{(j^c)}) $。现在最后一个断言清楚了。 \nobreak\hspace{0.35em}■

\item

对 $ \alpha \in \bigwedge^r V^* $ 且 $ 0 \le r \le m $，我们有 $ \ast\ast\alpha = \operatorname{sign}(\mathbf{b})(-1)^{r(m-r)}\alpha $。

\proofhead
与 (c) 的证明类似，由
\[
\mathbf{b}_r(\ast\varepsilon^{(j^c)}, \varepsilon^{(j)})\omega = \varepsilon^{(j^c)} \wedge \varepsilon^{(j)} = (-1)^{r(m-r)}\varepsilon^{(j)} \wedge \varepsilon^{(j^c)} = s(j)(-1)^{r(m-r)}\omega ,
\]

得到 $ \ast\varepsilon^{(j^c)} = s(j)(-1)^{r(m-r)}\mathbf{b}_r(\varepsilon^{(j)}, \varepsilon^{(j)})\varepsilon^{(j)} $。因此由 (c) 得到

\[
\ast(\ast\varepsilon^{(j)}) = \ast\bigl(s(j)\mathbf{b}_{m-r}(\varepsilon^{(j^c)}, \varepsilon^{(j^c)})\varepsilon^{(j^c)}\bigr) = s(j)^2(-1)^{r(m-r)}\mathbf{b}_r(\varepsilon^{(j)}, \varepsilon^{(j)})\mathbf{b}_{m-r}(\varepsilon^{(j^c)}, \varepsilon^{(j^c)})\varepsilon^{(j)} ,
\]

由此得到断言。 \nobreak\hspace{0.35em}■

\item

对 $ \alpha, \beta \in \bigwedge^r V^* $，我们有

\[
\alpha \wedge \ast\beta = \beta \wedge \ast\alpha = \operatorname{sign}(\mathbf{b})\, \mathbf{b}_r(\alpha, \beta)\omega .
\]

\proofhead
对例 2.17(f) 的证明作显然的修改即得。 \nobreak\hspace{0.35em}■
\end{items}

这些思想的一个重要应用是\textbf{闵可夫斯基空间} $ \mathbb{R}^{4}_{1,3} := (\mathbb{R}^4, (\cdot \mid \cdot)_{1,3}) $，即狭义相对论的“时空”，配有\textbf{闵可夫斯基度量}

\[
(x \mid y)_{1,3} := x_0 y_0 - x_1 y_1 - x_2 y_2 - x_3 y_3 .
\]

（在相对论中，“第 0 个坐标”是时间。）我们将在后面详细讨论这一点。

非退化对称双线性形式 $ \mathbf{b} $ 也称为\textbf{不定内积}；相应地，$ (V, \mathbf{b}) $ 称为\textbf{不定内积空间}。

\end{remark}

\needspace{3\baselineskip}\subsection*{张量}

为完整起见，下面我们简要介绍一般张量的概念，后面几节中我们会遇到它。设 $ r, s \in \mathbb{N} $。一个 $ (r + s) $ 线性映射

\[
\gamma: \underbrace{V^* \times \cdots \times V^*}_{r} \times \underbrace{V \times \cdots \times V}_{s} \to \mathbb{R}
\]

称为 $ V $ 上型为 $ (r, s) $ 的\textbf{张量}，或 $ (r, s) $-张量。特别地，$ \gamma $ 是 $ r $ 阶\textbf{逆变}、$ s $ 阶\textbf{协变}的（或称 $ r $ 逆变、$ s $ 协变）。我们用 $ T^r_s(V) $ 表示 $ V $ 上全体 $ (r, s) $-张量构成的赋范\footnote{见定理 VII.4.2。}向量空间。

对 $ \gamma_1 \in T^{r_1}_{s_1}(V) $ 与 $ \gamma_2 \in T^{r_2}_{s_2}(V) $，\textbf{张量积} $ \gamma_1 \otimes \gamma_2 $ 由下式定义：

\[
\gamma_1 \otimes \gamma_2(\alpha^1, \ldots, \alpha^{r_1}, \beta^1, \ldots, \beta^{r_2}, v_1, \ldots, v_{s_1}, w_1, \ldots, w_{s_2}) := \gamma_1(\alpha^1, \ldots, \alpha^{r_1}, v_1, \ldots, v_{s_1})\gamma_2(\beta^1, \ldots, \beta^{r_2}, w_1, \ldots, w_{s_2}) ,
\]

其中 $ \alpha^1, \ldots, \alpha^{r_1}, \beta^1, \ldots, \beta^{r_2} \in V^* $ 且 $ v_1, \ldots, v_{s_1}, w_1, \ldots, w_{s_2} \in V $。

下面如常，我们利用注 2.8 的典范同构 $ \kappa $ 把 $ V^{\ast\ast} $ 与 $ V $ 等同起来。

\begin{remark}

\begin{items}
\item

$ T^1_0(V) = V $，$ T^0_1(V) = V^* $，且 $ T^0_2(V) = \mathcal{L}^2(V, \mathbb{R}) $。

\item

对 $ \gamma \in T^1_1(V) $，恰有一个 $ C \in \mathcal{L}(V) $ 使得

\begin{equation}
\gamma(v^*, v) = \langle v^*, Cv \rangle \qquad \text{对 } v \in V,\ v^* \in V^* .\tag{2.26}
\end{equation}

映射

\[
T^1_1(V) \to \mathcal{L}(V), \qquad \gamma \mapsto C
\]

是等距同构。

\proofhead
对 $ v \in V $，映射 $ \gamma(\cdot, v) $ 属于 $ V^{\ast\ast} = V $。因为 $ \gamma $ 是双线性的，我们有
\[
C := (v \mapsto \gamma(\cdot, v)) \in \mathcal{L}(V) ,
\]

且 $ \langle v^*, Cv \rangle = \gamma(v^*, v) $ 对 $ (v, v^*) \in V \times V^* $。反之，每个 $ C \in \mathcal{L}(V) $ 都借助 (2.26) 定义一个 $ \gamma \in T^1_1(V) $。最后一个断言现在清楚了。 \nobreak\hspace{0.35em}■

\item

张量积是双线性的且结合的。

\item

令 $ m := \dim(V) $，有 $ \dim(T^r_s(V)) = m^{r+s} $。若 $ (e_1, \ldots, e_m) $ 是 $ V $ 的基，$ (\varepsilon^1, \ldots, \varepsilon^m) $ 是其对偶基，则

\[
\{ e_{j_1} \otimes \cdots \otimes e_{j_r} \otimes \varepsilon^{k_1} \otimes \cdots \otimes \varepsilon^{k_s} ;\ j_i, k_i \in \{1, \ldots, m\} \}
\]

是 $ T^r_s(V) $ 的基。

\proofhead
我们将这个简单的证明留给你。\nobreak\hspace{0.35em}■
\item

$ \bigwedge^{r}V^{*} $ 是 $ T_{r}^{0}(V) $ 的向量子空间。

\item

对偶配对 $ \langle\cdot,\cdot\rangle: V^{*}\times V \to \mathbb{R} $ 是 $ V $ 上的 $ (1,1) $ 张量。\nobreak\hspace{0.35em}■

\end{items}

\end{remark}

\exerciseshead

\begin{exercises}
\item

对 $ T \in \mathcal{L}^{r}(V,\mathbb{R}) $，\textbf{交错化算子} $ \operatorname{Alt}(T) $ 由下式定义：

\[
\operatorname{Alt}(T)(v_{1},\ldots,v_{r}) := \frac{1}{r!}\sum_{\sigma\in S_{r}}\operatorname{sign}(\sigma)T(v_{\sigma(1)},\ldots,v_{\sigma(r)})
\]

其中 $ v_{1},\ldots,v_{r} \in V $。证明：

\begin{items}
\item

$ \operatorname{Alt} \in \mathcal{L}(\mathcal{L}^{r}(V,\mathbb{R}),\bigwedge^{r}V^{*}) $；

\item

$ \operatorname{Alt}^{2} = \operatorname{Alt} $。

\end{items}

\item

对 $ S \in \mathcal{L}^{s}(V,\mathbb{R}) $ 与 $ T \in \mathcal{L}^{t}(V,\mathbb{R}) $，定义 $ S\otimes T \in \mathcal{L}^{s+t}(V,\mathbb{R}) $ 为

\[
S\otimes T(v_{1},\ldots,v_{s},v_{s+1},\ldots,v_{s+t}) := S(v_{1},\ldots,v_{s})T(v_{s+1},\ldots,v_{s+t}),
\]

其中 $ v_{1},\ldots,v_{s+t} \in V $。证明：对 $ \alpha \in \bigwedge^{r}V^{*} $ 与 $ \beta \in \bigwedge^{s}V^{*} $，

\[
\alpha\wedge\beta = \frac{(r+s)!}{r!\,s!}\operatorname{Alt}(\alpha\otimes\beta).
\]

在习题 3–8 中，设 $ (V,(\cdot\mid\cdot),\mathcal{O}r) $ 是定向内积空间，$ \omega $ 是它的体积元，$ \Theta $ 是里斯同构。

\item

设 $ \dim(V) = 3 $。则 $ V $ 上的\textbf{向量积或叉积} $ \times $ 由下式定义\footnote{见注 VIII.2.14。}：

\[
\times: V\times V \to V, \quad (v,w)\mapsto v\times w := \Theta^{-1}\omega(v,w,\cdot).
\]

证明以下各项：

\begin{items}
\item

对 $ u,v,w \in V $，有 $ (v\times w\mid u) = \omega(v,w,u) $。

\item

向量积是双线性的且交错的。

\item

向量 $ v\times w $ 不为零当且仅当 $ v $ 与 $ w $ 线性无关。

\item

若 $ v $ 与 $ w $ 线性无关，则 $ (v,w,v\times w) $ 是 $ V $ 的正定向基。

\item

向量 $ v\times w $ 与 $ v $ 和 $ w $ 都正交。

\item

对 $ v,w \in V\setminus\{0\} $，有

\[
|v\times w| = \sqrt{|v|^{2}|w|^{2} - (v\mid w)^{2}} = |v|\,|w|\sin\varphi,
\]

其中 $ \varphi \in [0,\pi] $ 是向量 $ v $ 与 $ w $ 之间的（非定向）夹角。

\item

设 $ (e_{1},e_{2},e_{3}) $ 是 $ V $ 的正定向标准正交基。则对 $ v = \sum_{j}\xi^{j}e_{j} $ 与 $ w = \sum_{j}\eta^{j}e_{j} $，有

\[
v\times w = (\xi^{2}\eta^{3} - \xi^{3}\eta^{2})e_{1} + (\xi^{3}\eta^{1} - \xi^{1}\eta^{3})e_{2} + (\xi^{1}\eta^{2} - \xi^{2}\eta^{1})e_{3}.
\]
\item

（格拉斯曼恒等式）$ v_{1}\times(v_{2}\times v_{3}) = (v_{1}\mid v_{3})v_{2} - (v_{1}\mid v_{2})v_{3} $。

\item

向量积不是结合的。

\item

$ (v_{1}\times v_{2})\times(v_{3}\times v_{4}) = \omega(v_{1},v_{2},v_{4})v_{3} - \omega(v_{1},v_{2},v_{3})v_{4} $。

\item

（雅可比恒等式）$ v_{1}\times(v_{2}\times v_{3}) + v_{2}\times(v_{3}\times v_{1}) + v_{3}\times(v_{1}\times v_{2}) = 0 $。
（提示：(f) 回忆命题 2.13 与 (a)。(h) 向量积由它在基 $ (e_{1},e_{2},e_{3}) $ 上的值确定。）

\end{items}

\item

对 $ 0 \le r \le m $，验证以下公式：

\begin{items}
\item

对 $ \alpha \in \bigwedge^{r}V^{*} $ 与 $ \beta \in \bigwedge^{m-r}V^{*} $，有 $ (\ast\alpha\mid\beta)_{m-r} = (-1)^{r(m-r)}(\alpha\mid\ast\beta)_{r} $。

\item

对 $ \alpha,\beta \in \bigwedge^{r}V^{*} $，有 $ (\ast\alpha)\wedge\beta = (\ast\beta)\wedge\alpha $。

\item

对 $ v,w \in V $，有 $ \ast(\Theta v\wedge\ast\Theta w) = (v\mid w) $。

\end{items}

\item

设 $ (b_{1},\ldots,b_{m}) $ 是 $ V $ 的正定向基，其对偶基为 $ (\beta^{1},\ldots,\beta^{m}) $。证明：

\begin{items}
\item

对 $ 1 \le j,k \le m $，有 $ \beta^{j}\wedge\ast\beta^{k} = g^{jk}\sqrt{G}\,\beta^{1}\wedge\cdots\wedge\beta^{m} $。

\item

对 $ 1 \le j \le m $，有 $ \ast\beta^{j} = \sum_{k=1}^{m}(-1)^{k-1}g^{jk}\sqrt{G}\,\beta^{1}\wedge\cdots\wedge\widehat{\beta^{k}}\wedge\cdots\wedge\beta^{m} $。若 $ V $ 是三维的，则证明

\[
\ast(\beta^{j}\wedge\beta^{k}) = \frac{1}{\sqrt{G}}\operatorname{sign}(j,k,\ell)\sum_{i=1}^{3}g_{\ell i}\beta^{i} = \frac{1}{G}\operatorname{sign}(j,k,\ell)\Theta b_{\ell}
\]

对 $ (j,k,\ell) \in \mathbb{S}_{3} $。

\end{items}

\item

在 $ \dim(V) = 3 $ 的情形，证明对 $ v,w \in V $ 有 $ v\times w = \Theta^{-1}(\ast(\Theta v\wedge\Theta w)) $。

\item

设 $ (b_{1},b_{2},b_{3}) $ 是 $ V $ 的正定向基，其对偶基为 $ (\beta^{1},\beta^{2},\beta^{3}) $。证明

\[
b_{j}\times b_{k} = \sqrt{G}\operatorname{sign}(j,k,\ell)\sum_{i=1}^{3}g^{\ell i}b_{i} = \sqrt{G}\operatorname{sign}(j,k,\ell)\Theta^{-1}\beta^{\ell}
\]

对 $ (j,k,\ell) \in \mathbb{S}_{3} $。

\item

设 $ (W,(\cdot\mid\cdot)_{W},\mathcal{O}r(W)) $ 是定向内积空间，$ A \in \mathcal{L}(V,W) $ 是保定向的等距（即 $ A^{*}\omega_{W} = \omega_{V} $）。证明下面的图对 $ 0 \le r \le m $ 交换：

\bookfigure[0.40]{figures/f84f8b993464fae9.png}

\item

对不定内积空间叙述并证明习题 4 与习题 5 的断言。

\item

对 $ k \in \mathbb{N} $，设 $ \mathbb{K}_{k}[X] $ 是 $ \mathbb{K} $ 上次数 $ \le k $ 的多项式的向量空间。证明 $ \mathbb{K}[X] = \bigoplus_{k\ge 0}\mathbb{K}_{k}[X] $ 关于多项式的通常乘法，即关于 §I.8 的卷积，是分次交换代数。

\item

设 $ (\varepsilon^{0},\varepsilon^{1},\varepsilon^{2},\varepsilon^{3}) $ 是 $ \mathbb{R}^{4} $ 的标准基的对偶基。对 $ c,E_{j},H_{j} \in \mathbb{R} $，置

\begin{equation*}
\begin{aligned}
\alpha &:= (E_{1}\varepsilon^{1} + E_{2}\varepsilon^{2} + E_{3}\varepsilon^{3})\wedge c\varepsilon^{0} + (H_{1}\varepsilon^{2}\wedge\varepsilon^{3} + H_{2}\varepsilon^{3}\wedge\varepsilon^{1} + H_{3}\varepsilon^{1}\wedge\varepsilon^{2}), \\ \beta &:= -(H_{1}\varepsilon^{1} + H_{2}\varepsilon^{2} + H_{3}\varepsilon^{3})\wedge c\varepsilon^{0} + (E_{1}\varepsilon^{2}\wedge\varepsilon^{3} + E_{2}\varepsilon^{3}\wedge\varepsilon^{1} + E_{3}\varepsilon^{1}\wedge\varepsilon^{2})
\end{aligned}
\end{equation*}

并计算关于 $ (\cdot\mid\cdot)_{1,3} $ 的 $ \ast\alpha $ 与 $ \ast\beta $。

\end{exercises}

\section{微分形式的局部理论}

在第 VIII.3 节中，我们对 1 次微分形式即普法夫形式已有诸多了解，并发展了一套构成线积分理论基础的演算。现在我们把这一想法推广到更多维的情形。作为第一步（本节即专门讨论这一步），我们引入欧氏空间开集上任意次数的微分形式，并给出这种“局部”情形下微分形式的演算。在随后的各节中，我们再考虑一般情形，即流形上的微分形式。

$ \mathbb{R}^{m} $ 的开子集 $ X $ 上的 $ r $ 次微分形式，不过是这样一个集合：对每个 $ x \in X $，它由切空间 $ T_{x}X $ 上的一个交错 $ r $ 形式组成。正因如此，本节第一部分实际上只是对第 2 节所给线性代数结果的重新表述。我们不再陈述新定理，而是用注和例来解释这些定义。当我们引入微分形式上的一种运算——外微分——时，分析才开始发挥作用。外微分用到分析中的概念，超出了线性代数的范围。

在本节中，$ X $ 是 $ \mathbb{R}^{m} $ 中的开集，且 $ \mathbb{K} = \mathbb{R} $。

\needspace{3\baselineskip}\subsection*{定义与基表示}

对 $ x \in X $，余切空间 $ T_{x}^{*}X = \{x\} \times (\mathbb{R}^{m})^{*} $ 是与切空间 $ T_{x}X = \{x\} \times \mathbb{R}^{m} $ 对偶的空间。因此，外积

\begin{equation}
\bigwedge^{r} T_{x}^{*}X = \{x\} \times \bigwedge^{r}(\mathbb{R}^{m})^{*} \qquad \text{对 } r \in \mathbb{N}\tag{3.1}
\end{equation}

与格拉斯曼代数

\[
\bigwedge T_{x}^{*}X = \{x\} \times \bigwedge(\mathbb{R}^{m})^{*}
\]

在 $ T_{x}^{*}X $ 上都有良好定义。我们可以通过如下定义 $ X $ 上的\textbf{交错 $ r $ 形式丛}来推广切丛与余切丛：

\[
\bigwedge^{r} T^{*}X := \bigcup_{x \in X} \bigwedge^{r} T_{x}^{*}X = X \times \bigwedge^{r}(\mathbb{R}^{m})^{*}
\]

并通过如下定义 $ X $ 的\textbf{格拉斯曼丛}：

\[
\bigwedge T^{*}X := \bigcup_{x \in X} \bigwedge T_{x}^{*}X = X \times \bigwedge(\mathbb{R}^{m})^{*}.
\]

一个映射

\[
\boldsymbol{\alpha}: X \to \bigwedge^{r} T^{*}X \qquad \text{满足 } \boldsymbol{\alpha}(x) \in \bigwedge^{r} T_{x}^{*}X \text{ 且 } x \in X,
\]

即格拉斯曼丛的一个\textbf{截面}\footnote{我们使用了“向量丛”理论的语言。此处不作详细展开（但可参见例如 [Con93]、[Dar94] 或 [HR72]），尽管它会带来各种想法的统一。}，称为 $ X $ 上的 \textbf{$ r $ 次微分形式}（简称 \textbf{$ r $ 形式}）。由 (3.1)，$ X $ 上每个 $ r $ 形式都有唯一表示

\[
\boldsymbol{\alpha}(x) = (x, \alpha(x)) \qquad \text{对 } x \in X
\]

其 \textbf{$ r $ 余向量部分}（简称余向量部分）为

\[
\alpha: X \to \bigwedge^{r}(\mathbb{R}^{m})^{*}.
\]

设 $ k \in \mathbb{N} \cup \{\infty\} $。$ r $ 形式 $ \boldsymbol{\alpha} $ 属于 \textbf{$ C^{k} $ 类}（或称 \textbf{$ k $ 次连续可微}\footnote{自然，我们把 $ C^{0} $ 类的 $ r $ 形式称为连续的。}，当 $ k = \infty $ 时称\textbf{光滑}），如果其余向量部分具有这一性质，即如果

\begin{equation}
\alpha \in C^{k}\big(X, \bigwedge^{r}(\mathbb{R}^{m})^{*}\big).\tag{3.2}
\end{equation}

这个定义是有意义的，因为根据注 2.2(a)，$ \bigwedge^{r}(\mathbb{R}^{m})^{*} $ 是 $ \mathcal{L}^{r}(\mathbb{R}^{m},\mathbb{R}) $ 的（闭）向量子空间。

为简单起见，并为专注于理论的本质方面，我们几乎只考虑光滑 $ r $ 形式与光滑向量场。$ C^{k} $ 情形只在注中简要提及，其验证留给你们。

我们用 $ \Omega^{r}(X) $ 表示 $ X $ 上所有光滑 $ r $ 形式的集合。为简便起见，我们记

\[
\mathcal{E}(X) := C^{\infty}(X) \quad \text{与} \quad \mathcal{V}(X) := \mathcal{V}^{\infty}(X).
\]

若 $ \boldsymbol{v}_{1},\ldots,\boldsymbol{v}_{r} $ 是 $ X $ 上的向量场，其相应的向量部分为 $ v_{1},\ldots,v_{r} $，即若 $ \boldsymbol{v}_{j}(x) = (x,v_{j}(x)) $ 对 $ x \in X $ 与 $ 1 \le j \le r $ 成立，则我们定义

\begin{equation}
\boldsymbol{\alpha}(v_{1},\ldots,v_{r})(x) := \boldsymbol{\alpha}(x)\big(\boldsymbol{v}_{1}(x),\ldots,\boldsymbol{v}_{r}(x)\big) \qquad \text{对 } x \in X.\tag{3.3}
\end{equation}

于是由 (VIII.3.1) 推出

\[
\boldsymbol{\alpha}(x)\big(\boldsymbol{v}_{1}(x),\ldots,\boldsymbol{v}_{r}(x)\big) = \alpha(x)\big(v_{1}(x),\ldots,v_{r}(x)\big) \qquad \text{对 } x \in X,
\]

即

\begin{equation}
\boldsymbol{\alpha}(\boldsymbol{v}_{1},\ldots,\boldsymbol{v}_{r}) = \alpha(v_{1},\ldots,v_{r}).\tag{3.4}
\end{equation}

这表明，在不引起误解的情况下，我们可以把 $ r $ 形式 $ \boldsymbol{\alpha} $ 与其余向量部分 $ \alpha $ 等同起来，把向量场 $ \boldsymbol{v} $ 与其向量部分 $ v $ 等同起来。因此，从现在起我们将用普通字体（非粗体）书写微分形式与向量场。在每种情形下，你都能毫不费力地判断一个符号描述的是形式还是它的余向量部分（或者它指的是向量场还是其向量部分）。

加法

\[
\Omega^{r}(X) \times \Omega^{r}(X) \to \Omega^{r}(X), \quad (\alpha,\beta) \mapsto \alpha+\beta
\]

与外积

\[
\wedge: \Omega^{r}(X) \times \Omega^{s}(X) \to \Omega^{r+s}(X), \quad (\alpha,\beta) \mapsto \alpha \wedge \beta
\]

是逐点进行的：

\[
(\alpha+\beta)(x) := \alpha(x)+\beta(x) \quad \text{与} \quad (\alpha \wedge \beta)(x) := \alpha(x) \wedge \beta(x) \qquad \text{对 } x \in X.
\]

这些映射显然是良定义的。

\begin{remark}

\begin{items}
\item

$ \Omega^{0}(X) = \mathcal{E}(X) $。

\item

$ \Omega^{1}(X) = \Omega_{(\infty)}(X) $，即 $ X $ 的光滑 1 形式是 $ X $ 上的 $ C^{\infty} $ 普法夫形式。

\item

当 $ r > m $ 时 $ \Omega^{r}(X) = \{0\} $。

\item

对 $ 0 \le r \le m $，$ \Omega^{r}(X) $ 是无限维实向量空间，也是维数为 $ \binom{m}{r} $ 的自由 $ \mathcal{E}(X) $ 模（关于逐点乘法）。$ \Omega^{r}(X) $ 的一个模基由

\begin{equation}
\{ dx^{(j)} := dx^{j_{1}} \wedge \cdots \wedge dx^{j_{m}} ;\ (j) \in \mathbb{J}_{r} \}\tag{3.5}
\end{equation}

给出。

\proofhead
由于 (a) 以及 $ \mathbb{R} $ 与 $ \mathcal{E}(X) $ 的子环\footnote{$ \mathbf{1}(x) = 1 $ 对 $ x \in X $。} $ \mathbb{R}\mathbf{1} $ 的典范等同，我们对 $ \alpha \in \mathbb{R} $ 与 $ \beta \in \Omega(X) $ 有 $ \alpha \wedge \beta = \alpha\beta $。第一个断言随即由注 2.2(a) 与例 I.12.3(e) 推出。
根据注 VIII.3.3，$ (dx^{1}(x),\ldots,dx^{m}(x)) $ 是与 $ T_{x}X $ 的典范基 $ ((e_{1})_{x},\ldots,(e_{m})_{x}) $ 对偶的基。于是其余断言由命题 2.3 推出。\nobreak\hspace{0.35em}■
\item

$ X $ 上的 $ r $ 形式 $ \alpha $ 属于 $ C^{k} $ 类，当且仅当 $ \mathcal{V}^{k}(X) $ 中的每个 $ r $ 元组 $ v_{1},\ldots,v_{r} $ 都满足

\begin{equation}
\alpha(v_{1},\ldots,v_{r}) \in C^{k}(X).\tag{3.6}
\end{equation}

这成立当且仅当典范基表示\footnote{由命题 2.3 可知（正如在 (d) 的证明中那样），(3.5) 是 $ X $ 上所有 $ r $ 形式构成的 $ \mathbb{R}^{X} $ 模的基。}的系数 $ a_{(j)} $

\begin{equation}
\alpha = \sum_{(j) \in \mathbb{J}_{r}} a_{(j)} dx^{(j)}\tag{3.7}
\end{equation}

满足关系

\begin{equation}
a_{(j)} \in C^{k}(X) \qquad \text{对 } (j) \in \mathbb{J}_{r}.\tag{3.8}
\end{equation}

\proofhead
当 $ \alpha $ 属于 $ C^{k} $ 类时，由 (3.2) 与推论 VII.4.7 容易推出 (3.6) 成立。又因为 (2.4) 蕴含
\begin{equation}
a_{(j)} = \alpha(e_{j_{1}},\ldots,e_{j_{r}}),\tag{3.9}
\end{equation}

故由 (3.6) 得到 (3.8)。若 (3.8) 满足，则由 (3.7) 以及基形式 $ dx^{(j)} $ 的常值性可知 $ \alpha $ 属于 $ C^{k} $ 类。\nobreak\hspace{0.35em}■

\item

外积\footnote{有时我们用“楔积”而不用“外积”。}是双线性的、结合的和分次反交换的。因此

\[
\Omega(X) := \bigoplus_{r \ge 0} \Omega^{r}(X)
\]

是（无限维的）结合分次反交换代数（关于乘积 $ \wedge $）。此外，$ \Omega(X) $ 是维数为 $ 2^{m} $ 的自由 $ \mathcal{E}(X) $ 模（关于逐点乘法）；我们称之为 $ X $ 上的\textbf{微分形式模}。

\proofhead
这些都是定理 2.7 与 (d) 的简单推论。\nobreak\hspace{0.35em}■
\item

每个 $ \alpha \in \Omega^{r}(X) $ 都是 $ \mathcal{V}(X) $ 上的交错 $ r $ 形式。

\proofhead
这由定义 (3.3) 立即推出。\nobreak\hspace{0.35em}■
\item

（正则性）对 $ k \in \mathbb{N} $，设 $ \Omega^{r}_{(k)}(X) $ 为 $ X $ 上 $ C^{k} $ 类 $ r $ 形式的集合。则把 $ \mathcal{E}(X) $ 处处换成 $ C^{k}(X) $ 后，前述各断言对 $ \Omega^{r}_{(k)}(X) $ 仍类似成立。\nobreak\hspace{0.35em}■

\end{items}

在下文中，我们一般不再声明 $ \alpha \in \Omega^{r}(X) $ 的典范基表示 (3.7) 的系数 $ a_{(j)} $ 属于 $ \mathcal{E}(X) $。这将被视为自明的。

\end{remark}

\begin{example}

\begin{items}
\item

如我们已知，每个普法夫形式 $ \alpha \in \Omega^{1}(X) $ 都有典范基表示

\[
\alpha = \sum_{j=1}^{m} a_{j} dx^{j}.
\]
\item

对 $ \alpha \in \Omega^{m-1}(X) $，基表示具有形式

\[
\alpha = \sum_{j=1}^{m} (-1)^{j-1} a_{j} dx^{1} \wedge \cdots \wedge \widehat{dx^{j}} \wedge \cdots \wedge dx^{m}.
\]
\item

当 $ m = 3 $ 时，任意 $ \alpha \in \Omega^{2}(X) $ 都有基表示\footnote{注意指标的循环置换。}

\[
\alpha = a_{1} dx^{2} \wedge dx^{3} + a_{2} dx^{3} \wedge dx^{1} + a_{3} dx^{1} \wedge dx^{2}.
\]

\proofhead
因为 $ dx^{3} \wedge dx^{2} = -dx^{2} \wedge dx^{3} $，故这由 (b) 推出。\nobreak\hspace{0.35em}■
\item

每个 $ \alpha \in \Omega^{m}(X) $ 都具有形式 $ a\,dx^{1} \wedge \cdots \wedge dx^{m} $，其中 $ a \in \mathcal{E}(X) $。

\item

对 $ m = 3 $，楔积

\[
\alpha = a_{1}\,dx^{1} + a_{2}\,dx^{2} + a_{3}\,dx^{3} \quad \text{与} \quad \beta = b_{1}\,dx^{1} + b_{2}\,dx^{2} + b_{3}\,dx^{3}
\]

是

\[
\alpha \wedge \beta = (a_{2}b_{3} - a_{3}b_{2})\,dx^{2} \wedge dx^{3} + (a_{3}b_{1} - a_{1}b_{3})\,dx^{3} \wedge dx^{1} + (a_{1}b_{2} - a_{2}b_{1})\,dx^{1} \wedge dx^{2}.
\]

\proofhead
这由注 3.1(f) 推出。\nobreak\hspace{0.35em}■
\end{items}

\end{example}

\needspace{3\baselineskip}\subsection*{拉回}

设 $ Y $ 是 $ \mathbb{R}^{n} $ 中的开集，$ \varphi \in C^{\infty}(X, Y) $。作为普法夫形式拉回的一种推广，我们引入微分形式经 $ \varphi $ 的\textbf{拉回}。它是映射

\begin{equation}
\varphi^{*}: \Omega(Y) \to \Omega(X)\tag{3.10}
\end{equation}

定义如下：

\begin{equation}
(\varphi^{*}\beta)(x) := (T_{x}\varphi)^{*}\beta(\varphi(x)) \quad \text{对 } x \in X \text{ 和 } \beta \in \Omega(Y).\tag{3.11}
\end{equation}

若 $ \beta \in \Omega^{r}(Y) $，则由于 $ T_{x}\varphi \in \mathcal{L}(T_{x}X, T_{\varphi(x)}Y) $ 且由注 2.9(a)，$ (T_{x}\varphi)^{*}\beta(\varphi(x)) $ 与 $ \beta(\varphi(x)) \in \bigwedge^{r}T_{\varphi(x)}^{*}Y $ 都落在 $ \bigwedge^{r}T_{x}^{*}X $ 中。由 $ T_{x}\varphi = (\varphi(x), \partial\varphi(x)) $ 以及

\[
\partial\varphi \in C^{\infty}(X, \mathcal{L}(\mathbb{R}^{m}, \mathbb{R}^{n}))
\]

并且 (3.4) 蕴含

\[
\varphi^{*}\beta(v_{1}, \ldots, v_{r}) = (\beta \circ \varphi)((\partial\varphi)v_{1}, \ldots, (\partial\varphi)v_{r}) \quad \text{对 } v_{1}, \ldots, v_{r} \in \mathcal{V}(X),
\]

由注 3.1(e) 可见 $ \varphi^{*}\beta $ 属于 $ \Omega^{r}(X) $。因此 (3.10) 通过 (3.11) 得到良定义。

\begin{remark}

\begin{items}
\item

映射 (3.10) 是 $ \mathbb{R} $ 线性的，且满足

\[
(\psi \circ \varphi)^{*} = \varphi^{*} \circ \psi^{*} \quad \text{与} \quad (\mathrm{id}_{X})^{*} = \mathrm{id}_{\Omega(X)},
\]

即拉回反变地作用。它也与外积相容，即

\[
\varphi^{*}(\alpha \wedge \beta) = \varphi^{*}\alpha \wedge \varphi^{*}\beta \quad \text{对 } \alpha, \beta \in \Omega(Y).
\]

因此 $ \varphi^{*} $ 是从 $ \Omega(Y) $ 到 $ \Omega(X) $ 的代数同态。

\proofhead
这由注 2.9 以及注 VII.10.2(b) 中给出的链式法则推出。\nobreak\hspace{0.35em}■
\item

（正则性）拉回也可对 $ \varphi \in C^{k+1}(X, Y) $ 自然地定义。若 $ 1 \le r \le m $，则一个 $ C^{k+1} $ 类的 $ r $ 形式一般只成为一个 $ C^{k} $ 类的 $ r $ 形式，而一个 $ C^{k} $ 类的 $ r $ 形式仍留在同一类中。当 $ r = 0 $ 时，拉回保持正则性。\nobreak\hspace{0.35em}■

\end{items}

\end{remark}

\begin{example}

设 $ (x^{1}, \ldots, x^{m}) $ 与 $ (y^{1}, \ldots, y^{n}) $ 分别是 $ X $ 与 $ Y $ 的欧几里得坐标。

\begin{items}
\item
\[
\varphi^{*}dy^{j} = d\varphi^{j} = \sum_{k=1}^{m} \partial_{k}\varphi^{j}\,dx^{k} \quad \text{且} \quad 1 \le j \le n.
\]

\proofhead
见例 VIII.3.14(a)。\nobreak\hspace{0.35em}■
\item

对

\[
\beta = \sum_{(j) \in \mathbb{J}_{r}} b_{(j)}\,dy^{(j)} \in \Omega^{r}(Y),
\]

我们有

\[
\varphi^{*}\beta = \sum_{(j) \in \mathbb{J}_{r}} (\varphi^{*}b_{(j)})\,d\varphi^{(j)}.
\]

\proofhead
这是 (a) 和注 3.3(a) 的推论。\nobreak\hspace{0.35em}■
\item

当 $ m = n $ 时，我们有

\[
\varphi^{*}(dy^{1} \wedge \cdots \wedge dy^{m}) = d\varphi^{1} \wedge \cdots \wedge d\varphi^{m} = (\det\partial\varphi)\,dx^{1} \wedge \cdots \wedge dx^{m}.
\]

\proofhead
第一个等号由 (b) 推出。因为对 $ x \in X $ 有 $ \det T_{x}\varphi = \det\partial\varphi(x) $，所以论断由命题 2.10 以及基形式 $ dx^{1} \wedge \cdots \wedge dx^{m} $ 在 $ X $ 上的常性推出。\nobreak\hspace{0.35em}■
\item

设 $ m = 2 $、$ n = 3 $，并设 $ (u, v) $ 与 $ (x, y, z) $ 分别是 $ X $ 与 $ Y $ 的欧几里得坐标。则\footnote{见注 VII.7.9。}

\[
\varphi^{*}(a\,dy \wedge dz + b\,dz \wedge dx + c\,dx \wedge dy) = \left[a \circ \varphi\,\frac{\partial(\varphi^{2}, \varphi^{3})}{\partial(u, v)} + b \circ \varphi\,\frac{\partial(\varphi^{3}, \varphi^{1})}{\partial(u, v)} + c \circ \varphi\,\frac{\partial(\varphi^{1}, \varphi^{2})}{\partial(u, v)}\right] du \wedge dv.
\]

\proofhead
因为对 $ 1 \le j \le 3 $ 有 $ d\varphi^{j} = \varphi_{u}^{j}\,du + \varphi_{v}^{j}\,dv $，并且
\begin{equation*}
\frac{\partial(\varphi^{2}, \varphi^{3})}{\partial(u, v)} = \det\begin{bmatrix} \varphi_{u}^{2} & \varphi_{v}^{2} \\ \varphi_{u}^{3} & \varphi_{v}^{3} \end{bmatrix} = \varphi_{u}^{2}\varphi_{v}^{3} - \varphi_{u}^{3}\varphi_{v}^{2}
\end{equation*}

等等，所以该论断由 (b) 和例 3.2(e) 推出。\nobreak\hspace{0.35em}■

\item

（平面极坐标）设

\[
f_{2}: \mathbb{R}^{2} \to \mathbb{R}^{2}, \quad (r, \varphi) \mapsto (x, y) := (r\cos\varphi, r\sin\varphi)
\]

为极坐标映射。则

\[
f_{2}^{*}(dx \wedge dy) = r\,dr \wedge d\varphi.
\]

\proofhead
这由 (c) 和例 X.8.7 推出。\nobreak\hspace{0.35em}■
\item

（球坐标）对球坐标映射

\[
f_{3}: \mathbb{R}^{3} \to \mathbb{R}^{3}, \quad (r, \varphi, \vartheta) \mapsto (x, y, z) := (r\cos\varphi\sin\vartheta, r\sin\varphi\sin\vartheta, r\cos\vartheta),
\]

我们有

\[
f_{3}^{*}(dx \wedge dy \wedge dz) = -r^{2}\sin\vartheta\,dr \wedge d\varphi \wedge d\vartheta.
\]

\proofhead
引理 X.8.8 和 (c)。\nobreak\hspace{0.35em}■
\item

（m 维极坐标）设

\[
f_{m}: \mathbb{R}^{m} \to \mathbb{R}^{m}, \quad (r, \varphi, \vartheta_{1}, \ldots, \vartheta_{m-2}) \mapsto (x^{1}, \ldots, x^{m})
\]

为 $ m $ 维极坐标映射 (X.8.17)。则

\[
f_{m}^{*}dx^{1} \wedge \cdots \wedge dx^{m} = (-1)^{m}r^{m-1}w_{m}(\vartheta)\,dr \wedge d\varphi \wedge d\vartheta_{1} \wedge \cdots \wedge d\vartheta_{m-2},
\]

其中 $ w_{m}(\vartheta) := \sin\vartheta_{1}\sin^{2}\vartheta_{2} \cdots \sin^{m-2}\vartheta_{m-2} $。

\proofhead
这由引理 X.8.8 推出。\nobreak\hspace{0.35em}■
\item

（柱坐标）设

\[
f: \mathbb{R}^{3} \to \mathbb{R}^{3}, \quad (r, \varphi, z) \mapsto (x, y, z) := (r\cos\varphi, r\sin\varphi, z)
\]

为柱坐标映射。则

\[
f^{*}(dx \wedge dy \wedge dz) = r\,dr \wedge d\varphi \wedge dz.
\]

\proofhead
例 VII.9.11(c) 和 (c)。\nobreak\hspace{0.35em}■
\item

若 $ \varphi $ 是常值映射，则对 $ r \ge 1 $ 的 $ \alpha \in \Omega^{r}(Y) $ 有 $ \varphi^{*}\alpha = 0 $。

\proofhead
因为对 $ 1 \le j \le n $ 有 $ d\varphi^{j} = 0 $，所以该论断是 (b) 的推论。\nobreak\hspace{0.35em}■
\item

设 $ m \le n $，并设 $ i: \mathbb{R}^{m} \hookrightarrow \mathbb{R}^{n} $ 是把 $ \mathbb{R}^{m} $ 与 $ \mathbb{R}^{m} \times \{0\} \subset \mathbb{R}^{n} $ 等同起来的自然嵌入。又设 $ Y $ 是 $ \mathbb{R}^{n} $ 中的开集，满足

\[
Y \cap (\mathbb{R}^{m} \times \{0\}) \supset i(X).
\]

注意 $ X $ 是 $ Y $ 的 $ m $ 维子流形。

\end{items}

\bookfigure[0.50]{figures/9616091fd992dd52.png}

对 $ \alpha \in \Omega^{r}(Y) $，定义 $ \alpha \mid X $，即 $ \alpha $ 在 $ X $ 上的\textbf{限制}，为

\[
(\alpha \mid X)(x) := \alpha(x, 0) \mid (T_{x}X)^{r} \quad \text{对 } x \in X.
\]

换言之，当 $ \alpha $ 具有基表示

\[
\alpha = \sum_{(j) \in \mathbb{J}_{r}^{n}} a_{(j)}\,dx^{(j)},
\]

由此推出

\begin{equation*}
\begin{aligned}
(\alpha \mid X)(x) = \sum_{\substack{(j) \in \mathbb{J}_{r}^{n} \\ j_{r} \le m}} a_{(j)}(x, 0)\,dx^{(j)} \quad \text{对 } x \in X.
\end{aligned}
\end{equation*}

则 $ i^{*}\alpha = \alpha \mid X $。

\proofhead
由于 $ i^{*} $ 与限制映射
\[
\Omega^{r}(Y) \to \Omega^{r}(X), \quad \alpha \mapsto \alpha \mid X
\]

的线性性，只需考虑 $ \alpha = a\,dx^{(j)} $（$ (j) \in \mathbb{J}_{r}^{n} $）的情形。于是由 (b) 推出

\[
i^{*}\alpha = (i^{*}a)\,di^{(j)}.
\]

由 $ (i^{*}a)(x) = a(x, 0) $ 以及 $ i^{k} = \mathrm{pr}_{k}i = 0 $（$ m+1 \le k \le n $，其中 $ \mathrm{pr}_{k}: \mathbb{R}^{n} \to \mathbb{R} $ 是标准投影），我们有 $ di^{(j)} = 0 $（$ j_{r} > m $）。对 $ j_{r} \le m $，可得 $ di^{(j)} = dx^{(j)} $。现在论断显而易见。\nobreak\hspace{0.35em}■

(k) 设 $ (q, p) \in \mathbb{R}^{m} \times \mathbb{R}^{m} = \mathbb{R}^{2m} $ 是 $ \mathbb{R}^{2m} $ 的任意一点。我们通过

\[
\sigma := \sum_{j=1}^{m} dp^{j} \wedge dq^{j}
\]

定义 $ \mathbb{R}^{2m} $ 上的（标准）\textbf{辛形式}。
我们用 $ \mathrm{Sp}(2m) $ 表示所有满足 $ S^{*}\sigma = \sigma $ 的 $ S \in \mathcal{L}(\mathbb{R}^{2m}) $ 的集合。则 $ \mathrm{Sp}(2m) $ 是 $ \mathcal{L}\mathrm{aut}(\mathbb{R}^{2m}) $ 的子群，即辛群。任何 $ S \in \mathrm{Sp}(2m) $ 都满足 $ \det(S) = 1 $。

\proofhead
我们通过 $ \alpha := \sigma \wedge \cdots \wedge \sigma $（$ m $ 个因子）定义 $ \alpha \in \Omega^{2m}(\mathbb{R}^{2m}) $。则存在 $ a \in \mathbb{R}^{\times} $ 使 $ \alpha = a\omega $，其中 $ \omega $ 表示 $ \mathbb{R}^{2m} $ 的体积元。现设 $ S \in \mathrm{Sp}(2m) $。则由 $ S^{*}\sigma = \sigma $ 和注 3.3(a) 推出
\[
S^{*}\alpha = S^{*}\sigma \wedge \cdots \wedge S^{*}\sigma = \sigma \wedge \cdots \wedge \sigma = \alpha.
\]

因为 $ S^{*}\alpha = S^{*}(a\omega) = aS^{*}\omega $，并且由 (c) 得

\[
\alpha = S^{*}\alpha = a\det(S)\omega = \det(S)\alpha,
\]

因此 $ \det(S) = 1 $。我们把 $ \mathrm{Sp}(2m) $ 是 $ \mathcal{L}\mathrm{aut}(\mathbb{R}^{2m}) $ 的子群这一证明留作练习。\nobreak\hspace{0.35em}■

\end{example}

\needspace{3\baselineskip}\subsection*{外微分}

在 §VIII.3 中我们看到，函数 $ f \in \mathcal{E}(X) = \Omega^{0}(X) $ 的微分 $ df $ 是光滑普法夫形式，从而是 $ \Omega^{1}(X) $ 的元素。显然 $ d: \Omega^{0}(X) \to \Omega^{1}(X) $ 是线性的。此外，由命题 VIII.3.12 我们知道 $ d $ 与拉回可交换。下面的定理表明，$ d $ 可延拓为从微分形式模 $ \Omega(X) $ 到自身的 $ \mathbb{R} $ 线性映射；该映射同样与拉回可交换。

\qyxffntextdef{1}{有时外微分也称为\textbf{嘉当导数}。}
\begin{theorem}

恰有一个映射

\[
d: \Omega(X) \to \Omega(X),
\]

即\textbf{外微分}\qyxffnmark{1}，具有性质

\begin{subitems}
\item

–(iv)：

\begin{items}
\item

$ d $ 是 $ \mathbb{R} $ 线性的，并把 $ \Omega^{r}(X) $ 映到 $ \Omega^{r+1}(X) $。

\end{items}

\item

$ d $ 满足\textbf{乘积法则}

\[
d(\alpha \wedge \beta) = d\alpha \wedge \beta + (-1)^{r}\alpha \wedge d\beta \quad \text{对 } \alpha \in \Omega^{r}(X) \text{ 和 } \beta \in \Omega(X).
\]
\item

$ d^{2} := d \circ d = 0 $。

\item

对 $ f \in \mathcal{E}(X) $，外微分 $ df $ 等于 $ f $ 的微分。
若 $ Y $ 是 $ \mathbb{R}^{n} $ 中的开集且 $ \varphi \in C^{\infty}(X, Y) $，则

\begin{equation}
d \circ \varphi^{*} = \varphi^{*} \circ d,\tag{3.12}
\end{equation}

即外微分与拉回可交换。

\end{subitems}

证明

\begin{items}
\item

（唯一性）对

\begin{equation}
\alpha = \sum_{(j) \in \mathbb{J}_{r}} a_{(j)}\,dx^{(j)} \in \Omega^{r}(X),\tag{3.13}
\end{equation}

由 (i)–(iv) 容易推出

\begin{equation}
d\alpha = \sum_{(j) \in \mathbb{J}_{r}} da_{(j)} \wedge dx^{(j)} \in \Omega^{r+1}(X).\tag{3.14}
\end{equation}

这表明至多有一个映射能满足性质 (i)–(iv)。

\item

（存在性）对按 (3.13) 展开的 $ \alpha \in \Omega^{r}(X) $，我们用 (3.14) 定义 $ d\alpha $。于是 $ d $ 显然满足要求 (i) 与 (iv)。
为证 (ii)，注意 (i) 意味着我们只需考虑 $ \alpha = a\,dx^{(j)} $ 与 $ \beta = b\,dx^{(k)} $（$ (j) \in \mathbb{J}_{r} $，$ (k) \in \mathbb{J}_{s} $）的情形。于是由 (3.14)、外积的性质以及推论 VII.3.8 中的普通乘积法则推出

\begin{equation*}
\begin{aligned}
d(\alpha \wedge \beta) &= d(ab\,dx^{(j)} \wedge dx^{(k)}) = d(ab) \wedge dx^{(j)} \wedge dx^{(k)} \\ &= da \wedge dx^{(j)} \wedge b\,dx^{(k)} + (-1)^{r}a\,dx^{(j)} \wedge db \wedge dx^{(k)} \\ &= d(a\,dx^{(j)}) \wedge b\,dx^{(k)} + (-1)^{r}a\,dx^{(j)} \wedge d(b\,dx^{(k)}) \\ &= d\alpha \wedge \beta + (-1)^{r}\alpha \wedge d\beta,
\end{aligned}
\end{equation*}

正如所愿。
为证 (iii)，我们可以利用 $ d $ 的线性性再次限制到 $ \alpha = a\,dx^{(j)} $（$ (j) \in \mathbb{J}_{r} $）的情形。于是由 (3.14) 和 (ii) 推出

\[
d(d\alpha) = d(da \wedge dx^{(j)}) = d^{2}a \wedge dx^{(j)} - da \wedge d(dx^{(j)}).
\]

对 $ d(dx^{(j)}) $ 反复应用乘积法则 (ii)，我们看到：只要能够证明对 $ a \in \Omega^{0}(X) = \mathcal{E}(X) $ 有 $ d^{2}a = 0 $，论断便成立。
因此设 $ a \in \mathcal{E}(X) $。则我们可以利用 (i)、(ii) 与 (iv) 导出关系

\begin{equation*}
\begin{aligned}
d(da) &= d\left(\sum_{k=1}^{m} \partial_{k}a\,dx^{k}\right) = \sum_{k=1}^{m} d(\partial_{k}a) \wedge dx^{k} \\ &= \sum_{j,k=1}^{m} \partial_{j}\partial_{k}a\,dx^{j} \wedge dx^{k} = \sum_{1 \le j < k \le m} (\partial_{j}\partial_{k}a - \partial_{k}\partial_{j}a)\,dx^{j} \wedge dx^{k} = 0,
\end{aligned}
\end{equation*}

其中最后一个等号由施瓦茨定理（推论 VII.5.5）推出。因此 (iii) 得到满足。

\item

设 $ \varphi \in C^{\infty}(X, Y) $，$ (j) \in \mathbb{J}_{r}^{n} $。令 $ \beta = b\,dy^{(j)} \in \Omega^{r}(Y) $。则根据例 3.4(b)，我们有

\begin{equation}
\varphi^{*}\beta = \varphi^{*}b\,d\varphi^{(j)} \in \Omega^{r}(X).\tag{3.15}
\end{equation}

由 (3.14) 以及注 3.3(a) 中说明的拉回性质，我们得到

\[
\varphi^{*}d\beta = \varphi^{*}(db \wedge dy^{(j)}) = \varphi^{*}db \wedge \varphi^{*}dy^{(j)} = \varphi^{*}db \wedge d\varphi^{(j)}.
\]

命题 VIII.3.12 蕴含 $ \varphi^{*}db = d(\varphi^{*}b) $。因此利用 (i)、(iii) 和 (3.15) 得到

\begin{equation*}
\begin{aligned}
\varphi^{*}d\beta &= d(\varphi^{*}b) \wedge d\varphi^{(j)} = d(\varphi^{*}b) \wedge d\varphi^{(j)} + (-1)^{1}\varphi^{*}b \wedge d(d\varphi^{(j)}) \\ &= d(\varphi^{*}b \wedge d\varphi^{(j)}) = d(\varphi^{*}\beta).
\end{aligned}
\end{equation*}

现在 (3.12) 由 $ \varphi^{*} $ 与 $ d $ 的线性性以及注 3.1(e) 推出。\nobreak\hspace{0.35em}■

\end{items}

\end{theorem}

\begin{remark}

\begin{items}
\item

对 $ \alpha = \sum_{(j) \in \mathbb{J}_{r}} a_{(j)}dx^{(j)} \in \Omega^{r}(X) $，我们有

\[
d\alpha = \sum_{(j) \in \mathbb{J}_{r}} da_{(j)} \wedge dx^{(j)}.
\]

\proofhead
这就是陈述 (3.13)、(3.14)。\nobreak\hspace{0.35em}■
\item

对 $ \varphi \in C^{\infty}(X, Y) $ 与 $ r \in \mathbb{N} $，图表

\begin{equation*}
\begin{array}{ccc} \Omega^{r}(Y) & \xrightarrow{d} & \Omega^{r+1}(Y) \\ \varphi^{\ast}\downarrow & & \downarrow\varphi^{\ast} \\ \Omega^{r}(X) & \xrightarrow{d} & \Omega^{r+1}(X) \end{array}
\end{equation*}

可交换。

\proofhead
这就是 (3.12)。\nobreak\hspace{0.35em}■
\item

（正则性）若 $ \alpha $ 是 $ C^{k+1} $ 类的 $ r $ 形式，则 $ d\alpha $ 显然是 $ C^{k} $ 类的 $ (r+1) $ 形式。然而，对 $ \alpha = a\,dx^{(j)} $，其中 $ (j) \in \mathbb{J}_{r}^{m} $，有

\[
d\alpha = da \wedge dx^{(j)} = \sum_{i} \partial_{i} a\, dx^{i} \wedge dx^{(j)},
\]

其中我们只对满足 $ i \neq j_{k} $（$ 1 \le k \le r $）的下标 $ i \in \{1, \ldots, m\} $ 求和，因为对其余下标有 $ dx^{i} \wedge dx^{(j)} = 0 $。因此存在 $ C^{k} $ 类的 $ r $ 形式 $ \alpha $，使得 $ d\alpha $ 也属于 $ C^{k} $ 类。\nobreak\hspace{0.35em}■

\end{items}

\end{remark}

\begin{example}

\begin{items}
\item

对 $ \alpha = \sum_{j=1}^{m} a_{j}\, dx^{j} \in \Omega^{1}(X) $，有

\[
d\alpha = \sum_{1 \le j < k \le m} (\partial_{j} a_{k} - \partial_{k} a_{j})\, dx^{j} \wedge dx^{k}.
\]
\item

对 $ \alpha = \sum_{j=1}^{m} (-1)^{j-1} a_{j}\, dx^{1} \wedge \cdots \wedge \widehat{dx^{j}} \wedge \cdots \wedge dx^{m} \in \Omega^{m-1}(X) $，得到

\[
d\alpha = \left( \sum_{j=1}^{m} \partial_{j} a_{j} \right) dx^{1} \wedge \cdots \wedge dx^{m}.
\]
\item

对 $ \alpha \in \Omega^{m}(X) $，$ d\alpha = 0 $。\nobreak\hspace{0.35em}■

\end{items}

\end{example}

\needspace{3\baselineskip}\subsection*{庞加莱引理}

若 $ d\alpha = 0 $，则称微分形式 $ \alpha \in \Omega(X) $ 为\textbf{闭的}。若存在\textbf{原函数} $ \beta \in \Omega(X) $ 使得\footnote{说一个形式是恰当的，意味着它的次数至少为 1。} $ d\beta = \alpha $，则称它是\textbf{恰当的}。

\begin{examplenote}

\begin{items}
\item

例 3.7(a) 表明，$ \alpha = \sum_{j=1}^{m} a_{j}\, dx^{j} $ 是闭的，当且仅当 $ \partial_{j} a_{k} = \partial_{k} a_{j} $（$ 1 \le j, k \le m $）。因此，在普法夫形式的情形，这个推广了的闭性概念归结为第 VIII.3 节的定义。

\item

每个恰当形式都是闭形式。

\proofhead
这由 $ d^{2} = 0 $ 推出。\nobreak\hspace{0.35em}■
\item

$ X $ 上的每个 $ m $ 形式都是闭形式。

\proofhead
见例 3.7(c)。\nobreak\hspace{0.35em}■
\item

（正则性）闭形式的定义显然对 $ C^{1} $ 类形式有意义；恰当形式的概念对连续微分形式有意义。\nobreak\hspace{0.35em}■

\end{items}

在定理 VIII.3.8 中我们已经看到，若 $ X $ 是星形的，则每个闭普法夫形式都是恰当的。在下文中，我们将证明庞加莱的这个“引理”在一般情形下也成立。

令 $ I := [0,1] $，并令 $ t $ 为 $ I $ 中的一般点。对 $ \ell \in \{0,1\} $，单射

\[
i_{\ell} : X \to I \times X ,\quad x \mapsto (\ell, x)
\]

是光滑的。显然，$ i_{0} $ 与 $ i_{1} $ 分别把 $ X $ 与 $ X $ 上的柱体 $ I \times X $ 的“底” $ \{0\} \times X $ 和“顶” $ \{1\} \times X $ 等同起来。因此\footnote{因为偏导数 $ \partial_{t} $ 定义在 $ I $ 上，所以微分形式在 $ I \times X $ 上如何定义是清楚的。注意 $ I \times X $ 是带边流形，并见第 4 节。}

\[
i_{\ell}^{\ast} : \Omega^{r}(I \times X) \to \Omega^{r}(X)
\]

有定义。对 $ \alpha \in \Omega(I \times X) $，形式 $ i_{0}^{\ast}\alpha $（或 $ i_{1}^{\ast}\alpha $）是 $ \alpha $ 到 $ X $ 的限制。它可通过把 $ \alpha $ 的典范基表示中的系数里的 $ (t,x) $ 换为 $ (0,x) $（或 $ (1,x) $），并删去所有出现 $ dt $ 的项而得到（见例 3.4(j)）。

我们定义线性映射

\[
K : \Omega^{r+1}(I \times X) \to \Omega^{r}(X)
\]

为

\begin{equation}
K\alpha := \sum_{(j)\in\mathbb{J}_{r}} \int_{0}^{1} a_{(j)}(t, \cdot)\, dt\, dx^{(j)}\tag{3.16}
\end{equation}

其中

\begin{equation}
\alpha = \sum_{(j)\in\mathbb{J}_{r}} a_{(j)}\, dt \wedge dx^{(j)} + \sum_{(k)\in\mathbb{J}_{r+1}} b_{(k)}\, dx^{(k)}.\tag{3.17}
\end{equation}
\end{examplenote}

\begin{lemma}

$ K $ 是良定义的，并且满足

\begin{equation}
K \circ d + d \circ K = i_{1}^{\ast} - i_{0}^{\ast}.\tag{3.18}
\end{equation}

\proofhead
关于含参积分可微性的定理（定理 X.3.18）容易推出，由 (3.16) 对 (3.17) 中的 $ \alpha $ 所定义的 $ K\alpha $ 属于 $ \Omega^{r}(X) $。显然映射 $ K $ 也是线性的。

为证明 (3.18)，只需考虑 $ \alpha = a\,dt \wedge dx^{(j)} $ 与 $ \alpha = b\,dx^{(k)} $ 这两种情形，其中 $ (j) \in \mathbb{J}_{r} $，$ (k) \in \mathbb{J}_{r+1} $。

\begin{subitems}
\item

令 $ \alpha = a\,dt \wedge dx^{(j)} $。则 $ i_{0}^{\ast}\alpha = i_{1}^{\ast}\alpha = 0 $。我们还得到

\begin{equation*}
\begin{aligned}
K d\alpha &= K(da \wedge dt \wedge dx^{(j)}) = K\left( \sum_{\ell=1}^{m} \partial_{x^{\ell}} a\, dx^{\ell} \wedge dt \wedge dx^{(j)} \right) \\ &= -\sum_{\ell=1}^{m} \int_{0}^{1} \partial_{x^{\ell}} a(t, \cdot)\, dt\, dx^{\ell} \wedge dx^{(j)},
\end{aligned}
\end{equation*}

其中我们用了 $ dt \wedge dt \wedge dx^{(j)} = 0 $。另一方面，定理 X.3.18 给出

\[
d(K\alpha) = d\left( \int_{0}^{1} a(t, \cdot)\, dt\, dx^{(j)} \right) = \sum_{\ell=1}^{m} \int_{0}^{1} \partial_{x^{\ell}} a(t, \cdot)\, dt\, dx^{\ell} \wedge dx^{(j)}.
\]

这就证明了本情形的断言。

\item

令 $ \alpha = b\,dx^{(k)} $，其中 $ (k) \in \mathbb{J}_{r+1} $。则 $ K\alpha = 0 $，因此 $ dK\alpha = 0 $。此外，我们求得

\[
d\alpha = \partial_{t} b\, dt \wedge dx^{(k)} + \sum_{\ell=1}^{m} \partial_{x^{\ell}} b\, dx^{\ell} \wedge dx^{(k)}
\]

以及

\[
K d\alpha = \int_{0}^{1} \partial_{t} b(\tau, \cdot)\, d\tau\, dx^{(k)} = (b(1, \cdot) - b(0, \cdot))\, dx^{(k)} = i_{1}^{\ast}\alpha - i_{0}^{\ast}\alpha,
\]

所以本情形的断言也成立。\nobreak\hspace{0.35em}■

\end{subitems}

\end{lemma}

设 $ M $ 与 $ N $ 是流形。若存在映射\footnote{注意 $ I \times M $ 是带边流形。} $ h \in C^{\infty}(I \times M, N) $，即一个\textbf{同伦}，使得 $ h(j, \cdot) = f_{j} $（$ j = 0, 1 $），则称两个映射 $ f_{0}, f_{1} \in C^{\infty}(M, N) $ 在 $ N $ 中\textbf{同伦}。若映射 $ f \in C^{\infty}(M, N) $ 在 $ N $ 中同伦于常值映射，则称它在 $ N $ 中是\textbf{零同伦的}。最后，若从 $ M $ 到 $ M $ 的恒等映射是零同伦的，则称 $ M $ 是\textbf{可缩的}。

\begin{remark}

\begin{items}
\item

语句“$ f_{1} $ 在 $ N $ 中同伦于 $ f_{2} $”在 $ C^{\infty}(M, N) $ 中（或更一般地，在 $ C^{k}(M, N) $ 中）定义了一个等价关系。

\item

（连续）同伦的概念显然推广了环路同伦的想法（见第 VIII.4 节）。

\item

每个星形开集都是可缩的。

\proofhead
设 $ X $ 关于 $ x_{0} \in X $ 是星形的。则
\[
h : I \times X \to X ,\quad (t, x) \mapsto x_{0} + t(x - x_{0})
\]

显然是同伦，且 $ h(0, \cdot) = x_{0} $，$ h(1, \cdot) = \mathrm{id}_{X} $。\nobreak\hspace{0.35em}■

\item

（正则性）对 $ k \in \mathbb{N}^{\times} $，若出现的所有函数都属于 $ C^{k} $ 类，则上述定义对 $ C^{k} $ 流形 $ M $ 与 $ N $ 有意义。若 $ M $ 与 $ N $ 是拓扑空间且所考虑的所有函数都连续，则这些定义也有意义。\nobreak\hspace{0.35em}■

\end{items}

我们现在可以容易地证明广义庞加莱引理。

\end{remark}

\begin{theorem}

（庞加莱引理）若 $ X $ 可缩，则 $ X $ 上每个闭微分形式都是恰当的。

\proofhead
设 $ \alpha \in \Omega^{r+1}(X) $ 是闭的。因为 $ X $ 可缩，存在 $ h \in C^{\infty}(I \times X, X) $ 使得 $ h(1, \cdot) = \mathrm{id}_{X} $，且对某个 $ p \in X $ 有 $ h(0, \cdot) = p $。因为 $ \alpha $ 是闭的，$ h^{\ast}\alpha \in \Omega^{r+1}(I \times X) $ 也是闭的，因为 $ d \circ h^{\ast} = h^{\ast} \circ d $。因此由引理 3.9 推出
\[
d(Kh^{\ast}\alpha) = i_{1}^{\ast}h^{\ast}\alpha - i_{0}^{\ast}h^{\ast}\alpha = (h \circ i_{1})^{\ast}\alpha = \alpha.
\]

这是因为 $ h \circ i_{1} = \mathrm{id}_{X} $，且根据例 3.4(i)，$ i_{0}^{\ast}h^{\ast}\alpha = (h \circ i_{0})^{\ast}\alpha $ 是零形式。\nobreak\hspace{0.35em}■

\end{theorem}

我们应当指出，庞加莱引理的证明给出了一个构造给定闭微分形式的原函数的显式程序。当 $ X $ 是星形时情形特别简单，此时我们可以不失一般性地（通过作适当的平移）假设 $ X $ 关于 $ 0 $ 是星形的。

\begin{corollary}

设 $ X $ 关于 $ 0 $ 是星形的。设 $ r \in \mathbb{N}^{\times} $，并设

\[
\alpha = \sum_{(j)\in\mathbb{J}_{r}} a_{(j)}\, dx^{(j)} \in \Omega^{r}(X)
\]

是闭的。又令

\begin{equation}
\beta := \sum_{(j)\in\mathbb{J}_{r}} \sum_{k=1}^{r} (-1)^{k-1} \int_{0}^{1} t^{r-1} a_{(j)}(tx)\, dt\, x^{j_{k}} dx^{j_{1}} \wedge \cdots \wedge \widehat{dx^{j_{k}}} \wedge \cdots \wedge dx^{j_{r}}.\tag{3.19}
\end{equation}

则 $ \beta $ 属于 $ \Omega^{r-1}(X) $，且 $ d\beta = \alpha $。

\proofhead
在本情形中，$ h(t, x) := tx $（$ (t, x) \in I \times X $）定义了“$ X $ 到 $ 0 $ 的收缩”。由 $ dh^{j} = x^{j}\, dt + t\, dx^{j} $ 与例 3.4(b) 推出
\begin{equation*}
\begin{aligned}
h^{\ast}\alpha(t, x) &= \sum_{(j)\in\mathbb{J}_{r}} a_{(j)}(tx) t^{r}\, dx^{(j)} \\ &\quad + \sum_{(j)\in\mathbb{J}_{r}} \sum_{k=1}^{r} (-1)^{k-1} a_{(j)}(tx) t^{r-1} x^{j_{k}}\, dt \wedge dx^{j_{1}} \wedge \cdots \wedge \widehat{dx^{j_{k}}} \wedge \cdots \wedge dx^{j_{r}},
\end{aligned}
\end{equation*}

因为那些 $ dt $ 至少出现两次的项为零。由 (3.16) 与 (3.17) 推出 $ \beta = Kh^{\ast}\alpha $，而断言现在由庞加莱引理的证明推出。\nobreak\hspace{0.35em}■

\end{corollary}

\begin{remark}

\begin{items}
\item

当 $ r = 1 $ 时，即当 $ \alpha $ 是普法夫形式时，$ \beta $ 的公式与 (VIII.3.4) 中的情形相同。

\item

令 $ m = 3 $，且

\[
\alpha = a_{1}\, dx^{2} \wedge dx^{3} + a_{2}\, dx^{3} \wedge dx^{1} + a_{3}\, dx^{1} \wedge dx^{2} \in \Omega^{2}(X).
\]

则求 $ \beta = \sum_{j=1}^{3} b_{j}\, dx^{j} $ 使 $ d\beta = \alpha $ 的问题，等价于求三个函数 $ b_{1}, b_{2}, b_{3} \in \mathcal{E}(X) $ 使之满足 $ X $ 中的偏微分方程组

\begin{equation}
\begin{aligned}
\partial_{1} b_{2} - \partial_{2} b_{1} &= a_{3}, \\ \partial_{2} b_{3} - \partial_{3} b_{2} &= a_{1}, \\ \partial_{3} b_{1} - \partial_{1} b_{3} &= a_{2}
\end{aligned}\tag{3.20}
\end{equation}

的问题。然后，对给定的 $ a_{j} \in \mathcal{E}(X) $，

\begin{equation}
\partial_{1} a_{1} + \partial_{2} a_{2} + \partial_{3} a_{3} = 0\tag{3.21}
\end{equation}

是 (3.20) 有解所必需的。若 $ X $ 可缩（例如 $ X = \mathbb{R}^{3} $），则 (3.21) 也是充分的。

\proofhead
由例 3.7(a)，(3.20) 等价于 $ d\beta = \alpha $。例 3.7(b) 表明 (3.21) 等价于 $ d\alpha = 0 $。现在断言由 $ d^{2} = 0 $ 与庞加莱引理推出。\nobreak\hspace{0.35em}■
\end{items}

由推论 3.12 特别推出，在星形区域的情形，(3.20) 可用公式 (3.19) 通过求积解出。在一般情形，方程 $ d\beta = \alpha $ 显然也可以重新表述为等价的偏微分方程组。

当然，(3.20) 没有唯一解，因为可以给 $ \beta $ 加上一个闭形式，也就是说，可以加上把 (3.20) 右端置零所得齐次方程组的任意解 $ (b_{1}, b_{2}, b_{3}) $。

\end{remark}

\needspace{3\baselineskip}\subsection*{张量}

设 $ r, s \in \mathbb{N} $。对 $ x \in X $，令

\begin{equation}
T_{s}^{r}(T_{x}X) := \{x\} \times T_{s}^{r}(\mathbb{R}^{m})\tag{3.22}
\end{equation}

并称 $ \gamma \in T_{s}^{r}(T_{x}X) $ 为 $ T_{x}X $ 上的 $ r $ 逆变、$ s $ 协变的\textbf{张量}，或类型 $ (r,s) $ 的\textbf{张量}。$ X $ 上 $ (r,s) $ 张量的\textbf{丛}定义为

\[
T_{s}^{r}(X) := \bigcup_{x\in X} T_{s}^{r}(T_{x}X) = X \times T_{s}^{r}(\mathbb{R}^{m}).
\]

映射

\[
\gamma : X \to T_{s}^{r}(X) \quad \text{且} \quad \gamma(x) \in T_{s}^{r}(T_{x}X),
\]

即张量丛 $ T_{s}^{r}(X) $ 的一个截面，称为 $ X $ 上的 $ (r,s) $ \textbf{张量（场）}或类型 $ (r,s) $ 的\textbf{张量}。由 (3.22)，$ X $ 上每个 $ (r,s) $ 张量 $ \gamma $ 都有唯一表示

\[
\gamma(x) = (x, \gamma(x)) \quad \text{对 } x \in X,
\]

其\textbf{主部}\footnote{当 $ s = 0 $ 时，我们称之为向量部分；当 $ r = 0 $ 时，我们称之为余向量部分。张量把向量与余向量结合在一起，因此这样的术语不再可行。}为

\[
\gamma : X \to T_{s}^{r}(\mathbb{R}^{m}).
\]

设 $ k \in \mathbb{N} \cup \{\infty\} $。若 $ (r,s) $ 张量 $ \gamma $ 的主部具有这一性质，即若\footnote{照例，我们把 $ T_{x}\mathbb{R}^{m} $ 和 $ T_{x}^{\ast}\mathbb{R}^{m} $ 与 $ \mathbb{R}^{m} $ 等同起来。}

\[
\gamma \in C^{k}\left(X, \mathcal{L}^{r+s}(\mathbb{R}^{m}, \mathbb{R})\right).
\]

则 $ \gamma $ 属于 $ C^{k} $ 类（或当 $ k = \infty $ 时，称 $ \gamma $ 为 $ k $ 次连续可微或光滑）。

我们用

\[
\mathcal{T}_{s}^{r}(X)
\]

表示 $ X $ 上所有光滑 $ (r,s) $ 张量的集合。

若 $ \boldsymbol{\alpha}_{1}, \ldots, \boldsymbol{\alpha}_{r} $ 是 $ X $ 上的普法夫形式，$ \boldsymbol{v}_{1}, \ldots, \boldsymbol{v}_{s} $ 是 $ X $ 上的向量场，其对应的主部分别为 $ \alpha_{1}, \ldots, \alpha_{r} $ 与 $ v_{1}, \ldots, v_{s} $，则我们令

\[
\gamma(\boldsymbol{\alpha}_{1}, \ldots, \boldsymbol{\alpha}_{r}, \boldsymbol{v}_{1}, \ldots, \boldsymbol{v}_{s})(x) := \gamma(x)\left(\alpha_{1}(x), \ldots, \alpha_{r}(x), v_{1}(x), \ldots, v_{s}(x)\right)
\]

对 $ x \in X $。（这显然与 (3.4) 一致。）由于这些理由，我们可以沿用以前处理向量场与微分形式时所用的记号约定，也就是说，把张量与其主部等同起来，并且从现在起使用普通字体而不是粗体。

\textbf{加法}

\[
T_{s}^{r}(X) \times T_{s}^{r}(X) \to T_{s}^{r}(X) ,\quad (\gamma, \delta) \mapsto \gamma + \delta,
\]

\textbf{乘以函数}

\[
\mathcal{E}(X) \times T_{s}^{r}(X) \to T_{s}^{r}(X), ,\quad (f, \gamma) \mapsto f\gamma
\]

\textbf{以及张量积}

\begin{equation}
T_{s_{1}}^{r_{1}}(X) \times T_{s_{2}}^{r_{2}}(X) \to T_{s_{1}+s_{2}}^{r_{1}+r_{2}}(X) ,\quad (\gamma, \delta) \mapsto \gamma \otimes \delta\tag{3.23}
\end{equation}

也将逐点定义：

\[
(\gamma + \delta)(x) := \gamma(x) + \delta(x), \quad (f\gamma)(x) := f(x)\gamma(x), \quad (\gamma \otimes \delta)(x) := \gamma(x) \otimes \delta(x) .
\]

下面的注是注 2.20 与链式法则的简单推论。我们把详细的证明留给你们作为练习。

\begin{remark}

\begin{items}
\item

$ \mathcal{T}_0^1(X) = \mathcal{V}(X) $ 且 $ \mathcal{T}_1^0(X) = \Omega^1(X) $。又有

\[
\mathcal{T}_2^0(X) = C^\infty(X, \mathcal{L}^2(\mathbb{R}^m)) ,
\]

其中我们使用了张量与其主部的典范等同。

\item

张量积是 $ \mathcal{E}(X) $-双线性的且结合的。

\item

$ \mathcal{T}_s^r(X) $ 是无限维 $ \mathbb{R} $-向量空间，也是 $ m^{r+s} $ 维的 $ \mathcal{E}(X) $-模。取 $ \mathbb{R}^m $ 的典范基 $ (\partial/\partial x^1, \ldots, \partial/\partial x^m) $，

\begin{equation}
\left\{ \frac{\partial}{\partial x^{j_1}} \otimes \cdots \otimes \frac{\partial}{\partial x^{j_r}} \otimes dx^{k_1} \otimes \cdots \otimes dx^{k_s} ;\ j_i, k_i \in \{1, \ldots, m\} \right\}\tag{3.24}
\end{equation}

是 $ \mathcal{T}_s^r(X) $ 的一个模基。

\item

$ X $ 上的 $ (r, s) $-张量 $ \gamma $ 属于 $ \mathcal{T}_s^r(X) $ 当且仅当 $ \Omega^1(X) $ 中的每个 $ r $ 元组 $ \alpha_1, \ldots, \alpha_r $ 与 $ \mathcal{V}(X) $ 中的每个 $ s $ 元组 $ v_1, \ldots, v_s $ 都满足

\[
\gamma(\alpha_1, \ldots, \alpha_r, v_1, \ldots, v_s) \in \mathcal{E}(X) .
\]

这当且仅当 $ \gamma $ 在基 (3.24) 下的系数属于 $ \mathcal{E}(X) $。

\item

（正则性）上述定义与断言显然有类似结论，它们对 $ C^k $ 类张量仍然成立。\nobreak\hspace{0.35em}■

\end{items}

\end{remark}

\exerciseshead

\begin{exercises}
\item

设 $ \alpha, \beta \in \Omega(\mathbb{R}^4) $ 由

\[
\alpha := dx^1 + x^2 dx^2 \quad \text{与} \quad \beta := \sin(x^2)\, dx^1 \wedge dx^3 + \cos(x^3)\, dx^2 \wedge dx^4 ,
\]

给出，并定义 $ h \in C^\infty(\mathbb{R}^4, \mathbb{R}^4) $ 为 $ h(x) := (x^1, x^2, x^3 x^4, x^4) $。计算：

\begin{subitems}
\item

$ \gamma := \alpha \wedge \beta $；

\item

$ h^*\gamma $；

\item

$ h^*\gamma(0)(e_1, e_2, e_3 + e_4) $，其中 $ (e_1, e_2, e_3, e_4) $ 是 $ \mathbb{R}^4 $ 中的标准基；

\item

$ d\alpha, d\beta, d\gamma, d(h^*\gamma) $。

\end{subitems}

\item

设 $ f_3: \mathbb{R}^3 \to \mathbb{R}^3 $，$ (r, \varphi, \vartheta) \mapsto (x, y, z) $ 是球坐标映射。计算

\begin{items}
\item

$ f_3^* dx, f_3^* dy, f_3^* dz $；

\item

$ f_3^*(dy \wedge dz) $；

\item

$ f_3^* dx \wedge f_3^*(dy \wedge dz) $。

\end{items}

\item

一个简单的热力学系统（例如理想气体）由其体积 $ V $ 与温度 $ T $ 刻画（这里 $ V, T \in \mathbb{R} $）。这样的系统的状态则由压强 $ p := p(V, T) $ 与内能 $ E := E(V, T) $ 描述。根据热力学第二定律，该系统还有另一个状态函数 $ S := S(V, T) $，即熵，其微分由下式给出：

\[
dS := \frac{dE + p\,dV}{T} \quad \text{对 } T > 0 .
\]

证明以下事实：

\begin{items}
\item

$ E $ 与 $ p $ 满足关系

\[
\frac{\partial E}{\partial V} = T\frac{\partial p}{\partial T} - p .
\]
\item

满足状态方程 $ pV = RT $ 的理想气体（其中 $ R \in \mathbb{R} $ 是普适气体常数），其内能与体积无关，即 $ E = E(T) $。

\item

对范德瓦尔斯气体，其状态方程为

\begin{equation}
\left( p + \frac{a}{V^2} \right)(V - b) = cT \quad \text{对 } a, b, c \in \mathbb{R}^\times ,\tag{3.25}
\end{equation}

其内能确实依赖于体积。
（提示：(a) $ d^2 = 0 $。(c) 由 (3.25) 推出 $ T\,\partial p/\partial T = p + a/V^2 $。）
注 在物理学文献中，当 1 形式 $ \alpha $ 不是恰当形式时，$ d\alpha $ 常写作 $ \delta\alpha $。

\end{items}

\item

若存在 $ \alpha_1, \ldots, \alpha_r \in \Omega^1(X) $ 使得

\[
\alpha = \alpha_1 \wedge \alpha_2 \wedge \cdots \wedge \alpha_r ,
\]

则称 $ r $ 形式 $ \alpha \in \Omega^r(X) $ 是\textbf{可分解的}。

设 $ \alpha, \beta \in \Omega^r(X) $ 可分解。计算 $ (\alpha + \beta) \wedge (\alpha + \beta) $。

\item

设 $ \alpha = \sum_{j \le k} a_{jk} dx^j \wedge dx^k \in \Omega^2(X) $。证明 $ \alpha $ 可分解当且仅当

\[
a_{ij} a_{k\ell} + a_{jk} a_{i\ell} + a_{ki} a_{j\ell} = 0 \quad \text{对 } 1 \le i, j, k, \ell \le n ,
\]

其中 $ a_{jk} := -a_{kj} $（$ j \ge k $）。

\item

设 $ \alpha = \sum_{i \le j} a_{ij} dx^i \wedge dx^j \in \Omega^2(X) $。证明

\[
d\alpha = \sum_{i<j<k} \left( \frac{\partial a_{ij}}{\partial x^k} + \frac{\partial a_{jk}}{\partial x^i} + \frac{\partial a_{ki}}{\partial x^j} \right) dx^i \wedge dx^j \wedge dx^k .
\]
\item

计算下列外微分：

\begin{items}
\item

$ d\alpha \wedge \beta - \alpha \wedge d\beta $ 与

\item

$ d\alpha \wedge \beta \wedge \gamma + \alpha \wedge d\beta \wedge \gamma + \alpha \wedge \beta \wedge d\gamma $，其中 (b) 中 $ \alpha $ 与 $ \beta $ 为偶次。

\end{items}

\item

求 $ d\alpha $，若

\[
\alpha := \sum_{j=1}^{m} (-1)^{j-1} \frac{x^j}{|x|^m} dx^1 \wedge \cdots \wedge \widehat{dx^j} \wedge \cdots \wedge dx^m \in \Omega^{m-1}(\mathbb{R}^m \setminus \{0\}) .
\]
\item

设 $ \alpha := 2xz\,dy \wedge dz + dz \wedge dx - (z^2 + e^x)\,dx \wedge dy \in \Omega^2(\mathbb{R}^3) $。证明 $ \alpha $ 是恰当的，并求出一个原函数。

\item

设 $ \omega \in \Omega^2(X) $ 非退化。证明

\[
\Theta_\omega : \mathcal{V}(X) \to \Omega^1(X), \quad v \mapsto \omega(v, \cdot)
\]

是 $ \mathcal{E}(X) $-模同构。

\item

证明以下三个断言：

\begin{items}
\item

辛形式 $ \sigma \in \Omega^2(\mathbb{R}^{2m}) $ 非退化且闭。

\item

$ m $ 重外积 $ \sigma^m := \sigma \wedge \cdots \wedge \sigma \in \Omega^{2m}(\mathbb{R}^{2m}) $ 满足 $ \sigma^m \neq 0 $。

\item

根据习题 10 与 (b)，\textbf{辛梯度} $ \operatorname{sgrad} f := \Theta_\sigma^{-1} df \in \mathcal{V}(\mathbb{R}^{2m}) $ 对每个 $ f \in \mathcal{E}(\mathbb{R}^{2m}) $ 有定义。在坐标 $ (q, p) \in \mathbb{R}^m \times \mathbb{R}^m $ 中计算 $ \operatorname{sgrad} f $。

\end{items}

\item

若 $ \sigma $ 是 $ \mathbb{R}^{2m} $ 上的辛形式，则

\[
\{\cdot, \cdot\} : \mathcal{E}(\mathbb{R}^{2m}) \times \mathcal{E}(\mathbb{R}^{2m}) \to \mathcal{E}(\mathbb{R}^{2m}), \quad (f, g) \mapsto \sigma(\operatorname{sgrad} f, \operatorname{sgrad} g)
\]

称为\textbf{泊松括号}。对 $ f, g, h \in \mathcal{E}(\mathbb{R}^{2m}) $ 与 $ c \in \mathbb{R} $ 证明

\begin{subitems}
\item

在局部坐标 $ (q_1, \ldots, q_m, p_1, \ldots, p_m) $ 中，泊松括号为

\[
\{f, g\} = \sum_{j=1}^{m} \left( \frac{\partial f}{\partial p_j} \frac{\partial g}{\partial q_j} - \frac{\partial f}{\partial q_j} \frac{\partial g}{\partial p_j} \right);
\]
\item

$ \{f, cg + h\} = c\{f, g\} + \{f, h\} $；

\item

$ \{f, g\} = -\{g, f\} $；

\item

$ \{f, \{g, h\}\} + \{g, \{h, f\}\} + \{h, \{f, g\}\} = 0 $（雅可比恒等式）；

\item

$ \{f, gh\} = g\{f, h\} + h\{f, g\} $；

\item

$ \operatorname{sgrad}\{f, g\} = (\operatorname{sgrad} f \mid \operatorname{sgrad} g)_{\mathbb{R}^{2m}} $。

\end{subitems}

\item

证明泊松括号与 $ \mathbb{R}^{2m} $ 上的辛形式 $ \sigma $ 由关系

\[
df \wedge dg \wedge \sigma^{m-1} = \frac{1}{m} \{f, g\}\sigma^m .
\]

相联系。

\end{exercises}

\section{向量场与微分形式}

本节致力于微分形式的整体理论，即流形上的微分形式。第一部分本质上只是局部理论的简单移植，它要求我们关注正则性问题。借助一个关于单位分解的定理，我们进而可以把外微分这一重要概念推广到流形的情形，并证明我们为局部理论所建立的法则仍然适用。

整体理论引出了一个重要的新思想：流形的可定向性。我们给出刻画这一核心概念的多种方式，并考察大量例子。为了给流形上的积分理论作准备，我们给出许多重要流形的体积元的显式表示。

本节中，

- $ M $ 是 $ m $ 维的，而 $ N $ 是 $ n $ 维流形；
- $ r \in \mathbb{N} $。

\needspace{3\baselineskip}\subsection*{向量场}

所谓 $ M $ 上的\textbf{向量场} $ v $，是指一个映射

\[
v : M \to TM \qquad \text{其中 } v(p) \in T_{p}M \text{ 对 } p \in M ,
\]

即切丛的一个截面。若 $ v $ 是 $ M $ 上的向量场，则我们可以用从 $ M $ 到 $ N $ 的微分同胚把它“移植”过去。于是，对 $ \varphi \in \operatorname{Diff}^{1}(M,N) $，我们通过令

\[
\varphi_{*}v(q) := (T_{\varphi^{-1}(q)}\varphi)v(\varphi^{-1}(q)) , \quad q \in N .
\]

来定义 $ v $ 沿 $ \varphi $ 的\textbf{前推} $ \varphi_{*}v $。

\bookfigure[0.78]{figures/7f8d942111026200.png}

因此 $ \varphi_{*}v $ 是 $ N $ 上的向量场。对于 $ M $ 上的函数，双射 $ \psi : M \to N $ 的\textbf{前推}是如下对应：

\[
\psi_{*} : \mathbb{R}^{M} \to \mathbb{R}^{N} , \quad a \mapsto \psi_{*}a := a \circ \psi^{-1} .
\]

\begin{remark}

\begin{items}
\item

对于函数，前推 $ \psi_{*} $ 显然与拉回 $ \psi^{-1} $ 相同：$ \psi_{*} = (\psi^{-1})^{*} $。然而要注意，与拉回不同，前推只对双射有定义。特别地，必有 $ \dim(M) = \dim(N) $。\footnote{见习题 VII.10.9。}

\item

设 $ \varphi \in \operatorname{Diff}^{1}(M,N) $。则

\[
\varphi_{*}(a+b) = \varphi_{*}a + \varphi_{*}b , \quad \varphi_{*}(v+w) = \varphi_{*}v + \varphi_{*}w ,
\]

以及

\[
\varphi_{*}(av) = \varphi_{*}a\,\varphi_{*}v
\]

对 $ a, b \in \mathbb{R}^{M} $ 与 $ M $ 上的向量场 $ v $ 和 $ w $ 成立。

\item

设 $ \varphi \in \operatorname{Diff}(M,N) $ 与 $ \psi \in \operatorname{Diff}(N,L) $，其中 $ L $ 是另一个流形。则

\begin{equation}
(\psi \circ \varphi)_{*} = \psi_{*} \circ \varphi_{*} \quad \text{且} \quad (\operatorname{id}_{M})_{*} = \operatorname{id}_{\mathcal{F}(M)}\tag{4.1}
\end{equation}

对 $ \mathcal{F}(M) := \mathcal{E}(M) $ 或 $ \mathcal{F}(M) := \mathcal{V}(M) $ 成立。法则 (4.1) 意味着前推是\textbf{协变地}作用的。

\proofhead
该陈述对函数的前推是显然的。对于向量场，(4.1) 由注 VII.10.9(b) 的链式法则与注 1.14(c) 推出。\nobreak\hspace{0.35em}■
\end{items}

设 $ k \in \mathbb{N} \cup \{\infty\} $。$ M $ 上的向量场 $ v $ 属于类 $ C^{k} $（即它是 $ k $ 次\textbf{连续可微}的，或在 $ k = \infty $ 的情形是\textbf{光滑}的），若 $ M $ 的每一点 $ p $ 都有一个 $ p $ 附近的坐标卡 $ (\varphi,U) $，使得\footnote{见第 VIII.3 节。那里所说的对 $ \mathbb{H}^{m} $ 的开子集同样成立。} $ \varphi_{*}v \in \mathcal{V}^{k}(\varphi(U)) $。我们把 $ M $ 上所有 $ C^{k} $ 类向量场的集合记作 $ \mathcal{V}^{k}(M) $。为记号简便，我们置

\[
\mathcal{V}(M) := \mathcal{V}^{\infty}(M) \quad \text{且} \quad \mathcal{E}(M) := C^{\infty}(M) .
\]
\end{remark}

\begin{remark}

\begin{items}
\item

$ C^{k} $ 向量场的定义与坐标无关。若 $ v $ 是 $ C^{k} $ 向量场，且 $ (\psi,V) $ 是 $ M $ 的任意坐标卡，则 $ \psi_{*}v $ 属于类 $ C^{k} $。

\proofhead
因此设 $ (\psi,V) $ 是 $ M $ 的坐标卡。则我们需要证明 $ \psi_{*}v $ 属于类 $ C^{k} $。每个 $ q \in V $ 都有 $ M $ 的一个包含它的坐标卡 $ (\varphi,U) $，使得 $ \varphi_{*}v \in \mathcal{V}^{k}(\varphi(U)) $。于是由 (4.1) 推出 $ \psi_{*}v = (\psi \circ \varphi^{-1})_{*}\varphi_{*}v $。因为
\[
\psi \circ \varphi^{-1} \in \operatorname{Diff}(\varphi(U \cap V), \psi(U \cap V))
\]

且 $ \varphi_{*}v \in \mathcal{V}^{k}(\varphi(U)) $，我们得到 $ \psi_{*}v \in \mathcal{V}^{k}(\psi(U \cap V)) $。由于这对每个 $ q \in V $ 都成立，且可微性是局部性质，我们推出 $ \psi_{*}v \in \mathcal{V}^{k}(\psi(V)) $。\nobreak\hspace{0.35em}■

\item

逐点定义的运算

\[
\mathcal{V}(M) \times \mathcal{V}(M) \to \mathcal{V}(M) , \quad (v,w) \mapsto v + w
\]

以及

\[
\mathcal{E}(M) \times \mathcal{V}(M) \to \mathcal{V}(M) , \quad (a,v) \mapsto av
\]

使 $ \mathcal{V}(M) $ 成为 $ \mathcal{E}(M) $-模。特别地，$ \mathcal{E}(M) $ 与 $ \mathcal{V}(M) $ 是（无限维的）$ \mathbb{R} $-向量空间。
若 $ \varphi \in \operatorname{Diff}(M,N) $，则 $ \varphi_{*} $ 是从 $ \mathcal{E}(M) $ 到 $ \mathcal{E}(N) $ 以及从 $ \mathcal{V}(M) $ 到 $ \mathcal{V}(N) $ 的模同构。

\proofhead
由 (4.1) 推出
\[
\operatorname{id}_{M} = (\varphi^{-1} \circ \varphi)_{*} = (\varphi^{-1})_{*}\varphi_{*} \quad \text{且} \quad \operatorname{id}_{N} = (\varphi \circ \varphi^{-1})_{*} = \varphi_{*}(\varphi^{-1})_{*} .
\]

因此 $ \varphi_{*} $ 是双射，且 $ \varphi_{*}^{-1} = (\varphi^{-1})_{*} $。其余断言是注 4.1(b) 与 $ \mathbb{H}^{m} $ 的开子集上向量场的性质（见第 VIII.3 节）的简单推论。\nobreak\hspace{0.35em}■

\item

设 $ X_{0} $ 与 $ X_{1} $ 在 $ \mathbb{R}^{m} $ 中开，并设 $ \varphi \in \operatorname{Diff}(X_{0},X_{1}) $。又对 $ j = 0,1 $，用 $ \Theta_{j} : \mathcal{V}(X_{j}) \to \Omega^{1}(X_{j}) $ 表示注 VIII.3.3(g) 中定义的那个典范模同构。则

\[
(\varphi^{-1})^{*} \circ \Theta_{0} = \Theta_{1} \circ \varphi_{*} ,
\]

也就是说，下图交换。

\end{items}

\bookfigure[0.45]{figures/0d8eb2bc0c958241.png}

\begin{items}
\item[（d）]

（正则性）设 $ k \in \mathbb{N} $。对 $ \varphi \in \operatorname{Diff}^{k+1}(M,N) $ 与 $ 0 \le \ell \le k $，前推 $ \varphi_{*} $ 把 $ C^{\ell}(M) $ 映到 $ C^{\ell}(N) $，把 $ \mathcal{V}^{\ell}(M) $ 映到 $ \mathcal{V}^{\ell}(N) $，但这一陈述（去掉不等式）对 $ \ell = k+1 $ 不成立。
若 $ M $ 是 $ C^{k+1} $ 流形，则 $ C^{\ell}(M) $-模 $ C^{\ell}(M) $ 与 $ \mathcal{V}^{\ell}(M) $ 对 $ 0 \le \ell \le k $ 有定义；然而\footnote{平凡情形除外。}模 $ C^{k+1}(M) $ 与 $ \mathcal{V}^{k+1}(M) $ 则没有定义。

\proofhead
这是因为切映射“失去”一阶导数。\nobreak\hspace{0.35em}■
\end{items}

\end{remark}

\needspace{3\baselineskip}\subsection*{局部基表示}

设 $ (\varphi,U) $ 是 $ M $ 在 $ p $ 附近的坐标卡。则我们用

\[
\partial_{j}|_{p} = \left.\frac{\partial}{\partial x^{j}}\right|_{p} \in T_{p}M \quad \text{对 } 1 \le j \le m
\]

表示 $ T_{p}M $ 的对应于局部坐标 $ \varphi = (x^{1},\ldots,x^{m}) $ 的基向量。换言之，$ \partial_{j}|_{p} $ 是坐标路径 $ t \mapsto \varphi^{-1}(\varphi(p) + te_{j}) $ 在点\footnote{若 $ p $ 在 $ M $ 的内部。若 $ p $ 是边界点，我们必须把 $ \varphi $ 取作 $ M $ 在 $ p $ 附近的 $ \mathbb{R}^{\overline{m}} $ 的子流形坐标卡。} $ p $ 处的切向量，即

\begin{equation}
\partial_{j}|_{p} := (T_{p}\varphi)^{-1}(\varphi(p), e_{j}) , \quad \text{对 } 1 \le j \le m ,\tag{4.2}
\end{equation}

其中 $ (e_{1},\ldots,e_{m}) $ 是 $ \mathbb{R}^{m} $ 的典范基。

\bookfigure[0.70]{figures/c8a0db9f82bab82c.png}

\begin{remark}

\begin{items}
\item

设 $ i_{M} : M \hookrightarrow \mathbb{R}^{\overline{m}} $，并设 $ g_{\varphi} := i_{M} \circ \varphi^{-1} : \varphi(U) \to \mathbb{R}^{\overline{m}} $ 是属于 $ \varphi $ 的参数化。则

\[
(T_{p}i_{M})\partial_{j}|_{p} = (p, \partial_{j}g_{\varphi}(\varphi(p))) \in T_{p}\mathbb{R}^{\overline{m}} \quad \text{对 } 1 \le j \le m .
\]

这意味着，若我们在典范单射

\[
T_{p}i_{M} : T_{p}M \to T_{p}\mathbb{R}^{\overline{m}} ,
\]

下把 $ \partial_{j}|_{p} \in T_{p}M $ 与它在 $ T_{p}\mathbb{R}^{\overline{m}} $ 中的像等同，则我们得到 $ \partial_{j}|_{p} = (p, \partial_{j}g_{\varphi}(\varphi(p))) $。

\proofhead
由例 VII.10.9(b) 与注 1.14(c)，我们得到
\[
T_{\varphi(p)}g_{\varphi} = T_{\varphi(p)}(i_{M} \circ \varphi^{-1}) = T_{p}i_{M} \circ T_{\varphi(p)}(\varphi^{-1}) = T_{p}i_{M} \circ (T_{p}\varphi)^{-1} .
\]

于是由 (4.2) 推出

\[
(T_{p}i_{M})\partial_{j}|_{p} = (T_{\varphi(p)}g_{\varphi})(\varphi(p), e_{j}) = (p, \partial g_{\varphi}(\varphi(p))e_{j}) = (p, \partial_{j}g_{\varphi}(\varphi(p))) .\quad\blacksquare
\]
\item

映射

\[
\partial_{j} = \frac{\partial}{\partial x^{j}} : U \to TU , \quad p \mapsto \partial_{j}|_{p} \quad \text{对 } 1 \le j \le m
\]

是 $ U $ 上的光滑向量场。

\proofhead
这是显然的，因为对 $ p \in U $ 与 $ 1 \le j \le m $，有
\[
(\varphi_{*}\partial_{j})(\varphi(p)) = (T_{p}\varphi)(T_{p}\varphi)^{-1}(\varphi(p), e_{j}) = (\varphi(p), e_{j}) .\quad\blacksquare
\]
\item

对 $ p \in U $，我们有 $ T_{p}M $ 的基 $ (\partial_{1}|_{p},\ldots,\partial_{m}|_{p}) $ 以及 $ \mathcal{V}(U) $ 的模基 $ (\partial_{1},\ldots,\partial_{m}) $。$ U $ 上的向量场 $ v $ 属于 $ \mathcal{V}(U) $ 当且仅当基表示

\[
v = \sum_{j=1}^{m} v^{j}\partial_{j}
\]

的系数 $ v^{j} $ 全都属于 $ \mathcal{E}(U) $。

\proofhead
第一个断言由注 VII.10.5 以及边界点处切空间的定义推出。第二个断言是
\[
\varphi_{\ast}v = \varphi_{\ast}\left(\sum_{j=1}^{m} v^{j}\partial_{j}\right) = \sum_{j=1}^{m}(\varphi_{\ast}v^{j})\varphi_{\ast}\partial_{j},
\]

以及 (b)、注 VIII.3.3(c) 的推论。\nobreak\hspace{0.35em}■

\item

（正则性）设 $ k \in \mathbb{N} $，并设 $ M $ 是 $ C^{k+1} $ 流形。此时 $ (\partial_{1},\ldots,\partial_{m}) $ 是 $ \mathcal{V}^{k}(U) $ 的 $ C^{k}(U) $ 模基。$ U $ 上的向量场 $ v $ 属于 $ \mathcal{V}^{k}(U) $ 当且仅当它关于该基表示的系数属于 $ C^{k}(U) $。\nobreak\hspace{0.35em}■

\end{items}

\end{remark}

\needspace{3\baselineskip}\subsection*{微分形式}

为推广 $ \mathbb{R}^{m} $ 的开子集 $ X $ 的余切空间 $ T_{p}^{\ast}X $ 与余切丛 $ T^{\ast}X $，我们现在定义 $ M $ 在点 $ p $ 处的\textbf{余切空间}

\[
T_{p}^{\ast}M := (T_{p}M)^{\ast} = \mathcal{L}(T_{p}M,\mathbb{R}).
\]

我们定义 $ M $ 的\textbf{余切丛}

\[
T^{\ast}M := \bigcup_{p\in M} T_{p}^{\ast}M.
\]

我们记

\[
\langle\cdot,\cdot\rangle_{p} : T_{p}^{\ast}M\times T_{p}M\to\mathbb{R} \qquad \text{对 } p\in M
\]

为对偶配对\footnote{参见第 VIII.3 节。}，并同样把

\[
\langle\cdot,\cdot\rangle : T^{\ast}M\times TM\to\mathcal{E}(M), \qquad (\alpha,v)\mapsto[p\mapsto\langle\alpha(p),v(p)\rangle_{p}]
\]

称为\textbf{对偶配对}。

由于 $ T_{p}M $ 是 $ m $ 维向量空间，$ T_{p}^{\ast}M $ 也是。因此，对 $ r\in\mathbb{N} $ 与 $ p\in M $，$ T_{p}^{\ast}M $ 的 $ r $ 重外积 $ \bigwedge^{r}T_{p}^{\ast}M $ 与 格拉斯曼代数

\[
\bigwedge T_{p}^{\ast}M = \bigoplus_{r\ge 0}\bigwedge^{r}T_{p}^{\ast}M
\]

都得以定义。

为扩展上一节引入的概念，我们定义 $ M $ 上的\textbf{交错 r 形式丛}

\[
\bigwedge^{r}T^{\ast}M := \bigcup_{p\in M}\bigwedge^{r}T_{p}^{\ast}M.
\]

我们定义 $ M $ 的\textbf{格拉斯曼丛}

\[
\bigwedge T^{\ast}M := \bigcup_{p\in M}\bigwedge T_{p}^{\ast}M.
\]

于是 $ M $ 上的\textbf{微分形式}是一个映射

\[
\alpha : M\to\bigwedge T^{\ast}M \qquad \text{其中 } \alpha(p)\in\bigwedge T_{p}^{\ast}M \text{ 对 } p\in M,
\]

也就是说，它是格拉斯曼丛的一个截面。若 $ \alpha(M)\subset\bigwedge^{r}T^{\ast}M $，则它有\textbf{次数 r}（或称它是一个 \textbf{r 形式}）。有时我们称 1 形式为\textbf{普法夫形式}。

若 $ \alpha $ 与 $ \beta $ 是 $ M $ 上的微分形式，则\textbf{和} $ \alpha+\beta $ 与\textbf{外积}\footnote{符号 $ \wedge $ 也称为楔积。} $ \alpha\wedge\beta $ 逐点定义：

\[
(\alpha+\beta)(p) := \alpha(p)+\beta(p) \quad \text{且 } \alpha\wedge\beta(p) := \alpha(p)\wedge\beta(p) \qquad \text{对 } p\in M.
\]

若 $ \alpha $ 是 $ M $ 上的 r 形式，则它对向量场的作用也逐点定义：

\[
\alpha(v_{1},\ldots,v_{r})(p) := \alpha(p)\big(v_{1}(p),\ldots,v_{r}(p)\big) \qquad \text{对 } p\in M \text{ 且 } v_{1},\ldots,v_{r}\in\mathcal{V}(M).
\]

最后，设 $ \varphi\in C^{1}(M,N) $，$ \beta $ 是 $ N $ 上的微分形式。则 $ \beta $ 沿 $ \varphi $ 的\textbf{拉回}同样逐点定义：

\[
\varphi^{\ast}\beta(p) := (T_{p}\varphi)^{\ast}\beta\big(\varphi(p)\big) \qquad \text{对 } p\in M.
\]

显然 $ \varphi^{\ast}\beta $ 是 $ M $ 上的微分形式，即 $ \beta $ 沿 $ \varphi $ 的拉回。若 $ \varphi $ 是从 $ M $ 到 $ N $ 的 $ C^{1} $ 微分同胚，则

\[
\varphi_{\ast}\alpha := (\varphi^{-1})^{\ast}\alpha
\]

是 $ M $ 上微分形式 $ \alpha $ 的\textbf{前推}。

设 $ k\in\mathbb{N}\cup\{\infty\} $。$ M $ 上的微分形式 $ \alpha $ 属于\textbf{类 $ C^{k} $}（或称它为 $ k $ 次\textbf{连续可微}，\footnote{当然，我们把类 $ C^{0} $ 的微分形式说成连续的。} 或在 $ k=\infty $ 的情形称为\textbf{光滑}），如果 $ M $ 的每一点都有一个坐标卡 $ (\varphi,U) $，使得 $ \varphi_{\ast}\alpha $ 是 $ \varphi(U) $ 上类 $ C^{k} $ 的微分形式。我们把 $ M $ 上所有类 $ C^{k} $ 的 r 形式之集记为

\[
\Omega_{(k)}^{r}(M),
\]

而

\[
\Omega^{r}(M) := \Omega_{(\infty)}^{r}(M)
\]

是 $ M $ 上所有光滑 r 形式之集。最后，

\[
\Omega(M) := \Omega_{(\infty)}(M)
\]

是 $ M $ 上所有光滑微分形式之集。

仿照我们对向量场的处理，我们将一般地把注意力限制在光滑微分形式的研究上。我们把它留给你去证明：我们所证明的关于光滑形式的所有结论，只要按情形对 $ k $ 加以限制，对类 $ C^{k} $ 的形式也类似地成立。

\begin{remark}

\begin{items}
\item

上述微分形式的可微性概念与坐标无关。
若 $ \alpha $ 是 $ M $ 上类 $ C^{k} $ 的 r 形式，$ (\psi,V) $ 是 $ M $ 上的坐标卡，则 $ \psi_{\ast}\alpha $ 是 $ \psi(V) $ 上类 $ C^{k} $ 的 r 形式。

\proofhead
$ M $ 的每一点 $ p $ 都有一个坐标卡 $ (\varphi,U) $ 使 $ \varphi_{\ast}\alpha\in\Omega_{(k)}^{r}(\varphi(U)) $。由注 3.3(a) 以及前推的逐点定义可得
\[
\psi_{\ast}\alpha = (\psi\circ\varphi^{-1})_{\ast}\varphi_{\ast}\alpha.
\]

至此，仿照注 4.2(a) 的证明即可得到结论。\nobreak\hspace{0.35em}■

\item

$ \Omega(M) $ 与 $ \Omega^{r}(M) $ 是 $ \mathcal{E}(M) $ 模，因而特别是 $ \mathbb{R} $ 向量空间。又

\[
\Omega(M) = \bigoplus_{r\ge 0}\Omega^{r}(M).
\]

外积是 $ \mathbb{R} $ 双线性的、结合的与分次反交换的，也就是说，它满足以下规则：

\begin{subitems}
\item

映射

\[
\Omega^{r}(M)\times\Omega^{s}(M)\to\Omega^{r+s}(M), \quad (\alpha,\beta)\mapsto\alpha\wedge\beta
\]

是良定义的且 $ \mathbb{R} $ 双线性的。

\item

对 $ \alpha,\beta,\gamma\in\Omega(M) $ 有 $ \alpha\wedge(\beta\wedge\gamma) = (\alpha\wedge\beta)\wedge\gamma $。

\item

对 $ \alpha\in\Omega^{r}(M) $ 与 $ \beta\in\Omega^{s}(M) $ 有 $ \alpha\wedge\beta = (-1)^{rs}\beta\wedge\alpha $。

\proofhead
这由光滑性的定义、$ \wedge $ 的逐点定义以及定理 2.7 推出。\nobreak\hspace{0.35em}■
\end{subitems}

\item

每个 $ \alpha\in\Omega^{r}(M) $ 都是 $ \mathcal{V}(M) $ 上的交错 r 形式。

\item

$ \Omega^{0}(M) = \mathcal{E}(M) $，且对 $ r>m $ 有 $ \Omega^{r}(M) = \{0\} $。

\item

对 $ h\in C^{\infty}(M,N) $，拉回 $ h^{\ast} : \Omega(N)\to\Omega(M) $ 是代数同态，即

\[
h^{\ast}(\alpha+\beta) = h^{\ast}\alpha+h^{\ast}\beta, \quad h^{\ast}(\alpha\wedge\beta) = h^{\ast}\alpha\wedge h^{\ast}\beta
\]

对 $ \alpha,\beta\in\Omega(N) $。若 $ \alpha\in\Omega^{r}(N) $，则 $ h^{\ast}\alpha $ 属于 $ \Omega^{r}(M) $。又

\[
(k\circ h)^{\ast} = h^{\ast}\circ k^{\ast} \quad \text{且 } (\operatorname{id}_{M})^{\ast} = \operatorname{id}_{\Omega(M)}.
\]

若 $ h $ 是微分同胚，则 $ h^{\ast} $ 是双射，且 $ (h^{\ast})^{-1} = (h^{-1})^{\ast} = h_{\ast} $。

\proofhead
简单的验证留给你。\nobreak\hspace{0.35em}■
\item

设 $ M $ 是 $ N $ 的子流形，$ i : M\hookrightarrow N $ 是自然嵌入。\footnote{在这种情形下，我们总假设 $ N $ 没有边界。} 则对 $ \alpha\in\Omega^{r}(N) $，

\[
\alpha\,|\,M := i^{\ast}\alpha\in\Omega^{r}(M)
\]

是 $ \alpha $ 到 $ M $ 的\textbf{限制}。\footnote{参见例 3.4(j)。} 设 $ p\in M $。由于切空间 $ T_{p}M $ 可视为 $ T_{p}N $ 的向量子空间，我们有 $ (\alpha\,|\,M)(p) = \alpha(p)\,|\,(T_{p}M)^{r} $。\nobreak\hspace{0.35em}■

\end{items}

\end{remark}

\needspace{3\baselineskip}\subsection*{局部表示}

设 $ f\in C^{1}(M) := C^{1}(M,\mathbb{R}) $。如同第 VII.10 节，我们通过

\[
df(p) := \operatorname{pr}\circ T_{p}f \qquad \text{对 } p\in M,
\]

定义 $ f $ 的\textbf{微分} $ df $，其中

\[
\operatorname{pr} := \operatorname{pr}_{2} : T_{f(p)}\mathbb{R} = \{f(p)\}\times\mathbb{R}\to\mathbb{R}
\]

是标准投影。

设 $ (\varphi,U) $ 是 $ p\in M $ 处的坐标卡。则由 $ df(p) $ 与 $ \partial_{j}|_{p} $ 的定义以及注 VII.10.9(b) 与 1.14(c) 的链式法则可得，对 $ 1\le j\le m $，

\begin{equation*}
\begin{aligned}
\langle df(p),\partial_{j}|_{p}\rangle_{p} &= \langle df(p),(T_{\varphi(p)}\varphi^{-1})(\varphi(p),e_{j})\rangle_{p}\\ &= \operatorname{pr}\circ T_{p}f\circ T_{\varphi(p)}\varphi^{-1}(\varphi(p),e_{j})\\ &= \operatorname{pr}\circ T_{\varphi(p)}(f\circ\varphi^{-1})(\varphi(p),e_{j})\\ &= \partial(f\circ\varphi^{-1})(\varphi(p))e_{j}\\ &= \partial_{j}(f\circ\varphi^{-1})(\varphi(p))\\ &= \partial_{j}(\varphi_{\ast}f)(\varphi(p))
\end{aligned}
\end{equation*}

以缩写

\begin{equation}
\partial_{j}f(p) := \frac{\partial f}{\partial x^{j}}(p) := \partial_{j}(f\circ\varphi^{-1})(\varphi(p)) = \partial_{j}(\varphi_{\ast}f)(\varphi(p))\tag{4.3}
\end{equation}

表示，对 $ 1\le j\le m $ 与 $ p\in U $，我们于是有

\begin{equation}
\langle df(p),\partial_{j}|_{p}\rangle_{p} = \partial_{j}f(p) \quad \text{对 } 1\le j\le m, \quad p\in U.\tag{4.4}
\end{equation}

因此

\begin{equation}
\langle df,\partial_{j}\rangle = \partial_{j}f \quad \text{对 } 1\le j\le m.\tag{4.5}
\end{equation}

注意，$ M $ 上通常的偏导数 $ \partial_{j}f $（在注 VII.2.7(a) 的意义下）在 $ M $ 不是“平坦的”，即不是 $ \mathbb{R}^{m} $ 的开子集时没有定义。由于流形上函数的导数只能通过局部表示来定义，(4.5) 中的 $ \partial_{j}f $ 除非解释为由 $ \varphi $ “下推”到参数区域 $ \varphi(U) $ 上的函数的偏导数，即 (4.3) 中出现的 $ \varphi_{\ast}f $ 的偏导数，否则没有意义。这就排除了实践中任何误解。记号 $ \partial f/\partial x^{j} $ 的优点是它给出了 $ f $ 被局部写出的“坐标的名称” $ (x^{1},\ldots,x^{m}) = \varphi $。

在第 VII.2 节中，对 $ \mathbb{R}^{m} $ 的开子集的情形，我们把偏导数 $ \partial_{j}f(p) $ 定义为第 $ j $ 个坐标单位向量 $ e_{j} $ 在（全）导数 $ \partial f(p) $（即 $ f $ 在 $ p $ 处的线性化）下的像。由于 $ df(p) $ 正是切映射 $ T_{p}f $ 的切部分，从而是“$ f $ 在点 $ p $ 处的线性化”，并且 $ \partial_{j}|_{p} $ 是 $ T_{p}M $ 的第 $ j $ 个坐标基向量，(4.4) 表明 $ \partial_{j}f(p) $ 是这些坐标向量在 $ f $ 的切映射下的像的切部分。因此 (4.3) 确实是把偏导数概念推广到流形上函数的正确方式。

最后，显然当 $ M $ 在 $ \mathbb{R}^{m} $ 中开且 $ \varphi $ 表示平凡坐标卡 $ \operatorname{id}_{M} $ 时，(4.3) 与经典的偏导数一致。

\begin{remark}

设 $ (\varphi,U) $ 是 $ M $ 的坐标卡。

\begin{items}
\item

对 $ f\in\mathcal{E}(M) = \Omega^{0}(M) $，微分 $ df $ 属于 $ \Omega^{1}(M) $。映射

\[
d : \Omega^{0}(M)\to\Omega^{1}(M), \quad f\mapsto df
\]

是 $ \mathbb{R} $ 线性的。

\item

设 $ (x^{1},\ldots,x^{m}) = \varphi $ 是 $ \varphi $ 在 $ U $ 上诱导的局部坐标，于是

\[
x^{j} := \operatorname{pr}_{j}\circ\varphi\in\mathcal{E}(U) \quad \text{对 } 1\le j\le m,
\]

其中 $ \operatorname{pr}_{j} : \mathbb{R}^{m}\to\mathbb{R} $ 是标准投影。则 $ \Omega^{1}(U) $ 是 $ m $ 维自由 $ \mathcal{E}(U) $ 模，且 $ (dx^{1},\ldots,dx^{m}) $ 是一组模基，满足

\begin{equation}
\left\langle dx^{j},\frac{\partial}{\partial x^{k}}\right\rangle = \delta_{k}^{j} \quad \text{对 } 1\le j,k\le m\tag{4.6}
\end{equation}

并且它是 $ \mathcal{V}(U) $ 的基 $ (\partial/\partial x^{1},\ldots,\partial/\partial x^{m}) $ 的对偶基。基表示

\begin{equation}
v = \sum_{j=1}^{m}v^{j}\frac{\partial}{\partial x^{j}}\in\mathcal{V}(U) \quad \text{且 } \alpha = \sum_{j=1}^{m}a_{j}\,dx^{j}\in\Omega^{1}(U)\tag{4.7}
\end{equation}

要求有关系

\begin{equation}
v^{j} = \langle dx^{j},v\rangle\in\mathcal{E}(U) \quad \text{且 } a_{j} = \left\langle\alpha,\frac{\partial}{\partial x^{j}}\right\rangle\in\mathcal{E}(U)\tag{4.8}
\end{equation}

对 $ 1\le j\le m $。特别地，对 $ f\in\mathcal{E}(U) $ 有

\[
df = \sum_{j=1}^{m}\frac{\partial f}{\partial x^{j}}dx^{j}\in\Omega^{1}(U).
\]

\proofhead
(4.3) 与 (4.5) 蕴含
\[
\langle dx^{j},\partial_{k}\rangle = \partial_{k}x^{j} = \varphi^{\ast}\partial_{k}(\varphi_{\ast}x^{j}) = \varphi^{\ast}\partial_{k}[(\operatorname{pr}_{j}\circ\varphi)\circ\varphi^{-1}] = \varphi^{\ast}\partial_{k}\operatorname{pr}_{j} = \delta_{k}^{j}
\]

从而得到 (4.6)。对具有 (4.7) 中表示的 $ v $，我们得到

\begin{equation}
\langle dx^{j}(p),v(p)\rangle_{p} = \sum_{k=1}^{m}v^{k}(p)\left\langle dx^{j}(p),\frac{\partial}{\partial x^{k}}\Big|_{p}\right\rangle_{p} = \sum_{k=1}^{m}v^{k}(p)\delta_{k}^{j} = v^{j}(p)\tag{4.9}
\end{equation}

对 $ p\in U $ 与 $ 1\le j\le m $，因为 $ dx^{j}(p) $ 是 $ T_{p}M = T_{p}U $ 上的线性形式。因此 (4.7) 的第一部分与注 4.3(c) 蕴含 (4.8) 的第一个断言。
对于 $ dx^{j} $ 经 $ \varphi $ 的前推，我们应用注 VII.10.9(b) 与 1.14(c) 以及 (4.2) 与 (4.6) 得到

\begin{equation}
\begin{aligned}
\langle(\varphi_{\ast}dx^{j})(\varphi(p)),(\varphi(p),e_{k})\rangle_{\varphi(p)} &= \langle dx^{j}(p),(T_{\varphi(p)}\varphi^{-1})(\varphi(p),e_{k})\rangle_{p}\\ &= \langle dx^{j}(p),(T_{p}\varphi)^{-1}(\varphi(p),e_{k})\rangle_{p}\\ &= \left\langle dx^{j}(p),\frac{\partial}{\partial x^{k}}\Big|_{p}\right\rangle_{p} = \delta_{k}^{j}.
\end{aligned}\tag{4.10}
\end{equation}

这表明 $ (\varphi_{\ast}dx^{1},\ldots,\varphi_{\ast}dx^{m}) $ 在每一点 $ \varphi(p)\in\varphi(U) $ 处是 $ T_{\varphi(p)}\varphi(U) $ 的标准基的对偶基。特别地，$ \varphi_{\ast}dx^{j} $ 的余向量部分在 $ \varphi(U) $ 上是常值的。
注 4.4(e) 保证 $ \varphi_{\ast} $ 是从 $ \Omega^{1}(U) $ 到 $ \Omega^{1}(\varphi(U)) $ 的向量空间同构。由此，(4.10) 以及命题 2.3，我们断定每个 $ \alpha\in\Omega^{1}(U) $ 都有形如 (4.7) 的表示，其系数是 $ U $ 上的实值函数 $ a_{j} $。由于

\begin{equation}
\varphi_{\ast}\alpha = \sum_{j=1}^{m}(\varphi_{\ast}a_{j})\,\varphi_{\ast}dx^{j}\tag{4.11}
\end{equation}

并且由于 1 形式 $ \varphi_{\ast}dx^{j} $ 的余向量部分在 $ \varphi(U) $ 上是常值的，由注 3.1(e) 可知 $ \alpha $ 属于 $ \Omega^{1}(U) $ 当且仅当 $ a_{j}\in\mathcal{E}(U) $ 对 $ 1\le j\le m $。最后，$ a_{j} = \langle\alpha,\partial_{j}\rangle $ 由类似于 (4.9) 的计算推出。\nobreak\hspace{0.35em}■

\item

对 $ r\in\mathbb{N} $，$ \Omega^{r}(U) $ 是 $ \binom{m}{r} $ 维自由 $ \mathcal{E}(U) $ 模，且

\begin{equation}
\left\{ dx^{(j)} = dx^{j_{1}}\wedge\cdots\wedge dx^{j_{r}} ;\ (j) = (j_{1},\ldots,j_{r})\in\mathbb{J}_{r} \right\}\tag{4.12}
\end{equation}

为一组基。$ U $ 上的 r 形式 $ \alpha $ 有唯一确定的\textbf{局部坐标下的基表示}

\begin{equation}
\alpha = \sum_{(j)\in\mathbb{J}_{r}}a_{(j)}\,dx^{(j)}\tag{4.13}
\end{equation}

其系数为

\begin{equation}
a_{(j)} = \alpha\left(\frac{\partial}{\partial x^{j_{1}}},\ldots,\frac{\partial}{\partial x^{j_{r}}}\right) \quad \text{对 } (j)\in\mathbb{J}_{r}.\tag{4.14}
\end{equation}

若 $ k\in\mathbb{N}\cup\{\infty\} $，则 $ \alpha $ 在 $ U $ 上属于类 $ C^{k} $ 当且仅当 $ a_{(j)}\in C^{k}(U) $ 对 $ (j)\in\mathbb{J}_{r} $。

\proofhead
由 (4.10) 以及注 4.4(e) 中给出的拉回性质 $ (\varphi^{-1})^{\ast} = \varphi_{\ast} $ 可得
\begin{equation}
\varphi_{\ast}dx^{(j)} = \varepsilon^{(j)} \quad \text{对 } (j)\in\mathbb{J}_{r},\tag{4.15}
\end{equation}

其中 $ (\varepsilon^{1},\ldots,\varepsilon^{m}) $ 表示 $ T_{\varphi(p)}\varphi(U) $ 的标准基的对偶基，$ p\in U $。因为 $ \varphi_{\ast} $ 是从 $ \Omega^{r}(U) $ 到 $ \Omega^{r}(\varphi(U)) $ 的向量空间同构，由命题 2.3 与 (4.2) 可知 $ U $ 上每个 r 形式 $ \alpha $ 都有形如 (4.13) 的唯一表示，其系数由 (4.14) 给出。因为

\[
\varphi_{\ast}\alpha = \sum_{(j)\in\mathbb{J}_{r}}(\varphi_{\ast}a_{(j)})\varphi_{\ast}dx^{(j)}
\]

并且由 (4.15)，类 $ C^{k} $ 的 r 形式的可微性定义蕴含 $ \alpha $ 属于类 $ C^{k} $ 当且仅当 $ a_{(j)} $ 落在 $ C^{k}(U) $ 中。\nobreak\hspace{0.35em}■

\item

注意我们只证明了 $ \mathcal{V}(U) $ 与 $ \Omega(U) $ 是自由模，而对 $ \mathcal{V}(M) $ 与 $ \Omega(M) $ 没有作这样的断言。事实上，相应的断言在整体情形下一般并不成立，即对不能由单个坐标卡描述的流形不成立。例如，已知\footnote{由 Bott、Kervaire 与 Milnor 的工作。} $ n $ 维球面不支持 $ n $ 个（非平凡的）线性无关向量场（即 $ \mathcal{V}(S^{n}) $ 不是 $ n $ 维自由模），除非 $ n = 0,1,3 $ 或 $ 7 $。

\item

（正则性）若 $ k\in\mathbb{N} $，则当 $ M $ 是 $ C^{k+1} $ 流形时，(c) 的结论仍然成立。\nobreak\hspace{0.35em}■

\end{items}

属于坐标卡 $ \varphi $ 的 $ U $ 上的局部坐标 $ x^{1},\ldots,x^{m} $ 是 $ U $ 上的光滑函数；即它们是映射 $ \operatorname{pr}_{j}\circ\varphi\in\mathcal{E}(U) $，$ 1\le j\le m $。另一方面，我们也用 $ (x^{1},\ldots,x^{m}) $ 表示 $ \varphi(U) $ 的一般点，即 $ \mathbb{R}^{m} $ 的坐标也称为 $ x^{1},\ldots,x^{m} $。对两个不同对象使用同一记号是故意的。只要从上下文可以清楚哪个解释是正确的，这就大大简化了（局部）坐标下的计算。例如，表达式

\begin{equation}
\alpha = \sum_{(j)\in\mathbb{J}_{r}}a_{(j)}\,dx^{(j)}\tag{4.16}
\end{equation}

（若没有其他说明，这在实践中是通常的）有两种含义。第一，我们可以像前几节那样，把 (4.16) 视为 $ \overline{\mathbb{H}}^{m} $ 的开子集 $ X = \varphi(U) $ 上 r 形式的基表示。或者，我们可以把 (4.16) 解释为 $ U $ 上 r 形式关于相应坐标卡中局部坐标的基表示。这就是我们在此采取的立场。在第一种情形，$ a_{(j)} $ 是 $ X $ 上的函数，而 $ dx^{(j)} $ 是 $ \mathbb{R}^{m} $ 的常值基形式。在第二种情形，$ a_{(j)} $ 是 $ U\subset M $ 上的函数，而 $ dx^{(j)} $ 是“生活”在 $ U $ 上的依赖于位置的 r 形式。由于 (4.15)，为了从第二种解释过渡到第一种，我们必须用 $ \varphi $ 把系数函数 $ a_{(j)} = a_{(j)}(p) $ “下推”到参数区域。也就是说，$ a_{(j)} $ 必须解释为 $ \varphi_{\ast}a_{(j)} = a_{(j)}\circ\varphi^{-1} $，并且我们必须认为 $ a_{(j)} = a_{(j)}(x) $，$ x\in X $。

\end{remark}

\begin{example}

\begin{items}
\item

记 $ \mathbb{R}^{m+1} $ 中 $ m $ 维球面 $ S^{m} $ 的上、下半球面分别为 $ S_{+}^{m} $ 与 $ S_{-}^{m} $。也就是说，设

\[
S_{\pm}^{m} := \{x\in\mathbb{R}^{m+1} ;\ |x| = 1,\ \pm x^{m+1} > 0\}.
\]

又设

\[
\varphi_{\pm} : S_{\pm}^{m}\to\mathbb{B}^{m}, \quad x\mapsto x' := (x^{1},\ldots,x^{m})
\]

是 $ \mathbb{B}^{m} = \mathbb{B}^{m}\times\{0\} $ 到与 $ x^{m+1} $ 轴正交的超平面上的投影。则 $ (\varphi_{+},S_{+}^{m}) $ 与 $ (\varphi_{-},S_{-}^{m}) $ 是 $ S^{m} $ 的坐标卡。对

\[
\alpha := \sum_{j=1}^{m+1}(-1)^{j-1}x^{j}\,dx^{1}\wedge\cdots\wedge\widehat{dx^{j}}\wedge\cdots\wedge dx^{m+1}\in\Omega^{m}(\mathbb{R}^{m+1}),
\]

在由 $ \varphi_{\pm} $ 诱导的局部坐标下，到 $ S^{m} $ 的限制读作

\[
\alpha\,|\,S_{\pm}^{m} = \pm\frac{(-1)^{m}}{\sqrt{1-|x'|^{2}}}\,dx^{1}\wedge\cdots\wedge dx^{m}.
\]

\proofhead
设 $ g_{\pm}(x') := (x',\pm\sqrt{1-|x'|^{2}}) $，$ x'\in\mathbb{B}^{m} $。则 $ g_{\pm} $ 是光滑的，并且是属于 $ \varphi_{\pm} $ 的、把半球面 $ S_{\pm}^{m} $ 表示为 $ \mathbb{B}^{m} $ 上图像的参数化。又 $ g_{\pm} = i\circ\varphi_{\pm}^{-1} $，其中 $ i : S^{m}\hookrightarrow\mathbb{R}^{m+1} $。因此 $ (\varphi_{\pm},S_{\pm}^{m}) $ 是 $ S^{m} $ 的坐标卡。对这些坐标卡，我们得到
\begin{equation*}
\begin{aligned}
(\varphi_{\pm})_{\ast}(\alpha\,|\,S_{\pm}^{m}) &= (\varphi_{\pm}^{-1})^{\ast}\circ i^{\ast}\alpha = g_{\pm}^{\ast}\alpha\\ &= \sum_{j=1}^{m+1}(-1)^{j-1}g_{\pm}^{j}\,dg_{\pm}^{1}\wedge\cdots\wedge\widehat{dg_{\pm}^{j}}\wedge\cdots\wedge dg_{\pm}^{m+1}\\ &= \sum_{j=1}^{m}(-1)^{j-1}x^{j}\,dx^{1}\wedge\cdots\wedge\widehat{dx^{j}}\wedge\cdots\wedge dx^{m}\wedge\sum_{k=1}^{m}\frac{\mp x^{k}\,dx^{k}}{\sqrt{1-|x'|^{2}}}\\ &\quad \pm(-1)^{m}\sqrt{1-|x'|^{2}}\,dx^{1}\wedge\cdots\wedge dx^{m}\\ &= \pm\frac{(-1)^{m}}{\sqrt{1-|x'|^{2}}}\left[-\sum_{j=1}^{m}(-1)^{m+j-1+m-j}(x^{j})^{2}+1-|x'|^{2}\right]dx^{1}\wedge\cdots\wedge dx^{m}.
\end{aligned}
\end{equation*}

由方括号中的表达式化为 1 即得结论。\nobreak\hspace{0.35em}■

\item

令 $ \omega_{S^{1}} := (x\,dy - y\,dx)|S^{1} $，并令

\[
g_{1}: (0,2\pi) \to S^{1}\setminus\{(1,0)\}, \qquad t \mapsto (\cos t, \sin t)
\]

为 $ S^{1}\setminus\{(1,0)\} $ 的一个参数化。那么关于由坐标卡 $ (\varphi, U) $（其中 $ \varphi := g_{1}^{-1} $，$ U := S^{1}\setminus\{(1,0)\} $）诱导的局部坐标，我们有

\[
\omega_{S^{1}}|U = dt .
\]

\proofhead
这由 $ \varphi_{\ast}\omega_{S^{1}} = (g_{1}^{1}\dot{g}_{1}^{2} - g_{1}^{2}\dot{g}_{1}^{1})\,dt $ 推出。\nobreak\hspace{0.35em}■
\item

令 $ U := S^{2}\setminus H_{3} $ 为 2 球面 $ S^{2} $ 减去它与半平面 $ H_{3} := \mathbb{R}^{+}\times\{0\}\times\mathbb{R} $ 相交的半圆。\footnote{见例 VII.9.11(b)。}又令

\[
(0,2\pi)\times(0,\pi) \to U, \qquad (\varphi,\vartheta) \mapsto (\cos\varphi\sin\vartheta, \sin\varphi\sin\vartheta, \cos\vartheta)
\]

为 $ U $ 用球坐标给出的参数化。最后，令

\[
\alpha := x\,dy\wedge dz + y\,dz\wedge dx + z\,dx\wedge dy \in \Omega^{2}(\mathbb{R}^{3}) .
\]

则形式 $ \omega_{S^{2}} := \alpha|S^{2} \in \Omega^{2}(S^{2}) $ 关于局部坐标 $ (\varphi,\vartheta) $ 有表示

\[
\omega_{S^{2}}|U = -\sin\vartheta\,d\varphi\wedge d\vartheta .
\]

\proofhead
经过简单计算，\footnote{注意（例如）$ \partial(x,y)/\partial(\varphi,\vartheta) $ 是由 (VII.9.3) 去掉最后一行所得的矩阵的行列式。}我们由例 3.4(d) 得到此式。\nobreak\hspace{0.35em}■
\end{items}

\end{example}

\needspace{3\baselineskip}\subsection*{坐标变换}

为了高效地进行具体计算，选择最适合问题的坐标是重要的。例如，当我们要描述旋转对称的问题时，就使用极坐标，正如我们在第 X.8 节处理积分理论时所做的那样。

由于给定的问题通常已经在某个坐标系中描述，我们必须能够不费周折地换到另一个坐标系。这一背景为下文关于向量场与普法夫形式的变换定理作了铺垫。

设 $ (\varphi, U) $ 与 $ (\psi, V) $ 是 $ M $ 的坐标卡，且 $ U\cap V\neq\varnothing $。设 $ \varphi = (x^{1}, \ldots, x^{m}) $，$ \psi = (y^{1}, \ldots, y^{m}) $。在 $ U\cap V $ 上，我们可以把 $ y^{j} $ 看作局部坐标 $ x = (x^{1}, \ldots, x^{m}) $ 的函数；也可以把 $ x^{j} $ 看作 $ y = (y^{1}, \ldots, y^{m}) $ 的函数。这里通常而方便的做法是不引入新记号，而直接写 $ y = y(x) $ 与 $ x = x(y) $。显然映射 $ y(\cdot) $ 是从 $ U\cap V $ 到自身的微分同胚，即\textbf{坐标变换}，我们也把它记为 $ x\mapsto y $。其逆映射是 $ x(\cdot) $，即坐标变换 $ y\mapsto x $。然而，我们也可以把 $ x $［或 $ y $］看成 $ \mathbb{H}^{m} $ 中 $ X := \varphi(U\cap V) $［或 $ Y := \psi(U\cap V) $］里的一般点。这时坐标变换 $ x\mapsto y $ 不过是转移函数 $ \psi\circ\varphi^{-1}\in\mathrm{Diff}(X,Y) $。从上下文总能看清应在这两种解释中选哪一种。

在下面的公式中，我们把从上下文判断 $ x^{j} $ 表示自变量还是表示函数 $ x^{j}(\cdot) $ 这件事留给你。这种双重含义在实践中几乎不成为问题，它是有意采用的，因为这有助于把公式写成更直观、更易记的形式。

\begin{proposition}

对坐标变换 $ x\mapsto y $，我们有

\[
\frac{\partial}{\partial y^{j}} = \sum_{k=1}^{m}\frac{\partial x^{k}}{\partial y^{j}}\frac{\partial}{\partial x^{k}} \qquad \text{与} \qquad dy^{j} = \sum_{k=1}^{m}\frac{\partial y^{j}}{\partial x^{k}}\,dx^{k}
\]

对 $ 1\le j\le m $。

\proofhead
由注 4.5(c) 可得
\[
\frac{\partial}{\partial y^{j}} = \sum_{k=1}^{m} v_{j}^{k}\frac{\partial}{\partial x^{k}} \qquad \text{与} \qquad v_{j}^{k} = \left\langle dx^{k}, \frac{\partial}{\partial y^{j}} \right\rangle \qquad \text{对 } 1\le j, k\le m .
\]

由 $ x = f(y) $ 与 (4.5)，我们得到

\begin{equation}
\left\langle dx^{k}, \frac{\partial}{\partial y^{j}} \right\rangle = \frac{\partial x^{k}}{\partial y^{j}} \qquad \text{对 } 1\le j, k\le m ,\tag{4.17}
\end{equation}

这证明了第一个断言。

类似地，我们有

\[
dy^{j} = \sum_{k=1}^{m} a_{k}\,dx^{k} \qquad \text{与} \qquad a_{k} = \left\langle dy^{j}, \frac{\partial}{\partial x^{k}} \right\rangle = \frac{\partial y^{j}}{\partial x^{k}}
\]

对 $ 1\le j, k\le m $，这证明了第二个断言。\nobreak\hspace{0.35em}■

\end{proposition}

\begin{corollary}

\begin{items}
\item

坐标变换 $ x\mapsto y $ 的雅可比矩阵满足

\[
\left[\frac{\partial y^{j}}{\partial x^{k}}\right] = \left[\frac{\partial x^{j}}{\partial y^{k}}\right]^{-1} .
\]
\item
\[
dy^{1}\wedge\cdots\wedge dy^{m} = \frac{\partial(y^{1}, \ldots, y^{m})}{\partial(x^{1}, \ldots, x^{m})}\,dx^{1}\wedge\cdots\wedge dx^{m} .
\]

\proofhead
（a）因为
\[
y(\cdot) = \psi\circ\varphi^{-1}\in\mathrm{Diff}(X,Y) \qquad \text{与} \qquad y(\cdot)^{-1} = x(\cdot) = \varphi\circ\psi^{-1}\in\mathrm{Diff}(Y,X) ,
\]

所以断言是显然的。

\end{items}

（b）这可由例 3.4(c)、注 4.5(e) 之后的讨论以及下述事实推出：

\[
\frac{\partial(y^{1}, \ldots, y^{m})}{\partial(x^{1}, \ldots, x^{m})}
\]

是坐标变换 $ x\mapsto y $ 的雅可比行列式（见注 VII.7.9(a)）。\nobreak\hspace{0.35em}■

\end{corollary}

\begin{example}

\begin{items}
\item

（平面极坐标）利用极坐标变换

\[
V_{2} \to \mathbb{R}^{2}, \qquad (r,\varphi) \mapsto (x,y) := (r\cos\varphi, r\sin\varphi)
\]

其中 $ V_{2} := (0,\infty)\times(0,2\pi) $，我们有

\[
\frac{\partial}{\partial r} = \frac{\partial x}{\partial r}\frac{\partial}{\partial x} + \frac{\partial y}{\partial r}\frac{\partial}{\partial y} = \cos\varphi\,\frac{\partial}{\partial x} + \sin\varphi\,\frac{\partial}{\partial y}
\]

与

\[
\frac{\partial}{\partial\varphi} = \frac{\partial x}{\partial\varphi}\frac{\partial}{\partial x} + \frac{\partial y}{\partial\varphi}\frac{\partial}{\partial y} = -r\sin\varphi\,\frac{\partial}{\partial x} + r\cos\varphi\,\frac{\partial}{\partial y} .
\]
\item

（球坐标）令 $ V_{3} := (0,\infty)\times(0,2\pi)\times(0,\pi) $。利用球坐标变换

\[
V_{3} \to \mathbb{R}^{3}, \qquad (r,\varphi,\vartheta) \mapsto (x,y,z) = (r\cos\varphi\sin\vartheta, r\sin\varphi\sin\vartheta, r\cos\vartheta) ,
\]

我们得到

\begin{equation*}
\begin{aligned}
\frac{\partial}{\partial r} &= \cos\varphi\sin\vartheta\,\frac{\partial}{\partial x} + \sin\varphi\sin\vartheta\,\frac{\partial}{\partial y} + \cos\vartheta\,\frac{\partial}{\partial z} \\ \frac{\partial}{\partial\varphi} &= -r\sin\varphi\sin\vartheta\,\frac{\partial}{\partial x} + r\cos\varphi\sin\vartheta\,\frac{\partial}{\partial y} \\ \frac{\partial}{\partial\vartheta} &= r\cos\varphi\cos\vartheta\,\frac{\partial}{\partial x} + r\sin\varphi\cos\vartheta\,\frac{\partial}{\partial y} - r\sin\vartheta\,\frac{\partial}{\partial z} .
\end{aligned}
\end{equation*}
\item

（柱坐标）令 $ X := (0,\infty)\times(0,2\pi)\times\mathbb{R} $。对柱坐标变换

\[
X \to \mathbb{R}^{3}, \qquad (r,\varphi,\zeta) \mapsto (x,y,z) := (r\cos\varphi, r\sin\varphi, \zeta) ,
\]

我们得到

\[
\frac{\partial}{\partial r} = \cos\varphi\,\frac{\partial}{\partial x} + \sin\varphi\,\frac{\partial}{\partial y}, \quad \frac{\partial}{\partial\varphi} = -r\sin\varphi\,\frac{\partial}{\partial x} + r\cos\varphi\,\frac{\partial}{\partial y}, \quad \frac{\partial}{\partial\zeta} = \frac{\partial}{\partial z} .\quad\blacksquare
\]
\end{items}

\end{example}

\needspace{3\baselineskip}\subsection*{外微分}

下一个定理表明，外微分可以推广到在流形上整体地定义。

\begin{theorem}

恰有一个映射

\[
d: \Omega(M) \to \Omega(M) ,
\]

即\textbf{外}（或\textbf{嘉当}）微分，它具有以下四条性质：

\begin{subitems}
\item

$ d $ 是 $ \mathbb{R} $-线性的，并把 $ \Omega^{r}(M) $ 映入 $ \Omega^{r+1}(M) $。

\item

$ d $ 满足\textbf{乘积法则}

\[
d(\alpha\wedge\beta) = d\alpha\wedge\beta + (-1)^{r}\alpha\wedge d\beta \qquad \text{对 } \alpha\in\Omega^{r}(M) \text{ 与 } \beta\in\Omega(M) .
\]
\item

$ d^{2} = d\circ d = 0 $。

\item

$ f\in\mathcal{E}(M) = \Omega^{0}(M) $ 的微分 $ df $ 与 $ f $ 的微分相同。
又

\begin{equation}
d\circ h^{\ast} = h^{\ast}\circ d\tag{4.18}
\end{equation}

对 $ h\in C^{\infty}(M,N) $。

\proofhead
（a）（存在性）设 $ (\varphi, U) $ 是 $ M $ 的坐标卡。由定理 3.5，恰有一个映射 $ d: \Omega(\varphi(U))\to\Omega(\varphi(U)) $ 具有性质 (i)–(iv)。我们通过要求下图交换来定义 $ d_{U}: \Omega(U)\to\Omega(U) $：
\begin{equation}
\begin{array}{ccc} \Omega(U) & \xrightarrow{d_{U}} & \Omega(U) \\ \varphi_{\ast}\downarrow & & \uparrow\varphi^{\ast} \\ \Omega(\varphi(U)) & \xrightarrow{d} & \Omega(\varphi(U)) \end{array}\tag{4.19}
\end{equation}

等价地，我们令 $ d_{U} := \varphi^{\ast}\circ d\circ\varphi_{\ast} $。由注 4.4(e) 可知 $ \varphi_{\ast} $ 是代数同构，且 $ (\varphi_{\ast})^{-1} = (\varphi^{-1})_{\ast} $。由此与 (4.19)，我们容易验证 $ d_{U} $ 具有性质 (i)–(iv) 且被唯一确定。
设 $ (\psi, V) $ 是 $ M $ 的另一个坐标卡，且 $ U\cap V\neq\varnothing $。则由 $ \varphi = (\varphi\circ\psi^{-1})\circ\psi $、拉回的性质与 (3.12) 推出

\begin{equation}
\begin{aligned}
d_{U} &= \varphi^{\ast}\circ d\circ\varphi_{\ast} = \psi^{\ast}\circ(\varphi\circ\psi^{-1})^{\ast}\circ d\circ(\varphi\circ\psi^{-1})_{\ast}\circ\psi_{\ast} \\ &= \psi^{\ast}\circ d\circ\psi_{\ast} = d_{V}
\end{aligned}\tag{4.20}
\end{equation}

（当然是在 $ U\cap V $ 上）。因此 $ d_{U} $ 与所选取的特殊坐标无关。
设 $ \{(\varphi_{\kappa}, U_{\kappa}) ; \kappa\in K\} $ 是 $ M $ 的图册，并设 $ i_{\kappa}: U_{\kappa}\hookrightarrow M $ 为自然嵌入。则我们如下定义 $ d: \Omega(M)\to\Omega(M) $：

\[
d\alpha(p) := d_{U_{\kappa}}[(i_{\kappa})^{\ast}\alpha](p) \qquad \text{对 } \alpha\in\Omega(M) ,
\]

其中 $ \kappa\in K $ 的选取使得 $ p $ 位于 $ U_{\kappa} $ 中。由 (4.20)，这个定义是有意义的，并且显然 $ d $ 具有性质 (i)–(iv)。

\end{subitems}

\begin{items}
\item[（b）]

（唯一性）设 $ \alpha\in\Omega^{r}(M) $，$ p\in M $。又设 $ (\varphi, U) $ 是 $ p $ 附近的坐标卡。由注 4.5(c)，$ \alpha|U $ 在局部坐标下可以写成

\[
\alpha|U = \sum_{(j)\in\mathbb{J}_{r}} a_{(j)}\,dx^{(j)}
\]

其中 $ a_{(j)}\in\mathcal{E}(U) $。现在由 (4.19) 以及注 3.6(a) 与 4.4(e) 推出

\begin{equation}
\begin{aligned}
d_{U}(\alpha|U) &= \varphi^{\ast}d\varphi_{\ast}(\alpha|U) = \varphi^{\ast}\sum_{(j)\in\mathbb{J}_{r}} d(\varphi_{\ast}a_{(j)})\wedge\varphi_{\ast}dx^{(j)} \\ &= \sum_{(j)\in\mathbb{J}_{r}} d_{U}a_{(j)}\wedge dx^{(j)} = \sum_{(j)\in\mathbb{J}_{r}} da_{(j)}\wedge dx^{(j)} ,
\end{aligned}\tag{4.21}
\end{equation}

因为 $ d_{U}a_{(j)} $ 是 $ a_{(j)}\in\mathcal{E}(U) $ 的微分。
设 $ V $ 是 $ p $ 的开邻域，且 $ V\subset\subset U $。则 $ \varphi(V)\subset\subset\varphi(U) $。因此注 1.21(a) 蕴含存在 $ \widetilde{\chi}\in\mathcal{D}(\varphi(U)) $ 使得 $ \widetilde{\chi}|\varphi(V) = 1 $。对

\begin{equation*}
\chi := \begin{cases} \varphi^{\ast}\widetilde{\chi} & \text{在 } U \text{ 上} , \\ 0 & \text{在 } M\setminus U \text{ 上} , \end{cases}
\end{equation*}

我们有 $ \chi\in\mathcal{E}(M) $ 且 $ \chi|V = 1 $。这意味着

\[
b_{(j)} := \chi a_{(j)} \quad \text{对 } (j)\in\mathbb{J}_{r} \qquad \text{与} \qquad \xi^{j} := \chi x^{j} \quad \text{对 } 1\le j\le m
\]

都属于 $ \mathcal{E}(M) $。因此微分 $ d\xi^{j}\in\Omega^{1}(M) $ 有定义，从而

\[
\beta := \sum_{(j)\in\mathbb{J}_{r}} b_{(j)}\,d\xi^{(j)}
\]

也有定义，且属于 $ \Omega(M) $。
现在设 $ \widetilde{d} $ 是从 $ \Omega(M) $ 到自身的、满足 (i)–(iv) 的映射。则我们容易得到

\[
\widetilde{d}\beta = \sum_{(j)\in\mathbb{J}_{r}} db_{(j)}\wedge d\xi^{(j)} .
\]

对 $ a\in\mathcal{E}(U) $，乘积法则给出

\[
d(\chi a) = a\,d\chi + \chi\,da
\]

（见推论 VII.3.8 以及切映射的定义）。因为 $ \chi|V = 1 $，我们可以利用自然嵌入 $ i: V\hookrightarrow M $ 得出

\[
\left\langle i^{\ast}d(\chi a)(q), v(q)\right\rangle_{q} = \left\langle d(\chi a)(q), v(q)\right\rangle_{q} = \left\langle da(q), v(q)\right\rangle_{q} \qquad \text{对 } q\in V ,
\]

对 $ v\in\mathcal{V}(M) $。即 $ d(\chi a)|V = da|V $。由此与 (4.21) 推出 $ \beta|V = \alpha|V $ 以及

\begin{equation}
\widetilde{d}\beta|V = d_{V}(\alpha|V) .\tag{4.22}
\end{equation}

因为 $ d_{V} $ 是唯一的，且每个 $ p\in M $ 都有使 (4.22) 成立的开坐标邻域 $ V $，所以我们看到 $ \widetilde{d} = d $。

\item[（c）]

为证明 (4.18)，我们可以利用前面的工作化归到局部情形，即可以假设 $ M = U $。于是断言由 (4.19)、(3.12) 与定理 3.5 推出。\nobreak\hspace{0.35em}■

\end{items}

\end{theorem}

\begin{remark}

\begin{items}
\item

设

\[
\alpha|U = \sum_{(j)\in\mathbb{J}_{r}} a_{(j)}\,dx^{(j)}
\]

是 $ \alpha\in\Omega^{r}(M) $ 在坐标卡 $ (\varphi, U) $ 的局部坐标下的表示。则

\[
d(\alpha|U) = \sum_{(j)\in\mathbb{J}_{r}} da_{(j)}\wedge dx^{(j)} .
\]

\proofhead
这由 (4.21) 推出。\nobreak\hspace{0.35em}■
\item

（正则性）对 $ k\in\mathbb{N} $，映射

\[
d: \Omega^{r}_{(k+1)}(M) \to \Omega^{r+1}_{(k)}(M) \qquad \text{对 } r\in\mathbb{N}
\]

是有定义且 $ \mathbb{R} $-线性的。当 $ M $ 是 $ C^{k+2} $ 流形时，这仍然成立。\nobreak\hspace{0.35em}■

\end{items}

\end{remark}

\needspace{3\baselineskip}\subsection*{闭形式与恰当形式}

与局部理论中一样，若 $ d\alpha = 0 $，则称 $ \alpha\in\Omega(M) $ 是\textbf{闭的}。若存在 $ \beta\in\Omega(M) $，即一个\textbf{原函数}，使得 $ d\beta = \alpha $，则称它是\textbf{恰当的}。

\begin{examplenote}

\begin{items}
\item

因为 $ d^{2} = 0 $，每个恰当形式都是闭形式。

\item

$ M $ 上的每个 $ m $ 形式都是闭形式。

\proofhead
这是因为 $ \Omega^{m+1}(M) = \{0\} $。\nobreak\hspace{0.35em}■
\item

（庞加莱引理）设 $ r\in\mathbb{N}^{\times} $，且 $ \alpha\in\Omega^{r}(M) $ 是闭形式。则 $ \alpha $ 是局部恰当的，即每个 $ p\in M $ 都有开邻域 $ U $ 与 $ \beta\in\Omega^{r-1}(U) $，使得 $ d\beta = \alpha|U $。

\proofhead
设 $ (\varphi, U) $ 是 $ p $ 附近的坐标卡，其中 $ \varphi(U) $ 是星形的。因为 $ d\alpha = 0 $ 且 $ d\varphi_{\ast}\alpha = \varphi_{\ast}d\alpha $，所以形式 $ \varphi_{\ast}\alpha\in\Omega^{r}(\varphi(U)) $。由于 $ \varphi(U) $ 是可缩的，由庞加莱引理（定理 3.11）存在 $ \beta_{0}\in\Omega^{r-1}(\varphi(U)) $ 使得 $ d\beta_{0} = \varphi_{\ast}\alpha $。令 $ \beta := \varphi^{\ast}\beta_{0}\in\Omega^{r-1}(U) $，我们便有 $ d\beta = \varphi^{\ast}d\beta_{0} = \varphi^{\ast}\varphi_{\ast}\alpha = \alpha|U $。\nobreak\hspace{0.35em}■
\end{items}

\end{examplenote}

\needspace{3\baselineskip}\subsection*{缩并}

设 $ \alpha\in\Omega^{r+1}(M) $，$ v\in\mathcal{V}(M) $。则 $ \alpha $ 关于 $ v $ 的\textbf{缩并} $ v\lrcorner\alpha $ 由

\[
v\lrcorner\alpha(v_{1}, \ldots, v_{r}) := \alpha(v, v_{1}, \ldots, v_{r}) \qquad \text{对 } v_{j}\in\mathcal{V}(M) \text{ 与 } 1\le j\le r .
\]

定义。我们有时把 $ v\lrcorner\cdot $ 写成 $ i_{v} $，并称 $ i_{v}\alpha $ 为 $ v $ 与 $ \alpha $ 的\textbf{内乘}。我们容易验证 $ v\lrcorner\alpha $ 属于 $ \Omega^{r}(M) $。为了完备性并避开一个麻烦的特殊情形，我们直接令

\[
v\lrcorner\alpha := 0 \qquad \text{对 } \alpha\in\Omega^{0}(M) .
\]

\begin{examplenote}

\begin{items}
\item

若 $ \varphi: M\to N $ 是微分同胚，则

\[
v\lrcorner(\varphi^{\ast}\alpha) = \varphi^{\ast}(\varphi_{\ast}v\lrcorner\alpha)
\]

对 $ \alpha\in\Omega(N) $ 与 $ v\in\mathcal{V}(M) $。特别地，对每个 $ r $，下图交换：

\begin{equation*}
\begin{array}{ccc} \Omega^{r+1}(M) & \xleftarrow{\varphi^{\ast}} & \Omega^{r+1}(N) \\ v\lrcorner\downarrow & & \downarrow\varphi_{\ast}v\lrcorner \\ \Omega^{r}(M) & \xleftarrow{\varphi^{\ast}} & \Omega^{r}(N) \end{array}
\end{equation*}

\proofhead
若 $ \alpha $ 是零形式，则断言显然成立。因此我们可以假设 $ \alpha\in\Omega^{r+1}(N) $。于是对 $ p\in M $ 与 $ v_{1}, \ldots, v_{r}\in T_{p}M $，我们得到
\begin{equation*}
\begin{aligned}
v\lrcorner(\varphi^{\ast}\alpha)(p)(v_{1}, \ldots, v_{r}) &= (\varphi^{\ast}\alpha)(p)(v(p), v_{1}, \ldots, v_{r}) \\ &= \alpha(\varphi(p))((T_{p}\varphi)v(p), (T_{p}\varphi)v_{1}, \ldots, (T_{p}\varphi)v_{r}) \\ &= \alpha(\varphi(p))(\varphi_{\ast}v(\varphi(p)), (T_{p}\varphi)v_{1}, \ldots, (T_{p}\varphi)v_{r}) \\ &= (\varphi_{\ast}v\lrcorner\alpha)(\varphi(p))((T_{p}\varphi)v_{1}, \ldots, (T_{p}\varphi)v_{r}) \\ &= \varphi^{\ast}(\varphi_{\ast}v\lrcorner\alpha)(p)(v_{1}, \ldots, v_{r}) ,
\end{aligned}
\end{equation*}

这证明了断言。\nobreak\hspace{0.35em}■

\item

设 $ X $ 是 $ \mathbb{H}^{m} $ 中的开集，$ \omega := dx^{1}\wedge\cdots\wedge dx^{m} $。对 $ v = \sum_{j=1}^{m} v^{j}\partial_{j} $，我们有

\[
v\lrcorner\omega = \sum_{j=1}^{m}(-1)^{j-1}v^{j}\,dx^{1}\wedge\cdots\wedge\widehat{dx^{j}}\wedge\cdots\wedge dx^{m} .
\]

\proofhead
我们令 $ v_{1} := v $。则对 $ v_{2}, \ldots, v_{m}\in\mathcal{V}(X) $，我们有
\[
(v_{1}\lrcorner\omega)(v_{2}, \ldots, v_{m}) = \omega(v_{1}, \ldots, v_{m}) = \det\left[\langle dx^{j}, v_{k}\rangle\right] .
\]

把这个行列式按第一列展开，我们发现它的值为

\[
\sum_{j=1}^{m}(-1)^{j+1}\langle dx^{j}, v_{1}\rangle\det(A_{j}) ,
\]

其中 $ A_{j} $ 是从 $ [\langle dx^{j}, v_{k}\rangle] $ 中划去第一列与第 $ j $ 行所得的矩阵。由此推出

\[
\det(A_{j}) = dx^{1}\wedge\cdots\wedge\widehat{dx^{j}}\wedge\cdots\wedge dx^{m}(v_{2}, \ldots, v_{m}) .
\]

现在断言随之成立，因为

\[
\langle dx^{j}, v_{1}\rangle = \sum_{k=1}^{m} v^{k}\langle dx^{j}, \partial_{k}\rangle = v^{j} .\quad\blacksquare
\]
\item

设

\[
\rho: \mathbb{R}^{m+1}\setminus\{0\} \to S^{m}, \qquad x\mapsto x/|x|
\]

为 $ \mathbb{R}^{m+1} $ 中 $ m $ 球面上的\textbf{径向收缩}\footnote{若 $ X $ 是拓扑空间，$ A $ 是 $ X $ 的子集，则连续映射 $ \rho: X\to A $ 若对 $ a\in A $ 满足 $ \rho(a) = a $，就称为 $ X $ 到 $ A $ 上的收缩。若存在 $ X $ 到 $ A $ 上的收缩，则称 $ A $ 是 $ X $ 的收缩核。}。又令

\[
\alpha := \sum_{j=1}^{m+1}(-1)^{j-1}x^{j}\,dx^{1}\wedge\cdots\wedge\widehat{dx^{j}}\wedge\cdots\wedge dx^{m+1}
\]

以及

\[
\omega_{S^{m}} := \alpha|S^{m} .
\]

于是令 $ r(x) := |x| $（$ x\in\mathbb{R}^{m+1} $），我们有

\[
\rho^{\ast}\omega_{S^{m}} = \frac{1}{r^{m+1}}\alpha = \sum_{j=1}^{m+1}(-1)^{j-1}\frac{x^{j}}{|x|^{m+1}}\,dx^{1}\wedge\cdots\wedge\widehat{dx^{j}}\wedge\cdots\wedge dx^{m+1} ,
\]

且 $ \rho^{\ast}\omega_{S^{m}} $ 是闭形式。

\proofhead
因为 $ \rho\in C^{\infty}(\mathbb{R}^{m+1}\setminus\{0\}, \mathbb{R}^{m+1}) $ 且 $ \mathrm{im}(\rho) = S^{m} $，我们知道 $ \rho $ 是从 $ \mathbb{R}^{m+1}\setminus\{0\} $ 到 $ S^{m} $ 的光滑映射。因此 $ \rho^{\ast}\omega_{S^{m}}\in\Omega^{m}(\mathbb{R}^{m+1}\setminus\{0\}) $ 有定义。由注 4.12(b) 以及 $ d(\rho^{\ast}\omega_{S^{m}}) = \rho^{\ast}d\omega_{S^{m}} = 0 $，它是闭形式。
为证明 $ \rho^{\ast}\omega_{S^{m}} = r^{-(m+1)}\alpha $，我们必须验证：对每个 $ p\in\mathbb{R}^{m+1}\setminus\{0\} $，两边在取自 $ T_{p}\mathbb{R}^{m+1} $ 的一组基向量的每个 $ m $ 元组上取值相同。因此设 $ p\in\mathbb{R}^{m+1}\setminus\{0\} $。$ T_{p}\mathbb{R}^{m+1} $ 的一组基由向量 $ \{(p)_{p}, (v_{1})_{p}, \ldots, (v_{m})_{p}\} $ 给出，其中
$ \{ (v_1)_p, \ldots, (v_m)_p \} $ 是 $ T_{p}(r(p)S^{m}) $ 的一组基。若 $ m $ 元组 $ (w_1)_p, \ldots, (w_m)_p $ 含有向量 $ (p)_p $，则令 $ \omega := dx^{1} \wedge \cdots \wedge dx^{m+1} $，我们可由 (b) 得到
\[
\alpha(p)((w_1)_p,\ldots,(w_m)_p) = ((p)_p \lrcorner \omega)((w_1)_p,\ldots,(w_m)_p) = \omega((p)_p,(w_1)_p,\ldots,(w_m)_p) = 0 ,
\]

因为有两个分量相同。因此 $ r^{-(m+1)}(p)\alpha(p) $ 也在这个 $ m $ 元组上为零。我们还得到

\[
\rho^{*}\omega_{S^{m}}(p)((w_1)_p,\ldots,(w_m)_p) = \omega_{S^{m}}(\rho(p))((T_{p}\rho)(w_1)_p,\ldots,(T_{p}\rho)(w_m)_p)
\]

其中

\[
(T_{p}\rho)(w_j)_p = (\rho(p), \partial\rho(p)w_j) ,
\]

这里按命题 VII.2.5，有

\[
\partial\rho(p)w_j = \partial_{t}\rho(p + tw_j)\big|_{t=0} \qquad \text{对 } 1 \le j \le m .
\]

因为 $ \rho(p+tp) = \rho(p) $（$ t \in (-1,1) $），特别地推出 $ (T_{p}\rho)(p)_p = 0 $。因此当 $ m $ 元组 $ (w_1)_p, \ldots, (w_m)_p $ 含有向量 $ (p)_p $ 时，$ \rho^{*}\omega_{S^{m}}(p)((w_1)_p,\ldots,(w_m)_p) $ 也为零。
剩下要证明

\begin{equation}
\rho^{*}\omega_{S^{m}}(p)((v_1)_p,\ldots,(v_m)_p) = \frac{1}{r(p)^{m+1}} \alpha(p)((v_1)_p,\ldots,(v_m)_p) .\tag{4.23}
\end{equation}

对 $ (v)_p \in T_{p}(r(p)S^{m}) $，定理 VII.10.6 给出 $ \varepsilon > 0 $ 与 $ \gamma \in C^{1}((-\varepsilon,\varepsilon), r(p)S^{m}) $，使得 $ \gamma(0) = p $ 且 $ \dot{\gamma}(0) = v $。现在我们利用 $ \rho \circ \gamma(t) = \gamma(t)/r(p) $ 得到

\[
\partial\rho(p)v = (\rho \circ \gamma)'(0) = v/r(p) .
\]

由此推出

\[
\rho^{*}\omega_{S^{m}}(p)((v_1)_p,\ldots,(v_m)_p) = r(p)^{-m}\alpha(\rho(p))((v_1)_p,\ldots,(v_m)_p) ,
\]

这蕴含 (4.23)，从而完成了证明。\nobreak\hspace{0.35em}■

\item

（正则性）设 $ k \in \mathbb{N} $，并设 $ M $ 是 $ C^{k+1} $ 流形。对 $ \alpha \in \Omega^{r+1}_{(k)}(M) $ 与 $ v \in \mathcal{V}^{k}(M) $，缩并 $ v \lrcorner \alpha $ 属于 $ \Omega^{r}_{(k)}(M) $。\nobreak\hspace{0.35em}■

\end{items}

\end{examplenote}

\needspace{3\baselineskip}\subsection*{可定向性}

正如我们在 §2 中学到的，$ T_{p}M $ 可通过选取体积形式 $ \alpha(p) \in \Lambda^{m}T_{p}^{*}M $ 来定向。由此得到 $ M $ 上的一个 $ m $ 形式 $ \alpha $，满足 $ \alpha(p) \neq 0 $（$ p \in M $）。反之，每个映射 $ p \mapsto \alpha(p) \in \Lambda^{m}T_{p}^{*}M $，只要 $ \alpha(p) \neq 0 $（$ p \in M $），就在每个 $ T_{p}M $ 上诱导一个定向。然而这样的 $ \alpha $ 一般不会连续。直观地说，这意味着切空间的定向并不“协调”，即切空间在从一点移动到下一点时可以“翻转”。为避免这一点，我们还要求 $ \alpha $ 是光滑的（更确切地说，具有流形的正则性所允许的正则性）。

流形 $ M $ 称为\textbf{可定向的}，若存在 $ \alpha \in \Omega^{m}(M) $ 使 $ \alpha(p) \neq 0 $ 对每个 $ p \in M $ 成立；这样的 $ m $ 形式 $ \alpha $ 称为 $ M $ 上的\textbf{体积形式}。

\begin{remark}

\begin{items}
\item

若 $ M $ 可定向，则 $ \Omega^{m}(M) $ 是一维 $ \mathcal{E}(M) $ 模。

\proofhead
设 $ \alpha $ 是 $ M $ 上的体积形式，并设 $ \beta \in \Omega^{m}(M) $。因为 $ \dim \bigwedge^{m}T_{p}^{*}M = 1 $（$ p \in M $），存在 $ f: M \to \mathbb{R} $ 使 $ \beta = f\alpha $。我们必须证明 $ f $ 是光滑的。在局部坐标下，有
\[
\alpha \mid U = a\,dx^{1} \wedge \cdots \wedge dx^{m} \qquad \text{且} \qquad \beta \mid U = b\,dx^{1} \wedge \cdots \wedge dx^{m}
\]

其中 $ a, b \in \mathcal{E}(U) $ 且 $ a(p) \neq 0 $（$ p \in U $）。由此推出 $ \beta \mid U = f\alpha \mid U $，其中 $ f := b/a $ 属于 $ \mathcal{E}(U) $。\nobreak\hspace{0.35em}■

\item

（正则性）设 $ k \in \mathbb{N} $，$ M $ 是 $ C^{k+1} $ 流形。则 $ M $ 可定向当且仅当存在 $ \alpha \in \Omega^{m}_{(k)}(M) $ 使 $ \alpha(p) \neq 0 $（$ p \in M $）。这又当且仅当 $ C^{k}(M) $ 模 $ \Omega^{m}_{(k)}(M) $ 是一维的。\nobreak\hspace{0.35em}■

\end{items}

下一个命题表明，也可以用坐标卡来刻画流形的可定向性。

若 $ X $ 与 $ Y $ 在 $ \mathbb{H}^{m} $ 中开，则称 $ \varphi \in \operatorname{Diff}(X,Y) $ 是\textbf{保定向的}[或\textbf{反定向的}]，若对每个 $ x \in X $ 有 $ \det \partial\varphi(x) > 0 $[或 $ \det \partial\varphi(x) < 0 $]，也就是说，若对每个 $ x \in X $，$ \partial\varphi(x) \in \mathcal{L}(\mathbb{R}^{m}) $ 是保定向的[或反定向的]自同构。流形 $ M $ 的图册称为\textbf{定向的}，若其所有转移函数都是保定向的。

\end{remark}

\begin{proposition}

维数 $ \ge 2 $ 的流形可定向当且仅当它具有定向图册。

证明

\begin{items}
\item

设 $ M $ 可定向，$ \alpha \in \Omega^{m}(M) $ 是体积形式。此外，设 $ \{ (\varphi_{\kappa}, U_{\kappa}) ; \kappa \in K \} $ 是 $ M $ 的图册。则在 $ X_{\kappa} := \varphi_{\kappa}(U_{\kappa}) \subset \mathbb{H}^{m} $ 上 $ (\varphi_{\kappa})_{\ast}\alpha = a_{\kappa}\,dx^{1} \wedge \cdots \wedge dx^{m} $，其中 $ a_{\kappa}(x) \neq 0 $（$ x \in X_{\kappa} $）。因为必要时可作坐标变换 $ x \mapsto (-x^{1}, x^{2}, \ldots, x^{m}) $，我们可以假定对某个 $ x_{\kappa} \in X_{\kappa} $，$ a_{\kappa}(x_{\kappa}) $ 严格为正。又因为可以假定 $ U_{\kappa} $ 从而 $ X_{\kappa} $ 是连通的，由介值定理（定理 III.4.7）推出 $ a_{\kappa}(x) > 0 $ 对所有 $ x \in X_{\kappa} $ 与每个 $ \kappa \in K $ 成立。
现设 $ (\varphi_{\kappa}, U_{\kappa}) $ 与 $ (\varphi_{\lambda}, U_{\lambda}) $ 是满足 $ U_{\kappa} \cap U_{\lambda} \neq \varnothing $ 的局部坐标卡。又设 $ \varphi_{\kappa} = (x^{1}, \ldots, x^{m}) $，$ \varphi_{\lambda} = (y^{1}, \ldots, y^{m}) $。则得到

\begin{equation}
\begin{aligned}
&(\varphi_{\lambda} \circ \varphi_{\kappa}^{-1})^{*}(a_{\lambda}\,dy^{1} \wedge \cdots \wedge dy^{m} \mid \varphi_{\lambda}(U_{\kappa} \cap U_{\lambda})) \\ &= (\varphi_{\kappa})_{\ast}\varphi_{\lambda}^{*}(a_{\lambda}\,dy^{1} \wedge \cdots \wedge dy^{m} \mid \varphi_{\lambda}(U_{\kappa} \cap U_{\lambda})) \\ &= (\varphi_{\kappa})_{\ast}\alpha \mid (U_{\kappa} \cap U_{\lambda}) = a_{\kappa}\,dx^{1} \wedge \cdots \wedge dx^{m} .
\end{aligned}\tag{4.24}
\end{equation}

由例 3.4(c)，有

\[
(\varphi_{\lambda} \circ \varphi_{\kappa}^{-1})^{*}dy^{1} \wedge \cdots \wedge dy^{m} = \det \partial(\varphi_{\lambda} \circ \varphi_{\kappa}^{-1})\,dx^{1} \wedge \cdots \wedge dx^{m} .
\]

与 (4.24) 比较，我们看到

\[
(\varphi_{\lambda} \circ \varphi_{\kappa}^{-1})^{*}a_{\lambda}(x) \det \partial(\varphi_{\lambda} \circ \varphi_{\kappa}^{-1})(x) = a_{\kappa}(x) > 0 \qquad \text{对 } x \in \varphi_{\kappa}(U_{\kappa} \cap U_{\lambda}) .
\]

因为 $ a_{\lambda} $ 为正，推出 $ M $ 具有定向图册。

\item

设 $ \{ (\varphi_{\kappa}, U_{\kappa}) ; \kappa \in K \} $ 是定向图册。命题 1.20 保证存在从属于 $ M $ 的覆盖 $ \{ U_{\kappa} ; \kappa \in K \} $ 的光滑单位分解 $ \{ \pi_{\kappa} ; \kappa \in K \} $。对 $ \kappa \in K $，定义 $ \alpha_{\kappa} \in \Omega^{m}(U_{\kappa}) $ 为

\begin{equation*}
\alpha_{\kappa} := \begin{cases} \pi_{\kappa}\varphi_{\kappa}^{*}\,dx^{1} \wedge \cdots \wedge dx^{m} & \text{在 } U_{\kappa} \text{ 中} , \\ 0 & \text{其他情形} . \end{cases}
\end{equation*}

容易验证，定义

\[
\alpha := \sum_{\kappa \in K} \alpha_{\kappa} \in \Omega^{m}(M)
\]

是有意义的。我们必须证明 $ \alpha(p) \neq 0 $（$ p \in M $）。
设 $ p \in M $，选取 $ \kappa \in K $ 使 $ \pi_{\kappa}(p) > 0 $。对满足 $ \lambda \neq \kappa $ 且 $ U_{\kappa} \cap U_{\lambda} \neq \varnothing $ 的 $ \lambda \in K $，如在 (a) 中那样，推出

\begin{equation*}
\begin{aligned}
\alpha_{\lambda} = \pi_{\lambda}\varphi_{\lambda}^{*}\,dy^{1} \wedge \cdots \wedge dy^{m} &= \pi_{\lambda}\varphi_{\kappa}^{*}(\varphi_{\lambda} \circ \varphi_{\kappa}^{-1})^{*}\,dy^{1} \wedge \cdots \wedge dy^{m} \\ &= \pi_{\lambda}\big(\varphi_{\kappa}^{*}\det\big(\partial(\varphi_{\lambda} \circ \varphi_{\kappa}^{-1})\big)\big)\varphi_{\kappa}^{*}\,dx^{1} \wedge \cdots \wedge dx^{m} .
\end{aligned}
\end{equation*}

由此得到

\begin{equation*}
\begin{aligned}
\alpha(p) = \Big( \pi_{\kappa}(p) + \sum_{\substack{\lambda \in K \\ \lambda \neq \kappa}} \pi_{\lambda}(p)\det\big(\partial(\varphi_{\lambda} \circ \varphi_{\kappa}^{-1})\big)(\varphi_{\kappa}(p)) \Big)\varphi_{\kappa}^{*}\,dx^{1} \wedge \cdots \wedge dx^{m}(p) ,
\end{aligned}
\end{equation*}

其中只有有限多个被加项不为零。因为 $ \pi_{\lambda}(p) \ge 0 $，且转移函数保定向，我们看到 $ \alpha(p) \neq 0 $。因此 $ \alpha $ 是体积形式，$ M $ 可定向。\nobreak\hspace{0.35em}■

\end{items}

\end{proposition}

设 $ M $ 可定向。则称两个体积形式 $ \alpha, \beta \in \Omega^{m}(M) $ 是\textbf{等价的}，若存在 $ f \in \mathcal{E}(M) $ 使 $ f(p) > 0 $（$ p \in M $）且 $ \alpha = f\beta $。这显然是 $ M $ 上所有体积形式之集上的等价关系。关于此关系的每个等价类称为 $ M $ 上的一个\textbf{定向}。给定 $ M $ 的定向 $ \mathcal{O}r := \mathcal{O}r(M) $，则称 $ (M, \mathcal{O}r) $ 为\textbf{定向流形}。若 $ M $ 的定向由上下文清楚，我们可以把 $ (M, \mathcal{O}r) $ 简写为 $ M $。

若 $ \alpha \in \mathcal{O}r $，则 $ -\alpha $ 是不属于 $ \mathcal{O}r $ 的体积形式。我们把相应的等价类记作 $ -\mathcal{O}r $，并称它为\textbf{与 $ \mathcal{O}r $ 相反的定向}。显然 $ -\mathcal{O}r $ 与它的具体代表无关。

\begin{remark}

\begin{items}
\item

可定向流形连通当且仅当它恰好有两个定向。

\proofhead
设 $ M $ 连通，$ \alpha $ 与 $ \beta $ 是两个体积形式。由注 4.14(a)，存在 $ f \in \mathcal{E}(M) $ 使 $ \alpha = f\beta $。因为 $ \alpha $ 处处不为零，我们有 $ f(p) \neq 0 $（$ p \in M $）。因为 $ M $ 连通，介值定理（见定理 III.4.7）蕴含对每个 $ p \in M $，或者 $ f(p) > 0 $，或者 $ f(p) < 0 $。因此 $ \alpha $ 或者等价于 $ \beta $，或者等价于 $ -\beta $。所以 $ M $ 恰有两个定向。
现设 $ M $ 连通。命题 III.4.2 保证存在 $ M $ 的非空、既开又闭的真子集 $ X $。对 $ M $ 上的体积形式 $ \alpha $，我们置
\begin{equation*}
\beta(p) := \begin{cases} \alpha(p) & \text{若 } p \in X , \\ -\alpha(p) & \text{若 } p \in M \setminus X . \end{cases}
\end{equation*}

则 $ \beta $ 显然是体积形式，且 $ \beta \notin \mathcal{O}r \cup (-\mathcal{O}r) $，其中 $ \mathcal{O}r $ 是 $ \alpha $ 的等价类。因此 $ M $ 有两个以上的定向。\nobreak\hspace{0.35em}■

\item

设 $ M = (M, \mathcal{O}r) $ 是定向流形。称 $ M $ 的坐标卡 $ (\varphi, U) $ 是\textbf{正向的（正定向的）}，若对 $ \alpha \in \mathcal{O}r $，$ \varphi_{\ast}(\alpha \mid U) $ 等价于 $ m $ 形式 $ dx^{1} \wedge \cdots \wedge dx^{m} \mid \varphi(U) $。否则称它是\textbf{负向的（负定向的）}。$ M $ 具有仅由正向坐标卡组成的图册，即定向图册。

\proofhead
对 $ \beta \in \mathcal{O}r $，有 $ \alpha = f\beta $，其中 $ f \in \mathcal{E}(M) $ 且 $ f(p) > 0 $（$ p \in M $）。以 $ i: U \hookrightarrow M $ 表示包含映射，由此推出
\[
\varphi_{\ast}\alpha \mid U = \varphi_{\ast}i^{\ast}\alpha = \varphi_{\ast}i^{\ast}(f\beta) = (\varphi_{\ast}i^{\ast}f)(\varphi_{\ast}i^{\ast}\beta) = g\varphi_{\ast}(\beta \mid U) ,
\]

其中 $ g := f \circ \varphi^{-1} \in \mathcal{E}(\varphi(U)) $ 且 $ g(x) > 0 $（$ x \in \varphi(U) $）。这表明该定义与所选的代表无关。确实存在由正向坐标卡组成的图册，这一点已在命题 4.15 的证明的 (a) 部分中说明。\nobreak\hspace{0.35em}■

\item

设 $ M $ 已定向。则 $ (\varphi, U) $ 是正向坐标卡当且仅当 $ (\partial_1|_p, \ldots, \partial_m|_p) $ 是 $ T_{p}M $ 的正向基（$ p \in U $）。

\proofhead
对 $ \alpha \in \Omega^{m}(M) $，注 4.5(c) 说局部坐标下的基表示为
\[
\alpha \mid U = a\,dx^{1} \wedge \cdots \wedge dx^{m}
\]

其中 $ a(p) = \alpha(p)(\partial_1|_p, \ldots, \partial_m|_p) $（$ p \in U $）。现在断言清楚了。\nobreak\hspace{0.35em}■

\end{items}

\end{remark}

\begin{example}

\begin{items}
\item

可定向流形 $ M $ 的每个开子集 $ U $ 本身可定向。\footnote{我们约定空集是可定向的。}

\proofhead
对 $ \alpha \in \mathcal{O}r(M) $，限制 $ \alpha \mid U $ 是 $ U $ 上的体积形式。\nobreak\hspace{0.35em}■
\item

若 $ M $ 与 $ N $ 可定向，且这两个流形中有一个没有边界，则乘积流形\footnote{见习题 VII.9.4 与习题 3。为什么我们假定这两个流形中有一个没有边界？} $ M \times N $ 可定向。

\proofhead
若 $ \{ (\varphi_{\kappa}, U_{\kappa}) ; \kappa \in K \} $ 与 $ \{ (\psi_{\lambda}, V_{\lambda}) ; \lambda \in L \} $ 分别是 $ M $ 与 $ N $ 的定向图册，则容易看出 $ \{ \varphi_{\kappa} \times \psi_{\lambda} ; (\kappa, \lambda) \in K \times L \} $ 是 $ M \times N $ 的定向图册，其中
\[
\varphi_{\kappa} \times \psi_{\lambda}(p,q) := (\varphi_{\kappa}(p), \psi_{\lambda}(q)) \in \mathbb{R}^{m} \times \mathbb{R}^{n} \qquad \text{对 } (p,q) \in U_{\kappa} \times V_{\lambda} ,
\]

因为可以不失一般性地假定 $ M $ 与 $ N $ 都至少是一维的，断言由命题 4.15 推出。\nobreak\hspace{0.35em}■

\item

任何可由单个坐标卡描述的流形（即具有只含一个坐标卡的图册的流形）都是可定向的。

\proofhead
这是平凡的（见命题 4.15 的证明的第一部分）。\nobreak\hspace{0.35em}■
\item

（图像）设 $ X $ 在 $ \mathbb{R}^{m} $ 中开，$ f \in C^{\infty}(X, \mathbb{R}^{n}) $。则 $ \operatorname{graph}(f) $ 是 $ \mathbb{R}^{m+n} $ 的 $ m $ 维可定向子流形。

\proofhead
由命题 VII.9.2，我们知道 $ \operatorname{graph}(f) $ 是 $ \mathbb{R}^{m+n} $ 的 $ m $ 维子流形。该结果的证明表明
\[
\varphi: \operatorname{graph}(f) \to X , \quad (x,f(x)) \mapsto x
\]

是描述 $ \operatorname{graph}(f) $ 的坐标卡。因此断言由 (c) 推出。\nobreak\hspace{0.35em}■

\item

（正则映射的纤维）设 $ X $ 在 $ \mathbb{R}^{m} $ 中开，$ \ell \in \{ 0, \ldots, m-1 \} $。又设 $ q $ 是 $ f \in C^{\infty}(X, \mathbb{R}^{m-\ell}) $ 的正则值。则 $ X $ 的 $ \ell $ 维子流形 $ f^{-1}(q) $ 是可定向的。

\proofhead
令 $ \omega := dx^{1} \wedge \cdots \wedge dx^{m} \mid X $ 以及
\[
\nabla f^{k} := \sum_{j=1}^{m} \partial_j f^{k} \frac{\partial}{\partial x^{j}} \in \mathcal{V}(X) \qquad \text{对 } 1 \le k \le m - \ell .
\]

采用注 VII.10.11(a) 的记号，我们有 $ \nabla f^{k}(p) = \nabla_{p}f^{k} $（$ p \in X $）。我们可以假定 $ L := f^{-1}(q) $ 非空。于是

\[
\alpha := \nabla f^{1} \lrcorner \big( \nabla f^{2} \lrcorner ( \cdots \lrcorner ( \nabla f^{m-\ell} \lrcorner\, \omega ) \cdots ) \big) \mid L \in \Omega^{\ell}(L) .
\]

命题 VII.10.13 保证 $ \nabla f^{1}(p), \ldots, \nabla f^{m-\ell}(p) $ 线性无关。因此

\[
\alpha(p) = \omega(\nabla f^{m-\ell}(p), \ldots, \nabla f^{1}(p), \ldots) \neq 0 \qquad \text{对 } p \in L ,
\]

也就是说，$ \alpha $ 是 $ L $ 上的体积形式。\nobreak\hspace{0.35em}■

\item

若 $ M $ 与 $ N $ 微分同胚，则 $ M $ 可定向当且仅当 $ N $ 可定向。

\proofhead
设 $ f \in \operatorname{Diff}(M,N) $，并设 $ (\varphi, U) $ 是 $ M $ 的坐标卡。则 $ \psi := \varphi \circ f^{-1} $ 是 $ N $ 的坐标卡，其中 $ V := f(U) = \operatorname{dom}(\psi) $。因为 $ M $ 与 $ N $ 微分同胚，$ m = n $。现设 $ \beta \in \Omega^{m}(N) $ 是 $ N $ 上的体积形式。局部坐标 $ (y^{1}, \ldots, y^{m}) = \psi $ 给出 $ \beta $ 的表示 $ \beta \mid V = b\,dy^{1} \wedge \cdots \wedge dy^{m} $，其中 $ b(q) \neq 0 $（$ q \in V $）。由此推出
\[
f^{*}(\beta \mid V) = (f^{*}b)f^{*}(dy^{1} \wedge \cdots \wedge dy^{m}) = b \circ f\,dx^{1} \wedge \cdots \wedge dx^{m} ,
\]

因为 $ f^{*}y^{j} = \operatorname{pr}_j \circ \psi \circ f = \operatorname{pr}_j \circ \varphi = x^{j} $，其中 $ (x^{1}, \ldots, x^{m}) = \varphi $。因为 $ b \circ f(p) \neq 0 $（$ p \in U $），我们看到 $ f^{*}\beta $ 是 $ M $ 上的体积形式。现在断言立刻可得。\nobreak\hspace{0.35em}■

\item

每个一维流形都是可定向的。

\proofhead
我们可以假定流形 $ M $ 连通，因为只需证明每个连通分支可定向。于是由定理 1.18，$ M $ 微分同胚于区间 $ J $ 或 $ S^{1} $。因为 $ J $ 与 $ S^{1} $ 可定向（其中 $ S^{1} $ 的可定向性例如由 (e) 推出），断言由 (f) 蕴含。\nobreak\hspace{0.35em}■
\item

（超曲面）$ \mathbb{R}^{m+1} $ 中的超曲面 $ M $ 可定向当且仅当 $ M $ 上存在光滑的\textbf{单位法向量场}，即存在 $ \nu \in C^{\infty}(M, \mathbb{R}^{m+1}) $，使得 $ |\nu(p)| = 1 $ 且 $ \boldsymbol{\nu}(p) = (p, \nu(p)) \in T_{p}^{\perp}M $（$ p \in M $）。

\proofhead
若 $ \nu $ 是 $ M $ 上的单位法向量场，则 $ (\boldsymbol{\nu} \lrcorner\, dx^{1} \wedge \cdots \wedge dx^{m+1}) \mid M $ 是 $ M $ 上的体积形式。因此 $ M $ 可定向。
设 $ M $ 可定向。若 $ (\varphi, U) $ 是正向坐标卡，$ \varphi = (x^{1}, \ldots, x^{m}) $，则因为 $ \dim(T_{p}^{\perp}M) = 1 $，对每个 $ p \in U $ 恰有一个 $ \boldsymbol{\nu}(p) = (p, \nu(p)) \in T_{p}^{\perp}M $ 满足 $ |\nu(p)| = 1 $，使得
\[
\big( \boldsymbol{\nu}(p), \frac{\partial}{\partial x^{1}}\big|_p, \ldots, \frac{\partial}{\partial x^{m}}\big|_p \big)
\]

是 $ T_{p}\mathbb{R}^{m+1} $ 的正向基。通过缩小 $ U $，我们可以假定存在 $ \mathbb{R}^{m+1} $ 的开集 $ \widetilde{U} $ 与 $ \widetilde{V} $，满足 $ U = \widetilde{U} \cap M $，并存在 $ \Phi \in \operatorname{Diff}(\widetilde{U}, \widetilde{V}) $ 使 $ U = f^{-1}(0) $，其中 $ f := \Phi^{m+1} \in \mathcal{E}(\widetilde{U}) $。因为 $ \nabla f(p) \neq 0 $（$ p \in \widetilde{U} $），推出 $ f $ 是正则的。因此由命题 VII.10.13 推出

\[
\nu(p) = \varepsilon \nabla f(p)/|\nabla f(p)| \qquad \text{对 } p \in U
\]

其中 $ \varepsilon \in \{ \pm 1 \} $。这表明 $ \nu $ 是光滑的。
现设 $ (\psi, V) $ 是第二个正向坐标卡，$ U \cap V \neq \varnothing $，$ \psi = (y^{1}, \ldots, y^{m}) $，并设 $ \boldsymbol{\mu}(q) = (q, \mu(q)) \in T_{q}^{\perp}M $ 满足 $ \mu \in C^{\infty}(V, \mathbb{R}^{m+1}) $ 且 $ |\mu(q)| = 1 $（$ q \in V $）。又设 $ (\boldsymbol{\mu}(q), \frac{\partial}{\partial y^{1}}|_q, \ldots, \frac{\partial}{\partial y^{m}}|_q) $ 是 $ T_{q}\mathbb{R}^{m+1} $ 的正向基（$ q \in V $）。因为两个基

\[
\big( \frac{\partial}{\partial x^{1}}\big|_p, \ldots, \frac{\partial}{\partial x^{m}}\big|_p \big) \qquad \text{与} \qquad \big( \frac{\partial}{\partial y^{1}}\big|_p, \ldots, \frac{\partial}{\partial y^{m}}\big|_p \big)
\]

对 $ p \in U \cap V $ 具有相同的定向，推出 $ \mu(p) = \nu(p) $（$ p \in U \cap V $）。于是单位法向量场的存在性由 $ M $ 的定向图册的存在性推出。\nobreak\hspace{0.35em}■

\item

（莫比乌斯带）设 $ R > 0 $，定义

\[
f: [-\pi, \pi) \times (-1,1) \to \mathbb{R}^{3}
\]

为

\[
f(\theta,t) := \Big( \big( R + t\cos\frac{\theta}{2} \big)\cos\theta, \big( R + t\cos\frac{\theta}{2} \big)\sin\theta, t\sin\frac{\theta}{2} \Big) .
\]

则 $ f $ 的像 $ M $ 是不可定向曲面，即\textbf{莫比乌斯带}。直观上，映射 $ f $ 的作用如下：因为

\[
f(\pm\pi, t) = (-R, 0, \pm t) ,
\]

它把矩形 $ [-\pi, \pi] \times (-1,1) $ 的末端 $ \{ \pi \} \times (-1,1) $ 相对于起点 $ \{ -\pi \} \times (-1,1) $ 扭转 180 度。这两个末端随后被粘合在一起。

\end{items}

\bookfigure[0.50]{figures/d36a9d513a56e42a.png}

将 $ f $ 表示为

\[
f(\theta,t) = R(\cos\theta, \sin\theta, 0) + tg(\theta)
\]

其中 $ g(\theta) := (\cos(\theta/2)\cos\theta, \sin(\theta/2)\sin\theta, \sin(\theta/2)) $，我们就得到 $ f $ 的参数化的一种解释：以角速度 1 运动的点描出 $ (x,y) $ 平面内以 0 为圆心、以 $ R $ 为半径的圆；这描述第一项。长度为 2 的杆的中点（沿其长度方向）固定在该点上，并允许它同时以角速度 $ 1/2 $ 绕自身中点旋转，使得转一圈后其方向反转；这就是第二项。

\proofhead
关于 $ M $ 是光滑曲面的证明留给你。

对 $ -\pi \le \theta \le \pi $，有

\begin{equation*}
\begin{aligned}
v_1(\theta) &:= \partial_1 f(\theta,0) = R(-\sin\theta, \cos\theta, 0) , \\ v_2(\theta) &:= \partial_2 f(\theta,0) = (\cos(\theta/2)\cos\theta, \cos(\theta/2)\sin\theta, \sin(\theta/2)) .
\end{aligned}
\end{equation*}

由此推出，对每个 $ \theta \in [-\pi,\pi) $，附着在 $ p(\theta) := f(\theta,0) $ 上的向量 $ v_1(\theta), v_2(\theta) $ 构成 $ T_{p(\theta)}M $ 的一组基。因此附着在 $ p(\theta) $ 上的向量

\[
n(\theta) := (-v_1(\theta) \times v_2(\theta))/R = (-\cos\theta\sin(\theta/2), -\sin\theta\sin(\theta/2), \cos(\theta/2))
\]

对 $ -\pi \le \theta < \pi $ 是单位法向量。特别地，$ n(0) = e_3 $。

我们假定 $ \nu: M \to \mathbb{R}^{3} $ 是单位法向量场，且 $ \nu(p(0)) = e_3 $。则因为 $ T_{p(\theta)}^{\perp}M $ 是连续且一维的，推出向量 $ \nu(p(\theta)) $ 与 $ n(\theta) $ 在 $ -\pi \le \theta < \pi $ 上重合。由此并结合关系 $ p(-\pi) = p(\pi) $，当 $ \theta \to \pi $ 时得到

\[
-e_1 = n(-\pi) = \nu(p(-\pi)) = \nu(p(\pi)) = n(\pi) = e_1 ,
\]

这是不可能的。因此 $ M $ 上不存在光滑（甚至连续）的单位法向量场；由 (h)，这表明 $ M $ 不可定向。\nobreak\hspace{0.35em}■

\begin{items}
\item[（j）]

（正则性）作明显的修改后，上述结论对 $ C^{1} $ 流形仍然成立。\nobreak\hspace{0.35em}■

\end{items}

\end{example}

\needspace{3\baselineskip}\subsection*{张量场}

设 $ r, s \in \mathbb{N} $。则按 §2，由 $ r $ 逆变与 $ s $ 协变张量组成的向量空间 $ T_{s}^{r}(T_{p}M) $ 在 $ T_{p}M $ 上是良好定义的。因此 $ M $ 上的 \textbf{(r,s)-张量丛}，

\[
T_{s}^{r}(M) := \bigcup_{p \in M} T_{s}^{r}(T_{p}M) ,
\]

也是良好定义的。$ M $ 上的一个 \textbf{(r,s)-张量}（更确切地说，一个 $ r $ 逆变与 $ s $ 协变张量）是这个丛的一个截面，即它是一个映射

\[
\gamma: M \to T_{s}^{r}(M) \qquad \text{其中 } \gamma(p) \in T_{s}^{r}(T_{p}M) \text{ 对 } p \in M .
\]

若 $ \gamma $ 与 $ \delta $ 是 $ M $ 上的 $ (r,s) $ 张量，$ f \in \mathbb{R}^{M} $，则\textbf{和} $ \gamma + \delta $、\textbf{与函数的乘积} $ f\gamma $ 以及\textbf{张量积} $ \gamma \otimes \delta $ 仍逐点定义为

\[
(\gamma + \delta)(p) := \gamma(p) + \delta(p) , \quad (f\gamma)(p) := f(p)\gamma(p) , \quad \gamma \otimes \delta(p) := \gamma(p) \otimes \delta(p)
\]

（$ p \in M $）。同样，$ \gamma \in T_{s}^{r}(M) $ 对普法夫形式的 $ r $ 元组 $ \alpha_1, \ldots, \alpha_r $ 与向量场的 $ s $ 元组 $ v_1, \ldots, v_s $ 的作用逐点定义为

\[
\gamma(\alpha_1, \ldots, \alpha_r, v_1, \ldots, v_s)(p) := \gamma(p)(\alpha_1(p), \ldots, \alpha_r(p), v_1(p), \ldots, v_s(p)) \qquad \text{对 } p \in M .
\]

最后，设 $ \varphi \in \operatorname{Diff}^{1}(M,N) $。则定义 $ \gamma \in T_{s}^{r}(M) $ 经 $ \varphi $ 的\textbf{前推}为

\[
(\varphi_{\ast}\gamma)(\alpha_1, \ldots, \alpha_r, v_1, \ldots, v_s) := (\gamma \circ \varphi^{-1})(\varphi^{\ast}\alpha_1, \ldots, \varphi^{\ast}\alpha_r, \varphi^{\ast}v_1, \ldots, \varphi^{\ast}v_s) ,
\]

其中 $ \alpha_1, \ldots, \alpha_r $ 是 $ N $ 上的普法夫形式，$ v_1, \ldots, v_s $ 是 $ N $ 上的向量场，并置

\begin{equation}
\varphi^{\ast}v := (\varphi^{-1})_{\ast}v\tag{4.25}
\end{equation}

其中 $ v $ 是 $ N $ 上的向量场。自然，$ \varphi^{\ast}\gamma := (\varphi^{-1})_{\ast}\gamma $ 就是 $ \gamma \in T_{s}^{r}(N) $ 的\textbf{拉回}。

设 $ k \in \mathbb{N} \cup \{ \infty \} $。则 $ (r,s) $ 张量 $ \gamma $ 属于类 $ C^{k} $（或者说，是 $ k $ 次\textbf{连续可微}的，或当 $ k = \infty $ 时是\textbf{光滑}的），若 $ M $ 的每个点都有坐标卡 $ (\varphi, U) $，使得 $ \varphi_{\ast}\gamma $ 是 $ \varphi(U) $ 上类 $ C^{k} $ 的 $ (r,s) $ 张量。我们把 $ M $ 上所有光滑 $ (r,s) $ 张量之集记作

\[
\mathcal{T}_{s}^{r}(M) .
\]

下列注的证明可由 §2 与 §3\footnote{若先建立（初等的）向量丛理论，这种重复是可以避免的。} 中关于微分形式的相应证明直接搬来。因此我们把这些证明留给你。

\begin{remark}

\begin{items}
\item

可微性的定义与坐标无关。

\item

$ T_{s}^{r}(M) $ 是 $ \mathcal{E}(M) $ 模。张量积映射

\[
\otimes: T_{s_1}^{r_1}(M) \times T_{s_2}^{r_2}(M) \to T_{s_1+s_2}^{r_1+r_2}(M) , \quad (\gamma, \delta) \mapsto \gamma \otimes \delta
\]

是 $ \mathcal{E}(M) $ 双线性的且结合的。

\item

$ M $ 上的 $ (r,s) $ 张量 $ \gamma $ 光滑当且仅当

\[
\gamma(\alpha_1, \ldots, \alpha_r, v_1, \ldots, v_s) \in \mathcal{E}(M)
\]

对所有 $ v_1, \ldots, v_s \in \mathcal{V}(M) $ 与 $ \alpha_1, \ldots, \alpha_r \in \Omega^{1}(M) $ 成立。

\item

设 $ (\varphi, U) $ 是 $ M $ 的坐标卡。则

\begin{equation}
\left\{ \frac{\partial}{\partial x^{j_{1}}} \otimes \cdots \otimes \frac{\partial}{\partial x^{j_{r}}} \otimes dx^{k_{1}} \otimes \cdots \otimes dx^{k_{s}} ;\ j_{i}, k_{i} \in \{1, \ldots, m\} \right\}\tag{4.26}
\end{equation}

是 $ \mathcal{T}_s^r(M) $ 的一组模基。于是 $ \gamma \in \mathcal{T}_s^r(U) $ 当且仅当 $ \gamma $ 在基表示 (4.26) 中的系数是光滑的。

\item

设 $ \varphi \in \operatorname{Diff}(M, N) $。则 $ \varphi_* $ 将模 $ \mathcal{T}_s^r(M) $ 映到 $ \mathcal{T}_s^r(N) $，并按协变方式作用为

\[
(\psi \circ \varphi)_* = \psi_* \circ \varphi_* \quad \text{且} \quad (\operatorname{id}_M)_* = \operatorname{id}_{\mathcal{T}_s^r(M)} .
\]

类似地，$ \varphi^*(\mathcal{T}_s^r(N)) = \mathcal{T}_s^r(M) $，且 $ \varphi^* $ 按反变方式作用。最后，$ \varphi_* $ 从而 $ \varphi^* $ 也与张量积映射相容，即

\[
\varphi_*(\gamma \otimes \delta) = \varphi_*\gamma \otimes \varphi_*\delta .
\]
\item

对 $ f \in C^\infty(M, N) $ 与 $ \gamma \in \mathcal{T}_s^0(N) $，$ \gamma $ 经 $ f $ 的\textbf{拉回} $ f^*\gamma $ 由下式确定：

\[
f^*\gamma(v_1, \ldots, v_s) := (\gamma \circ f)((Tf)v_1, \ldots, (Tf)v_s) \quad \text{对 } v_1, \ldots, v_s \in \mathcal{V}(M)
\]

其中 $ ((Tf)v)(p) := (T_p f)v(p) $ 对 $ p \in M $。于是映射

\[
f^*: \mathcal{T}_s^0(N) \to \mathcal{T}_s^0(M), \quad \gamma \mapsto f^*\gamma
\]

是良定义的且是 $ \mathbb{R} $ 线性的，按反变方式作用，并与张量积映射相容。对 $ f \in \operatorname{Diff}(M, N) $，它与先前定义的拉回相同。

\item

$ \mathcal{T}_0^1(M) = \mathcal{V}(M) $，$ \mathcal{T}_1^0(M) = \Omega^1(M) $，且对偶配对 $ \langle \cdot, \cdot \rangle $ 是 $ M $ 上的一个 $ (1, 1) $ 张量。

\item

（正则性）设 $ k \in \mathbb{N} $。则当 $ M $ 是 $ C^{k+1} $ 流形且把 $ C^\infty $ 换成 $ C^k $ 时，上述论断同样成立。\nobreak\hspace{0.35em}■

\end{items}

\end{remark}

\exerciseshead

\begin{exercises}
\item

设 $ N $ 是无边流形的子流形。证明（在典范等同下）对 $ k \in \mathbb{N} \cup \{\infty\} $ 有 $ \mathcal{V}^k(N) \subset \mathcal{V}^k(M) $。

\item

设 $ \alpha \in \Omega^r(M) $、$ \beta \in \Omega(M) $，并设 $ v \in \mathcal{V}(M) $。证明

\[
v \lrcorner (\alpha \wedge \beta) = (v \lrcorner \alpha) \wedge \beta + (-1)^r \alpha \wedge (v \lrcorner \beta) .
\]
\item

验证例 4.17(b) 的证明中所作的陈述。

\item

对 $ \alpha \in \Omega^1(M) $ 与 $ v \in \mathcal{V}(M) $，在局部坐标中计算 $ d\langle \alpha, v \rangle $。

\end{exercises}

\section{黎曼度量}

我们已从第 VII.10 节知道，$ \mathbb{R}^{\overline{m}} $ 上的欧几里得内积 $ (\cdot|\cdot) $ 通过限制到子流形 $ M $ 的切空间 $ T_p M $ 上，可以定义另一个内积。这就给出了在 $ T_p M $ 上度量长度与角度的方法。所以，例如，我们可以通过验证切空间 $ T_p\Gamma_1 $ 与 $ T_p\Gamma_2 $ 在 $ T_p M $ 中彼此正交，来确定 $ M $ 上的两条曲线 $ \Gamma_1 $ 与 $ \Gamma_2 $ 是否在点 $ p $ 处正交相交。

$ \mathbb{R}^{\overline{m}} $ 的欧几里得结构在 $ M $ 上——确切地说，在 $ M $ 的切丛上——诱导出一个欧几里得结构，这是流形上积分理论的基础，我们将在下一章讨论该理论。在本节中，我们探讨 $ M $ 上存在欧几里得结构的若干推论，并研究几个例子。我们还引入霍奇星算子与余微分，它们对更深入地研究微分形式理论具有重要意义——特别地，这些概念在（理论）物理学中很重要。

为便于引入材料，我们首先考虑由 $ (\cdot|\cdot) $ 在 $ M $ 上诱导的欧几里得结构的情形。然而，后面将会清楚，所有抽象定理在本质上更一般的框架，即黎曼几何的框架中仍然成立。由于这些事实具有重大的理论与实际意义，我们将引入（伪）黎曼度量的概念，它构成我们后续讨论的一般框架。

本节自始至终作如下假设：

- $ M $ 是 $ \mathbb{R}^{\overline{m}} $ 的 $ m $ 维子流形；$ N $ 是 $ \mathbb{R}^{\overline{n}} $ 的 $ n $ 维子流形。
- 除非另有说明，指标 $ i, j, k, l $ 总取 $ 1 $ 到 $ m $ 的值，而 $ \sum_{j} $ 表示 $ j $ 从 $ 1 $ 到 $ m $ 求和。

\needspace{3\baselineskip}\subsection*{体积元}

设 $ M $ 已定向。则 $ \mathcal{O}r $ 在每个切空间 $ T_p M $ 上诱导一个定向。又 $ T_p M $ 是内积空间，其内积 $ (\cdot|\cdot)_p $ 由周围空间 $ \mathbb{R}^{\overline{m}} $ 的欧几里得标量积诱导。于是由注 2.12(b)，$ T_p M $ 上存在唯一的体积元 $ \omega_p $。因此

\[
\omega_M(p) := \omega_p \qquad \text{对 } p \in M
\]

在 $ M $ 上定义一个 $ m $ 形式，即 $ M $ 上的\textbf{体积元}。

\begin{proposition}

设 $ M $ 已定向。则 $ \omega_M $ 属于 $ \mathcal{O}r(M) $。若 $ (\varphi, U) $ 是正向坐标卡，且 $ \varphi = (x^1, \ldots, x^m) $，则

\begin{equation}
\omega_M|U = \sqrt{G}\, dx^1 \wedge \cdots \wedge dx^m ,\tag{5.1}
\end{equation}

其中 $ G := \det[g_{jk}] \in \mathcal{E}(U) $ 是\textbf{格拉姆行列式}，且

\[
g_{jk}(p) := (\partial_j|_p | \partial_k|_p)_p \qquad \text{对 } 1 \le j, k \le m \text{ 及 } p \in U .
\]

设 $ g_\varphi := i \circ \varphi^{-1} \in C^\infty(\varphi(U), \mathbb{R}^{\overline{m}}) $ 是属于 $ \varphi $ 的参数化，其中 $ i : M \hookrightarrow \mathbb{R}^{\overline{m}} $，则

\begin{equation}
\varphi_{*}g_{jk}(x) = (\partial_j g_\varphi(x) | \partial_k g_\varphi(x)) \qquad \text{对 } 1 \le j, k \le m \text{ 及 } x \in \varphi(U) .\tag{5.2}
\end{equation}

\proofhead
因为 $ (\partial_1|_p, \ldots, \partial_m|_p) $ 是 $ T_p M $ 的正向基，故由命题 2.13 推出
\[
\omega_p = \sqrt{G(p)}\, dx^1 \wedge \cdots \wedge dx^m(p) \qquad \text{对 } p \in U .
\]

因此 (5.1) 成立。因为

\[
\varphi_{*}(\omega_M|U) = \varphi_{*}\sqrt{G}\, dx^1 \wedge \cdots \wedge dx^m \,|\, \varphi(U)
\]

且 $ \varphi_{*}\sqrt{G} = \sqrt{G \circ \varphi^{-1}} $，故由注 4.3(a) 推出 (5.2) 成立。由于标量积与行列式函数是光滑的（见命题 VII.4.6 与习题 VII.4.2），且 $ \sqrt{G(p)} > 0 $ 对 $ p \in U $ 成立，链式法则给出 $ \varphi_{*}\sqrt{G} \in \mathcal{E}(\varphi(U)) $。因此 $ \omega_M|U $ 是光滑的，从而证明了 $ \omega_M \in \mathcal{O}r $。 \nobreak\hspace{0.35em}■

\end{proposition}

\begin{remark}

（正则性） 以显然的方式修改本命题的陈述，我们发现当 $ M $ 是 $ C^1 $ 流形时它仍然成立。 \nobreak\hspace{0.35em}■

\end{remark}

\begin{example}

\begin{items}
\item

（$ \overline{\mathbb{H}}^m $ 中的开集） 设 $ X $ 是 $ \overline{\mathbb{H}}^m $ 的非空开子集。则 $ X $ 被赋予一个\textbf{自然定向}，关于它，每个切空间 $ T_p X = T_p\mathbb{R}^m $（$ p \in X $）都自然地定向，也就是说，该定向使典范基 $ ((e_1)_p, \ldots, (e_m)_p) $ 为正向。于是 $ X $ 的体积元由

\[
\omega_X = dx^1 \wedge \cdots \wedge dx^m \,|\, X
\]

给出。平凡坐标卡 $ (\operatorname{id}_X, X) $ 是正向的。

\item

（正则映射的纤维） 设 $ X $ 在 $ \mathbb{R}^m $ 中开，$ q $ 是 $ f \in \mathcal{E}(X) $ 的正则值，且 $ M := f^{-1}(q) \ne \varnothing $。我们按如下方式给超曲面 $ M $ 赋予\textbf{由 $ \nabla f $ 诱导的定向 $ \mathcal{O}r(M, \nabla f) $}：对每个 $ p \in M $，$ T_p M $ 的基 $ (v_1, \ldots, v_{m-1}) $ 为正向，当且仅当 $ T_p X = T_p\mathbb{R}^m $ 的基

\[
(\nabla f(p), v_1, \ldots, v_{m-1})
\]

为正向，其中 $ \nabla f = \sum_{k} \partial_k f\, \partial/\partial x^k $。以

\[
\nu := \nabla f/|\nabla f|
\]

表示 $ M $ 的单位法向量场，则 $ (M, \mathcal{O}r(M, \nabla f)) $ 的体积元由

\[
\omega_M := (\nu \lrcorner \omega_X)|M
\]

给出。

\end{items}

\bookfigure[0.30]{figures/cc7dc4ff5f99c1be.png}

若 $ m = 3 $，则 $ T_p M $ 的基 $ (v_1(p), v_2(p)) $ 为正向，当且仅当三个向量 $ (v_1, v_2, \nu(p)) $ 构成 $ T_p\mathbb{R}^3 $ 的“右手基”。这里 $ (w_1, w_2, w_3) $ 是\textbf{右手基}，是指可以伸出右手的拇指、食指与中指（中指弯向掌心），使这三根手指指向这三个向量的方向（且顺序相同）。这称为\textbf{右手定则}。

\proofhead
因为 $ q $ 是正则点，$ \nabla f(q) \ne 0 $ 对 $ q \in M $ 成立。由正则值定理，$ M $ 是 $ X $ 中的光滑超曲面。例 4.17(e) 的证明表明 $ \omega_M $ 是光滑体积形式。现在一切都很清楚了。 \nobreak\hspace{0.35em}■

\begin{items}
\item[（c）]

（球面） 对 $ m \in \mathbb{N} $，$ \mathbb{R}^{m+1} $ 中的 $ m $ 维球面 $ S^m $ 由外向单位法向量场

\[
\nu(x) := (x, x) \in T_x\mathbb{R}^{m+1}
\]

\textbf{典范地}定向。

\end{items}

\bookfigure[0.30]{figures/b814773309dd3233.png}

若 $ m = 0 $，则 $ S^0 $ 由两个点 $ \{\pm 1\} \subset \mathbb{R} $ 组成，而在 $ 1 $ 处[或在 $ -1 $ 处]的外向单位法向量场由 $ (1,1) \in T_1\mathbb{R} $[或 $ (-1,-1) \in T_{-1}\mathbb{R} $]给出。\footnote{见例 1.17(a)。}

当 $ m = 1 $ 时，$ S^1 $ 的典范定向与注 VIII.5.8 中给出的相同。因此，“沿 $ S^1 $ 行进的方向就是正方向”恰好在行进为逆时针时成立。此时，$ \nu $ 与弗雷内二标架意义下的负单位法向量 $ -\mathbf{n} $ 重合。

典范定向的 $ m $ 维球面的体积元是 $ m $ 形式\footnote{这为例 4.6 与例 4.13(c) 中所用的记号提供了依据。}：

\begin{equation*}
\begin{aligned}
\omega_{S^m} &= (\nu \lrcorner \omega_{\mathbb{R}^{m+1}})|S^m \\ &= \sum_{j=1}^{m+1} (-1)^{j-1} x^j dx^1 \wedge \cdots \wedge \widehat{dx^j} \wedge \cdots \wedge dx^{m+1} \,|\, S^m .
\end{aligned}
\end{equation*}

坐标卡 $ (\varphi_{\pm}, S_{\pm}^m) $ 描述上半[或下半]半球面 $ S_{\pm}^m $，它沿 $ x^{m+1} $ 轴投影到 $ \mathbb{B}^m \times \{0\} $ 上；当 $ m $ 为偶数[或奇数]时，该坐标卡是正向的；当 $ m $ 为奇数[或偶数]时，它是反向的。

$ S^1 $ 的球坐标卡是正向的，而 $ S^2 $ 的球坐标卡是反向的。

\proofhead
$ \omega_{S^m} $ 的公式是 (b) 的特例。关于 $ S^m $ 的各种坐标卡的陈述由例 4.6(a)–(b) 推出。 \nobreak\hspace{0.35em}■

\begin{items}
\item[（d）]

（图像） 设 $ X $ 在 $ \overline{\mathbb{H}}^m $ 中开，$ f \in C^\infty(X, \mathbb{R}^n) $。则图像 $ M := \operatorname{graph}(f) $ 的\textbf{自然定向}是使自然坐标卡 $ (\varphi, M) $ 为正的那个定向，其中

\[
\varphi : M \to \mathbb{R}^m , \quad (x, f(x)) \mapsto x .
\]

在 $ n = 1 $ 的情形，体积元 $ \omega_M $ 有局部表示

\[
\omega_M|M = \sqrt{1 + |\nabla f|^2}\, dx^1 \wedge \cdots \wedge dx^m
\]

其中 $ \nabla f = \sum_{j=1}^{m} \partial_j f\, \partial_j $。

\proofhead
因为 $ g_\varphi(x) = (x, f(x)) $ 对 $ x \in X = \varphi(M) $ 成立，故由注 4.3(a) 推出
\[
\partial_j|_p = (p, (e_j, \partial_j f(x))) \in T_p\mathbb{R}^{m+1} \qquad \text{对 } p = (x, f(x)) \in M \text{ 及 } 1 \le j \le m ,
\]

其中 $ T_p M $ 与 $ T_p\mathbb{R}^{m+1} $ 的向量子空间 $ (T_p i_M)(T_p M) $ 被典范地等同起来。置 $ d_j := \partial_j f $，我们便得到 $ g_{jk} = \delta_{jk} + d_j d_k $。
设 $ D_m := [\varphi_{*}g_{jk}] $。则由此推出

\begin{equation*}
\begin{aligned}
G &= \det D_m = \det\begin{bmatrix} 1+d_1^2 & d_1 d_2 & \cdots & d_1 d_m \\ d_2 d_1 & 1+d_2^2 & \cdots & d_2 d_m \\ \vdots & \vdots & \ddots & \vdots \\ d_m d_1 & d_m d_2 & \cdots & 1+d_m^2 \end{bmatrix} \\ &= \det\begin{bmatrix} D_{m-1} & \begin{matrix} 0 \\ \vdots \\ 0 \end{matrix} \\ \begin{matrix} d_m d_1 & \cdots & d_m d_{m-1} \end{matrix} & 1 \end{bmatrix} + d_m^2 \det\begin{bmatrix} D_{m-1} & \begin{matrix} d_1 \\ \vdots \\ d_{m-1} \end{matrix} \\ \begin{matrix} d_1 & \cdots & d_{m-1} \end{matrix} & 1 \end{bmatrix} .
\end{aligned}
\end{equation*}

为计算最后一个行列式，对 $ 1 \le j \le m-1 $，从第 $ j $ 列减去 $ d_j $ 倍最后一列，于是发现其值为 $ 1 $。这就给出递推公式

\[
\det D_m = \det D_{m-1} + d_m^2 .
\]

因为 $ \det D_1 = 1 + d_1^2 $，递推给出

\[
G = \det D_m = 1 + d_1^2 + \cdots + d_m^2 = 1 + |\nabla f|^2 ,
\]

从而断言得证。 \nobreak\hspace{0.35em}■

\item[（e）]

（曲线） 设 $ J $ 是 $ \mathbb{R} $ 中的完全区间，$ \gamma : J \to \mathbb{R}^m $ 是光滑嵌入。则 $ M := \gamma(J) $ 是 $ \mathbb{R}^m $ 中的嵌入曲线。又设 $ M $ 由 $ \gamma $ \textbf{定向}，即设 $ (\gamma(t), \dot{\gamma}(t)) $ 是 $ T_{\gamma(t)}M $ 的正向基（$ t \in J $）。最后，设 $ \varphi : M \to \mathbb{R} $ 是属于 $ \gamma $ 的 $ M $ 的坐标卡，且 $ \gamma = i_M \circ \varphi^{-1} $。则 $ \omega_M = |\dot{\gamma}|\, dt $。

\proofhead
这是命题 5.1 的直接推论。 \nobreak\hspace{0.35em}■
\item[（f）]

（参数曲面） 设 $ X $ 在 $ \overline{\mathbb{H}}^2 $ 中开，$ h : X \to \mathbb{R}^n $ 是光滑嵌入。则 $ M := h(X) $ 是 $ \mathbb{R}^n $ 中的二维子流形，即 $ \mathbb{R}^n $ 中的一张\textbf{曲面}，它可由单个坐标卡描述。因此 $ M $ 是可定向的。所谓\textbf{由参数化 $ h $ 诱导的定向}，我们指的是使 $ (\partial_1 h(x), \partial_2 h(x)) $ 为 $ T_{h(x)}M $ 的正向基（对每个 $ x \in X $）的那个定向。
设 $ \varphi: M \to \mathbb{R}^{2} $，$ \varphi = (u, v) $ 为属于 $ h $ 的坐标卡，即 $ h = i_{M} \circ \varphi^{-1} $。采用经典记号

\[
\mathrm{E} := |\partial_{1}h|^{2} , \quad \mathrm{F} := (\partial_{1}h \mid \partial_{2}h) , \quad \mathrm{G} := |\partial_{2}h|^{2} ,
\]

则

\[
\omega_{M} = \sqrt{\mathrm{EG} - \mathrm{F}^{2}} du \wedge dv .
\]

\proofhead
这由 $ \varphi_{*}\mathrm{G} = \mathrm{EG} - \mathrm{F}^{2} $ 推出。\nobreak\hspace{0.35em}■
\item[（g）]

（边界）设 $ M $ 是带边的定向流形，$ \nu(p) $ 是 $ \partial M $ 在 $ p \in \partial M $ 处的向外（单位）法向量。若 $ (\nu(p), v_{1}, \ldots, v_{m-1}) $ 是 $ T_{p}M $ 的正基，则称 $ T_{p}\partial M $ 的基 $ (v_{1}, \ldots, v_{m-1}) $ 为正的。这又确定了 $ \partial M $ 上的一个定向，即由外法向量诱导的定向。$ \partial M $ 的体积元 $ \omega_{\partial M} $ 满足

\[
\omega_{\partial M} = (\nu \lrcorner \omega_{M}) \mid \partial M = i^{*}_{\partial M}(\nu \lrcorner \omega_{M}) ,
\]

其中 $ i_{\partial M}: \partial M \hookrightarrow M $ 是自然嵌入。
显然 (c) 是这种情形的特殊情形。还要注意，若 $ \partial M $ 能像 (b) 中那样表示为正则映射的纤维，则由外法向量诱导的定向不必与由 $ \nabla f $ 诱导的定向一致。

\proofhead
由定理 1.15 与注 1.16(a) 知，$ \nu $ 可以局部写成 $ \nu(p) = \nabla f(p)/|\nabla f(p)| $ 的形式，其中 $ f $ 是满足 $ \nabla f(p) \ne 0 $ 的光滑函数。这表明单位法向量场是光滑的。由此容易推出 $ (\nu \lrcorner \omega_{M}) \mid \partial M $ 属于 $ \Omega^{m-1}(\partial M) $。若 $ (v_{1}, \ldots, v_{m-1}) $ 是 $ T_{p}\partial M $ 的标准正交基，则 $ (\nu(p), v_{1}, \ldots, v_{m-1}) $ 是 $ T_{p}M $ 的标准正交基。当 $ (\nu(p), v_{1}, \ldots, v_{m-1}) $ 是 $ T_{p}M $ 的正标准正交基时，
\[
1 = \omega_{M}(\nu(p), v_{1}, \ldots, v_{m-1}) = (\nu \lrcorner \omega_{M})(p)(v_{1}, \ldots, v_{m-1}) .\quad\blacksquare
\]
\item[（h）]

设 $ m \ge 2 $。则 $ \partial\mathbb{H}^{m} = \mathbb{R}^{m-1} $ 由外法向量 $ \nu = -e_{m} $ 诱导的定向与 $ \mathbb{R}^{m-1} $ 的自然定向一致，当且仅当 $ m $ 为偶数。

\proofhead
这可以从下式读出：
\[
\det[\nu, e_{1}, \ldots, e_{m-1}] = (-1)^{m-1}\det[e_{1}, \ldots, e_{m-1}, -e_{m}] = (-1)^{m} .\quad\blacksquare
\]
\end{items}

\end{example}

\needspace{3\baselineskip}\subsection*{黎曼流形}

命题 5.1 的证明依赖于如下事实：每个切空间 $ T_{p}M $ 都自然地带有一个随 $ p \in M $ 可微变化的内积。这种情形出现得相当频繁，尽管 $ TM $ 上的内积常常以另一种方式生成。因此，更精确地考察这些问题是有益的。

$ M $ 上的\textbf{黎曼度量}是张量 $ g \in \mathcal{T}^{0}_{2}(M) $，使得对每个 $ p \in M $，$ g(p) $ 都是 $ T_{p}M $ 上的内积。此时 $ (M, g) $ 称为\textbf{黎曼流形}。

设 $ (M, g) $ 是黎曼流形。我们常把 $ M $ 的坐标卡 $ (\varphi, U) $（其中 $ \varphi = (x^{1}, \ldots, x^{m}) $）写作 $ ((x^{1}, \ldots, x^{m}), U) $。然后令

\begin{equation}
g_{jk} := g\left(\frac{\partial}{\partial x^{j}}, \frac{\partial}{\partial x^{k}}\right) \in \mathcal{E}(U)\tag{5.3}
\end{equation}

并置

\begin{equation}
[g^{jk}] := [g_{jk}]^{-1} \in C^{\infty}(U, \mathbb{R}^{m \times m}_{\mathrm{sym}}) \quad \text{且} \quad G := \det[g_{jk}] \in \mathcal{E}(U) .\tag{5.4}
\end{equation}

我们也把 $ g $ 称为（\textbf{第一}）基本张量。这里 $ [g_{jk}] $ 是 $ g $ 在局部坐标 $ (x^{1}, \ldots, x^{m}) $ 下的表示矩阵（或简称矩阵）；它也称为（\textbf{第一}）基本矩阵\footnote{见注 VII.10.3(b)。}。如前，$ G $ 称为\textbf{格拉姆行列式}。

\begin{remark}

\begin{items}
\item

若 $ g $ 是 $ M $ 上的黎曼度量，则映射

\begin{equation}
\mathcal{V}(M) \times \mathcal{V}(M) \to \mathcal{E}(M) , \quad (v, w) \mapsto g(v, w)\tag{5.5}
\end{equation}

是良定义的、双线性的、对称的，并且是\textbf{正}的，即

\begin{equation}
g(v, v) \ge 0 \quad \text{且} \quad g(v, v) = 0 \iff v = 0 .\tag{5.6}
\end{equation}

\proofhead
映射 (5.5) 是良定义的，这由注 4.18(c) 立即推出。其余断言是内积性质的直接推论。\nobreak\hspace{0.35em}■
\item

设 $ ((x^{1}, \ldots, x^{m}), U) $ 是黎曼流形 $ (M, g) $ 的坐标卡。则

\[
g \mid U = \sum_{j,k} g_{jk} dx^{j} \otimes dx^{k} .
\]

在此上下文中，我们通常把 $ dx^{j} \otimes dx^{k} $ 写作 $ dx^{j}dx^{k} $。

\proofhead
根据 (4.8)，任意 $ v \in \mathcal{V}(U) $ 都有基表示
\[
v = \sum_{j} \langle dx^{j}, v \rangle \frac{\partial}{\partial x^{j}} .
\]

该断言由此、映射 (5.5) 的双线性以及 $ dx^{j} \otimes dx^{k} $ 与 $ g_{jk} $ 的定义推出。\nobreak\hspace{0.35em}■

\item

设 $ ((x^{1}, \ldots, x^{m}), U) $ 是定向黎曼流形 $ (M, g) $ 的正坐标卡。则 $ M $ 的体积元 $ \omega_{M} $ 满足

\[
\omega_{M} \mid U = \sqrt{G} dx^{1} \wedge \cdots \wedge dx^{m} .
\]

\proofhead
这由命题 2.13 推出。\nobreak\hspace{0.35em}■
\item

设 $ g $ 是 $ M $ 上的黎曼度量，$ (x^{1}, \ldots, x^{m}) $ 与 $ (y^{1}, \ldots, y^{m}) $ 是 $ M $ 的开集 $ U $ 上的局部坐标。则

\[
g \mid U = \sum_{j,k} g_{jk} dx^{j} \otimes dx^{k} = \sum_{r,s} \overline{g}_{rs} dy^{r} \otimes dy^{s}
\]

其中

\[
\overline{g}_{rs} = \sum_{j,k} \frac{\partial x^{j}}{\partial y^{r}} \frac{\partial x^{k}}{\partial y^{s}} g_{jk} \quad \text{对 } 1 \le r, s \le m .
\]

\proofhead
因为
\[
dx^{j} = \sum_{r} \frac{\partial x^{j}}{\partial y^{r}} dy^{r} \quad \text{对 } 1 \le r \le m ,
\]

所以这是 (b) 的推论。\nobreak\hspace{0.35em}■

\item

若我们只要求 $ g \in \mathcal{T}^{0}_{2}(M) $ 使得双线性形式 $ g(p) $ 在 $ T_{p}M $ 上对每个 $ p \in M $ 都对称且非退化，则称 $ g $ 为\textbf{不定黎曼度量}，$ (M, g) $ 为\textbf{伪黎曼流形}。此时我们仍使用记号 (5.3) 与 (5.4)。则除 (5.6) 外，(a)、(b) 与 (d) 仍然成立。每个黎曼流形也是伪黎曼流形。

\item

设 $ (M, g) $ 是（伪）黎曼流形，$ W $ 是 $ M $ 中的开集。若 $ v_{1}, \ldots, v_{m} \in \mathcal{V}(W) $ 满足

\[
g(v_{j}, v_{k}) = \pm\delta_{jk} \quad \text{对 } 1 \le j, k \le m ,
\]

则称 $ (v_{1}, \ldots, v_{m}) $ 是 $ W $ 上的\textbf{标准正交标架}。当然，对黎曼流形有 $ g(v_{j}, v_{j}) = 1 $，$ 1 \le j \le m $。因此 $ W $ 上的标准正交标架 $ (v_{1}, \ldots, v_{m}) $ 是 $ W $ 上（光滑）向量场的 $ m $ 元组，它们在每一点 $ p \in W $ 都构成 $ T_{p}M $ 的（不定）内积 $ g(p) $ 下的标准正交基。这种标准正交标架一般并不存在，因为根据注 4.5(d)，一般找不到 $ m $ 个处处线性无关的向量场。
若 $ (\varphi, U) $ 是 $ M $ 的坐标卡，则 $ U $ 上存在标准正交标架。

\proofhead
基向量场 $ \partial_{1}, \ldots, \partial_{m} \in \mathcal{V}(U) $ 在每一点线性无关。由于 $ g $ 非退化，格拉姆–施密特正交化过程（例如见 [Art93, §§ 7.1 与 7.2]）便生成一个标准正交标架。细节留给你们。\nobreak\hspace{0.35em}■
\item

设 $ (M, g) $ 是定向伪黎曼流形。若 $ (\varphi, U) $ 是 $ M $ 的正坐标卡，$ \varphi = (x^{1}, \ldots, x^{m}) $，则令

\[
\omega_{M} \mid U := \sqrt{|G|} dx^{1} \wedge \cdots \wedge dx^{m} .
\]

这便在 $ M $ 上定义了体积形式 $ \omega_{M} \in \Omega^{m}(M) $，我们称之为 $ M $ 的\textbf{体积元}。$ U $ 上的每个正标准正交标架 $ (v_{1}, \ldots, v_{m}) $ 都满足

\[
\omega_{M}(v_{1}, \ldots, v_{m}) = 1 .
\]

\proofhead
我们先证明 $ \omega_{M} \in \Omega^{m}(M) $ 是良定义的。设 $ (e_{1}, \ldots, e_{m}) $ 是 $ U $ 上任一标准正交标架，$ (\varepsilon^{1}, \ldots, \varepsilon^{m}) $ 是其对偶标架；即 $ \varepsilon^{j} \in \Omega^{1}(U) $ 且 $ \langle \varepsilon^{j}, e_{k} \rangle = \delta^{j}_{k} $，$ 1 \le j, k \le m $。则由注 2.18(d) 推出
\[
\varepsilon^{1} \wedge \cdots \wedge \varepsilon^{m} = \sqrt{|G|} dx^{1} \wedge \cdots \wedge dx^{m} .
\]

由于这对 $ U $ 上每个正坐标系 $ (x^{1}, \ldots, x^{m}) $ 都成立，故 $ \omega_{M} \mid U \in \Omega^{m}(U) $ 是良定义的，且与局部坐标的特殊选取无关。现在设 $ \{(\varphi_{\alpha}, U_{\alpha}); \alpha \in A\} $ 是 $ M $ 的正图册，$ (v_{\alpha,1}, \ldots, v_{\alpha,m}) $ 是 $ U_{\alpha} $ 上的正标准正交标架，其对偶标架为 $ (\varepsilon^{1}_{\alpha}, \ldots, \varepsilon^{m}_{\alpha}) $。于是我们在 $ M $ 上定义 $ \omega_{M} $ 为 $ \omega_{M} \mid U_{\alpha} := \varepsilon^{1}_{\alpha} \wedge \cdots \wedge \varepsilon^{m}_{\alpha} $。由前面的讨论推出，$ \omega_{M} $ 是良定义的且属于 $ \Omega^{m}(M) $。最后一个断言现在显然成立。\nobreak\hspace{0.35em}■

\item

（正则性）设 $ k \in \mathbb{N} $，$ M $ 是 $ C^{k+1} $ 流形。则上面的定义与陈述在把 $ \mathcal{V}(M) $ 与 $ \mathcal{E}(M) $ 处处分别换成 $ \mathcal{V}^{k}(M) $ 与 $ C^{k}(M) $ 之后仍然成立。\nobreak\hspace{0.35em}■

\end{items}

设 $ (N, \tilde{g}) $ 是黎曼流形，$ f: M \to N $ 是浸入。则 $ f^{*}\tilde{g} $（$ \tilde{g} $ 经 $ f $ 的拉回）是 $ M $ 上的黎曼度量。若 $ M $ 是 $ N $ 的子流形，$ i: M \hookrightarrow N $ 是自然嵌入，则 $ i^{*}\tilde{g} $ 是\textbf{由 $ N $ 诱导的黎曼度量}（更确切地说，由 $ (N, \tilde{g}) $ 诱导）。

设 $ (M, g) $ 与 $ (N, \tilde{g}) $ 是黎曼流形。若 $ g = f^{*}\tilde{g} $，则称浸入 $ f: M \to N $ 为\textbf{等距}。若 $ f $ 是等距微分同胚，即既是等距又是微分同胚，则 $ M $ 与 $ N $ 是\textbf{等距同构}的。

\end{remark}

\begin{example}

\textbf{
}

\begin{items}
\item

设 $ (M, g) $ 是黎曼流形，$ (\varphi, U) $ 是坐标卡，$ \varphi = (x^{1}, \ldots, x^{m}) $。则 $ (U, g) $ 与 $ (\varphi(U), \varphi_{*}g) $ 等距同构，且

\[
\varphi_{*}g = \sum_{j,k} g_{jk} dx^{j} dx^{k} .
\]

\proofhead
这由基本矩阵的定义立即推出。\nobreak\hspace{0.35em}■
\item

$ \mathbb{R}^{\overline{m}} $ 是带有欧几里得度量 $ g_{\overline{m}} := (\,\cdot \mid \cdot\,) $ 的黎曼流形，该度量称为标准度量。因此 $ \mathbb{R}^{\overline{m}} $ 在 $ M $ 上诱导一个黎曼度量 $ g $，我们也称其为\textbf{标准度量}。它显然与 $ \mathbb{R}^{\overline{m}} $ 无关，其意义是：当 $ M $ 位于 $ \mathbb{R}^{n} $ 中时，$ \mathbb{R}^{n} $ 在 $ M $ 上诱导相同的度量。特别地，对 $ p \in M $ 有 $ g(p) = (\,\cdot \mid \cdot\,)_{p} $（采用我们一直使用的记号），即 $ T_{p}M $ 中由 $ (\,\cdot \mid \cdot\,) $ 诱导的内积。
若 $ (\varphi, U) $ 是 $ M $ 的坐标卡，$ \varphi = (x^{1}, \ldots, x^{m}) $，而 $ h := i \circ \varphi^{-1} $，其中 $ i: M \hookrightarrow \mathbb{R}^{\overline{m}} $ 是相应的参数化，则

\[
\varphi_{*}g = \sum_{j,k} (\partial_{j}h \mid \partial_{k}h) dx^{j} dx^{k} .
\]

换言之，第一基本矩阵 $ [g_{jk}] $ 在局部坐标 $ (x^{1}, \ldots, x^{m}) $ 下由下式给出：

\[
[(\partial_{j}h \mid \partial_{k}h)] \in C^{\infty}(\varphi(U), \mathbb{R}^{m \times m}) .
\]

这与注 5.4(b) 一致，也表明命题 5.1 是注 5.4(c) 的特殊情形。

\proofhead
由 $ g = i^{*}g_{\overline{m}} $ 推出 $ \varphi_{*}g = (\varphi^{-1})^{*}i^{*}g_{\overline{m}} = h^{*}g_{\overline{m}} $。设 $ (y^{1}, \ldots, y^{\overline{m}}) $ 是 $ \mathbb{R}^{\overline{m}} $ 的欧几里得坐标。则
\[
g_{\overline{m}} = \sum_{j=1}^{\overline{m}} (dy^{j})^{2} \quad \text{且} \quad h^{*}g_{\overline{m}} = \sum_{j=1}^{\overline{m}} h^{*}(dy^{j} \otimes dy^{j}) = \sum_{j=1}^{\overline{m}} dh^{j} \otimes dh^{j} ,
\]

它们容易由 $ (0, 2) $ 张量 $ dy^{j} \otimes dy^{j} $ 的拉回定义以及例 3.4(a) 推出。现在该断言容易由 $ (\alpha, \beta) \mapsto \alpha \otimes \beta $（$ \alpha, \beta \in \Omega^{1}(\varphi(U)) $）的双线性以及 $ dh^{j} = \sum_{k} \partial_{k}h^{j} dx^{k} $ 推出。\nobreak\hspace{0.35em}■

\item

（图像）设 $ X $ 是 $ \mathbb{H}^{m} $ 中的开集，$ f \in C^{\infty}(X, \mathbb{R}^{n}) $。设 $ M $ 是 $ f $ 的图像，并令

\[
\varphi: M \to \mathbb{R}^{m} , \quad (x, f(x)) \mapsto x
\]

为自然坐标卡 $ (\varphi, M) $。则 $ M $ 的标准度量 $ g $ 满足

\[
g = \sum_{j} (dx^{j})^{2} + \sum_{j,k} (\partial_{j}f \mid \partial_{k}f) dx^{j} dx^{k} .
\]

特别地，在曲面（$ m = 2 $）的情形，

\[
g = (1 + |\partial_{1}f|^{2})(dx)^{2} + 2(\partial_{1}f \mid \partial_{2}f) dxdy + (1 + |\partial_{2}f|^{2})(dy)^{2} .
\]

\proofhead
因为 $ g_{jk} = \delta_{jk} + (\partial_{j}f \mid \partial_{k}f) $，$ 1 \le j, k \le m $，这由 (b) 推出。\nobreak\hspace{0.35em}■
\item

（参数化曲面）设 $ X $ 是 $ \mathbb{H}^{2} $ 中的开集，$ h: X \to \mathbb{R}^{n} $ 是嵌入。则曲面 $ M := h(X) $ 的标准度量为

\[
g = \mathrm{E}(du)^{2} + 2\mathrm{F} dudv + \mathrm{G}(dv)^{2} ,
\]

其中我们使用了例 5.3(f) 的记号。

\item

（平面极坐标）令 $ V_{2} := (0, \infty) \times (0, 2\pi) $。则极坐标映射

\[
f_{2}: V_{2} \to \mathbb{R}^{2} , \quad (r, \varphi) \mapsto (x, y) := (r \cos \varphi, r \sin \varphi)
\]

是嵌入，$ M := f_{2}(V_{2}) = \mathbb{R}^{2} \setminus (\mathbb{R}^{+} \times \{0\}) $，且

\[
g_{2} \mid M = (dx)^{2} + (dy)^{2} = (dr)^{2} + r^{2}(d\varphi)^{2} .
\]

\proofhead
这容易由 (d) 推出。\nobreak\hspace{0.35em}■
\item

（圆周坐标）关于参数化

\[
h: (0, 2\pi) \to \mathbb{R}^{2} , \quad t \mapsto (\cos t, \sin t)
\]

在 $ S^{1} \setminus \{(1, 0)\} $ 上，圆的标准度量满足 $ g_{S^{1}} = (dt)^{2} $。

\proofhead
因为 $ |\partial h| = 1 $，这由 (b) 推出。\nobreak\hspace{0.35em}■
\item

（$ m $ 维极坐标）设 $ m \ge 3 $，令\footnote{我们使用 (X.8.11)–(X.8.24) 的记号。}

\[
f_{m}: V_{m} \to \mathbb{R}^{m} , \quad (r, \varphi, \vartheta_{1}, \ldots, \vartheta_{m-2}) \mapsto (x^{1}, x^{2}, x^{3}, \ldots, x^{m})
\]

为（限制在 $ V_{m} $ 上的）极坐标映射 (X.8.17)。则 $ f_{m} $ 是 $ \mathbb{R}^{m} \setminus H_{m-1} $ 的参数化，且

\[
\sum_{j=1}^{m} (dx^{j})^{2} = (dr)^{2} + r^{2}\left[ a_{m,0}(d\varphi)^{2} + \sum_{k=1}^{m-2} a_{m,k}(d\vartheta_{k})^{2} \right]
\]

其中

\[
a_{m,k} := \prod_{i=k+1}^{m-2} \sin^{2}\vartheta_{i} \quad \text{对 } 0 \le k \le m-3 \quad \text{且} \quad a_{m,m-2} := 1 .
\]

特别地，球坐标满足（$ m = 3 $）

\[
(dx)^{2} + (dy)^{2} + (dz)^{2} = (dr)^{2} + r^{2}\left[\sin^{2}\vartheta (d\varphi)^{2} + (d\vartheta)^{2}\right] .
\]

\proofhead
设 $ y = (r, z) \in \mathbb{R} \times \mathbb{R}^{m-1} $，由 (X.8.14) 读出
\begin{equation}
\partial_{1}f_{m}(y) = h_{m-1}(z) \quad \text{且} \quad \partial_{j}f_{m}(y) = r\partial_{j-1}h_{m-1}(z) \quad \text{对 } 2 \le j \le m .\tag{5.7}
\end{equation}

因此 (X.8.13) 蕴含

\begin{equation}
|\partial_{1}f_{m}|^{2} = 1 .\tag{5.8}
\end{equation}

对 $ |h_{m-1}|^{2} = 1 $ 求导得 $ (h_{m-1} \mid \partial_{k}h_{m-1}) = 0 $，$ 1 \le k \le m-1 $。于是由 (5.7) 推出

\begin{equation}
(\partial_{1}f_{m}(y) \mid \partial_{k}f_{m}(y)) = r(h_{m-1}(z) \mid \partial_{k-1}h_{m-1}(z)) = 0 \quad \text{对 } 2 \le k \le m .\tag{5.9}
\end{equation}

由 (5.7) 我们还得到

\begin{equation}
(\partial_{j}f_{m}(y) \mid \partial_{k}f_{m}(y)) = r^{2}(\partial_{j-1}h_{m-1}(z) \mid \partial_{k-1}h_{m-1}(z)) \quad \text{对 } 2 \le j, k \le m .\tag{5.10}
\end{equation}

以 $ z = (z', z_{m-1}) \in \mathbb{R}^{m-2} \times \mathbb{R} $ 代入递推公式 (X.8.12) 得到

\begin{equation}
\partial_{j}h_{m-1}(z) = \left(\partial_{j}h_{m-2}(z') \sin z_{m-1}, 0\right) \quad \text{对 } 1 \le j \le m-2 ,\tag{5.11}
\end{equation}

以及

\[
\partial_{m-1}h_{m-1}(z) = \left(h_{m-2}(z') \cos z_{m-1}, -\sin z_{m-1}\right) .
\]

由此及 (X.8.13) 推出

\[
|\partial_{m-1}h_{m-1}(z)|^{2} = |h_{m-2}(z')|^{2}\cos^{2} z_{m-1} + \sin^{2} z_{m-1} = 1 ,
\]

并且与上面类似，我们有

\[
(\partial_{j}h_{m-1}(z) \mid \partial_{m-1}h_{m-1}(z)) = \sin z_{m-1}\cos z_{m-1}(h_{m-2}(z') \mid \partial_{j}h_{m-2}(z')) = 0
\]

对 $ 1 \le j \le m-2 $。由 (5.7)，这证明了

\begin{equation}
|\partial_m f_m(z)|^2 = r^2 \quad \text{且} \quad (\partial_j f_m \mid \partial_m f_m) = 0 \quad \text{对 } 2 \le j \le m-1.\tag{5.12}
\end{equation}

最后，(5.11) 蕴含

\[
(\partial_j h_{m-1}(z) \mid \partial_k h_{m-1}(z)) = \sin^2 z_{m-1} (\partial_j h_{m-2}(z') \mid \partial_k h_{m-2}(z')) \quad \text{对 } 1 \le j \le m-2.
\]

因此 (5.7) 与 (5.10) 给出递推公式

\begin{equation}
(\partial_j f_m(y) \mid \partial_k f_m(y)) = \sin^2 z_{m-1} (\partial_j f_{m-1}(y') \mid \partial_k f_{m-1}(y'))\tag{5.13}
\end{equation}

对 $ 2 \le j,k \le m-1 $，其中 $ y = (y', y_m) \in \mathbb{R}^{m-1} \times \mathbb{R} $。因为 (5.8) 与 (5.12) 对所有 $ m \ge 3 $ 成立，对 (5.13) 作归纳给出

\begin{equation}
|\partial_j f_m|^2 = r^2 a_{m,j-2} \quad \text{对 } 2 \le j \le m-1,\tag{5.14}
\end{equation}

以及

\begin{equation}
(\partial_j f_m \mid \partial_k f_m) = 0 \quad \text{对 } 2 \le j,k \le m-1 \text{ 且 } j \ne k.\tag{5.15}
\end{equation}

现在论断由 (5.8)、(5.9)、(5.12)、(5.14)、(5.15) 与 (b) 推出。\nobreak\hspace{0.35em}■

\item

（m 维球坐标）** 对 $ m \ge 2 $，令

\[
h_m : W_m \to \mathbb{R}^{m+1}, \quad (\varphi, \vartheta_1, \ldots, \vartheta_{m-1}) \mapsto (y^1, y^2, \ldots, y^{m+1}),
\]

其中 $ W_m := (0, 2\pi) \times (0, \pi)^{m-1} $，且

\begin{equation*}
\begin{aligned}
y^1 &= \cos\varphi \sin\vartheta_1 \sin\vartheta_2 \cdots \sin\vartheta_{m-1}, \\ y^2 &= \sin\varphi \sin\vartheta_1 \sin\vartheta_2 \cdots \sin\vartheta_{m-1}, \\ y^3 &= \cos\vartheta_1 \sin\vartheta_2 \cdots \sin\vartheta_{m-1}, \\ & \vdots \\ y^m &= \cos\vartheta_{m-2} \sin\vartheta_{m-1}, \\ y^{m+1} &= \cos\vartheta_{m-1}
\end{aligned}
\end{equation*}

是（$ m $ 维）球坐标。\footnote{见例 VII.9.11(b)。} 于是 $ h_m $ 是 $ m $ 球面的开子集 $ U_m := S^m \setminus H_m $ 的参数化。$ S^m $ 的标准度量 $ g_{S^m} $ 满足

\[
g_{S^m} = a_{m+1,0}(d\varphi)^2 + \sum_{k=1}^{m-1} a_{m+1,k}(d\vartheta_k)^2.
\]

在 2 球面（取 $ \vartheta := \vartheta_1 $）的情形，这成为

\[
g_{S^2} = \sin^2\vartheta (d\varphi)^2 + (d\vartheta)^2.
\]

\proofhead
因为 $ h_m = f_{m+1}(1, \cdot) $，所以论断是 (g) 的简单推论。\nobreak\hspace{0.35em}■
\item

（闵可夫斯基度量）** 我们把 $ \mathbb{R}^4 $ 的欧几里得坐标记作 $ (t, x, y, z) $ 或 $ (x^0, x^1, x^2, x^3) $，并置 $ \mathbb{R}^4_{1,3} := (\mathbb{R}^4, (\cdot\mid\cdot)_{1,3}) $，其中闵可夫斯基度量为

\[
(\cdot\mid\cdot)_{1,3} = (dt)^2 - (dx)^2 - (dy)^2 - (dz)^2 = (dx^0)^2 - \sum_{j=1}^{3}(dx^j)^2.
\]

则 $ \mathbb{R}^4_{1,3} $ 是伪黎曼流形，即（狭义）相对论的\textbf{时空}或\textbf{闵可夫斯基空间}。
对 $ v = (v^0, \ldots, v^3) \subset \mathbb{R}^4_{1,3} $，我们称

\[
|v|^2_{1,3} := (v\mid v)_{1,3} = (v^0)^2 - \sum_{j=1}^{3}(v^j)^2
\]

为向量 $ v $ 的\textbf{闵可夫斯基范数}。具有正闵可夫斯基范数的向量称为\textbf{类时}的；具有负范数的称为\textbf{类空}的。闵可夫斯基范数为零的向量称为\textbf{类光}的；在 $ \mathbb{R}^4_{1,3} $ 中，类光向量构成一个（双）锥，即\textbf{光锥} $ \mathcal{L}_{1,3} $。

\end{items}

\bookfigure[0.30]{figures/371fe58bc2982ef7.png}

\textbf{
}

\begin{items}
\item[（j）]

（伪球坐标）** 令 $ V_{1,3} := \mathbb{R} \times V_3 $，且

\[
f_{1,3} : V_{1,3} \to \mathbb{R}^4, \quad (\rho, \chi, \varphi, \vartheta) \mapsto (x^0, x^1, x^2, x^3)
\]

其中

\begin{equation*}
\begin{aligned}
x^0 &= \rho \cosh\chi, \\ x^1 &= \rho \sinh\chi \cos\varphi \sin\vartheta, \\ x^2 &= \rho \sinh\chi \sin\varphi \sin\vartheta, \\ x^3 &= \rho \sinh\chi \cos\vartheta;
\end{aligned}
\end{equation*}

这就是\textbf{伪球坐标映射}。于是 $ f_{1,3} $ 是从 $ V_{1,3}\setminus\{0\} $ 到光锥的\textbf{内部}

\[
\mathring{\mathcal{L}}_{1,3} := \{ x \in \mathbb{R}^4 ; |x|^2_{1,3} > 0 \}
\]

的光滑微分同胚，且

\[
(\cdot\mid\cdot)_{1,3} = (d\rho)^2 - \rho^2\Bigl[ (d\chi)^2 + \sinh^2\chi \sin^2\vartheta (d\varphi)^2 + \sinh^2\chi (d\vartheta)^2 \Bigr].
\]

\proofhead
这容易由 $ \sinh $ 与 $ \cosh $ 的性质（见习题 III.6.5 与 IV.2.5）以及注 5.4(d) 推出。\nobreak\hspace{0.35em}■
\item[（k）]

（双曲空间）** 为推广闵可夫斯基空间，我们置

\[
(\cdot\mid\cdot)_{1,m} := (dx^0)^2 - \sum_{j=1}^{m}(dx^j)^2
\]

对 $ n \in \mathbb{N}^\times $。则 $ \mathbb{R}^{m+1}_{1,m} := (\mathbb{R}^{m+1}, (\cdot\mid\cdot)_{1,m}) $ 是 $ m $ 维伪黎曼流形。
令

\[
M^m := \{ (x^0, x) \in \mathbb{R} \times \mathbb{R}^m ; (x^0)^2 - |x|^2 = 1, x^0 > 0 \},
\]

即 $ M^m $ 是 $ m $ 维双叶双曲面的上连通分支

\[
K_1 := \{ x \in \mathbb{R}^{m+1} ; (Ax\mid x) = 1 \}, \quad \text{其中 } A := \operatorname{diag}(1, -1, \ldots, -1)
\]

（见例 1.17(b)）。又令 $ i : M^m \hookrightarrow \mathbb{R}^{m+1} $ 为标准嵌入，并令

\[
g_{H^m} := -i^*(\cdot\mid\cdot)_{1,m}.
\]

则

\[
H^m := (M^m, g_{H^m})
\]

是 $ m $ 维黎曼流形，即 $ m $ 维\textbf{双曲空间}。若 $ N := (N, g) $ 与 $ H^m $ 等距同构，则称 $ N $ 是 $ H^m $ 的一个\textbf{模型}。特别地，若给 $ \mathbb{R}^m $ 配上度量

\[
\frac{(dr)^2}{1+r^2} + r^2 g_{S^{m-1}},
\]

写于“极坐标” $ (r, \sigma) \in \mathbb{R}^+ \times S^{m-1} $ 之下，则 $ \mathbb{R}^m $ 是 $ H^m $ 的一个模型。

\proofhead
对 $ u : \mathbb{R}^m \to \mathbb{R}^{m+1} $，$ x \mapsto \sqrt{1+|x|^2} $，有 $ M^m = \operatorname{graph}(u) $。因此
\[
\varphi : M^m \to \mathbb{R}^m, \quad (h(x), x) \mapsto x
\]

是从 $ \mathbb{R}^{m+1} $ 中的超曲面 $ M^m $ 到 $ \mathbb{R}^m $ 的微分同胚。因此我们只需证明由 $ -(\cdot\mid\cdot)_{1,m} $ 在 $ M $ 上诱导的双线性形式 $ g_{H^m}(0) $ 是正定的，且 $ \varphi_* g_{H^m} $ 具有所指出的形式，因为由此可以看出，对每个 $ p \in M \setminus \{\varphi^{-1}(0)\} $，$ g_{H^m}(p) $ 都是正定的。
对 $ x \in \mathbb{R}^m $ 令 $ h(x) := (u(x), x) $，则 $ h = i \circ \varphi^{-1} $，且

\[
\partial_j h = (\partial_j u, e_j) \quad \text{对 } 1 \le j \le m,
\]

其中 $ e_j $ 是 $ \mathbb{R}^m $ 的第 $ j $ 个标准基向量。因为 $ \partial_j u(x) = x^j/u(x) $，所以

\[
(\varphi_* g_{H^m})_{jk}(x) = (\partial_j h \mid \partial_k h)_{1,m}(x) = (\delta_{jk} - x^j x^k)/u^2(x) \quad \text{对 } x \in \mathbb{R}^m;
\]

特别地，$ (\varphi_* g_{H^m})(0) = \sum_j (dx^j)^2 $。
如同 (g) 中那样，令

\[
f_m : (0, \infty) \times W_{m-1} \to \mathbb{R}^m, \quad (r, \vartheta) \mapsto r h_{m-1}(\vartheta)
\]

为 $ m $ 维极坐标映射。则 $ \psi := f_m^{-1} \circ \varphi $ 是 $ M $ 的局部坐标卡，且 $ a := i \circ \psi^{-1} = h \circ f_m = f_m^* h $ 是相伴的参数化。这满足

\[
a(r, \vartheta) = (\sqrt{1+r^2}, r h_{m-1}(\vartheta)),
\]

因此

\[
\partial_r a(r, \vartheta) = \left( \frac{r}{\sqrt{1+r^2}}, h_{m-1}(\vartheta) \right), \quad \partial_{\vartheta^j} a(r, \vartheta) = (0, r\partial_j h_{m-1}(r, \vartheta))
\]

对 $ (r, \vartheta) \in (0, \infty) \times W_{m-1} $。因为 $ |h_{m-1}| = 1 $，我们得到

\[
\psi_* g_{H^m} = -a^*(\cdot\mid\cdot)_{1,m}
\]
\[
= r^2 \sum_{j,k} (\partial_j h_{m-1} \mid \partial_k h_{m-1}) dx^j dx^k + \left( 1 - \frac{r^2}{\sqrt{1+r^2}} \right)(dr)^2.
\]

现在论断由此推出，因为 (h)，并且 $ M^m \setminus \{\varphi^{-1}(0)\} $ 中仍然缺失的部分可以类似地通过绕 $ x^0 $ 轴旋转 $ M^m $ 来参数化。\nobreak\hspace{0.35em}■

\end{items}

\textbf{(l)（庞加莱模型）} 类比于球面到平面的球极投影，考虑\textbf{伪球面}

\[
S^2_{1,3} := \{ (t, x, y) \in \mathbb{R}^3 ; t^2 - x^2 - y^2 = 1 \}
\]

的球极投影。我们置 $ N := (1, 0, 0) $，即 $ S^2_{1,3} $ 的\textbf{北极}，并把\textbf{南极}定义为 $ S := (-1, 0, 0) $。于是\textbf{球极投影} $ s : M^2 \to \mathbb{R}^2 $ 在点 $ p \in M^2 $ 处的值 $ s(p) $ 定义为：从 $ S $ 到 $ p $ 的直线与 $ \mathbb{R}^3 $ 中平面 $ \mathbb{R}^2 \times \{0\} $ 的交点。

\bookfigure[0.40]{figures/d00e8c9752593bfb.png}

若 $ p \in M^2 $ 的（欧几里得）坐标是 $ (t, x, y) $，而 $ s(p) $ 的坐标是 $ (u, v) $，则从上面这张图可以看出

\[
\frac{x}{u} = \frac{t+1}{1} \quad \text{且} \quad \frac{y}{v} = \frac{t+1}{1}.
\]

因为 $ t^2 - x^2 - y^2 = 1 $，所以 $ t^2 - (u^2+v^2)(t+1)^2 = 1 $。由此算得

\[
t = \frac{1+u^2+v^2}{1-u^2-v^2}, \quad x = \frac{2u}{1-u^2-v^2}, \quad y = \frac{2v}{1-u^2-v^2}.
\]

这表明

\[
\pi : \mathbb{B}^2 \to M^2, \quad (u, v) \mapsto \left( \frac{1+u^2+v^2}{1-u^2-v^2}, \frac{2u}{1-u^2-v^2}, \frac{2v}{1-u^2-v^2} \right)
\]

是 $ M^2 $ 在 $ \mathbb{B}^2 $ 上的参数化。它满足

\[
\pi^* g_{H^2} = 4\frac{(du)^2+(dv)^2}{(1-u^2-v^2)^2}.
\]

因此

\[
\left( \mathbb{B}^2, 4\frac{(dx)^2+(dy)^2}{(1-x^2-y^2)^2} \right)
\]

是双曲平面的一个模型，即\textbf{庞加莱模型}。

\proofhead
$ \pi^* g_{H^2} $ 具有给定形式一事留作习题。\nobreak\hspace{0.35em}■

\textbf{(m)（罗巴切夫斯基模型）} 仿照 (h) 与 (j) 中的做法，我们可以用\textbf{伪球坐标}参数化 $ M^2 $：
\[
h_{1,2} : \mathbb{R}^+ \times [0, 2\pi) \to \mathbb{R}^3, \quad (\chi, \varphi) \mapsto (t, x, y)
\]

其中

\[
t = \cosh\chi, \quad x = \sinh\chi \cos\varphi, \quad y = \sinh\chi \sin\varphi.
\]

它们满足 $ h^*_{1,2} g_{H^3} = (d\chi)^2 + \sinh^2\chi (d\varphi)^2 $。因此

\[
(\mathbb{R}^+ \times [0, 2\pi), (d\chi)^2 + \sinh^2\chi (d\varphi)^2)
\]

是双曲平面 $ H^2 $ 的一个模型，即\textbf{罗巴切夫斯基模型}。

\proofhead
验证所给公式同样留作习题。\nobreak\hspace{0.35em}■

\textbf{(n)（一般伪黎曼度量）} 设 $ X $ 在 $ \overline{\mathbb{H}}^m $ 中开，且设 $ g_{jk} = g_{kj} \in \mathcal{E}(X) $ 对 $ 1 \le j,k \le m $ 成立，并且 $ \det[g_{jk}(x)] \ne 0 $ 对 $ x \in X $ 成立。则
\[
g := \sum_{j,k} g_{jk} dx^j dx^k
\]

定义了 $ X $ 上的一个伪黎曼度量。若矩阵 $ [g_{jk}(x)] $ 对每个 $ x \in X $ 都正定，则 $ g $ 是 $ X $ 上的黎曼度量。

现在设 $ \{ (\varphi_\alpha, U_\alpha) ; \alpha \in \mathrm{A} \} $ 是 $ M $ 的图册，

\[
g_{\alpha,jk} = g_{\alpha,kj} \in \mathcal{E}(\varphi_\alpha(U_\alpha)) \quad \text{对 } 1 \le j,k \le m,
\]

且 $ \det[g_{\alpha,jk}(x)] \ne 0 $ 对 $ x \in \varphi_\alpha(U_\alpha) $ 与 $ \alpha \in \mathrm{A} $ 成立。则存在唯一的 $ M $ 上的伪黎曼度量 $ g $，使得

\[
g\mid U_\alpha = g_\alpha := \sum_{j,k} g_{\alpha,jk} dx^j dx^k
\]

只要转移函数 $ h := \varphi_\beta^{-1} \circ \varphi_\alpha $ 满足

\[
g_{\beta,rs} = \sum_{j,k} \frac{\partial h^j}{\partial x^r} \frac{\partial h^k}{\partial x^s} g_{\alpha,jk}
\]

对满足 $ U_\alpha \cap U_\beta \ne \varnothing $ 的 $ \alpha, \beta \in \mathrm{A} $。

\proofhead
这是注 5.4(d) 的推论。\nobreak\hspace{0.35em}■

黎曼流形在每个切空间上都有欧几里得结构，从而可以度量长度与角度。这使得欧几里得几何中的许多概念得以推广。例如，我们在第 VIII.1 节中介绍了

曲线长度的一个公式。现在它可以自然地推广：黎曼流形 $ M $ 上的曲线 $ \gamma: I \to M $ 的长度为
\[
\int_{I} \sqrt{g(\dot\gamma(t), \dot\gamma(t))}\, dt ,
\]

其中 $ \dot\gamma(t) \in T_{\gamma(t)}M $ 是“速度向量”

\[
\dot\gamma(t) = (T_{t}\gamma)(t,1) \qquad \text{对 } t \in I ,
\]

它位于点 $ \gamma(t) $ 处。这里我们不再展开这一主题，因为所提出的问题最好在黎曼几何的框架中处理（不过参见习题 5）。

\end{example}

\needspace{3\baselineskip}\subsection*{霍奇星\footnote{本章其余部分初读时可以跳过。}}

设 $ (M,g) $ 是黎曼流形，$ \omega_{M} $ 是 $ M $ 的体积元。对 $ 0 \le r \le m $，我们定义双线性映射

\begin{equation}
(\cdot \mid \cdot)_{g,r}: \Omega^{r}(M) \times \Omega^{r}(M) \to \mathcal{E}(M)\tag{5.16}
\end{equation}

为

\begin{equation}
(\alpha \mid \beta)_{g,r}(p) := (\alpha(p) \mid \beta(p))_{g(p),r} \qquad \text{对 } p \in M \text{ 且 } \alpha, \beta \in \Omega^{r}(M) ,\tag{5.17}
\end{equation}

其中 $ (\cdot \mid \cdot)_{g(p),r} $ 表示 (2.14) 与 (2.15) 中引入的 $ \bigwedge^{r}T_{p}^{*}M $ 上的标量积。\textbf{霍奇星算子}（或简称霍奇星）

\begin{equation}
\ast: \Omega^{r}(M) \to \Omega^{m-r}(M) , \quad \alpha \mapsto \ast\alpha\tag{5.18}
\end{equation}

也是按点定义的：

\[
(\ast\alpha)(p) := \ast\alpha(p) \qquad \text{对 } p \in M \text{ 且 } \alpha \in \Omega(M) .
\]

\begin{remark}

\begin{items}
\item

映射 (5.16) 是良定义的、双线性的、对称的，且为正。

\proofhead
我们只需证明对 $ \alpha, \beta \in \Omega^{r}(M) $，$ (\alpha \mid \beta)_{g,r} $ 属于 $ \mathcal{E}(M) $，因为其余断言可由 $ (\cdot \mid \cdot)_{g(p),r} $ 的性质推出。因此设 $ (\varphi, U) $ 是 $ M $ 的正向坐标卡。根据注 5.4(f)，我们可以在 $ U $ 上选取定向标准正交标架 $ (v_{1},\ldots,v_{m}) $。设 $ (\eta^{1},\ldots,\eta^{m}) $ 是其对偶标架。则注 3.1(e) 推出
\begin{equation}
\alpha|U = \sum_{(j)\in\mathbb{J}_{r}} \alpha_{(j)}\eta^{j_{1}} \wedge \cdots \wedge \eta^{j_{r}}\tag{5.19}
\end{equation}

其中

\begin{equation}
\alpha_{(j)} = \alpha(v_{j_{1}},\ldots,v_{j_{r}}) \in \mathcal{E}(U) \qquad \text{对 } (j) \in \mathbb{J}_{r} .\tag{5.20}
\end{equation}

现在断言由 (2.14)、(2.15) 与注 2.15(c) 推出。\nobreak\hspace{0.35em}■

\item

星算子是良定义的 $ \mathcal{E}(M) $ 模同构，且

\begin{equation}
\ast\ast\alpha = (-1)^{r(m-r)}\alpha \qquad \text{对 } \alpha \in \Omega^{r}(M) .\tag{5.21}
\end{equation}

\proofhead
因为 (5.21) 由例 2.17(e) 与按点定义 (5.18) 推出，又因为 (5.21) 也表明星算子是双射，所以只需证明 $ \ast\alpha $ 是光滑的。于是设 $ (\varphi, U) $ 是 $ M $ 的正向坐标卡。如同 (a) 的证明，设 $ (v_{1},\ldots,v_{m}) $ 是 $ U $ 上的标准正交标架，并设 $ (\eta^{1},\ldots,\eta^{m}) $ 是其对偶标架。则由 (5.19)、(5.20) 以及例 2.17(d) 中 $ \ast\alpha $ 的显式表示可得 $ \ast\alpha|U \in \mathcal{E}(U) $。\nobreak\hspace{0.35em}■
\item

对 $ \alpha, \beta \in \Omega^{r}(M) $，有

\begin{equation}
\alpha \wedge \ast\beta = \beta \wedge \ast\alpha = (\alpha \mid \beta)_{g,r}\,\omega_{M} .\tag{5.22}
\end{equation}

\proofhead
这可由例 2.17(f) 以及所有涉及运算的按点定义直接推出。\nobreak\hspace{0.35em}■
\item

$ \ast 1 = \omega_{M} $ 且 $ \ast\omega_{M} = 1 $。

\item

（正则性）显然，当 $ M $ 是 $ C^{k+1} $ 流形并将 $ \Omega(M) $ 换为 $ \Omega_{(k)}(M) $ 时，上述断言仍然成立。\nobreak\hspace{0.35em}■

\end{items}

利用星算子的按点定义，我们可以把例 2.17 的其他公式搬到当前情形。下面的例子汇集了若干由此得到的法则。

设 $ ((x^{1},\ldots,x^{m}),U) $ 是 $ M $ 的坐标卡。当 $ (\partial_{1},\ldots,\partial_{m}) $ 是 $ U $ 上的标准正交标架时，我们说 $ (x^{1},\ldots,x^{m}) $ 是 $ U $ 上的\textbf{标准正交坐标}。若 $ \partial_{j} $ 不必是标准化的，即我们只知道当 $ j \ne k $ 时 $ g(\partial_{j},\partial_{k}) = 0 $，则称这些坐标是\textbf{正交的}。

\end{remark}

\begin{example}

在这些例子中，$ (x^{1},\ldots,x^{m}) $ 是 $ U \subset M $ 上的标准正交坐标，且 $ \alpha \in \Omega(U) $。

\begin{items}
\item

欧几里得坐标是标准正交坐标。极坐标、球坐标与伪球坐标是正交坐标。

\proofhead
这由例 5.5 推出。\nobreak\hspace{0.35em}■
\item

$ \ast\sum_{j} a_{j}\,dx^{j} = \sum_{j=1}^{m}(-1)^{j-1}a_{j}\,dx^{1} \wedge \cdots \wedge \widehat{dx^{j}} \wedge \cdots \wedge dx^{m} $。

\item

$ \ast\sum_{j}(-1)^{j-1}a_{j}\,dx^{1} \wedge \cdots \wedge \widehat{dx^{j}} \wedge \cdots \wedge dx^{m} = (-1)^{m-1}\sum_{j}a_{j}\,dx^{j} $。

\item

对 $ m = 3 $，有

\[
\ast d\Bigl(\sum_{j} a_{j}\,dx^{j}\Bigr) = (\partial_{2}a_{3} - \partial_{3}a_{2})\,dx^{1} + (\partial_{3}a_{1} - \partial_{1}a_{3})\,dx^{2} + (\partial_{1}a_{2} - \partial_{2}a_{1})\,dx^{3} .
\]

\proofhead
这由例 3.7(a)（以及 (4.16) 之后的注）推出，因为 (2.20) 蕴含关系
\[
\ast(dx^{2} \wedge dx^{3}) = dx^{1} , \quad \ast(dx^{3} \wedge dx^{1}) = dx^{2} , \quad \ast(dx^{1} \wedge dx^{2}) = dx^{3} .\quad\blacksquare
\]

当然，即使在非标准正交坐标下，我们也可以显式计算 $ \ast\sum_{(j)\in\mathbb{J}_{r}}a_{(j)}\,dx^{(j)} $。为简单起见，我们只考虑 1 形式的情形。

\end{items}

\end{example}

\begin{proposition}

设 $ ((x^{1},\ldots,x^{m}),U) $ 是 $ M $ 的正向坐标卡。则

\[
\ast dx^{j} = \sum_{k}(-1)^{k-1}g^{jk}\sqrt{G}\,dx^{1} \wedge \cdots \wedge \widehat{dx^{k}} \wedge \cdots \wedge dx^{m} .
\]

\proofhead
因为 $ \ast dx^{j} \in \Omega^{m-1}(U) $，例 3.2(b) 保证存在 $ a_{j\ell} \in \mathcal{E}(U) $，使得
\[
\ast dx^{j} = \sum_{\ell}(-1)^{\ell-1}a_{j\ell}\,dx^{1} \wedge \cdots \wedge \widehat{dx^{\ell}} \wedge \cdots \wedge dx^{m} .
\]

这给出

\begin{equation}
\begin{aligned}
dx^{k} \wedge \ast dx^{j} &= \sum_{\ell}(-1)^{\ell-1}a_{j\ell}\,dx^{k} \wedge dx^{1} \wedge \cdots \wedge \widehat{dx^{\ell}} \wedge \cdots \wedge dx^{m} \\ &= a_{jk}\,dx^{1} \wedge \cdots \wedge dx^{m} .
\end{aligned}\tag{5.23}
\end{equation}

由注 2.14(b) 得 $ (dx^{j} \mid dx^{k})_{g,1} = g^{jk} $。于是注 5.6(c) 给出

\begin{equation}
dx^{k} \wedge \ast dx^{j} = g^{kj}\omega_{M} = g^{jk}\sqrt{G}\,dx^{1} \wedge \cdots \wedge dx^{m} ,\tag{5.24}
\end{equation}

其中最后一个等式由注 5.4(c) 推出。现在断言由 (5.23) 与 (5.24) 推出。\nobreak\hspace{0.35em}■

\end{proposition}

\needspace{3\baselineskip}\subsection*{余微分}

设 $ (M,g) $ 是定向黎曼流形。为避免例外情形，我们置 $ \Omega^{-1}(M) := \{0\} $，于是由于 $ \Omega^{m+1}(M) = \{0\} $，我们也可以定义星运算 $ \ast: \Omega^{m+1}(M) \to \Omega^{-1}(M) $。借助（如此延拓的）星算子与外微分，我们对 $ 0 \le r \le m $ 定义\textbf{余微分}

\[
\delta: \Omega^{r}(M) \to \Omega^{r-1}(M)
\]

由\footnote{经常用规范化因子 $ (-1)^{m(r+1)+1} $ 代替 $ (-1)^{m(r+1)} $，特别是在微分几何中。我们这样选择的原因将在注 6.23(c) 中说明。}

\[
\delta\alpha := (-1)^{m(r+1)}\ast d\ast\alpha \qquad \text{对 } \alpha \in \Omega^{r}(M) .
\]

换言之，我们要求图表

\begin{equation*}
\begin{array}{ccc} \Omega^{r}(M) & \xrightarrow{\ast} & \Omega^{m-r}(M) \\ (-1)^{m(r+1)}\delta\downarrow & & \downarrow d \\ \Omega^{r-1}(M) & \xleftarrow{\ast} & \Omega^{m-r+1}(M) \end{array}
\end{equation*}

可交换。

下列注列出余微分的若干性质。

\begin{remark}

\begin{items}
\item

$ \delta^{2} = 0 $。

\proofhead
因为 $ \ast\ast\alpha = (-1)^{r(m-r)}\alpha $，我们有 $ \delta\delta\alpha = \pm\ast d\ast\ast d\ast\alpha = \pm\ast d^{2}\ast\alpha = 0 $，因为 $ d^{2} = 0 $。\nobreak\hspace{0.35em}■
\item

$ \ast\delta d = d\delta\ast $ 且 $ \ast d\delta = \delta d\ast $。

\proofhead
若 $ \alpha \in \Omega^{r}(M) $，则 $ d\alpha $ 属于 $ \Omega^{r+1}(M) $。因此
\[
\ast\delta d\alpha = (-1)^{m(r+2)}\ast\ast d\ast d\alpha = (-1)^{mr}\ast\ast d\ast d\alpha .
\]

因为 $ d\ast d\alpha \in \Omega^{m-r}(M) $，我们于是得到 $ \ast\delta d\alpha = (-1)^{-r^{2}}d\ast d\alpha $。类似地，

\[
d\delta\ast\alpha = (-1)^{m(m-r+1)}d\ast d\ast\ast\alpha = (-1)^{m(m+1)-r^{2}}d\ast d\alpha .
\]

因为 $ m(m+1) $ 是偶数，这证明了第一个断言。第二个断言类似可得。\nobreak\hspace{0.35em}■

\item

$ d\ast\delta = \delta\ast d = 0 $。

\proofhead
我们把这一简单证明留给你。\nobreak\hspace{0.35em}■
\item

对 $ \alpha \in \Omega^{r}(M) $，$ \ast\delta\alpha = (-1)^{r+1}d\ast\alpha $ 且 $ \delta(\ast\alpha) = (-1)^{r}\ast d\alpha $。

\proofhead
第一个断言由
\[
\ast\delta\alpha = (-1)^{m(r+1)}\ast\ast d\ast\alpha = (-1)^{mr+m}(-1)^{(m-r+1)(r-1)}d\ast\alpha = (-1)^{r+1}d\ast\alpha
\]

推出。第二个断言由类似计算推出。\nobreak\hspace{0.35em}■

\item

（正则性）由 $ \delta $ 的定义以及注 4.11(b) 与 5.6(e) 立即推出：对 $ 1 \le r \le m $ 与 $ k \in \mathbb{N}^{\times} $，$ \delta $ 是从 $ \Omega^{r}_{(k)} $ 到 $ \Omega^{r-1}_{(k-1)} $ 的 $ \mathbb{R} $ 线性映射。这对 $ C^{k+1} $ 流形仍然成立。\nobreak\hspace{0.35em}■

\end{items}

\end{remark}

\begin{example}

设 $ ((x^{1},\ldots,x^{m}),U) $ 是 $ M $ 的正向坐标卡。

\begin{items}
\item

对 $ \alpha = \sum_{j} a_{j}\,dx^{j} \in \Omega^{1}(U) $，有

\[
\delta\alpha = \frac{1}{\sqrt{G}}\sum_{j,k}\frac{\partial}{\partial x^{j}}\bigl(g^{jk}a_{k}\sqrt{G}\bigr) \in \mathcal{E}(U) .
\]

\proofhead
由命题 5.8 可得
\[
\ast\alpha = \sum_{j} a_{j}\sum_{k}(-1)^{k-1}g^{jk}\sqrt{G}\,dx^{1} \wedge \cdots \wedge \widehat{dx^{k}} \wedge \cdots \wedge dx^{m} .
\]

由此（因为 $ r = 1 $）我们推出

\begin{equation*}
\begin{aligned}
\delta\alpha = \ast d\ast\alpha &= \ast\sum_{j}\sum_{k}\sum_{\ell}(-1)^{k-1}\frac{\partial}{\partial x^{\ell}}\bigl(a_{j}g^{jk}\sqrt{G}\bigr)\,dx^{\ell} \wedge dx^{1} \wedge \cdots \wedge \widehat{dx^{k}} \wedge \cdots \wedge dx^{m} \\ &= \ast\sum_{j,k}\frac{\partial}{\partial x^{j}}\bigl(g^{jk}a_{k}\sqrt{G}\bigr)\,dx^{1} \wedge \cdots \wedge dx^{m} .
\end{aligned}
\end{equation*}

现在断言由

\begin{equation}
dx^{1} \wedge \cdots \wedge dx^{m} = \frac{1}{\sqrt{G}}\omega_{M}\tag{5.25}
\end{equation}

以及注 5.6(d) 推出。\nobreak\hspace{0.35em}■

\item

对标准正交坐标 $ (x^{1},\ldots,x^{m}) $，由 (a) 推出

\[
\delta\Bigl(\sum_{j}a_{j}\,dx^{j}\Bigr) = \sum_{j}\partial_{j}a_{j} .
\]
\item

对 $ a \in \mathcal{E}(M) $，$ \delta a = 0 $。

\item
\[
\delta(a\,dx^{1} \wedge \cdots \wedge dx^{m}) = \sum_{j,k}(-1)^{k-1}\frac{\partial}{\partial x^{j}}\Bigl(\frac{a}{\sqrt{G}}\Bigr)g^{jk}\sqrt{G}\,dx^{1} \wedge \cdots \wedge \widehat{dx^{k}} \wedge \cdots \wedge dx^{m} .
\]

\proofhead
利用 (5.25) 与注 5.6(d)，我们得到
\[
\ast(a\,dx^{1} \wedge \cdots \wedge dx^{m}) = a/\sqrt{G} .
\]

因此

\[
\ast d\ast(a\,dx^{1} \wedge \cdots \wedge dx^{m}) = \ast\sum_{j}\frac{\partial}{\partial x^{j}}\Bigl(\frac{a}{\sqrt{G}}\Bigr)dx^{j} .
\]

现在断言由命题 5.8 推出。\nobreak\hspace{0.35em}■

\item

在标准正交坐标下，有

\[
\delta\sum_{(j)\in\mathbb{J}_{r}}a_{(j)}\,dx^{(j)} = \sum_{(j)\in\mathbb{J}_{r}}\sum_{k=1}^{r}(-1)^{k-1}\partial_{j_{k}}a_{(j)}\,dx^{1} \wedge \cdots \wedge \widehat{dx^{j_{k}}} \wedge \cdots \wedge dx^{j_{r}} .
\]

\proofhead
由于线性性，只需考虑 $ \alpha = a\,dx^{(j)} $，其中 $ (j) \in \mathbb{J}_{r} $。由 (2.20) 与定理 4.10(ii)，我们得到
\[
d\ast\alpha = s(j)\,da \wedge dx^{(j^{c})} = s(j)\sum_{k=1}^{r}\partial_{j_{k}}a\,dx^{j_{k}} \wedge dx^{(j^{c})} .
\]

因此例 2.17(d) 蕴含

\[
\ast d\ast\alpha = s(j)\sum_{k=1}^{r}s\bigl(j_{k},(j^{c})\bigr)\partial_{j_{k}}a\,dx^{j_{1}} \wedge \cdots \wedge \widehat{dx^{j_{k}}} \wedge \cdots \wedge dx^{j_{r}}
\]

其中 $ s(j_{k},(j^{c})) := \operatorname{sign}\bigl(j_{k},(j^{c}),j_{1},\ldots,\widehat{j_{k}},\ldots,j_{r}\bigr) $。因为 $ (j^{c}) $ 由 $ m-r $ 个元素组成，由此推出

\begin{equation*}
\begin{aligned}
s(j_{k},(j^{c})) &= (-1)^{(m-r)(r-1)}\operatorname{sign}\bigl(j_{k},j_{1},\ldots,\widehat{j_{k}},\ldots,j_{r},(j^{c})\bigr) \\ &= (-1)^{(m-r)(r-1)+k-1}s(j) .
\end{aligned}
\end{equation*}

由于（mod 2）同余

\[
(m-r)(r-1)+k-1+m(r+1) \equiv k-r(r+1)-1 \equiv k-1 ,
\]

断言于是由 $ \delta $ 的定义推出。\nobreak\hspace{0.35em}■

\end{items}

\end{example}

\begin{remark}

\begin{items}
\item

通过作适当修改，上述星算子与余微分的性质可以推广到伪黎曼流形的情形。
更确切地说，设 $ (M,g) $ 是定向伪黎曼流形。我们可以给 $ T_{p}M $ 配备由 $ \mathbb{R}^{m} $ 的内积诱导的内积。由注 2.18(a) 推出，$ g(p) $ 在任意 $ p \in M $ 处的表示矩阵 $ g $ 在适当选择的基下是对角的，且其对角元为 $ \pm 1 $。现在 $ (-1)^{s} = \operatorname{sign}g(p) $ 由 $ g(p) $ 唯一确定，其中 $ s $ 表示负元素的个数。我们现在假设 $ \operatorname{sign}(g) = \operatorname{sign}g(p) $ 在 $ M $ 上是常数，即它与 $ p $ 无关。由注 5.4(f)（的证明）推出，若 $ M $ 可由单个坐标卡描述，则该假设成立。
在此假设下，通过 (2.25) 定义的星算子也可以按点定义。此时 (5.21) 与 (5.22) 必须分别换为

\[
\ast\ast\alpha = \operatorname{sign}(g)(-1)^{r(m-r)}\alpha \qquad \text{对 } \alpha \in \Omega^{r}(M) ,
\]

以及

\[
\alpha \wedge \ast\beta = \beta \wedge \ast\alpha = \operatorname{sign}(g)(\alpha \mid \beta)_{g,r}\,\omega_{M} \qquad \text{对 } \alpha, \beta \in \Omega^{r}(M) ,
\]

正如我们从注 2.19(d) 与 (e) 分别得知的那样。这里 $ \omega_{M} $ 是注 5.4(g) 中定义的 $ M $ 的体积元。注 2.19(c) 还蕴含

\[
\ast 1 = \operatorname{sign}(g)\omega_{M} \qquad \text{对 } \ast\omega_{M} = 1 .
\]

在此情形下，余微分定义为

\begin{equation}
\delta\alpha := \operatorname{sign}(g)(-1)^{m(r+1)}\ast d\ast\alpha \qquad \text{对 } \alpha \in \Omega^{r}(M) .\tag{5.26}
\end{equation}

我们容易验证，经过这些修改，注 5.6 的断言仍按原文成立。

\proofhead
这些断言由注 2.19 推出。\nobreak\hspace{0.35em}■
\end{items}

(b) 设 $ ((x^{1},\ldots,x^{m}),U) $ 是 $ M $ 的正向坐标卡。则

\[
\delta\sum_{j}a_{j}\,dx^{j} = \frac{1}{\sqrt{|G|}}\sum_{j,k}\frac{\partial}{\partial x^{j}}\bigl(g^{jk}a_{k}\sqrt{|G|}\bigr) \in \mathcal{E}(U) .
\]

\proofhead
这可由注 5.4(g) 推出，类似于例 5.10(a) 的证明。\nobreak\hspace{0.35em}■
\end{remark}

\begin{example}

我们现在考虑闵可夫斯基空间 $ \mathbb{R}^{4}_{1,3} $，其度量为 $ (\cdot \mid \cdot)_{1,3} $，因此 $ g := (dt)^{2} - (dx)^{2} - (dy)^{2} - (dz)^{2} $。

\begin{items}
\item

若 $ (i,j,k) $ 是 $ (1,2,3) $ 的循环置换，则令 $ (x^{1},x^{2},x^{3}) := (x,y,z) $，有

\[
\ast(dx^{i} \wedge dt) = dx^{j} \wedge dx^{k} \qquad \text{与} \qquad \ast(dx^{i} \wedge dx^{j}) = -dx^{k} \wedge dt .
\]

\proofhead
设 $ (e_{0},e_{1},e_{2},e_{3}) $ 是 $ \mathbb{R}^{4}_{1,3} $ 的标准基。则
\[
g(e_{0},e_{0}) = 1 \qquad \text{与} \qquad g(e_{j},e_{j}) = -1 \qquad \text{对 } 1 \le j \le 3 .
\]

这意味着

\[
(dt \mid dt)_{g,1} = 1 \qquad \text{与} \qquad (dx^{j} \mid dx^{j})_{g,1} = -1 .
\]

因此

\[
(dt \wedge dx^{j} \mid dt \wedge dx^{j})_{g,2} = -1 \qquad \text{与} \qquad (dx^{j} \wedge dx^{k} \mid dx^{j} \wedge dx^{k})_{g,2} = 1 \qquad \text{对 } 1 \le j < k \le 3 .
\]

现在断言由注 2.19(c) 推出。\nobreak\hspace{0.35em}■

\item

设 $ E_{j}, H_{j} \in \mathcal{E}(\mathbb{R}^{4}_{1,3}) $ 且

\begin{equation*}
\begin{aligned}
\alpha &:= (E_{1}dx^{1} + E_{2}dx^{2} + E_{3}dx^{3}) \wedge dt \\ &\quad + H_{1}dx^{2} \wedge dx^{3} + H_{2}dx^{3} \wedge dx^{1} + H_{3}dx^{1} \wedge dx^{2} .
\end{aligned}
\end{equation*}

则

\begin{equation*}
\begin{aligned}
\ast\alpha &= -(H_{1}dx^{1} + H_{2}dx^{2} + H_{3}dx^{3}) \wedge dt \\ &\quad + E_{1}dx^{2} \wedge dx^{3} + E_{2}dx^{3} \wedge dx^{1} + E_{3}dx^{1} \wedge dx^{2} .
\end{aligned}
\end{equation*}

\proofhead
这是 (a) 与星算子的 $ \mathcal{E}(\mathbb{R}^{4}_{1,3}) $ 线性性的直接推论。\nobreak\hspace{0.35em}■
\item

(b) 中的 $ \alpha $ 满足

\[
\delta\alpha = \sum_{j=1}^{3}\frac{\partial E_{j}}{\partial t}\,dx^{j} - \sum_{k=1}^{3}\frac{\partial E_{k}}{\partial x^{k}}\,dt + \sum_{(i,j,k)}\Bigl(\frac{\partial H_{i}}{\partial x^{j}} - \frac{\partial H_{j}}{\partial x^{i}}\Bigr)dx^{k} ,
\]

其中最后一项对所有 $ (1,2,3) $ 的循环置换求和。

\proofhead
由 (b) 可知
\[
\ast\alpha = -\sum_{i=1}^{3}H_{i}\,dx^{i} \wedge dt + \sum_{(i,j,k)}E_{i}\,dx^{j} \wedge dx^{k} .
\]

这意味着

\begin{equation*}
\begin{aligned}
d\ast\alpha &= -\sum_{i=1}^{3}\sum_{\substack{j=1 \\ j \ne i}}^{3}\frac{\partial H_{i}}{\partial x^{j}}\,dx^{j} \wedge dx^{i} \wedge dt \\ &\quad + \sum_{(i,j,k)}\Bigl(\frac{\partial E_{i}}{\partial t}\,dt \wedge dx^{j} \wedge dx^{k} + \frac{\partial E_{i}}{\partial x^{i}}\,dx^{i} \wedge dx^{j} \wedge dx^{k}\Bigr) .
\end{aligned}
\end{equation*}

由注 2.19(c)，我们得到

\[
\ast(dt \wedge dx^{i} \wedge dx^{j}) = -dx^{k} \qquad \text{且} \qquad \ast(dx^{i} \wedge dx^{j} \wedge dx^{k}) = dt .
\]

由此得到

\[
\ast d\ast\alpha = -\sum_{(i,j,k)} \left( \frac{\partial H_{i}}{\partial x^{j}} - \frac{\partial H_{j}}{\partial x^{i}} \right) dx^{k} - \sum_{i=1}^{3} \frac{\partial E_{i}}{\partial t} dx^{i} + \sum_{k=1}^{3} \frac{\partial E_{k}}{\partial x^{k}} dt .
\]

现在断言得证，因为 $ m = 4 $ 且 $ \operatorname{sign}(g) = -1 $。\nobreak\hspace{0.35em}■

\end{items}

\end{example}

\exerciseshead

\begin{exercises}
\item

设 $ (M_{j}, g_{j}) $（$ j = 1, 2 $）是伪黎曼流形，$ \partial M_{1} = \varnothing $，并记 $ \pi_{j} : M_{1} \times M_{2} \to M_{j} $ 为到 $ M_{j} $ 的典范投影。证明下列断言：

\begin{subitems}
\item

$ (M_{1} \times M_{2}, \pi_{1}^{\ast} g_{1} + \pi_{2}^{\ast} g_{2}) $ 是黎曼流形，即 $ M_{1} $ 与 $ M_{2} $ 的\textbf{积}。

\item

两个点 $ (p_{1}, p_{2}) $ 给出 $ M_{1} \times M_{2} $ 的子流形 $ M_{1} \times \{p_{2}\} $ 与 $ \{p_{1}\} \times M_{2} $。

\item

$ T_{(p_{1},p_{2})}(M_{1} \times M_{2}) = T_{(p_{1},p_{2})}(M_{1} \times \{p_{2}\}) \oplus T_{(p_{1},p_{2})}(\{p_{1}\} \times M_{2}) $。

\item

$ \omega_{M_{1} \times M_{2}} = \pi_{1}^{\ast}\omega_{M_{1}} \wedge \pi_{2}^{\ast}\omega_{M_{2}} $。

\end{subitems}

\item

设 $ M $ 是 $ \mathbb{R}^{m+1} $ 中的定向超曲面。若 $ \nu $ 是 $ M $ 的单位法向量，且使得对每个 $ p \in M $ 与 $ T_{p}M $ 的每个正基 $ (v_{1}, \ldots, v_{m}) $，$ (m+1) $ 元组 $ (\nu(p), v_{1}, \ldots, v_{m}) $ 都是 $ T_{p}\mathbb{R}^{m+1} $ 的正基，则我们称 $ \nu : M \to T\mathbb{R}^{m+1} $ 为\textbf{正单位法向量场}。

\begin{items}
\item

\proofhead
$ \nu $ 是良定义且唯一的。
\item

确定 $ \mathbb{R}^{3} $ 中下列曲面上的单位法向量：

\end{items}

(i) $ f $ 的图像，其中 $ X $ 是 $ \mathbb{R}^{2} $ 中的开集且 $ f \in \mathcal{E}(X) $，(ii) $ \mathbb{R} \times S^{1} $，(iii) $ S^{2} $，(iv) $ \mathbb{T}^{2}_{a,r} $。

（提示：(iv) 习题 VII.10.10 与例 VII.9.11(f)。）

\item

设 $ M $ 是 $ \mathbb{R}^{m+1} $ 中的定向超曲面，并配有标准度量。记 $ \nu $ 为 $ M $ 的正单位法向量。证明下列事实：

\begin{subitems}
\item

$ \nu $ 定义了从 $ M $ 到 $ S^{m} $ 的光滑映射，即\textbf{高斯映射}（也记作 $ \nu $）。

\item

对 $ p \in M $ 与 $ v \in T_{p}M $，有 $ ((T_{p}\nu)v \mid \nu(p))_{\mathbb{R}^{m+1}} = 0 $。因此 $ (T_{p}\nu)v $ 属于 $ T_{p}M $。

\item

映射

\[
L : M \to \bigcup_{p\in M} \mathcal{L}(T_{p}M) \in \mathcal{L}(T_{p}M) , \quad p \mapsto T_{p}\nu
\]

是良定义的。它称为 $ M $ 的\textbf{魏因加滕映射}。

\item

对 $ p \in M $ 与 $ v, w \in T_{p}M $，有

\[
g(p)(L(p)v, w) = g(p)(v, L(p)w) ,
\]

即 $ L(p) $ 在内积空间 $ (T_{p}M, g(p)) $ 上是对称的。由

\[
h(p)(v, w) := g(p)(L(p)v, w) \qquad \text{对 } p \in M \text{ 与 } v, w \in T_{p}M
\]

定义的张量 $ h \in \mathcal{T}_{2}^{0}(M) $ 称为 $ M $ 的\textbf{第二基本张量}。

\item

在局部坐标 $ (U, \varphi) $ 下，取自然嵌入 $ i : M \hookrightarrow \mathbb{R}^{m+1} $，并令 $ f := i \circ \varphi^{-1} $，有

\[
h_{jk} = (\partial_{j}\nu \mid \partial_{k}f) = -(\nu \mid \partial_{j}\partial_{k}f) ,
\]

其中 $ h_{jk} := h(\partial/\partial x^{j}, \partial/\partial x^{k}) $。

\end{subitems}

\item

计算 $ \mathbb{R}^{2} $、$ S^{2} $、$ \mathbb{R} \times S^{1} $ 与 $ \mathbb{T}^{2}_{a,r} $ 作为 $ \mathbb{R}^{3} $ 的子流形的第二基本形式。

\item

设 $ I $ 是 $ \mathbb{R} $ 中的紧区间，$ M $ 是黎曼流形。又设 $ \gamma \in C^{1}(I, M) $。令 $ i : M \hookrightarrow \mathbb{R}^{\overline{m}} $ 为自然嵌入，并置 $ \widetilde{\gamma} := i \circ \gamma $。于是 $ \widetilde{\gamma} $ 的长度 $ L(\widetilde{\gamma}) $ 按第 VIII.1 节定义。证明：若对 $ t \in I $ 有 $ \dot{\gamma}(t) := (T_{t}\gamma)(t, 1) $，则

\[
L(\widetilde{\gamma}) = \int_{I} \sqrt{g(\dot{\gamma}(t), \dot{\gamma}(t))} \, dt .
\]

当 $ L(\widetilde{\gamma}) = L(I) $ 时，我们称 $ \gamma $ 是\textbf{按弧长参数化}的。

\item

设 $ M $ 是 $ \mathbb{R}^{3} $ 中的定向曲面，$ \gamma \in C^{2}(I, M) $ 按弧长参数化。又记 $ \nu $ 为 $ M $ 的正单位法丛。则我们称

\[
\kappa_{g}(\gamma) := \det[\dot{\gamma}, \ddot{\gamma}, \nu]
\]

为 $ \gamma $ 在 $ M $ 中的\textbf{曲率}，或 $ \gamma $ 的\textbf{测地曲率}。

\begin{items}
\item

在欧氏情形 $ M = \mathbb{R}^{2} $ 中验证测地曲率与第 VIII.2 节中的（通常）曲率相同。

\item

设 $ M = S^{2} $，$ (x, y, z) $ 是 $ \mathbb{R}^{3} $ 中的欧氏坐标。又设对 $ z \in (-1, 1) $，$ \gamma_{z} $ 是 $ L_{z} := \sqrt{1 - |z|^{2}}\, S^{1} \times \{z\} $ 的按弧长参数化（见例 1.5(a)）。证明

\[
\kappa_{g}(\gamma_{z}) = \frac{z}{\sqrt{1 - |z|^{2}}} .
\]

因此测地曲率在圆 $ L_{z} $ 上为常数，并在赤道处为零。

\end{items}

\item

对 $ m = 2 $ 与 $ 3 $ 直接计算，证明例 4.13(c) 中的等式 $ \rho^{\ast}\omega_{S^{m}} = r^{-(m+1)}\alpha $。

\item

证明例 5.5(b) 的证明中所作的陈述。

\item

证明

\[
\{ z \in \mathbb{C} ; \operatorname{Im} z > 0 \} \to \mathbb{D} , \quad z \mapsto (1 + iz)/(1 - iz)
\]

给出从“复半平面的上半部分”到单位圆盘的一个微分同胚。然后用此映射证明

\[
\left( H^{2}, \frac{(dx)^{2} + (dy)^{2}}{y^{2}} \right)
\]

是双曲平面的一个模型，即\textbf{克莱因模型}。

\item

证明例 5.5(m) 中的罗巴切夫斯基平面是 $ H^{2} $ 的一个模型。

\item

对 $ \alpha \in \Omega^{r-1}(M) $ 与 $ \beta \in \Omega^{r}(M) $，证明

\[
d(\alpha \wedge \ast\beta) = d\alpha \wedge \ast\beta + \alpha \wedge \ast\delta\beta .
\]
\item

证明余微分不依赖于所在黎曼流形的定向。

\item

设 $ M $ 是定向的，$ (N, \widetilde{g}) $ 是另一个定向的 $ m $ 维黎曼流形，且 $ f : M \to N $ 是等距微分同胚。证明 $ f^{\ast}\omega_{N} = \pm\omega_{M} $。并证明 $ f^{\ast}\omega_{N} = \omega_{M} $ 当且仅当 $ f $ 是保定向的。

\item

设 $ M $ 与 $ N $ 如习题 13 所述，且 $ f : M \to N $ 是保定向的等距微分同胚。证明下面的图

\begin{equation*}
\begin{array}{ccc} \Omega^{r}(M) & \xrightarrow{\ast} & \Omega^{m-r}(M) \\ f^{\ast}\uparrow & & \uparrow f^{\ast} \\ \Omega^{r}(N) & \xrightarrow{\ast} & \Omega^{m-r}(N) \end{array}
\end{equation*}

对 $ 0 \le r \le m $ 交换。

\item

设 $ M $ 与 $ N $ 如习题 13 所述，且 $ f : M \to N $ 是等距微分同胚。\footnote{参见习题 12。}证明下面的图

\begin{equation*}
\begin{array}{ccc} \Omega^{r}(M) & \xrightarrow{\delta} & \Omega^{r-1}(M) \\ f^{\ast}\uparrow & & \uparrow f^{\ast} \\ \Omega^{r}(N) & \xrightarrow{\delta} & \Omega^{r-1}(N) \end{array}
\end{equation*}

对 $ 0 \le r \le m $ 交换。

\end{exercises}

\section{向量分析}

向量场与普法夫形式可以借助里斯同构互相转换。向量场具有直接的几何解释，而微分形式的演算在计算中极有价值。外积与外微分服从相对简单的规则，而这些规则本身代表了一套更复杂的、关于如何从一个局部坐标系过渡到另一个局部坐标系的法则。在本节中，我们将用里斯同构把微分形式的一些概念与定理转译成经典向量分析的语言。在此过程中，我们将了解向量场的散度与旋度，它们在物理学与偏微分方程理论中具有根本的重要性。

本节自始至终作如下假定：

- $ M $ 是 $ m $ 维流形；$ N $ 是 $ n $ 维流形。
- 指标 $ i, j, k, \ell $ 除另有说明外总是从 $ 1 $ 取到 $ m $，而 $ \sum_{j} $ 表示 $ j $ 从 $ 1 $ 到 $ m $ 求和。

\needspace{3\baselineskip}\subsection*{里斯同构}

设 $ g $ 是 $ M $ 上的伪黎曼度量。则我们如下定义\textbf{里斯同构} $ \Theta_{g} $：

\begin{equation}
\Theta_{g}: \mathcal{V}(M) \to \Omega^{1}(M), \quad v \mapsto \Theta_{g}v\tag{6.1}
\end{equation}

以及

\[
(\Theta_{g}v)(p) := \Theta_{g(p)}v(p) \quad \text{对 } p \in M,
\]

其中 $ \Theta_{g(p)}: T_{p}M \to T_{p}^{*}M $ 是 (2.12)（或注 2.18(b)）中的里斯同构，它由下式定义：

\[
\langle \Theta_{g(p)}u, w \rangle = g(p)(u, w) \quad \text{对 } u, w \in T_{p}M.
\]

在不致引起混淆时，我们可以用 $ \Theta $ 代替 $ \Theta_{g} $。

\begin{remark}

\begin{items}
\item

映射 (6.1) 是良定义的。

\proofhead
我们须证明，对 $ v \in \mathcal{V}(M) $，$ \Theta v $ 属于 $ \Omega^{1}(M) $。在局部坐标下，有
\[
v|U = \sum_{j} v^{j}\frac{\partial}{\partial x^{j}}
\]

其中 $ v^{j} \in \mathcal{E}(U) $。由此以及注 2.14(a) 与 2.18(b) 可得

\[
\Theta v(p) = \Theta_{g(p)}\sum_{j} v^{j}(p)\frac{\partial}{\partial x^{j}}\Big|_{p} = \sum_{j} v^{j}(p)\Theta_{g(p)}\frac{\partial}{\partial x^{j}}\Big|_{p} = \sum_{j} v^{j}(p)\sum_{k} g_{jk}(p)\,dx^{k}(p) = \sum_{j} a_{j}(p)\,dx^{j}(p),
\]

其中

\[
a_{k} := \sum_{j} g_{kj}v^{j} \in \mathcal{E}(U).
\]

现在由注 4.5(c) 与 5.4(e) 可得结论。\nobreak\hspace{0.35em}■

\item

在局部坐标下，

\begin{equation}
\Theta\Big(\sum_{j} v^{j}\frac{\partial}{\partial x^{j}}\Big) = \sum_{j} a_{j}\,dx^{j} \qquad \text{其中 } a_{j} := \sum_{k} g_{jk}v^{k}.\tag{6.2}
\end{equation}

我们常把 $ \Theta v $ 写成 $ v^{\flat} $ 或 $ g^{\flat}v $，因为如 (6.2) 所示，$ \Theta $ 实现的是“降低指标”（见注 2.14(d)）。

\proofhead
这已在 (a) 的证明中说明。\nobreak\hspace{0.35em}■
\item

映射 $ \Theta: \mathcal{V}(M) \to \Omega^{1}(M) $ 是 $ \mathcal{E}(M) $-模同构。

\proofhead
设 $ \alpha \in \Omega^{1}(M) $。则对 $ p \in M $ 有 $ \alpha(p) \in T_{p}^{*}M $。由第 2 节我们知道 $ \Theta_{g(p)} $ 是向量空间同构。因此 $ \Theta_{g(p)}^{-1}\alpha(p) \in T_{p}M $ 是良定义的。我们置
\[
(\overline{\Theta}_{g}\alpha)(p) := \Theta_{g(p)}^{-1}\alpha(p) \quad \text{对 } p \in M \text{ 与 } \alpha \in \Omega^{1}(M).
\]

在局部坐标下，由注 2.14(a) 与 2.18(b) 知

\[
\overline{\Theta}_{g}\alpha(p) = \Theta_{g(p)}^{-1}\sum_{j} a_{j}(p)\,dx^{j}(p) = \sum_{j} a_{j}(p)\Theta_{g(p)}^{-1}dx^{j}(p) = \sum_{j} a_{j}(p)\sum_{k} g^{jk}(p)\frac{\partial}{\partial x^{k}}\Big|_{p} = \sum_{j} v^{j}(p)\frac{\partial}{\partial x^{j}}\Big|_{p},
\]

其中

\[
v^{j} := \sum_{k} g^{jk}a_{k} \in \mathcal{E}(U).
\]

于是由注 4.3(c) 可知 $ \overline{\Theta}_{g}\alpha $ 属于 $ \mathcal{V}(M) $。由 $ \Theta_{g} $ 与 $ \overline{\Theta}_{g} $ 的定义立即得到 $ \Theta_{g}\overline{\Theta}_{g} = \operatorname{id}_{\Omega^{1}(M)} $ 与 $ \overline{\Theta}_{g}\Theta_{g} = \operatorname{id}_{\mathcal{V}(M)} $。因此 $ \Theta_{g} $ 是双射，且 $ \Theta_{g}^{-1} = \overline{\Theta}_{g} $。
最后我们看到，对 $ a \in \mathcal{E}(M) $ 与 $ v \in \mathcal{V}(M) $，有

\[
\Theta_{g}(av)(p) = \Theta_{g(p)}a(p)v(p) = a(p)\Theta_{g(p)}v(p) = (a\Theta_{g}v)(p) \quad \text{对 } p \in M.
\]

因此 $ \Theta $ 是 $ \mathcal{E}(M) $-模同构。\nobreak\hspace{0.35em}■

\item

在局部坐标下，

\begin{equation}
\Theta^{-1}\Big(\sum_{j} a_{j}\,dx^{j}\Big) = \sum_{j} v^{j}\frac{\partial}{\partial x^{j}} \qquad \text{其中 } v^{j} := \sum_{k} g^{jk}a_{k}.\tag{6.3}
\end{equation}

我们常把 $ \Theta^{-1}\alpha $ 写成 $ \alpha^{\sharp} $ 或 $ g^{\sharp}\alpha $，因为 $ \Theta^{-1} $ “提升指标”。

\proofhead
这已在 (c) 的证明中说明。\nobreak\hspace{0.35em}■
\item

（正交坐标）若 $ (x^{1}, \dots, x^{m}) $ 是正交坐标，即若

\[
g\Big(\frac{\partial}{\partial x^{j}}, \frac{\partial}{\partial x^{k}}\Big) = 0 \quad \text{对 } j \neq k,
\]

则 (6.2) 与 (6.3) 分别简化为

\[
\Theta v = \sum_{j} g_{jj}v^{j}\,dx^{j} \quad \text{与} \quad \Theta^{-1}\alpha = \sum_{j} g^{jj}a_{j}\frac{\partial}{\partial x^{j}}
\]

其中 $ v = \sum_{j} v^{j}\partial/\partial x^{j} $，$ \alpha = \sum_{j} a_{j}\,dx^{j} $。

\item

设 $ (N, \bar{g}) $ 是伪黎曼流形，又设 $ \varphi \in \operatorname{Diff}(M, N) $ 且对某个 $ \lambda \neq 0 $ 有 $ \varphi_{*}g = \lambda\bar{g} $。则下图交换。

\begin{equation*}
\begin{array}{ccc} \mathcal{V}(M) & \xleftarrow[\ \cong\ ]{\ \varphi^{*}\ } & \mathcal{V}(N) \\ \Theta_{M}\Big\downarrow\ \cong & & \cong\ \Big\downarrow\ \lambda\Theta_{N} \\ \Omega^{1}(M) & \xleftarrow[\ \cong\ ]{\ \varphi^{*}\ } & \Omega^{1}(N) \end{array}
\end{equation*}

因此 $ \Theta_{M}\varphi^{*} = \lambda\varphi^{*}\Theta_{N} $。

\proofhead
利用向量场与形式的前推和拉回的定义与性质（特别见 (4.25)），我们对 $ v, w \in \mathcal{V}(N) $ 得到
\[
\lambda\bar{g}(v, w) = \varphi_{*}g(v, w) = g(\varphi^{*}v, \varphi^{*}w) = \langle \Theta_{M}\varphi^{*}v, \varphi^{*}w \rangle_{M} = \langle \varphi_{*}\Theta_{M}\varphi^{*}v, w \rangle_{N} = \bar{g}(\Theta_{N}^{-1}\varphi_{*}\Theta_{M}\varphi^{*}v, w).
\]

由于 $ \bar{g} $ 非退化且为 $ \mathbb{R} $-线性的，由此可得

\[
\lambda v = \Theta_{N}^{-1}\varphi_{*}\Theta_{M}\varphi^{*}v \quad \text{对 } v \in \mathcal{V}(M),
\]

这就证明了结论。\nobreak\hspace{0.35em}■

\item

（正则性）设 $ k \in \mathbb{N} $ 且 $ M $ 是 $ C^{k+1} $ 流形。则当把光滑向量场、微分形式与函数分别换成 $ C^{k} $ 向量场、$ C^{k} $ 微分形式与 $ C^{k} $ 函数时，上述定义与陈述仍然成立。\nobreak\hspace{0.35em}■

\end{items}

\end{remark}

\begin{example}

\begin{items}
\item

（欧氏坐标）设 $ M $ 是 $ \mathbb{R}^{m} $ 中的开集。我们用 $ (x^{1}, \dots, x^{m}) $ 表示欧氏坐标，即 $ (\cdot|\cdot) = \sum_{j}(dx^{j})^{2} $。则

\[
\Theta\Big(\sum_{j} v^{j}\frac{\partial}{\partial x^{j}}\Big) = \sum_{j} a_{j}\,dx^{j} \qquad \text{其中 } a_{j} := v^{j}.
\]

最后这个赋值意味着，在欧氏情形下我们无需引入新记号；我们通常把 $ \sum_{j} v^{j}\partial/\partial x^{j} $ 在 $ \Theta $ 下的像写成 $ \sum_{j} v^{j}\,dx^{j} $。也就是说，$ \Theta $ 使我们可以把向量场 $ \sum_{j} v^{j}\partial/\partial x^{j} $ 的分量 $ v^{j} $ 看作普法夫形式 $ \sum_{j} v^{j}\,dx^{j} $ 的分量。

\proofhead
因为 $ g_{jk} = \delta_{jk} $，这由注 6.1(b) 即得。\nobreak\hspace{0.35em}■
\item

（球坐标）设 $ V_{3} := (0, \infty) \times (0, 2\pi) \times (0, \pi) $，并令

\[
f: V_{3} \to \mathbb{R}^{3}, \quad (r, \varphi, \vartheta) \mapsto (x, y, z)
\]

是例 VII.9.11(a) 中的球坐标变换。则关于标准度量，有

\[
\Theta\Big(v^{1}\frac{\partial}{\partial r} + v^{2}\frac{\partial}{\partial \varphi} + v^{3}\frac{\partial}{\partial \vartheta}\Big) = v^{1}\,dr + r^{2}\sin^{2}(\vartheta)v^{2}\,d\varphi + r^{2}v^{3}\,d\vartheta.
\]

\proofhead
这由注 6.1(b) 立得。\nobreak\hspace{0.35em}■
\item

（闵可夫斯基度量）在 $ \mathbb{R}^{4}_{1,3} $ 上，有

\[
\Theta_{g}\Big(\sum_{\mu=0}^{3} v^{\mu}\frac{\partial}{\partial x^{\mu}}\Big) = v^{0}\,dx^{0} - \sum_{j=1}^{3} v^{j}\,dx^{j}
\]

其中 $ g := (\cdot|\cdot)_{1,3} $。\nobreak\hspace{0.35em}■

\end{items}

\end{example}

\needspace{3\baselineskip}\subsection*{梯度}

若 $ f \in \mathcal{E}(M) $，则 $ df $ 属于 $ \Omega^{1}(M) $。因此

\[
\operatorname{grad}_{g} f := \Theta_{g}^{-1}df \in \mathcal{V}(M)
\]

是 $ M $ 上良定义的向量场，称为 $ f $ 在（伪）黎曼流形 $ (M, g) $ 上（或关于 $ g $）的\textbf{梯度}。在不致引起误解时，我们也可以把它写成 $ \operatorname{grad}_{M} f $ 或 $ \operatorname{grad} f $。因此 $ \operatorname{grad} f $ 由下图的可交换性定义。

\begin{equation}
\begin{array}{ccc} & \mathcal{E}(M) = \Omega^{0}(M) & \\ \operatorname{grad} & \swarrow \qquad \searrow & d \\ \mathcal{V}(M) & \xrightarrow[\ \cong\ ]{\ \Theta\ } & \Omega^{1}(M) \end{array}\tag{6.4}
\end{equation}

\begin{remark}

\begin{items}
\item

映射 $ \operatorname{grad}: \mathcal{E}(M) \to \mathcal{V}(M) $，$ f \mapsto \operatorname{grad} f $ 是 $ \mathbb{R} $-线性的。

\item

对 $ f \in \mathcal{E}(M) $，向量场 $ \operatorname{grad} f $ 由下述关系刻画：

\[
g(\operatorname{grad} f, w) = \langle df, w \rangle \quad \text{对 } w \in \mathcal{V}(M).
\]
\item

在局部坐标下，有

\begin{equation}
\operatorname{grad} f = \sum_{j}\Big(\sum_{k} g^{jk}\frac{\partial f}{\partial x^{k}}\Big)\frac{\partial}{\partial x^{j}}.\tag{6.5}
\end{equation}

\proofhead
因为由 (4.5) 与 (4.8) 知 $ df = \sum_{j}\partial f/\partial x^{j}\,dx^{j} $，所以由注 6.1(d) 即得结论。\nobreak\hspace{0.35em}■
\item

（正交坐标）在正交坐标下，(6.5) 简化为

\[
\operatorname{grad} f = \sum_{j} g^{jj}\frac{\partial f}{\partial x^{j}}\frac{\partial}{\partial x^{j}}.
\]

因为此时 $ g $ 具有形式

\begin{equation}
g = \sum_{j} g_{jj}(dx^{j})^{2},\tag{6.6}
\end{equation}

即因为基本矩阵是对角的，我们有 $ g^{jj} = 1/g_{jj} $。因此系数 $ g^{jj} $ 可以直接从表示式 (6.6) 读出。

\item

设 $ (N, \bar{g}) $ 是伪黎曼流形，又设 $ \varphi \in \operatorname{Diff}(M, N) $ 且对某个 $ \lambda \neq 0 $ 有 $ \varphi_{*}g = \lambda\bar{g} $。则下图交换。

\begin{equation*}
\begin{array}{ccc} \mathcal{E}(M) & \xrightarrow{\ \lambda\operatorname{grad}_{M}\ } & \mathcal{V}(M) \\ \varphi^{*}\Big\uparrow & & \Big\uparrow\ \varphi^{*} \\ \mathcal{E}(N) & \xrightarrow{\ \operatorname{grad}_{N}\ } & \mathcal{V}(N) \end{array}
\end{equation*}

因此 $ \operatorname{grad}_{M} \circ \varphi^{*} = \lambda^{-1}\varphi^{*} \circ \operatorname{grad}_{N} $。

\proofhead
因为由注 6.1(f) 可得关系 $ \lambda\Theta_{M}^{-1}\varphi^{*} = \varphi^{*}\Theta_{N}^{-1} $，所以对 $ f \in \mathcal{E}(N) $ 我们有
\[
\lambda\operatorname{grad}_{M}(\varphi^{*}f) = \lambda\Theta_{M}^{-1}d(\varphi^{*}f) = \lambda\Theta_{M}^{-1}\varphi^{*}df = \varphi^{*}\Theta_{N}^{-1}df = \varphi^{*}\operatorname{grad}_{N} f,
\]

这里我们用到了 (4.19)。\nobreak\hspace{0.35em}■

\item

（正则性）设 $ k \in \mathbb{N} $。对 $ f \in C^{k+1}(M) $，有 $ \operatorname{grad} f \in \mathcal{V}^{k}(M) $。这里只需假定 $ M $ 是 $ C^{k+1} $ 流形。\nobreak\hspace{0.35em}■

\end{items}

\end{remark}

\begin{example}

\begin{items}
\item

（欧氏坐标）设 $ M $ 是 $ \mathbb{R}^{m} $ 中的开集。用 $ (x^{1}, \dots, x^{m}) $ 表示欧氏坐标，则 $ g_{jk} = \delta_{jk} $，因此

\[
\operatorname{grad} f = \sum_{j}\frac{\partial f}{\partial x^{j}}\frac{\partial}{\partial x^{j}}.
\]

这一表示显然与命题 VII.2.16 中的表示一致。对于任意的局部黎曼度量，我们已在注 VII.2.17(c) 中验证了 (6.4)。

\item

（球坐标）设 $ V_{3} \to \mathbb{R}^{3} $，$ (r, \varphi, \vartheta) \mapsto (x, y, z) $ 是球坐标映射。在这些坐标下，关于标准度量的梯度为

\[
\operatorname{grad} f = \frac{\partial f}{\partial r}\frac{\partial}{\partial r} + \frac{1}{r^{2}\sin^{2}\vartheta}\frac{\partial f}{\partial \varphi}\frac{\partial}{\partial \varphi} + \frac{1}{r^{2}}\frac{\partial f}{\partial \vartheta}\frac{\partial}{\partial \vartheta}.
\]

\proofhead
因为球坐标是正交坐标，这由例 5.5(g) 即得。\nobreak\hspace{0.35em}■
\item

（球坐标）设 $ h_2: W_2 \to \mathbb{R}^3 $，$ (\varphi,\vartheta) \mapsto (x,y,z) $ 是 2 维球面 $ S^2 $ 的开子集 $ U_2 := S^2\setminus H_2 $ 的参数化。则对 $ f \in C^1(U_2,\mathbb{R}) $，有

\[
\operatorname{grad}_{S^2} f = \frac{1}{\sin^2\vartheta}\frac{\partial f}{\partial\varphi}\frac{\partial}{\partial\varphi} + \frac{\partial f}{\partial\vartheta}\frac{\partial}{\partial\vartheta}.
\]

\proofhead
这可以从例 5.5(h) 中 $ g_{S^2} $ 的表示读出。\nobreak\hspace{0.35em}■
\item

（闵可夫斯基度量）设 $ X $ 是 $ \mathbb{R}^4_{1,3} $ 中的开集，$ f \in C^1(X,\mathbb{R}) $。则在闵可夫斯基度量下，有

\[
\operatorname{grad} f = \frac{\partial f}{\partial t}\frac{\partial}{\partial t} - \frac{\partial f}{\partial x}\frac{\partial}{\partial x} - \frac{\partial f}{\partial y}\frac{\partial}{\partial y} - \frac{\partial f}{\partial z}\frac{\partial}{\partial z},
\]

正如我们从 $ (\cdot\mid\cdot)_{1,3} $ 的定义立即看到的。\nobreak\hspace{0.35em}■

\end{items}

\end{example}

\needspace{3\baselineskip}\subsection*{散度}

现在设 $ M $ 是定向的，且 $ \omega_M $ 表示 $ (M,g) $ 的体积元。则映射

\begin{equation}
\bullet\,\omega_{M}: \mathcal{E}(M) \to \Omega^{m}(M), \quad a \mapsto a\omega_{M}\tag{6.7}
\end{equation}

与

\begin{equation}
\lrcorner\,\omega_{M}: \mathcal{V}(M) \to \Omega^{m-1}(M), \quad v \mapsto v\lrcorner\,\omega_{M}\tag{6.8}
\end{equation}

是逐点定义的。

\qyxffnnumdef{1}
\begin{lemma}

映射 (6.7) 与 (6.8) 是良定义的 $ \mathcal{E}(M) $-模同构。若 $ ((x^1,\ldots,x^m),U) $ 是 $ M $ 的坐标卡，则

\begin{equation}
a\omega_{M}|U = \pm a\sqrt{|G|}\,dx^{1}\wedge\cdots\wedge dx^{m}\tag{6.9}
\end{equation}

与

\begin{equation}
\Bigl(\sum_{j}v^{j}\frac{\partial}{\partial x^{j}}\Bigr)\lrcorner\,\omega_{M} = \sum_{j}(-1)^{j-1}v^{j}\sqrt{|G|}\,dx^{1}\wedge\cdots\wedge\widehat{dx^{j}}\wedge\cdots\wedge dx^{m},\tag{6.10}
\end{equation}

其中当坐标卡为正定向时，(6.9) 中使用正号，否则使用负号。

证明

\begin{subitems}
\item

由 $ \bullet\,\omega_M $ 的逐点定义以及注 5.4(c) 与 (g)，(6.9) 的成立立即推出。由此以及注 4.5(c)，我们断定对 $ a \in \mathcal{E}(M) $ 有 $ a\omega_M \in \Omega^{m}(M) $。因此映射 (6.7) 是良定义的。它显然是 $ \mathcal{E}(M) $-线性的。由注 4.14(a)，每个 $ \alpha \in \Omega^{m}(M) $ 都有唯一的 $ a \in \mathcal{E}(M) $ 使 $ \alpha = a\omega_M $。因此 (6.7) 也是双射。

\item

若坐标卡是正的，则由注 4.13(b) 推出 (6.10) 的成立。否则我们替换\qyxffnmark{1} $ x^1 $ 为 $ -x^1 $。于是 $ v^1 $ 被 $ -v^1 $ 取代。这表明 (6.10) 与坐标卡的定向无关。
因为 $ \sqrt{|G|} \in \mathcal{E}(U) $，(6.10) 与注 4.5(c) 表明对 $ v \in \mathcal{V}(M) $ 有 $ v\lrcorner\,\omega_M \in \Omega^{m-1}(M) $。因此映射 (6.8) 是良定义的，且显然是 $ \mathcal{E}(M) $-线性的。
设 $ \alpha \in \Omega^{m-1}(M) $。则由例 3.2(b) 与注 4.5(c) 推出存在唯一的 $ a_j \in \mathcal{E}(U) $ 使

\[
\alpha|U = \sum_{j}(-1)^{j-1}a_j\,dx^{1}\wedge\cdots\wedge\widehat{dx^{j}}\wedge\cdots\wedge dx^{m}.
\]

于是 $ v^{j} := a_j/\sqrt{|G|} $ 属于 $ \mathcal{E}(U) $。因此

\[
v := \sum_{j}v^{j}\frac{\partial}{\partial x^{j}} \in \mathcal{V}(U),
\]

而 (6.10) 表明 $ (v\lrcorner\,\omega_M)|U = \alpha|U $。这蕴含映射 $ \lrcorner\,\omega_M $ 是满射。由于其单射性是显然的，我们看到它是从 $ \mathcal{V}(M) $ 到 $ \Omega^{m-1}(M) $ 的同构。\nobreak\hspace{0.35em}■

\end{subitems}

\end{lemma}
\qyxffntextof{1}{考虑对于带边一维流形的情形，这个证明应如何修改。}

\begin{remark}

\begin{items}
\item

设 $ (N,\tilde{g}) $ 是定向伪黎曼流形，并设 $ \varphi \in C^\infty(M,N) $ 满足对某个 $ \mu \ne 0 $ 有 $ \varphi^{*}\omega_N = \mu\omega_M $。则图

\begin{equation*}
\begin{array}{ccc} \mathcal{E}(M) & \xrightarrow{\mu(\bullet\,\omega_{M})} & \Omega^{m}(M) \\ \varphi^{*}\uparrow & & \uparrow\varphi^{*} \\ \mathcal{E}(N) & \xrightarrow{\bullet\,\omega_{N}} & \Omega^{n}(N) \end{array}
\end{equation*}

交换，即

\[
\mu(\varphi^{*}a)\cdot\omega_{M} = \varphi^{*}(a\cdot\omega_{N}), \quad \text{对 } a \in \mathcal{E}(N).
\]

\proofhead
这立即由（外）积在拉回下的行为推出。\nobreak\hspace{0.35em}■
\item

设 $ (N,\tilde{g}) $ 是定向伪黎曼流形，并设 $ \varphi $ 属于 $ \operatorname{Diff}(M,N) $ 且对某个 $ \mu \ne 0 $ 满足 $ \varphi^{*}\omega_N = \mu\omega_M $。则

\begin{equation*}
\begin{array}{ccc} \mathcal{V}(M) & \xrightarrow[\cong]{\mu(\lrcorner\,\omega_{M})} & \Omega^{m-1}(M) \\ \varphi^{*}\uparrow & & \uparrow\varphi^{*} \\ \mathcal{V}(N) & \xrightarrow[\cong]{\lrcorner\,\omega_{N}} & \Omega^{m-1}(N) \end{array}
\end{equation*}

是交换图，即对 $ v \in \mathcal{V}(N) $ 有 $ \mu((\varphi^{*}v)\lrcorner\,\omega_M) = \varphi^{*}(v\lrcorner\,\omega_N) $。

\proofhead
由注 4.13(a)，我们推出
\[
\mu((\varphi^{*}v)\lrcorner\,\omega_{M}) = \varphi^{*}v\lrcorner\,(\mu\omega_{M}) = \varphi^{*}v\lrcorner\,\varphi^{*}\omega_{N} = \varphi^{*}(\varphi_{*}\varphi^{*}v\lrcorner\,\omega_{N}) = \varphi^{*}(v\lrcorner\,\omega_{N})
\]

对 $ v \in \mathcal{V}(N) $ 成立。\nobreak\hspace{0.35em}■

\item

（正则性）设 $ k \in \mathbb{N} $。则显然

\[
\bullet\,\omega_{M}: C^{k}(M) \to \Omega^{m}_{(k)}(M)
\]

与

\[
\lrcorner\,\omega_{M}: \mathcal{V}^{k}(M) \to \Omega^{m-1}_{(k)}(M),
\]

且这些映射是 $ C^{k}(M) $-模同构。因此只需假设 $ M $ 是 $ C^{k+1} $ 流形。\nobreak\hspace{0.35em}■

\end{items}

借助同构 (6.7) 与 (6.8)，我们定义映射

\begin{equation}
\operatorname{div}_{g}: \mathcal{V}(M) \to \mathcal{E}(M), \quad v \mapsto \operatorname{div}_{g} v\tag{6.11}
\end{equation}

要求图

\begin{equation}
\begin{array}{ccc} \mathcal{V}(M) & \xrightarrow{\operatorname{div}_{g}} & \mathcal{E}(M) \\ \lrcorner\,\omega_{M}\downarrow\cong & & \cong\downarrow\bullet\,\omega_{M} \\ \Omega^{m-1}(M) & \xrightarrow{d} & \Omega^{m}(M) \end{array}\tag{6.12}
\end{equation}

交换。换言之，对 $ v \in \mathcal{V}(M) $，定向伪黎曼流形 $ (M,g) $（或关于 $ g $）上的向量场 $ v $ 的\textbf{散度} $ \operatorname{div}_g v $ 由关系

\begin{equation}
(\operatorname{div}_{g} v)\omega_{M} = d(v\lrcorner\,\omega_{M})\tag{6.13}
\end{equation}

定义。代替 $ \operatorname{div}_g $，我们也可以写 $ \operatorname{div}_M $，或者若不预期混淆就简单地写 $ \operatorname{div} $。

\end{remark}

\begin{remark}

\begin{items}
\item

映射 (6.11) 是 $ \mathbb{R} $-线性的。

\item

设 $ ((x^1,\ldots,x^m),U) $ 是 $ M $ 的坐标卡。对 $ v := \sum_j v^j\partial/\partial x^j \in \mathcal{V}(U) $，有

\begin{equation}
\operatorname{div} v = \frac{1}{\sqrt{|G|}}\sum_{j}\frac{\partial}{\partial x^{j}}\left(\sqrt{|G|}\,v^{j}\right).\tag{6.14}
\end{equation}

在正交坐标中，我们还有 $ \sqrt{|G|} = \sqrt{|g_{11}\cdot g_{22}\cdot\cdots\cdot g_{mm}|} $。

\proofhead
设 $ \varepsilon := 1 $ 若坐标卡是正的；若它是负的则用 $ \varepsilon := -1 $。由 (6.9)、(6.10) 与 (6.13)，我们（在 $ U $ 上）得到
\begin{equation*}
\begin{aligned}
\operatorname{div}(v)\omega_{M} &= d(v\lrcorner\,\omega_{M}) = \varepsilon d\Bigl(\sum_{j}(-1)^{j-1}v^{j}\sqrt{|G|}\,dx^{1}\wedge\cdots\wedge\widehat{dx^{j}}\wedge\cdots\wedge dx^{m}\Bigr) \\ &= \varepsilon\sum_{j,k}(-1)^{j-1}\frac{\partial(v^{j}\sqrt{|G|})}{\partial x^{k}}\,dx^{k}\wedge dx^{1}\wedge\cdots\wedge\widehat{dx^{j}}\wedge\cdots\wedge dx^{m} \\ &= \varepsilon\Bigl(\sum_{j}\frac{\partial(v^{j}\sqrt{|G|})}{\partial x^{j}}\Bigr)dx^{1}\wedge\cdots\wedge dx^{m} \\ &= \Bigl(\frac{1}{\sqrt{|G|}}\sum_{j}\frac{\partial(v^{j}\sqrt{|G|})}{\partial x^{j}}\Bigr)\omega_{M}
\end{aligned}
\end{equation*}

对 $ v \in \mathcal{V}(M) $ 成立。\nobreak\hspace{0.35em}■

\item

设 $ (N,\tilde{g}) $ 是定向伪黎曼流形，映射 $ \varphi \in \operatorname{Diff}(M,N) $ 对某个 $ \mu \ne 0 $ 满足 $ \varphi^{*}\omega_N = \mu\omega_M $。则

\begin{equation*}
\begin{array}{ccc} \mathcal{V}(M) & \xrightarrow{\operatorname{div}_{M}} & \mathcal{E}(M) \\ \varphi^{*}\uparrow & & \uparrow\varphi^{*} \\ \mathcal{V}(N) & \xrightarrow{\operatorname{div}_{N}} & \mathcal{E}(N) \end{array}
\end{equation*}

是交换图，即 $ \operatorname{div}_M\circ\varphi^{*} = \varphi^{*}\circ\operatorname{div}_N $。

\proofhead
由注 6.6(b) 与 (6.13)，利用 $ d\circ\varphi^{*} = \varphi^{*}\circ d $，我们得到
\begin{equation*}
\begin{aligned}
\mu\operatorname{div}_{M}(\varphi^{*}v)\omega_{M} &= \mu\,d(\varphi^{*}v\lrcorner\,\omega_{M}) = d\varphi^{*}(v\lrcorner\,\omega_{N}) = \varphi^{*}d(v\lrcorner\,\omega_{N}) \\ &= \varphi^{*}[(\operatorname{div}_{N}v)\omega_{N}] = \varphi^{*}(\operatorname{div}_{N}v)\varphi^{*}\omega_{N} = \mu\varphi^{*}(\operatorname{div}_{N}v)\omega_{M}
\end{aligned}
\end{equation*}

对 $ v \in \mathcal{V}(N) $ 成立。现在断言由引理 6.5 推出。\nobreak\hspace{0.35em}■

\item

（正则性）设 $ k \in \mathbb{N} $。则对 $ v \in \mathcal{V}^{k+1}(M) $ 有 $ \operatorname{div} v \in C^{k}(M) $，且映射

\[
\operatorname{div}: \mathcal{V}^{k+1}(M) \to C^{k}(M), \quad v \mapsto \operatorname{div} v
\]

是 $ \mathbb{R} $-线性的。因此只需假设 $ M $ 是 $ C^{k+2} $ 流形。

\proofhead
这是注 4.11(b) 与 6.6(c) 的推论。\nobreak\hspace{0.35em}■
\end{items}

正如我们将在下一章看到的，向量场的散度具有有趣的几何与物理解释。

\end{remark}

\begin{example}

\begin{items}
\item

（欧几里得坐标）设 $ U $ 是 $ \mathbb{R}^m $ 中的开集。记欧几里得坐标为 $ (x^1,\ldots,x^m) $，我们有

\[
\operatorname{div} v = \sum_{j}\frac{\partial v^{j}}{\partial x^{j}}
\]

对 $ v = \sum_j v^j\partial/\partial x^j $。当 $ ((x^1,\ldots,x^m),U) $ 是 $ (M,g) $ 上任何其他标准正交坐标时，该公式也成立。

\item

（平面极坐标）设 $ V_2 := (0,\infty)\times(0,2\pi) $，并设

\[
f_2: V_2 \to \mathbb{R}^{2}, \quad (r,\varphi) \mapsto (x,y) := (r\cos\varphi, r\sin\varphi)
\]

是平面极坐标映射。则关于标准度量，我们有

\[
\operatorname{div}\Bigl(v^{1}\frac{\partial}{\partial r} + v^{2}\frac{\partial}{\partial\varphi}\Bigr) = \frac{1}{r}\frac{\partial(rv^{1})}{\partial r} + \frac{\partial v^{2}}{\partial\varphi} = \frac{v^{1}}{r} + \frac{\partial v^{1}}{\partial r} + \frac{\partial v^{2}}{\partial\varphi}.
\]

\proofhead
这由 $ \sqrt{G} = r $ 推出，可以从例 5.5(e) 中给出的 $ g_2 $ 的表示读出。\nobreak\hspace{0.35em}■
\item

（球坐标）设 $ V_3 := (0,\infty)\times(0,2\pi)\times(0,\pi) $，并设

\[
f_3: V_3 \to \mathbb{R}^{3}, \quad (r,\varphi,\vartheta) \mapsto (x,y,z)
\]

是例 5.5(g) 的球坐标映射。关于标准度量 $ g_3 := (dx)^2 + (dy)^2 + (dz)^2 $，我们有

\begin{equation*}
\begin{aligned}
\operatorname{div}\Bigl(v^{1}\frac{\partial}{\partial r} + v^{2}\frac{\partial}{\partial\varphi} + v^{3}\frac{\partial}{\partial\vartheta}\Bigr) &= \frac{1}{r^{2}}\frac{\partial(r^{2}v^{1})}{\partial r} + \frac{\partial v^{2}}{\partial\varphi} + \frac{1}{\sin\vartheta}\frac{\partial(v^{3}\sin\vartheta)}{\partial\vartheta} \\ &= \frac{2}{r}v^{1} + \frac{\partial v^{1}}{\partial r} + \frac{\partial v^{2}}{\partial\varphi} + \cot(\vartheta)v^{3} + \frac{\partial v^{3}}{\partial\vartheta}.
\end{aligned}
\end{equation*}

\proofhead
例 5.5(g) 给出 $ \sqrt{|G|} = r^{2}\sin\vartheta $，这正是断言所需要的。\nobreak\hspace{0.35em}■
\item

（闵可夫斯基度量）设 $ M := \mathbb{R}^4_{1,3} $，$ g := (dt)^2 - (dx)^2 - (dy)^2 - (dz)^2 $。则

\[
\operatorname{div}\Bigl(v^{0}\frac{\partial}{\partial t} + v^{1}\frac{\partial}{\partial x} + v^{2}\frac{\partial}{\partial y} + v^{3}\frac{\partial}{\partial z}\Bigr) = \frac{\partial v^{0}}{\partial t} - \frac{\partial v^{1}}{\partial x} - \frac{\partial v^{2}}{\partial y} - \frac{\partial v^{3}}{\partial z}
\]

对 $ v^{j} \in \mathcal{E}(\mathbb{R}^4_{1,3}) $，$ 0 \le j \le 3 $。\nobreak\hspace{0.35em}■

\end{items}

\end{example}

\needspace{3\baselineskip}\subsection*{拉普拉斯–贝尔特拉米算子}

通过组合两个一阶微分算子 $ \operatorname{grad} $ 与 $ \operatorname{div} $，我们得到最重要的二阶微分算子，即\textbf{拉普拉斯–贝尔特拉米算子} $ \Delta_g $。

设 $ (M,g) $ 是定向伪黎曼流形。则我们通过

\[
\Delta_{g} := \operatorname{div}_{g}\operatorname{grad}_{g}
\]

定义 $ \Delta_g $，或等价地，通过要求图

\begin{equation*}
\begin{array}{ccc} \mathcal{E}(M) & \xrightarrow{\Delta_{g}} & \mathcal{E}(M) \\ \operatorname{grad}_{g}\searrow & & \nearrow\operatorname{div}_{g} \\ & \mathcal{V}(M) & \end{array}
\end{equation*}

交换来定义。代替 $ \Delta_g $，我们也可以写 $ \Delta_M $，或者若 $ g $ 从上下文清楚则简单地写 $ \Delta $。

\begin{remark}

\begin{items}
\item

映射 $ \Delta_M: \mathcal{E}(M) \to \mathcal{E}(M) $ 是 $ \mathbb{R} $-线性的。

\item

若 $ ((x^1,\ldots,x^m),U) $ 是 $ M $ 的坐标卡，则

\begin{equation}
\Delta_{M} f = \frac{1}{\sqrt{|G|}}\sum_{j,k}\frac{\partial}{\partial x^{j}}\left(\sqrt{|G|}\,g^{jk}\frac{\partial f}{\partial x^{k}}\right), \quad \text{对 } f \in \mathcal{E}(U).\tag{6.15}
\end{equation}

在正交坐标中，(6.15) 简化为

\begin{equation}
\Delta_{M} f = \frac{1}{\sqrt{|G|}}\sum_{j}\frac{\partial}{\partial x^{j}}\left(\sqrt{|G|}\,g^{jj}\frac{\partial f}{\partial x^{j}}\right), \quad \text{对 } f \in \mathcal{E}(U),\tag{6.16}
\end{equation}

其中 $ \sqrt{|G|} = \sqrt{|g_{11}\cdot g_{22}\cdot\cdots\cdot g_{mm}|} $。

\proofhead
这由注 6.3(c) 与 (d) 以及注 6.7(b) 推出。\nobreak\hspace{0.35em}■
\end{items}

\begin{items}
\item[（c）]

设 $ (N,\tilde{g}) $ 是定向伪黎曼流形。又设 $ \varphi \in \operatorname{Diff}(M,N) $，并设存在 $ \lambda \ne 0 $ 与 $ \mu \ne 0 $ 使 $ \varphi_{*}g = \lambda\tilde{g} $ 且 $ \varphi^{*}\omega_N = \mu\omega_M $。则图

\begin{equation*}
\begin{array}{ccc} \mathcal{E}(M) & \xrightarrow{\lambda\Delta_{M}} & \mathcal{E}(M) \\ \varphi^{*}\uparrow\cong & & \cong\uparrow\varphi^{*} \\ \mathcal{E}(N) & \xrightarrow{\Delta_{N}} & \mathcal{E}(N) \end{array}
\end{equation*}

交换：$ \lambda\Delta_M\circ\varphi^{*} = \varphi^{*}\circ\Delta_N $。

\proofhead
这是注 6.3(e) 与 6.7(c) 的推论。\nobreak\hspace{0.35em}■
\item[（d）]

（正则性）设 $ k \in \mathbb{N} $。则显然

\[
\Delta_{M}: C^{k+2}(M) \to C^{k}(M),
\]

且该映射是 $ \mathbb{R} $-线性的。这里只需假设 $ M $ 是 $ C^{k+2} $ 流形。\nobreak\hspace{0.35em}■

\end{items}

\end{remark}

\begin{example}

\begin{items}
\item

（欧几里得坐标）设 $ M $ 是 $ \mathbb{R}^m $ 中的开集，欧几里得坐标为 $ ((x^1,\ldots,x^m),M) $。则 $ \Delta_M $ 与（通常的）\textbf{m 维拉普拉斯算子}

\[
\Delta_{m} := \sum_{j}\partial_{j}^{2}.
\]

相同。参见习题 VII.5.3。

\item

（圆坐标）关于 $ S^1\setminus\{(1,0)\} $ 的参数化

\[
h: (0,2\pi) \to \mathbb{R}^{2}, \quad \varphi \mapsto (\cos\varphi, \sin\varphi)
\]

（以及标准度量），我们有 $ \Delta_{S^1} = \partial_{\varphi}^{2} $。

\proofhead
注 6.9(b) 与例 5.5(f)。\nobreak\hspace{0.35em}■
\item

（平面极坐标）在平面极坐标

\[
(0,\infty)\times(0,2\pi) \to \mathbb{R}^{2}, \quad (r,\varphi) \mapsto (r\cos\varphi, r\sin\varphi),
\]

中，拉普拉斯–贝尔特拉米算子（关于 $ \mathbb{R}^2 $ 的标准度量）为

\[
\Delta_{2} = \frac{1}{r}\partial_{r}(r\partial_{r}\,\cdot) + \frac{1}{r^{2}}\partial_{\varphi}^{2} = \partial_{r}^{2} + \frac{1}{r}\partial_{r} + \frac{1}{r^{2}}\partial_{\varphi}^{2} = \frac{1}{r^{2}}\left[(r\partial_{r})^{2} + \Delta_{S^1}\right].
\]

\proofhead
这由注 6.9(b)、例 5.5(e) 以及 (b) 推出。\nobreak\hspace{0.35em}■
\item

（m 维球坐标）对 $ m \ge 2 $，$ S^m $ 的拉普拉斯–贝尔特拉米算子（关于标准度量）在例 5.5(h) 的球坐标下具有形式

\[
\Delta_{S^{m}} = \frac{1}{\sin^{2}\vartheta_{1}\cdot\cdots\cdot\sin^{2}\vartheta_{m-1}}\frac{\partial^{2}}{\partial\varphi^{2}} + \sum_{k=1}^{m-1}\frac{1}{\sin^{k}\vartheta_{k}\sin^{2}\vartheta_{k+1}\cdot\cdots\cdot\sin^{2}\vartheta_{m-1}}\frac{\partial}{\partial\vartheta_{k}}\Bigl(\sin^{k}\vartheta_{k}\frac{\partial}{\partial\vartheta_{k}}\Bigr).
\]

特别地，

\[
\Delta_{S^{2}} = \frac{1}{\sin^{2}\vartheta}\partial_{\varphi}^{2} + \frac{1}{\sin\vartheta}\partial_{\vartheta}(\sin\vartheta\,\partial_{\vartheta}\,\cdot) = \frac{1}{\sin^{2}\vartheta}\partial_{\varphi}^{2} + \partial_{\vartheta}^{2} + \cot\vartheta\,\partial_{\vartheta}.
\]

\proofhead
由例 5.5(g) 与 (h) 推出
\[
G = \prod_{k=0}^{m-1}a_{m+1,k} = \prod_{k=0}^{m-2}\prod_{i=k+1}^{m-1}\sin^{2}\vartheta_{i}.
\]
\end{items}

\bookfigure[0.40]{figures/53ec4505d5088b03.png}

交换两个乘积的次序给出

\begin{equation}
G = \prod_{i=1}^{m-1}\sin^{2i}\vartheta_{i} = \left[w_{m+1}(\vartheta)\right]^{2},\tag{6.17}
\end{equation}

其中我们使用了命题 X.8.9 中引入的简写记号。

由球坐标的正交性，从已给出的例子还推出

\[
g^{jj} = \frac{1}{a_{m+1,j-1}} = \frac{1}{\prod_{i=j}^{m-1}\sin^{2}\vartheta_{i}}, \quad \text{对 } 1 \le j \le m.
\]

由此我们读出

\[
\sqrt{G}\,g^{jj} = \Bigl(\prod_{i=1,\,i\ne j-1}^{m-1}\sin^{i}\vartheta_{i}\prod_{k=j}^{m-1}\frac{1}{\sin^{2}\vartheta_{k}}\Bigr)\sin^{j-1}\vartheta_{j-1}
\]

对 $ 2 \le j \le m $。于是我们得到

\[
\frac{1}{\sqrt{G}}\frac{\partial}{\partial\vartheta_{j-1}}\Bigl(\sqrt{G}\,g^{jj}\frac{\partial}{\partial\vartheta_{j-1}}\Bigr) = \frac{1}{\sin^{j-1}\vartheta_{j-1}\prod_{i=j}^{m-1}\sin^{2}\vartheta_{i}}\frac{\partial}{\partial\vartheta_{j-1}}\Bigl(\sin^{j-1}\vartheta_{j-1}\frac{\partial}{\partial\vartheta_{j-1}}\Bigr)
\]

对 $ 2 \le j \le m $。现在断言清楚了。\nobreak\hspace{0.35em}■

\begin{items}
\item[（e）]

（m 维极坐标）在 $ m \ge 2 $ 的 $ m $ 维极坐标中，m 维拉普拉斯算子为

\begin{equation*}
\begin{aligned}
\Delta_{m} &= \frac{1}{r^{m-1}}\partial_{r}(r^{m-1}\partial_{r}\,\cdot) + \frac{1}{r^{2}}\Delta_{S^{m-1}} = \partial_{r}^{2} + \frac{m-1}{r}\partial_{r} + \frac{1}{r^{2}}\Delta_{S^{m-1}} \\ &= \frac{1}{r^{2}}\left[(r\partial_{r})^{2} + (m-2)r\partial_{r} + \Delta_{S^{m-1}}\right].
\end{aligned}
\end{equation*}

\proofhead
由例 5.5(g) 与 (h)，我们读出 $ g_m = (dr)^{2} + r^{2}g_{S^{m-1}} $。这便蕴含 $ G = r^{2(m-1)}G_{S^{m-1}} $。它还蕴含 $ g^{11} = 1 $ 以及
\[
g^{jj} = \frac{1}{r^{2}}g^{(j-1)(j-1)}_{S^{m-1}}, \quad \text{对 } 2 \le j \le m.
\]

现在由 (6.16) 以及坐标的正交性得到断言。\nobreak\hspace{0.35em}■

\item[（f）]

（闵可夫斯基度量）在标准正交坐标中，闵可夫斯基空间 $ \mathbb{R}^4_{1,3} $ 的拉普拉斯–贝尔特拉米算子具有形式 $ \partial_{t}^{2} - \Delta_{3} $，其中 $ \Delta_3 $ 是三维（欧几里得）拉普拉斯算子。也就是说，闵可夫斯基空间中的拉普拉斯–贝尔特拉米算子恰好就是波动算子。\footnote{参见习题 VII.5.10。}

\proofhead
这是 (6.16) 的直接推论。\nobreak\hspace{0.35em}■
\end{items}

在下一个命题中，我们列出向量分析中使用的微分算子的一些基本性质。这里以及下文中，我们用 $ (\cdot\mid\cdot)_M $ 表示 $ M $ 的伪黎曼度量。

\end{example}

\begin{proposition}

设 $ (M,(\cdot\mid\cdot)_M) $ 是定向伪黎曼流形，$ f, g \in \mathcal{E}(M) $，$ v, w \in \mathcal{V}(M) $。则

\begin{subitems}
\item

$ \operatorname{grad}(fg) = f\operatorname{grad} g + g\operatorname{grad} f $；

\item

$ \operatorname{div}(fv) = f\operatorname{div} v + (\operatorname{grad} f\mid v)_M $；

\item

$ \Delta(fg) = f\Delta g + 2(\operatorname{grad} f\mid\operatorname{grad} g)_M + g\Delta f $；

\item

$ f\Delta g - g\Delta f = \operatorname{div}(f\operatorname{grad} g) - \operatorname{div}(g\operatorname{grad} f) $。

\proofhead
(i) 因为 $ \Theta $ 是模同构，由 (6.4) 推出 (i) 等价于
\begin{equation}
d(fg) = f\,dg + g\,df.\tag{6.18}
\end{equation}

因为 (6.18) 是局部陈述，所以只需在局部坐标中证明该公式。此时，它是乘积法则的直接推论。

\end{subitems}

\begin{subitems}
\item[（ii）]

由 $ (fv)\lrcorner\,\omega_M = f(v\lrcorner\,\omega_M) = f\wedge(v\lrcorner\,\omega_M) $ 以及定理 4.10 的乘积法则推出

\begin{equation}
d((fv)\lrcorner\,\omega_{M}) = d(f\wedge(v\lrcorner\,\omega_{M})) = df\wedge(v\lrcorner\,\omega_{M}) + f\,d(v\lrcorner\,\omega_{M}).\tag{6.19}
\end{equation}

因为这也是局部陈述，我们可以使用局部表示。于是，取 $ v = \sum_j v^j\partial/\partial x^j $ 与正坐标卡，由 (6.9) 与 (6.10) 得到

\begin{equation}
\begin{aligned}
df\wedge(v\lrcorner\,\omega_{M}) &= \Bigl(\sum_{j}\frac{\partial f}{\partial x^{j}}dx^{j}\Bigr)\wedge\sum_{k}(-1)^{k-1}v^{k}\sqrt{|G|}\,dx^{1}\wedge\cdots\wedge\widehat{dx^{k}}\wedge\cdots\wedge dx^{m} \\ &= \Bigl(\sum_{j}\frac{\partial f}{\partial x^{j}}v^{j}\Bigr)\sqrt{|G|}\,dx^{1}\wedge\cdots\wedge dx^{m} = \sum_{j}\frac{\partial f}{\partial x^{j}}v^{j}\omega_{M}.
\end{aligned}\tag{6.20}
\end{equation}

然后由注 6.3(b) 与 (4.4) 读出

\begin{equation}
(\operatorname{grad} f\mid v)_M = \langle df, v\rangle = \sum_{j}\Bigl\langle df, \frac{\partial}{\partial x^{j}}\Bigr\rangle v^{j} = \sum_{j}\frac{\partial f}{\partial x^{j}}v^{j}.\tag{6.21}
\end{equation}

因此由 (6.19)–(6.21) 与定义 (6.13) 推出

\[
\operatorname{div}(fv)\omega_{M} = d((fv)\lrcorner\,\omega_{M}) = (\operatorname{grad} f\mid v)_M\omega_{M} + f\operatorname{div} v\,\omega_{M},
\]

这蕴含断言。

\item[（iii）]

这立即由 $ \Delta = \operatorname{div}\operatorname{grad} $ 以及 (i)、(ii) 得到。

\item[（iv）]

由 (ii) 推出

\begin{equation}
\operatorname{div}(f\operatorname{grad} g) = f\Delta g + (\operatorname{grad} f\mid\operatorname{grad} g)_M.\tag{6.22}
\end{equation}

交换 $ f $ 与 $ g $ 并从 (6.22) 中减去所得结果，便得到 (iv)。\nobreak\hspace{0.35em}■

\end{subitems}

\end{proposition}

\needspace{3\baselineskip}\subsection*{旋度}

设 $ (M, g) $ 是 3 维定向伪黎曼流形。于是我们通过要求下图交换来定义向量场 $ v \in \mathcal{V}(M) $ 的\textbf{旋度} $ \operatorname{curl} v $（有时写作 rot，是“rotation”的缩写）：

\[
\begin{tikzcd}[row sep=large, column sep=large] \mathcal{V}(M) \arrow[r, "\Theta", "\cong"'] \arrow[d, "\operatorname{curl}"'] & \Omega^{1}(M) \arrow[d, "d"] \\ \mathcal{V}(M) \arrow[r, "\lrcorner\,\omega_{M}", "\cong"'] & \Omega^{2}(M) \end{tikzcd} \tag{6.23}
\]

也就是要求

\begin{equation}
(\operatorname{curl} v) \lrcorner\, \omega_{M} = d(\Theta v) \qquad \text{对 } v \in \mathcal{V}(M) .\tag{6.24}
\end{equation}

这个定义显然只在 $ m = 3 $ 的情形下才可能。

\begin{remark}

\begin{items}
\item

映射 $ \operatorname{curl} : \mathcal{V}(M) \to \mathcal{V}(M) $，$ v \mapsto \operatorname{curl} v $ 是 $ \mathbb{R} $ 线性的。

\item

设 $ ((x^{1}, x^{2}, x^{3}), U) $ 是 $ M $ 的坐标卡。则

\[
\operatorname{curl} v = \frac{1}{\sqrt{|G|}} \sum_{i=1}^{3} \sum_{(j,k,\ell) \in S_{3}} \operatorname{sign}(j, k, \ell) \frac{\partial}{\partial x^{j}} (g_{ki} v^{i}) \frac{\partial}{\partial x^{\ell}}
\]

对 $ v = \sum_{j=1}^{3} v^{j} \partial/\partial x^{j} $。若坐标是正交的，这个表达式简化为

\begin{equation*}
\begin{aligned}
\operatorname{curl} v &= \frac{1}{\sqrt{|G|}} \sum_{(j,k,\ell) \in S_{3}} \operatorname{sign}(j, k, \ell) \frac{\partial}{\partial x^{j}} (g_{kk} v^{k}) \frac{\partial}{\partial x^{\ell}} \\ &= \frac{1}{\sqrt{|G|}} \left[ \bigl(\partial_{2}(g_{33} v^{3}) - \partial_{3}(g_{22} v^{2})\bigr) \frac{\partial}{\partial x^{1}} + \bigl(\partial_{3}(g_{11} v^{1}) - \partial_{1}(g_{33} v^{3})\bigr) \frac{\partial}{\partial x^{2}} + \bigl(\partial_{1}(g_{22} v^{2}) - \partial_{2}(g_{11} v^{1})\bigr) \frac{\partial}{\partial x^{3}} \right]
\end{aligned}
\end{equation*}

其中 $ \sqrt{|G|} = \sqrt{|g_{11} g_{22} g_{33}|} $。若坐标是正交规范的，这成为

\[
\operatorname{curl} v = (\partial_{2} v^{3} - \partial_{3} v^{2}) \frac{\partial}{\partial x^{1}} + (\partial_{3} v^{1} - \partial_{1} v^{3}) \frac{\partial}{\partial x^{2}} + (\partial_{1} v^{2} - \partial_{2} v^{1}) \frac{\partial}{\partial x^{3}} .
\]

\proofhead
注 6.1(b) 与外微分的性质给出
\[
d(\Theta v) = d \sum_{k} \Bigl( \sum_{i} g_{ki} v^{i} \Bigr) dx^{k} = \sum_{k} \sum_{j \neq k} \frac{\partial}{\partial x^{j}} \Bigl( \sum_{i} g_{ki} v^{i} \Bigr) dx^{j} \wedge dx^{k} .
\]

由 (6.10)，我们读出

\begin{equation}
\operatorname{curl} v \lrcorner\, \omega_{M} = \sqrt{|G|} \Bigl( (\operatorname{curl} v)^{1} dx^{2} \wedge dx^{3} + (\operatorname{curl} v)^{2} dx^{3} \wedge dx^{1} + (\operatorname{curl} v)^{3} dx^{1} \wedge dx^{2} \Bigr) .\tag{6.25}
\end{equation}

因此断言由 (6.24) 推出。\nobreak\hspace{0.35em}■

\item

（正则性）设 $ k \in \mathbb{N} $。则对 $ v \in \mathcal{V}^{k+1}(M) $ 有 $ \operatorname{curl} v \in \mathcal{V}^{k}(M) $。所以这里只需假设 $ M $ 是 $ C^{k+2} $ 流形。\nobreak\hspace{0.35em}■

\end{items}

在 $ m = 3 $ 的情形，算子 $ \operatorname{grad} $、$ \operatorname{div} $ 与 $ \operatorname{curl} $ 之间有重要的关系。下面的定理以图示方式把它们总结起来。

\end{remark}

\begin{theorem}

设 $ (M, g) $ 是 3 维定向（伪）黎曼流形。

\begin{subitems}
\item

下图交换。

\[
\begin{tikzcd}[row sep=large, column sep=large] \mathcal{E}(M) \arrow[r, "\operatorname{grad}"] \arrow[d, equal, "\cong"'] & \mathcal{V}(M) \arrow[r, "\operatorname{curl}"] \arrow[d, "\Theta_{M}", "\cong"'] & \mathcal{V}(M) \arrow[r, "\operatorname{div}"] \arrow[d, "\lrcorner\,\omega_{M}", "\cong"'] & \mathcal{E}(M) \arrow[d, "\cdot\,\omega_{M}", "\cong"'] \\ \Omega^{0}(M) \arrow[r, "d"] & \Omega^{1}(M) \arrow[r, "d"] & \Omega^{2}(M) \arrow[r, "d"] & \Omega^{3}(M) \end{tikzcd} \tag{6.26}
\]
\item

$ \operatorname{curl} \circ \operatorname{grad} = 0 $。

\item

$ \operatorname{div} \circ \operatorname{curl} = 0 $。

\proofhead
（i）立即由图 (6.4)、(6.12) 与 (6.23) 的交换性推出。（ii）与（iii）现在由 $ d^{2} = 0 $ 直接推出。\nobreak\hspace{0.35em}■
\end{subitems}

\end{theorem}

\begin{corollary}

设 $ X $ 在 $ \mathbb{R}^{3} $ 中开且可缩。又设 $ v $ 是 $ X $ 上的光滑向量场。

\begin{subitems}
\item

若 $ \operatorname{curl} v = 0 $，则存在 $ f \in \mathcal{E}(X) $ 使 $ v = \operatorname{grad} f $，即 $ v $ 的一个\textbf{势}。

\item

若 $ \operatorname{div} v = 0 $，则存在 $ w \in \mathcal{V}(X) $ 使 $ v = \operatorname{curl} w $，即 $ v $ 的一个\textbf{向量势}。

\proofhead
（i）由 (6.26) 我们得知 $ \operatorname{curl} v = 0 $ 等价于 $ d(\Theta_{X} v) = 0 $。因此 1 形式 $ \Theta_{M} v $ 是闭的，而庞加莱引理（定理 3.11）保证存在 $ f \in \Omega^{0}(X) = \mathcal{E}(X) $ 使 $ \Theta_{X} v = df $。由此推出 $ v = \Theta_{X}^{-1} df = \operatorname{grad} f $。
\end{subitems}

（ii）与（i）类似，由 $ \operatorname{div} v = 0 $ 推出 2 形式 $ v \lrcorner\, \omega_{X} $ 是闭的，从而再次由庞加莱引理是恰当的。于是存在 $ \alpha \in \Omega^{1}(X) $ 使 $ d\alpha = v \lrcorner\, \omega_{X} $。那么 $ w := \Theta_{X}^{-1} \alpha \in \mathcal{V}(X) $ 由 (6.26) 中间“回路”的交换性满足 $ \operatorname{curl} w = v $。\nobreak\hspace{0.35em}■

\end{corollary}

\begin{remark}

设 $ X $ 在 $ \mathbb{R}^{3} $ 中开。

\begin{items}
\item

在欧几里得坐标中，等式 $ \operatorname{curl} v = 0 $ 等价于可积性条件

\[
\partial_{j} v^{k} = \partial_{k} v^{j} \qquad \text{对 } 1 \le j, k \le 3 ,
\]

这可由注 6.12(b) 看出。因此推论 6.14(i) 是注 VIII.4.10(a) 的特例。

\item

（经典记号）在欧几里得坐标中，由例 6.4(a) 我们知道 $ \operatorname{grad} f $ 与命题 VII.2.16 中的 $ \nabla f $ 一致。物理与工程文献以及许多数学教材使用形式的 \textbf{nabla 向量}

\[
\nabla := \left( \frac{\partial}{\partial x}, \frac{\partial}{\partial y}, \frac{\partial}{\partial z} \right) .
\]

以 $ x \cdot y $ 记 $ \mathbb{R}^{3} $ 中的欧几里得内积、以 $ x \times y $ 记向量积，则 nabla 向量记号导致（形式的）关系

\[
\operatorname{div} v = \nabla \cdot v , \qquad \operatorname{curl} v = \nabla \times v , \qquad \Delta v = (\nabla \cdot \nabla) v =: \nabla^{2} v .
\]

这些容易由这些算子的相应局部表示以及注 VIII.2.14(d) 推出。特别地，向量 $ \operatorname{curl} v $ 的分量可通过沿第一行展开（形式的）行列式

\begin{equation*}
\begin{vmatrix} \vec{e}_{1} & \vec{e}_{2} & \vec{e}_{3} \\ \partial/\partial x & \partial/\partial y & \partial/\partial z \\ v^{1} & v^{2} & v^{3} \end{vmatrix}
\end{equation*}

得到。这里 $ \vec{e}_{1}, \vec{e}_{2}, \vec{e}_{3} $ 是 $ \mathbb{R}^{3} $ 的标准基向量，而 $ \partial/\partial x $、$ \partial/\partial y $、$ \partial/\partial z $ \textbf{不是}被解释为切向量，而是被解释为微分算子。
由于符号 $ \nabla $ 在黎曼几何的语境中有另一含义，在本书余下部分我们会很少使用 nabla 向量。

\end{items}

\bookfigure[0.35]{figures/d7f42909ad0f8ab5.png}

\begin{items}
\item[（c）]

（旋度的物理意义\footnote{向量场旋度的更深解释见第 XII.3 节。}）我们考虑一个绕固定轴以常（角）速度旋转的刚体。我们然后选取一个标准正交基 $ (\vec{e}_{1}, \vec{e}_{2}, \vec{e}_{3}) $ 与坐标原点，使 $ \vec{e}_{3} $ 指向旋转轴方向。又设 $ \omega $ 是\textbf{角速度}，即 $ \omega $ 是旋转刚体中固定在距旋转轴单位距离处的任一点 $ P $ 的速率。若 $ \vec{r} $ 是点 $ P $ 的半径向量，也就是点 $ P $ 在坐标系 $ (O; \vec{e}_{1}, \vec{e}_{2}, \vec{e}_{3}) $ 中的位置向量（见注 I.12.6 之后的陈述），且若 $ \theta $ 是 $ \vec{e}_{3} $ 与 $ \vec{r} $ 之间的夹角（在 $ \vec{e}_{3} $ 与 $ \vec{r} $ 张成的平面内），则 $ P $ 到旋转轴的距离 $ a $ 满足 $ a = |\vec{r}| \sin \theta $。因此点 $ P $ 的速度向量 $ \vec{v} $ 的模由

\[
|\vec{v}| = \omega a = \omega |\vec{r}| \sin \theta .
\]

给出。
记 $ \vec{w} := \omega \vec{e}_{3} $ 为“角速度向量”，并给它定向，使刚体绕它顺时针旋转。则由向量积的性质推出

\begin{equation}
\vec{v} = \vec{w} \times \vec{r} ,\tag{6.27}
\end{equation}

因为点 $ P $ 以常速率 $ \omega $ 在以 $ \vec{e}_{3} $ 轴为中心、且位于与 $ \vec{e}_{3} $ 轴正交的平面内的圆周上运动。\footnote{(6.27) 的形式证明留给你。}
设 $ (x, y, z) $ 是 $ P $ 关于 $ (O; \vec{e}_{1}, \vec{e}_{2}, \vec{e}_{3}) $ 的坐标。则

\[
\vec{r} = x \frac{\partial}{\partial x} + y \frac{\partial}{\partial y} + z \frac{\partial}{\partial z} \quad \text{且} \quad \vec{w} = \omega \frac{\partial}{\partial z} .
\]

因此

\[
\vec{v} = \vec{w} \times \vec{r} = -\omega y \frac{\partial}{\partial x} + \omega x \frac{\partial}{\partial y} .
\]

对向量场 $ \vec{v} $ 的旋度，我们求出 $ \operatorname{curl} \vec{v} = 2\omega\, \partial/\partial z = 2\vec{w} $。也就是说，对于绕固定轴旋转的刚体，速度向量的旋度是一个向量场，其元素平行于旋转轴，且绝对值为角速度的两倍。

\item[（d）]

（正则性）定理 6.13 与推论 6.14 的陈述可以在更弱的可微性假设下证明，这些假设容易由前面关于正则性的注推出。\nobreak\hspace{0.35em}■

\end{items}

\end{remark}

\needspace{3\baselineskip}\subsection*{李导数}

现在再设 $ M $ 是任意流形。对 $ f \in \mathcal{E}(M) $ 与 $ v \in \mathcal{V}(M) $，我们置

\[
L_{v} f := \langle df, v \rangle \in \mathcal{E}(M)
\]

并称 $ L_{v} f $ 为 $ f $ 关于 $ v $ 的\textbf{李导数}。

\begin{proposition}

\begin{subitems}
\item

映射 $ L_{v} : \mathcal{E}(M) \to \mathcal{E}(M) $，即关于 $ v $ 的李导数，具有性质
（α）$ L_{v} $ 是 $ \mathbb{R} $ 线性的；
（β）对 $ f, g \in \mathcal{E}(M) $ 有 $ L_{v}(fg) = L_{v}(f) g + f L_{v} g $。

\item

在局部坐标中，

\[
L_{v} f = \sum_{j} v^{j} \frac{\partial f}{\partial x^{j}} \quad \text{且} \quad v = \sum_{j} v^{j} \frac{\partial}{\partial x^{j}} .
\]

\proofhead
（i）立即由 $ d $ 的性质推出（见 (6.18)）。（ii）是 (4.4) 的推论。\nobreak\hspace{0.35em}■
\end{subitems}

\end{proposition}

\begin{remark}

\begin{items}
\item

命题 6.16(ii) 表明，李导数推广了第 VII.2 节的方向导数。

\item

设 $ A $ 是 $ \mathbb{R} $ 代数。若映射 $ D : A \to A $ 是 $ \mathbb{R} $ 线性的并且满足乘积法则

\[
D(ab) = (Da)b + a(Db) \qquad \text{对 } a, b \in A ,
\]

则称 $ D $ 是（$ A $ 的）\textbf{导子}。因此关于 $ v \in \mathcal{V}(M) $ 的李导数是代数 $ \mathcal{E}(M) $ 的导子。

\item

若 $ A $ 是有单位元 $ e $ 的代数、$ D $ 是 $ A $ 的导子，则 $ De = 0 $。

\proofhead
乘积法则给出
\[
De = D(ee) = (De)e + e(De) = De + De = 2De
\]

从而得证。\nobreak\hspace{0.35em}■

\end{items}

下一个定理表明，$ \mathcal{E}(M) $ 的每个导子都由一个李导数给出。

\end{remark}

\begin{theorem}

设 $ D $ 是 $ \mathcal{E}(M) $ 的导子。则恰有一个 $ v \in \mathcal{V}(M) $ 使 $ D = L_{v} $。

证明

\begin{subitems}
\item

我们先证明 $ D $ 是“局部算子”。设 $ U $ 是 $ p \in M $ 的开邻域、$ K $ 是 $ p $ 的紧邻域，且 $ K \subset\subset U $。注 1.21(a) 保证存在 $ \chi \in \mathcal{E}(M) $ 使 $ \chi|_{K} = 1 $ 且 $ \operatorname{supp}(\chi) \subset\subset U $。
设 $ f \in \mathcal{E}(M) $ 满足 $ f|_{U} = 0 $。则 $ f = f\chi + f(1 - \chi) = f(1 - \chi) $，从而

\[
Df(p) = Df(p)\bigl(1 - \chi(p)\bigr) + f(p) D(1 - \chi)(p) = 0 .
\]

由于这对每个 $ p \in U $ 成立，推出 $ D(f)|_{U} = 0 $。若 $ \chi_{1} \in \mathcal{E}(M) $ 是另一个满足 $ \operatorname{supp}(\chi_{1}) \subset\subset U $ 且在 $ p $ 的某邻域内恒等于 1 的函数，则对 $ f \in \mathcal{E}(U) $ 有 $ f\chi - f\chi_{1} \in \mathcal{E}(M) $ 在 $ p $ 的某邻域内为零。于是由上述结果推出 $ D(f\chi) = D(f\chi_{1}) $ 对 $ f \in \mathcal{E}(U) $ 成立。因此“$ D $ 到 $ U $ 的限制”由

\[
D_{U} f := D(f\chi) \qquad \text{对 } f \in \mathcal{E}(U)
\]

良定义，且与 $ \chi $ 的特殊选择无关。

\item

现在设 $ (\varphi, U) $ 是坐标卡，$ \varphi = (x^{1}, \ldots, x^{m}) $。我们可以假设 $ X := \varphi(U) $ 是凸的。对每个固定的 $ p \in U $，由积分形式的中值定理（定理 VII.3.10），取 $ a := \varphi(p) $，得到

\[
(\varphi_{*} f)(x) = (\varphi_{*} f)(a) + \sum_{j} (x^{j} - a^{j}) \widetilde{f}_{j}(x) \qquad \text{对 } x \in X ,
\]

其中我们置

\[
\widetilde{f}_{j}(x) := \int_{0}^{1} \partial_{j} f\bigl(a + t(x - a)\bigr) dt \qquad \text{对 } x \in X .
\]

因此

\[
f_{j} := \varphi^{*} \widetilde{f}_{j} \in \mathcal{E}(U) , \qquad f_{j}(p) = \frac{\partial f}{\partial x^{j}}(p) ,
\]

并且

\[
f(q) = f(p) + \sum_{j} \bigl(\varphi^{j}(q) - \varphi^{j}(p)\bigr) f_{j}(q) \qquad \text{对 } q \in U .
\]

由此、$ D $ 的性质以及注 6.17(c)，推出

\begin{equation}
Df(p) = \sum_{j} D\varphi^{j}(p) \frac{\partial f}{\partial x^{j}}(p) \qquad \text{对 } p \in U ,\tag{6.28}
\end{equation}

其中我们把 $ D_{U} $ 写成了 $ D $。

\item

设 $ (\psi, V) $ 是 $ p $ 附近的第二个坐标卡，$ \psi = (y^{1}, \ldots, y^{m}) $。现在定义转移函数 $ k := \psi \circ \varphi^{-1} $。那么，与（ii）类似，并由于我们可以假设 $ U = V $，我们有

\begin{equation}
k^{j}(x) = k^{j}(a) + \sum_{\ell} (x^{\ell} - a^{\ell}) k_{\ell}^{j}(x) \qquad \text{对 } x \in X ,\tag{6.29}
\end{equation}

其中 $ k_{\ell}^{j} \in \mathcal{E}(X) $ 且 $ k_{\ell}^{j}(a) = \partial_{\ell} k^{j}(a) $。因为 $ \varphi^{*} k = \psi $，把 $ \varphi^{*} $ 作用于 (6.29) 给出

\begin{equation}
\psi^{j}(q) = \psi^{j}(p) + \sum_{\ell} \bigl(\varphi^{\ell}(q) - \varphi^{\ell}(p)\bigr) h_{\ell}^{j}(q) \qquad \text{对 } q \in U ,\tag{6.30}
\end{equation}

其中 $ h_{\ell}^{j} := \varphi^{*} k_{\ell}^{j} \in \mathcal{E}(U) $ 且

\[
h_{\ell}^{j}(p) = (\varphi^{*} \partial_{\ell} k^{j})(p) = \frac{\partial y^{j}}{\partial x^{\ell}}(p) .
\]

这与 (6.30) 蕴涵

\begin{equation}
D\psi^{j}(p) = \sum_{k} D\varphi^{k}(p) \frac{\partial y^{j}}{\partial x^{k}}(p) \qquad \text{对 } p \in U \text{ 及 } 1 \le j, k \le m .\tag{6.31}
\end{equation}

现在我们置

\begin{equation}
v_{\varphi} := \sum_{j} D\varphi^{j} \frac{\partial}{\partial x^{j}} \quad \text{且} \quad v_{\psi} := \sum_{j} D\psi^{j} \frac{\partial}{\partial y^{j}} .\tag{6.32}
\end{equation}

于是由 (6.31) 与命题 4.7 推出

\[
v_{\psi} = \sum_{j} \sum_{k} D\varphi^{k} \frac{\partial y^{j}}{\partial x^{k}} \frac{\partial}{\partial y^{j}} = \sum_{k} D\varphi^{k} \frac{\partial}{\partial x^{k}} = v_{\varphi} .
\]

这表明 (6.32) 在 $ U $ 上定义一个与所选坐标无关的向量场 $ v_{U} \in \mathcal{V}(U) $。由 (6.28)、(6.31) 与命题 6.16(ii)，我们读出 $ D_{U} f = L_{v_{U}} f $ 对 $ f \in \mathcal{E}(U) $ 成立。

\item

现在设 $ \{ (\varphi_{\alpha}, U_{\alpha}) ; \alpha \in A \} $ 是 $ M $ 的图册。则由（iii）推出，对每个 $ \alpha \in A $ 存在 $ v_{\alpha} \in \mathcal{V}(U_{\alpha}) $ 使 $ D_{U_{\alpha}} f = L_{v_{\alpha}} f $ 对 $ f \in \mathcal{E}(U_{\alpha}) $ 成立。此外，（iii）中的观察表明恰有一个 $ v \in \mathcal{V}(M) $ 使 $ v|_{U_{\alpha}} = v_{\alpha} $ 对 $ \alpha \in A $ 成立。现在（i）与命题 6.16(ii) 给出 $ D = L_{v} $。

\item

设 $ v, w \in \mathcal{V}(M) $ 且 $ D = L_{v} $、$ D = L_{w} $。则 $ L_{v} f = L_{w} f $ 对每个 $ f \in \mathcal{E}(M) $ 成立。在任意局部坐标卡 $ ((x^{1}, \ldots, x^{m}), U) $ 中，我们于是有

\[
\sum_{j} (v^{j} - w^{j}) \frac{\partial f}{\partial x^{j}} = 0 \qquad \text{对 } f \in \mathcal{E}(U) .
\]

取 $ f := x^{k} $，我们得到 $ \partial f/\partial x^{j} = \delta_{j}^{k} $，因此 $ v^{k} - w^{k} = 0 $。这对 $ 1 \le k \le m $ 成立蕴涵 $ v|_{U} = w|_{U} $，从而 $ v = w $。于是完成。\nobreak\hspace{0.35em}■

\end{subitems}

\end{theorem}

\begin{lemma}

对 $ v, w \in \mathcal{V}(M) $，$ L_{v} L_{w} - L_{w} L_{v} $ 是 $ \mathcal{E}(M) $ 的导子。

\proofhead
显然 $ L_{v} L_{w} - L_{w} L_{v} $ 是 $ \mathcal{E}(M) $ 到自身的 $ \mathbb{R} $ 线性映射。对 $ f, g \in \mathcal{E}(M) $，由于 $ \mathcal{E}(M) $ 可交换，我们知道
\begin{equation*}
\begin{aligned}
L_{v} L_{w}(fg) &= L_{v}\bigl(L_{w}(f) g + f L_{w} g\bigr) \\ &= g L_{v} L_{w} f + L_{v} f L_{w} g + L_{v} g L_{w} f + f L_{v} L_{w} g .
\end{aligned}
\end{equation*}

断言现在显然。\nobreak\hspace{0.35em}■

\end{lemma}

设 $ v, w \in \mathcal{V}(M) $。则由定理 6.18 与引理 6.19 推出，恰有一个 $ M $ 上的光滑向量场 $ [v, w] $ 使

\begin{equation}
L_{[v, w]} = L_{v} L_{w} - L_{w} L_{v} .\tag{6.33}
\end{equation}

我们称 $ [v, w] $ 为 $ v $ 与 $ w $ 的\textbf{李括号}或\textbf{换位子}。

\begin{proposition}

\begin{subitems}
\item

映射 $ \mathcal{V}(M) \times \mathcal{V}(M) \to \mathcal{V}(M) $，$ (v, w) \mapsto [v, w] $ 具有性质
（α）（双线性）$ [\,\cdot\, , \cdot\,] $ 是 $ \mathbb{R} $ 双线性的。
（β）（斜对称性）$ [v, w] = -[w, v] $ 对 $ v, w \in \mathcal{V}(M) $。
(γ)（雅可比恒等式）$ u, v, w \in \mathcal{V}(M) $ 满足关系

\[
[u,[v,w]] + [v,[w,u]] + [w,[u,v]] = 0 .
\]
\item

在局部坐标下，

\begin{equation}
[v,w] = \sum_{j,k} \left( v^{k} \frac{\partial w^{j}}{\partial x^{k}} - w^{k} \frac{\partial v^{j}}{\partial x^{k}} \right) \frac{\partial}{\partial x^{j}}\tag{6.34}
\end{equation}

其中 $ v = \sum_{j} v^{j}\, \partial/\partial x^{j} $，$ w = \sum_{j} w^{j}\, \partial/\partial x^{j} $。

\proofhead
简单的证明留给你。 \nobreak\hspace{0.35em}■
\end{subitems}

\end{proposition}

\begin{remark}

\begin{items}
\item

设 $ M $ 是 $ \mathbb{R}^{m} $ 中的开集，$ (x^{1},\ldots,x^{m}) $ 是 $ M $ 上的欧几里得坐标。利用 nabla 向量 $ \nabla $，(6.34) 可以符号地写成直观形式

\[
[v,w] = (v\cdot\nabla)w - (w\cdot\nabla)v .
\]
\item

设 $ V $ 是向量空间，$ [\,\cdot\,,\cdot\,] : V\times V \to V $ 是具有命题 6.20(i) 中性质 (α)–(γ) 的映射。则 $ (V,[\,\cdot\,,\cdot\,]) $ 称为一个\textbf{李代数}。由于 (β)，“乘法” $ [\,\cdot\,,\cdot\,] $ 一般不是交换的。由 (β) 和 (γ) 推出

\[
[a,[b,c]] - [[a,b],c] = [[c,a],b] \quad \text{对 } a, b, c \in V .
\]

所以该乘法一般也不是结合的。于是李代数一般是不可交换、非结合的代数。\footnote{在平凡的可交换情形，即对 $ a, b \in V $ 有 $ [a,b] = 0 $ 时，它当然是可交换且结合的。} 因此 $ (\mathcal{V}(M),[\,\cdot\,,\cdot\,]) $ 是李代数。

\item

（正则性）设 $ k \in \mathbb{N} $，$ v, w \in \mathcal{V}^{k}(M) $。设 $ M $ 是 $ C^{k+1} $ 类流形。则 $ L_{v} $ 不是 $ \mathcal{E}^{k+1}(M) $ 上的导子，因为对 $ f \in \mathcal{E}^{k+1}(M) $，$ L_{v}f $ 一般只属于 $ \mathcal{E}^{k}(M) $。因此李括号也不能通过 (6.33) 来定义。我们只能选取局部坐标，并通过 (6.34) 对 $ v, w \in \mathcal{V}^{k}(M) $ 定义 $ [v,w] $。\footnote{你也许可以考虑一下，这样定义的 $ [v,w] $ 为什么在整个 $ M $ 上良定义。} 于是 $ [v,w] \in \mathcal{V}^{k-1}(M) $。 \nobreak\hspace{0.35em}■

\end{items}

\end{remark}

\needspace{3\baselineskip}\subsection*{霍奇–拉普拉斯算子}

在本节余下的部分，我们利用余微分和星算子，从向量分析中导出一些更重要的关系。

设 $ (M,(\cdot\mid\cdot)_{M}) $ 是定向伪黎曼流形。首先我们用余微分写出散度。

\begin{proposition}

下图

\bookfigure[0.32]{figures/6dac1f10f9f2db49.png}

是交换的，即 $ \operatorname{div} = \delta\circ\Theta $。

\proofhead
只需在局部证明该等式。设 $ ((x^{1},\ldots,x^{m}),U) $ 是局部坐标。则对 $ v = \sum_{j} v^{j}\, \partial/\partial x^{j} \in \mathcal{V}(U) $，注 6.1(b) 与 5.11(b) 蕴含
\[
\delta\Theta v = \delta\sum_{j}\left(\sum_{k} g_{jk}v^{k}\right) dx^{j} = \frac{1}{\sqrt{|G|}}\sum_{j}\frac{\partial}{\partial x^{j}}\left(\sqrt{|G|}\, v^{j}\right) .
\]

因此结论由 (6.14) 推出。 \nobreak\hspace{0.35em}■

\end{proposition}

利用外微分和余微分，我们对 $ 0 \le r \le m $ 用下式在 $ \Omega^{r}(M) $ 上定义一个 $ \mathbb{R} $ 线性映射

\begin{equation}
\Delta_{M} := d\delta + \delta d : \Omega^{r}(M) \to \Omega^{r}(M) .\tag{6.35}
\end{equation}

这就是\textbf{霍奇–拉普拉斯算子}。对 $ a \in \mathcal{E}(M) $，由 (6.4) 和命题 6.22 推出

\[
(d\delta + \delta d)a = \delta d a = \delta\Theta(\Theta^{-1}da) = \operatorname{div}\operatorname{grad} a .
\]

因此 $ \Omega^{0}(M) = \mathcal{E}(M) $ 上的霍奇–拉普拉斯算子与拉普拉斯–贝尔特拉米算子相同，这就解释了该记号。当 $ M $ 由上下文清楚时，我们写 $ \Delta $ 表示 $ \Delta_{M} $。 \nobreak\hspace{0.35em}■

\begin{remark}

\begin{items}
\item

$ \ast\Delta = \Delta\ast $。

\proofhead
注 5.9(b) 与 5.11(a) 给出
\[
\ast\Delta = \ast d\delta + \ast\delta d = \delta d\ast + d\delta\ast = \Delta\ast
\]

因此结论成立。 \nobreak\hspace{0.35em}■

\item

$ d\Delta = \Delta d = d\delta d $，且 $ \delta\Delta = \Delta\delta = \delta d\delta $。

\proofhead
由 $ d^{2} = 0 $ 得到
\[
d\Delta = dd\delta + d\delta d = d\delta d = d\delta d + \delta dd = \Delta d .
\]

第二个结论类似可得。 \nobreak\hspace{0.35em}■

\item

设 $ M $ 是 $ \mathbb{R}^{m} $ 中的开集，$ (x^{1},\ldots,x^{m}) $ 是欧几里得坐标。则对 $ 1 \le r \le m $ 有

\[
\Delta\left(\sum_{(j)\in\mathbb{J}_{r}} a_{(j)}\, dx^{(j)}\right) = \sum_{(j)\in\mathbb{J}_{r}} \Delta a_{(j)}\, dx^{(j)}
\]

\proofhead
由线性性，只需对满足 $ (j) \in \mathbb{J}_{r} $ 的 $ \alpha := a\, dx^{(j)} $ 证明该结论。利用例 5.10(e)，我们得到
\begin{equation*}
\begin{aligned}
d\delta\alpha &= d\left(\sum_{k=1}^{r}(-1)^{k-1}\partial_{j_{k}}a\, dx^{j_{1}}\wedge\cdots\wedge\widehat{dx^{j_{k}}}\wedge\cdots\wedge dx^{j_{r}}\right) \\ &= \sum_{k=1}^{r}(-1)^{k-1}\sum_{\ell=1}^{m}\partial_{\ell}\partial_{j_{k}}a\, dx^{\ell}\wedge dx^{j_{1}}\wedge\cdots\wedge\widehat{dx^{j_{k}}}\wedge\cdots\wedge dx^{j_{r}} \\ &= \sum_{k=1}^{r}\partial_{j_{k}}^{2}a\, dx^{(j)} + \sum_{k=1}^{r}(-1)^{k-1}\sum_{\ell\notin\{j_{1},\ldots,j_{r}\}}^{m}\partial_{\ell}\partial_{j_{k}}a\, dx^{\ell}\wedge dx^{j_{1}}\wedge\cdots\wedge\widehat{dx^{j_{k}}}\wedge\cdots\wedge dx^{j_{r}} .
\end{aligned}
\end{equation*}

类似地，我们得到

\begin{equation*}
\begin{aligned}
\delta d\alpha &= \delta\sum_{\substack{\ell=1\\ \ell\notin\{j_{1},\ldots,j_{r}\}}}^{m}\partial_{\ell}a\, dx^{\ell}\wedge dx^{(j)} = \sum_{\substack{\ell=1\\ \ell\notin\{j_{1},\ldots,j_{r}\}}}^{m}\partial_{\ell}^{2}a\, dx^{(j)} \\ &\quad - \sum_{\substack{\ell=1\\ \ell\notin\{j_{1},\ldots,j_{r}\}}}^{m}\sum_{k=1}^{r}(-1)^{k-1}\partial_{j_{k}}\partial_{\ell}a\, dx^{\ell}\wedge dx^{j_{1}}\wedge\cdots\wedge\widehat{dx^{j_{k}}}\wedge\cdots\wedge dx^{j_{r}} .
\end{aligned}
\end{equation*}

由此推出\footnote{我们之所以把 $ \delta $ 定义中的符号取为 $ (-1)^{m(r+1)} $，是为了让该公式取这种形式。对于几何中常用的符号约定，该公式为 $ (d\delta+\delta d)\alpha = -(\Delta a)\, dx^{(j)} $。}

\[
\Delta_{M}\alpha = (d\delta + \delta d)\alpha = \left(\sum_{k}\partial_{k}^{2}a\right) dx^{(j)} = (\Delta a)\, dx^{(j)}
\]

因此结论成立。 \nobreak\hspace{0.35em}■

\item

（正则性）显然，当 $ 0 \le r \le m $、$ k \in \mathbb{N} $ 且 $ k \ge 2 $ 时，$ \Delta_{M} $ 是从 $ \Omega^{r}_{(k)}(M) $ 到 $ \Omega^{r}_{(k-2)}(M) $ 的 $ \mathbb{R} $ 线性映射。此时只需假设 $ M $ 是 $ C^{k+2} $ 流形。 \nobreak\hspace{0.35em}■

\end{items}

最后，我们定义\textbf{向量场的拉普拉斯算子}，即 $ \vec{\Delta} $，定义为

\[
\vec{\Delta} := \vec{\Delta}_{M} := \Theta_{M}^{-1}\circ\Delta_{M}\circ\Theta_{M} : \mathcal{V}(M) \to \mathcal{V}(M)
\]

从而由该图的交换性

\bookfigure[0.36]{figures/b1bb2c6cd9394ca1.png}
\end{remark}

\begin{remark}

\begin{items}
\item

命题 6.22 与 (6.4) 给出 $ \vec{\Delta} = \operatorname{grad}\operatorname{div} + \Theta^{-1}\delta d\Theta $。

\item

设 $ M $ 是 $ \mathbb{R}^{m} $ 中的开集，$ (x^{1},\ldots,x^{m}) $ 是 $ M $ 上的欧几里得坐标。则

\[
\vec{\Delta}\left(\sum_{j} v^{j}\frac{\partial}{\partial x^{j}}\right) = \sum_{j} \Delta v^{j}\frac{\partial}{\partial x^{j}} .
\]

若按通常做法把向量场 $ v = \sum_{j} v^{j}\, \partial/\partial x^{j} $ 与 $ (v^{1},\ldots,v^{m}) $ 等同，则 $ \vec{\Delta}v $ 意味着拉普拉斯算子可以按分量作用

\[
\vec{\Delta}v = (\Delta v^{1},\ldots,\Delta v^{m}) .
\]

此时我们通常写 $ \Delta $，而不写 $ \vec{\Delta} $。

\proofhead
这由例 6.2(a) 与注 6.23(c) 推出。 \nobreak\hspace{0.35em}■
\item

（正则性）设 $ k \in \mathbb{N} $。则 $ \vec{\Delta} $ 把 $ \mathbb{R} $ 向量空间 $ \mathcal{V}^{k+2}(M) $ 线性地映入 $ \mathcal{V}^{k}(M) $。所以这里只需假设 $ M $ 是 $ C^{k+2} $ 流形。 \nobreak\hspace{0.35em}■

\end{items}

\end{remark}

\needspace{3\baselineskip}\subsection*{向量积与旋度}

在这最后一部分，我们导出旋度算子最重要的性质。

设 $ (M,(\cdot\mid\cdot)_{M}) $ 是具有体积元 $ \omega_{M} $ 的三维定向黎曼\footnote{为简单起见，这里我们将只限于讨论黎曼度量，因为它们在应用中最为重要。}流形。

在 $ \mathcal{V}(M) $ 上，我们定义\textbf{向量积}或\textbf{叉积}

\begin{equation}
\times : \mathcal{V}(M)\times\mathcal{V}(M) \to \mathcal{V}(M), \quad (v,w) \mapsto v\times w,\tag{6.36}
\end{equation}

为

\begin{equation}
v\times w := \Theta_{M}^{-1}\omega_{M}(v,w,\cdot) .\tag{6.37}
\end{equation}

显然该映射是良定义的。

\begin{remark}

\begin{items}
\item

设 $ (M,(\cdot\mid\cdot)_{M}) = (\mathbb{R}^{3},(\cdot\mid\cdot)) $。则在常向量场的情形，(6.37) 与第 VIII.2 节的定义一致。

\item

向量积是双线性的、交错的（斜对称的），并且满足

\begin{equation}
(u\mid v\times w)_{M} = \omega_{M}(u,v,w) \quad \text{对 } u, v, w \in \mathcal{V}(M) .\tag{6.38}
\end{equation}

对 $ p \in M $，向量积 $ (v\times w)(p) $ 关于 $ T_{p}M $ 的内积 $ (\cdot\mid\cdot)_{M}(p) $ 与 $ v(p) $ 和 $ w(p) $ 正交。令 $ |v|_{M} := \sqrt{(v\mid v)_{M}} $，我们有

\[
|v\times w|_{M} = \sqrt{|v|_{M}^{2}|w|_{M}^{2} - (v\mid w)_{M}^{2}} = |v|_{M}|w|_{M}\sin\varphi,
\]

其中对 $ p \in M $、$ v, w \in \mathcal{V}(M) $，$ \varphi(p) \in [0,\pi] $ 是向量 $ v(p) $ 与 $ w(p) $ 之间的夹角。
向量积满足\textbf{格拉斯曼恒等式}

\[
v_{1}\times(v_{2}\times v_{3}) = (v_{1}\mid v_{3})_{M}v_{2} - (v_{1}\mid v_{2})_{M}v_{3},
\]

\textbf{雅可比恒等式}

\[
v_{1}\times(v_{2}\times v_{3}) + v_{2}\times(v_{3}\times v_{1}) + v_{3}\times(v_{1}\times v_{2}) = 0
\]

以及关系

\[
(v_{1}\times v_{2})\times(v_{3}\times v_{4}) = \omega_{M}(v_{1},v_{2},v_{4})v_{3} - \omega_{M}(v_{1},v_{2},v_{3})v_{4}
\]

对 $ v_{1}, v_{2}, v_{3}, v_{4} \in \mathcal{V}(M) $ 成立。特别地，$ (\mathcal{V}(M),\times) $ 是李代数。

\proofhead
所有这些都容易化为已在习题 2.3 中证明的逐点陈述。 \nobreak\hspace{0.35em}■
\item

设 $ ((x^{1},x^{2},x^{3}),U) $ 是 $ M $ 上的正向标准正交坐标\footnote{即 $ (\partial/\partial x^{1},\partial/\partial x^{2},\partial/\partial x^{3}) $ 是正向标准正交标架。}。则向量场 $ v = \sum_{j} v^{j}\, \partial/\partial x^{j} $ 与 $ w = \sum_{j} w^{j}\, \partial/\partial x^{j} $ 的叉积具有形式

\[
v\times w = (v^{2}w^{3} - v^{3}w^{2})\frac{\partial}{\partial x^{1}} + (v^{3}w^{1} - v^{1}w^{3})\frac{\partial}{\partial x^{2}} + (v^{1}w^{2} - v^{2}w^{1})\frac{\partial}{\partial x^{3}} .
\]

\proofhead
习题 2.3。 \nobreak\hspace{0.35em}■
\item

（正则性）对 $ k \in \mathbb{N} $，上述陈述对 $ C^{k} $ 向量场仍成立，且只需假设 $ M $ 是 $ C^{k+1} $ 流形。 \nobreak\hspace{0.35em}■

\end{items}

下一个定理表明向量积与 1 形式的外积有何关系。

\end{remark}

\begin{proposition}

对 $ v, w \in \mathcal{V}(M) $，我们有 $ v\times w = \Theta^{-1}\ast(\Theta v\wedge\Theta w) $，即下图

\bookfigure[0.58]{figures/34ecebd4b33134bc.png}

交换。

\proofhead
只需在局部证明该等式，在那里我们可以选取正向标准正交坐标 $ ((x^{1},x^{2},x^{3}),U) $。若 $ (v^{1},v^{2},v^{3}) $ 与 $ (w^{1},w^{2},w^{3}) $ 是向量场 $ v, w \in \mathcal{V}(M) $ 的分量，则由注 6.1(e) 推出
\begin{equation*}
\begin{aligned}
\Theta v\wedge\Theta w &= \sum_{j} v^{j}\, dx^{j} \wedge \sum_{k} w^{k}\, dx^{k} \\ &= (v^{2}w^{3} - v^{3}w^{2})\, dx^{2}\wedge dx^{3} + (v^{3}w^{1} - v^{1}w^{3})\, dx^{3}\wedge dx^{1} \\ &\quad + (v^{1}w^{2} - v^{2}w^{1})\, dx^{1}\wedge dx^{2} .
\end{aligned}
\end{equation*}

由例 5.7(d) 的证明，我们知道

\begin{equation}
\ast(dx^{2}\wedge dx^{3}) = dx^{1}, \quad \ast(dx^{3}\wedge dx^{1}) = dx^{2}, \quad \ast(dx^{1}\wedge dx^{2}) = dx^{3} .\tag{6.39}
\end{equation}

现在结论由注 6.1(e) 与 6.25(c) 推出。 \nobreak\hspace{0.35em}■

\end{proposition}

下面我们导出旋度算子的一种表示。

\begin{proposition}

下图

\bookfigure[0.58]{figures/0b161536d972a961.png}

交换，即 $ \operatorname{curl} = \Theta^{-1}\ast d\Theta $。

\proofhead
同样只需在正向标准正交坐标 $ ((x^{1},x^{2},x^{3}),U) $ 中局部证明该等式。则对 $ v = \sum_{j=1}^{3} v^{j}\, \partial/\partial x^{j} $，利用注 6.1(e) 可得
\begin{equation*}
\begin{aligned}
d(\Theta v) &= d\left(\sum_{j} v^{j}\, dx^{j}\right) = \sum_{j,k}\frac{\partial v^{j}}{\partial x^{k}}\, dx^{k}\wedge dx^{j} \\ &= \left(\frac{\partial v^{3}}{\partial x^{2}} - \frac{\partial v^{2}}{\partial x^{3}}\right) dx^{2}\wedge dx^{3} + \left(\frac{\partial v^{1}}{\partial x^{3}} - \frac{\partial v^{3}}{\partial x^{1}}\right) dx^{3}\wedge dx^{1} \\ &\quad + \left(\frac{\partial v^{2}}{\partial x^{1}} - \frac{\partial v^{1}}{\partial x^{2}}\right) dx^{1}\wedge dx^{2} .
\end{aligned}
\end{equation*}

于是结论由 (6.39) 及注 6.1(e) 与 6.12(b) 推出。 \nobreak\hspace{0.35em}■

\end{proposition}

现在我们准备导出若干涉及三维向量场的重要微分恒等式。

\begin{proposition}

对 $ f \in \mathcal{E}(M) $ 与 $ v, w \in \mathcal{V}(M) $，

\begin{subitems}
\item

$ \operatorname{div}(v\times w) = (\operatorname{curl} v\mid w)_{M} - (v\mid\operatorname{curl} w)_{M} $；

\item

$ \operatorname{curl}(fv) = f\operatorname{curl} v + \operatorname{grad} f\times v $；

\item

$ \operatorname{curl}(v\times w) = (\operatorname{div} w)v - (\operatorname{div} v)w - [v,w] $；

\item

$ \operatorname{curl}(\operatorname{curl} v) = \operatorname{grad}\operatorname{div} v - \vec{\Delta}v $。

\proofhead
(i) 在注 5.9(d) 中取 $ m = 3 $，我们得到
\[
\ast\delta\alpha = (-1)^{m(r+1)}\ast\ast d\ast\alpha = d\ast\alpha \quad \text{对 } \alpha \in \Omega^{2}(M) .
\]

现在我们用命题 6.22 与 6.26 推出

\begin{equation*}
\begin{aligned}
\operatorname{div}(v\times w) &= \delta\Theta\left(\Theta^{-1}\ast(\Theta v\wedge\Theta w)\right) = \delta\ast(\Theta v\wedge\Theta w) \\ &= \ast d(\Theta v\wedge\Theta w) = \ast(d\Theta v\wedge\Theta w - \Theta v\wedge d\Theta w) .
\end{aligned}
\end{equation*}

由命题 6.27 推出 $ \Theta\operatorname{curl} = \ast d\Theta $。现在注 2.19(d) 蕴含 $ d\Theta = \ast\Theta\operatorname{curl} $，因为 $ m = 3 $ 且 $ r = 2 $。于是我们得到

\begin{equation*}
\begin{aligned}
\operatorname{div}(v\times w) &= \ast\left((\ast\Theta\operatorname{curl} v)\wedge\Theta w - \Theta v\wedge\ast\Theta\operatorname{curl} w\right) \\ &= \ast(\Theta w\wedge\ast\Theta\operatorname{curl} v - \Theta v\wedge\ast\Theta\operatorname{curl} w),
\end{aligned}
\end{equation*}

其中我们用到了 $ \ast\Theta\operatorname{curl} v \in \Omega^{2}(M) $。现在 (2.22)（取 $ r = 1 $）与 (2.13) 给出

\[
\operatorname{div}(v\times w) = \ast\left[(w\mid\operatorname{curl} v)_{M} - (v\mid\operatorname{curl} w)\right]\omega_{M} .
\]

结论现在由 $ \ast\omega_{M} = 1 $ 推出。

\end{subitems}

\begin{subitems}
\item[（ii）]

命题 6.27 给出

\begin{equation*}
\begin{aligned}
\operatorname{curl}(fv) &= \Theta^{-1}\ast d\Theta(fv) = \Theta^{-1}\ast d(f\Theta v) \\ &= \Theta^{-1}\ast(df\wedge\Theta v + f\, d\Theta v) \\ &= \Theta^{-1}\ast(\Theta\operatorname{grad} f\wedge\Theta v) + f\Theta^{-1}\ast d\Theta v \\ &= \operatorname{grad} f\times v + f\operatorname{curl} v .
\end{aligned}
\end{equation*}

这里我们还用到了命题 6.26 以及 $ d $ 的性质。

\item[（iii）]

只需在局部证明该陈述。我们可以使用正向标准正交坐标。于是经过一个简单的计算，结论由注 6.12(b) 与 6.25(c) 的局部表示以及命题 6.20 推出，我们把这个计算留给你。

\item[（iv）]

由命题 6.27 与 $ \delta $ 的定义推出

\begin{equation*}
\begin{aligned}
\operatorname{curl}\operatorname{curl} v &= \Theta^{-1}\ast d\Theta\Theta^{-1}\ast d\Theta v = \Theta^{-1}\ast d\ast d\Theta v \\ &= (-1)^{3(2+1)}\Theta^{-1}\delta d\Theta v = -\Theta^{-1}\delta d\Theta v .
\end{aligned}
\end{equation*}

现在结论由注 6.24(a) 推出。 \nobreak\hspace{0.35em}■

\end{subitems}

\end{proposition}

为了展示这套新演算的力量，我们利用余微分和星算子的性质证明了 (i)。当然，我们也可以在正向坐标卡的标准正交坐标中计算。换言之，我们可以假设 $ M $ 是 $ \mathbb{R}^{3} $ 中的开集，且 $ (\cdot\mid\cdot)_{M} $ 是标准度量 $ (\cdot\mid\cdot) $。利用 (6.38) 中的（形式）nabla 算子，我们得到

\begin{equation*}
\begin{aligned}
\nabla\cdot(v\times w) &= \det[\nabla, v, w] \\ &= \partial_{1}(v^{2}w^{3} - v^{3}w^{2}) + \partial_{2}(v^{3}w^{1} - v^{1}w^{3}) + \partial_{3}(v^{1}w^{2} - v^{2}w^{1})
\end{aligned}
\end{equation*}

这是把（形式）行列式按第一行展开得到的。利用乘积法则，我们容易看出最后一行与表达式 $ w\cdot\operatorname{curl} v - v\cdot\operatorname{curl} w $ 一致，正如 (i) 所述。

然而，以 nabla 算子进行的形式运算必须谨慎使用。例如，若用格拉斯曼恒等式形式地计算 $ \operatorname{curl}(v \times w) = \nabla \times (v \times w) $，便会得到错误的命题

\[
\nabla \times (v \times w) = (\nabla \cdot w)v - (\nabla \cdot v)w .
\]

错在哪里？

\exerciseshead

\begin{exercises}
\item

求拉普拉斯–贝尔特拉米算子关于下列各项的表示：

\begin{subitems}
\item

柱坐标 $ (0, 2\pi) \times \mathbb{R} \to \mathbb{R}^{3} $，$ (\varphi, z) \mapsto (\cos\varphi, \sin\varphi, z) $；

\item

例 VII.9.11(f) 中 2 环面 $ \mathrm{T}_{2,1}^{2} $ 的参数化

\[
(0, 2\pi)^{2} \to \mathbb{R}^{3}, \quad (\alpha, \beta) \mapsto ((2 + \cos\alpha)\cos\beta, (2 + \cos\alpha)\sin\beta, \sin\alpha)
\]
\item

当 $ X $ 是 $ \mathbb{R}^{2} $ 中的开集时，$ f \in \mathcal{E}(X) $ 的图像的参数化 $ X \to \mathbb{R}^{3} $，$ x \mapsto (x, f(x)) $。

\end{subitems}

\item

设 $ (M_{j}, g_{j}) $（$ j = 1, 2 $）是黎曼流形，且 $ \partial M_{1} = \varnothing $；并设 $ \pi_{j} $ 表示典范投影 $ M_{1} \times M_{2} \to M_{j} $。证明

\[
\Delta_{M_{1} \times M_{2}} = \pi_{1}^{*}\Delta_{M_{1}} + \pi_{2}^{*}\Delta_{M_{2}} .
\]
\item

设 $ M $ 与 $ N $ 是黎曼流形，$ f: M \to N $ 是等距微分同胚。则对 $ 0 \le r \le m $，证明下图交换：

\bookfigure[0.40]{figures/c7227e598eb977de.png}

\item

设 $ (M, g) $ 是伪黎曼流形。证明下图交换：

\bookfigure[0.35]{figures/b362788716af97fc.png}

并导出关系

\begin{subitems}
\item

$ \operatorname{div} = \ast d\ast\Theta $；

\item

$ \operatorname{curl} = \Theta^{-1}\ast d\Theta $（$ m = 3 $）；

\item

$ \Delta_{M} = \ast d\ast d $，其中 $ \Delta_{M} $ 是 $ M $ 的拉普拉斯–贝尔特拉米算子。

\end{subitems}

\item

设 $ \Omega $ 是 $ \mathbb{R}^{3} $ 中的开集。对 $ E, B, j \in C^{\infty}(\mathbb{R} \times \Omega, \mathbb{R}^{3}) $、$ \rho \in C^{\infty}(\mathbb{R} \times \Omega, \mathbb{R}) $ 及 $ c > 0 $，置

\[
F := \Theta_{e}E \wedge (c\,dt) + \ast(\Theta_{e}B \wedge (c\,dt)), \quad J := \Theta_{e}j - \rho\,dt \in \Omega(\mathbb{R}^{4}_{1,3}),
\]

其中 $ E $、$ B $、$ j $ 视为依赖时间的向量场，$ \rho $ 视为 $ \Omega $ 上依赖时间的函数，而 $ \Theta_{e}: \mathcal{V}(\mathbb{R}^{3}) \to \Omega^{1}(\mathbb{R}^{3}) $ 表示（欧氏）里斯同构。又设 $ dt $ 是 $ \Omega^{1}(\mathbb{R}^{4}_{1,3}) $ 中的第一个标准基向量。现证明以下事实：

\begin{items}
\item

下列命题等价：

\begin{subitems}
\item

$ dF = 0 $；

\item

$ \partial B/\partial t + c\operatorname{curl}E = 0 $ 且 $ \operatorname{div}B = 0 $。
（这些是齐次麦克斯韦方程。）也就是说，2 形式 $ F $ 是闭的，当且仅当向量场 $ E $ 与 $ B $ 满足麦克斯韦方程组中的这两个方程。

\end{subitems}

\item

下列命题等价：

\begin{subitems}
\item

$ dF = 4\pi J $；

\item

$ \partial E/\partial t - c\operatorname{curl}B = 4\pi j $ 且 $ \operatorname{div}E = 4\pi\rho $。
（这些是带源的麦克斯韦方程。）

\end{subitems}

\item

下列命题等价：

\begin{subitems}
\item

$ \Delta_{\mathbb{R}^{4}_{1,3}}F = 0 $；

\item

$ \partial B/\partial t + c\operatorname{curl}E = 0 $，$ \partial E/\partial t - c\operatorname{curl}B = 0 $，$ \operatorname{div}E = 0 $，$ \operatorname{div}B = 0 $。
因此，2 形式 $ F $ 是调和的，当且仅当向量场 $ E $ 与 $ B $ 满足齐次麦克斯韦方程。

\end{subitems}

\item

若 $ dF = 0 $，则 $ j $ 与 $ \rho $ 满足\textbf{连续性方程}

\[
\partial\rho/\partial t + \operatorname{div}j = 0 ,
\]

它也可写成 $ \operatorname{grad}_{\mathbb{R}^{4}_{1,3}}J = 0 $。

\item

下列两个命题等价：

\end{items}

\begin{subitems}
\item

$ F $ 是恰当的；

\item

存在\textbf{向量势} $ A \in C^{\infty}(\mathbb{R} \times \Omega, \mathbb{R}) $ 与\textbf{标量势} $ \Phi \in C^{\infty}(\mathbb{R} \times \Omega, \mathbb{R}) $，使得

\[
\operatorname{curl}A = B \quad \text{且} \quad -\partial A/\partial t - \operatorname{grad}\Phi = E .
\]
\end{subitems}

\item

设 $ X $ 是 $ \mathbb{R}^{3} $ 中的开集且可缩。又设 $ f, g \in \mathcal{E}(X) $。证明：

\begin{subitems}
\item

存在 $ v \in \mathcal{V}(X) $，使得 $ \operatorname{grad}f \times \operatorname{grad}g = \operatorname{curl}v $。

\item

若对 $ x \in X $ 有 $ f(x) \ne 0 $，则存在 $ h \in \mathcal{E}(X) $，使得 $ (\operatorname{grad}f)/f = \operatorname{grad}h $。

\end{subitems}

\item

证明：对 $ f \in \mathcal{E}(M) $ 有

\[
\vec{\Delta}_{M}(f\operatorname{grad}f) = \operatorname{grad}\operatorname{div}(f\operatorname{grad}f) = \Delta_{M}f\operatorname{grad}f + \operatorname{grad}|\operatorname{grad}f|_{M}^{2} + f\operatorname{grad}\Delta_{M}f ,
\]

其中 $ |v|_{M}^{2} := (v \mid v)_{M} $，$ v \in \mathcal{V}(M) $。

\item

证明 $ \alpha \in \Omega^{1}(M) $ 与 $ v, w \in \mathcal{V}(M) $ 满足

\[
d\alpha(v, w) = \mathcal{L}_{v}\langle \alpha, w \rangle - \mathcal{L}_{w}\langle \alpha, v \rangle - \langle \alpha, [v, w] \rangle .
\]
\item

设 $ M $ 与 $ N $ 是 $ m $ 维流形，并设 $ \varphi \in \operatorname{Diff}(M, N) $。证明

\[
\varphi_{*}[v, w] = [\varphi_{*}v, \varphi_{*}w] \quad \text{对 } v, w \in \mathcal{V}(M) .
\]
\item

设 $ T^{2} := S^{1} \times S^{2} \subset \mathbb{R}^{4} $，并设 $ \alpha, \beta \in \Omega^{1}(T^{2}) $ 满足

\[
\alpha := -x^{2}dx^{1} + x^{1}dx^{2}, \quad \beta := -x^{4}dx^{3} + x^{3}dx^{4} .
\]

证明 $ \Delta\alpha = \Delta\beta = 0 $。

\item

证明：对 $ H \in \mathcal{E}(\mathbb{R}^{2m}) $，向量场 $ \operatorname{sgrad}H \in \mathcal{V}(\mathbb{R}^{2m}) $ 的散度为零。

\end{exercises}
